\chapter{The C1 coefficient domain}
\label{ch:c1-short}

Let $K_c,K_t$ be separately labelled real closed extensions of
$\mathbb R$ and let $T=K_c\otimes_{\mathbb R}K_t$.
Source coefficients belong to $K_c$; independent testing coordinates
belong to $K_t$. A coefficient is bounded if it is a finite sum of
products of bounded elements of the two fields. Its residue is obtained
by taking the real residues in both factors. All sums of squares below
are finite sums in $T$.

We prove the complete C1 assertion under the native source hypotheses
of the contact-entry theorem: $d,m,C,A,\alpha>0$,
$\lambda_{\rm ev}=n=0$, $D,\beta\ge0$, $D+\beta>0$, the original
strict native gaps, and positive infinitesimal total scale. Normalize
$d=1$ by a pure-source square. The full target is the original
four-state quadratic parent in Chapter~\ref{ch:cover}, with its
homogeneous value used at projective testing endpoints.
No testing angle is identified with the source contact.

\section{Source coordinates and polynomial operations}
The contact equations determine the source coefficients. Their positive denominators and the lower bound for $E=1-x\eta h$ determine which rescalings remain bounded. We establish these facts before choosing the roots of the certificate.

\subsection{Finite polynomial operations used in the proof}
\label{sec:c1-finite-operations}

All coefficient arrays in this chapter have the following precise
meaning. If $\Pi$ is a displayed sum of retained positive squares,
$s$ is a vector of the remaining state coordinates and
$s=s_0+A\xi$ is the displayed change of state coordinates, put
\[
 \mathcal Q(\xi)=(F-\Pi)(s_0+A\xi),\qquad
 q_0=\mathcal Q(0),\quad q_i=[\xi_i]\mathcal Q,
 \quad H_{ij}={1\over2}\partial_i\partial_j\mathcal Q.
\]
Then exactly $\mathcal Q=q_0+\sum_i q_i\xi_i+\xi^TH\xi$.
Thus the scalar term, both linear state couplings and the entire state
quadratic form are computed from the same polynomial. Multiplying by the
displayed positive pure-source and pure-testing common denominators
makes these coefficient extractions polynomial. All later divisions
are divisions by the specified monic polynomials.

For a polynomial $f$ define its finite difference quotient by
\[
 \Delta_{z,a}f={f(z)-f(a)\over z-a}.
\]
It is the polynomial determined by
$(z^j-a^j)/(z-a)=\sum_{i=0}^{j-1}z^{j-1-i}a^i$.
More generally, a quotient/remainder by $z^r-a$ is determined uniquely
by repeatedly replacing $z^r$ by $a$ and collecting the difference.
Thus every quotient defined below is an explicit finite polynomial
operation, even where a source opening tends to zero. Weighted degree
is checked on the original polynomial before division; division by
a weighted-homogeneous monic polynomial preserves the filtration
with its stated degree shift. The support and weight bounds specified below ensure that the resulting
Gram coefficients remain bounded after rescaling.

\begin{lemma}[Auxiliary square allocation]
\label{lem:c1-square-allocation}
Let $r$ be roots with a bounded symmetric Gram whose real residue
$G_0$ is positive definite. Let $h=(h_1,\ldots,h_s)$ be additional
roots and suppose there are exact certificates
\[
 Q_j(r,h)-h_j^2\in\Sigma T^2,
\]
where every Gram coefficient of $Q_j$ is bounded with zero residue.
Suppose the remaining terms of the polynomial are a bounded quadratic
form in $(r,h)$ whose $r,r$ residue block is $G_0$. Then the complete
polynomial is a tensor sum of squares.
\end{lemma}
\begin{proof}
For a sufficiently large fixed real $M>0$, add
$M\sum_j(h_j^2-Q_j)$ to the remaining Gram. Its real residue is
\[
 \begin{pmatrix}G_0&B_0\\B_0^T&C_0+MI\end{pmatrix}\succ0.
\]
Existence of such an $M$ follows by a Schur complement of the finite
real matrix $G_0$. The strict positive-residue Gram lemma applies.
Restore $M\sum_j(Q_j-h_j^2)$ using the displayed exact certificates.
Only a finite real matrix was inverted in this proof.
\end{proof}

We also use the common product-relation allocation, the restricted
monic contact lift CC1--CC6, and the finite real two-block allocation
UA1--UA3. Their hypotheses are checked in the source cases below.
The arbitrary-scale noncommuting theorem is used only for the original
physical source, after its strong-event source conditions have been
proved. The generic LII and RR polynomial identities are used only
with the explicit C1 substitutions given below.

Whenever a proof below retains a nonnegative coefficient multiplying
a whole square, an unbounded positive excess stays outside its bounded
Gram. A strict leading block is obtained by clipping whole positive
squares at fixed real lower bounds, never by clipping individual
matrix entries. The real constrained-limit argument gives finite real
caps where a source opening makes a retained coefficient unbounded.
Every coefficient correction and both complete first state normals
are included before this argument is applied.

\subsection{The native source partition}
\label{sec:c1-source-partition}
\label{component:pfd0dd7bb10cf}

Use the native coordinates $x=k^2$, $0<k\le1$,
$\eta,\zeta>0$, $r=1-a\eta$, $\theta=t\zeta$, $h=1-\theta$,
$E=1-x\eta h>0$, $Q=1-x\zeta r>0$,
$M=1+xr-t>0$ and $J=1+xh-a>0$.
All satisfy $0<r,\theta<1$ and $0<a,t<2$.
Write I for positive infinitesimal, U for positive bounded with nonzero
residue and O for positive unbounded. Put $e=k\sqrt{\eta/\zeta}$.
The exhaustive source cases are
\[
\begin{array}{c|c|c|l}
k&\zeta&\text{other conditions}&\text{proof below}\\\hline
\mathrm U&\text{any}&\eta,\zeta:\mathrm I&\text{Section~\ref{sec:c1-uniform-projective-join}}\\
\mathrm U&\text{any}&\eta,\zeta\text{ not both I}&\text{Sections~\ref{component:pa0b854a04f49}, \ref{component:p42ee07a067c1}, \ref{component:pf7ffa921793c}}\\
\mathrm I&\mathrm U\text{ or }\mathrm O&\eta>0&\text{Section~\ref{sec:c1-L-partition}}\\
\mathrm I&\mathrm I&\eta:\mathrm I&\text{Section~\ref{component:pf5d2f33f694c}}\\
\mathrm I&\mathrm I&\eta:\mathrm O&\text{Section~\ref{component:p9afbe73899af}}\\
\mathrm I&\mathrm I&\eta:\mathrm U,e:\mathrm I&\text{Sections~\ref{component:pa7e744614f49}, \ref{component:p05a86be0046c}}\\
\mathrm I&\mathrm I&\eta,e,r:\mathrm U&\text{Section~\ref{component:p69995d829bb2}}\\
\mathrm I&\mathrm I&\eta,e:\mathrm U,r:\mathrm I&\text{Section~\ref{component:pccaa793fd0bf}}\\
\mathrm I&\mathrm I&\eta,r:\mathrm U,e:\mathrm O&\text{Section~\ref{component:p890eea5f538b}}\\
\mathrm I&\mathrm I&\eta:\mathrm U,e:\mathrm O,r:\mathrm I&\text{Section~\ref{component:pccaa793fd0bf}}
\end{array}
\]
The source identities below justify every listed implication, including
the ordinary UI/UH alternatives and the bound needed to split C1.1.
The last two centered-$A$ rows share the same bounded-reciprocal
construction. The source field alone determines this partition.

\subsubsection{The $L=x\zeta$ cases}
\label{sec:c1-L-partition}
For $k$ infinitesimal and $\zeta$ noninfinitesimal, put $L=x\zeta$.
When $L$ is unbounded, $Q>0$ forces $r$ infinitesimal and
$1-r=a\eta$, $a<2$ exclude infinitesimal $\eta$. Unit $\eta$
is treated in Section~\ref{component:pd4b92bf7de71}; unbounded
$\eta$ uses Theorem~\ref{thm:c1-unification-bounded-z-all-H}, with $z=1/L$ infinitesimal,
for $u,w$ infinitesimal. The unit-Alice results remain in
Sections~\ref{component:pb1104217f67b},~\ref{component:p909aa5765721}
and~\ref{component:pcc84202479d2}, and the complementary projective
tests remain in Section~\ref{component:p89fd3edb3edd}.
For unit $L$, infinitesimal and unit $\eta$ are covered in
Sections~\ref{component:p080aa62de16d} and~\ref{component:p686209b87229}.
Unbounded $\eta$ has $z$ a positive unit and uses
Theorem~\ref{thm:c1-unification-bounded-z-all-H}; its complementary
testing cases remain in Sections~\ref{component:pb0d9d28d0849}
and~\ref{component:pda4a21937f64}.

When $L$ is infinitesimal, unbounded $\eta$ gives
\[
 S=\sigma+L,\quad\sigma=1/\eta,\quad
 \kappa=x/S,\quad\nu=\sigma/S,\quad\ell=L/S,
 \quad \nu+\ell=1,\quad \zeta=\ell/\kappa.
\]
The identity $\kappa Q_s+\ell M+\nu E=1$, proved below, makes
$\kappa$ bounded and gives the five source cases
\[
\begin{array}{c|c|c|l}
\kappa&\nu&\ell&\text{proof}\\\hline
\mathrm U&\mathrm U&\mathrm U&\text{Sections~\ref{component:p3f02c51074be}, \ref{component:pd26b89cdb3f0}}\\
\mathrm U&\mathrm I&\mathrm U&\text{Section~\ref{component:p578719565e9d}}\\
\mathrm I&\mathrm U&\mathrm U&\text{Section~\ref{component:p2edb55cd1f69}}\\
\mathrm I&\mathrm U&\mathrm I&\text{Section~\ref{component:pb96c9bfcc12d}}\\
\mathrm I&\mathrm I&\mathrm U&\text{Sections~\ref{component:p80b0c5b9879f}, \ref{component:peaa56c1a3c7a}}
\end{array}
\]
In the last row $q_\nu=\nu/\ell$ and $\tau=h/q_\nu$ are
defined only there. BN proves every bounded $\tau$ and UN every
unbounded $\tau$. No separate restriction on $\beta/q_\nu^2$ is made.

If $L$ is infinitesimal and $\eta$ bounded, unbounded $\zeta$ with
unit $\eta$ is Section~\ref{component:p60eca5442075}. Small $\eta$
with unit or unbounded $\zeta$ uses
Sections~\ref{component:p024c821bfb59},~\ref{component:pc720c1e06aca}
and~\ref{component:pe3bbfd3a8151}.
The remaining $\eta,\zeta$ are units. Put
\[
 j=\eta(1+xh)-1+r=\eta J>0,\qquad j_0=\eta+r-1,
 \qquad j-j_0=\eta xh\in\operatorname{Inf}_c.
\]
Unit $j$ is Section~\ref{component:p3df6e6bdd8f4}; small $j$ with unit
$r$ is Section~\ref{component:p48da253a150f}; small $j,r$ is
Section~\ref{component:pbf6de1d7328a}.
Thus every positive source opening and each permitted literal
$D=0$ or $\beta=0$ has a case in the proof.


% Source formulas: p42ee07a067c1
\subsection{The native coefficient formulas and their uniform gap}

\label{component:p42ee07a067c1}

\label{p42ee07a067c1:sec:c1-global-native-cap}

For the physical C1 source, $\lambda=n=0$, $d>0$, $m,C,A,\alpha>0$,
$D,\beta\ge0$ with at most one zero, and the source has its actual
finite noncommuting stationary contact. Use the common global sign
entry and normalize $C=2$. The exact full-parent Alice reflection
permits $0<k\le1$. Put
\[
 x=k^2,\quad r=1-a\eta,\quad\theta=t\zeta,\quad
 h=1-\theta,\quad y=v_c^2=x\eta\zeta,
\]
\[
 K=1+x,\ R=1+xr,\ T=1+xh,\ G=h+\zeta,
 \quad M=R-t,\ J=T-a,\ E=1-x\eta h,\ Q=1-x\zeta r.
 \tag{NC1}
\]
The common exact signs give $0<r,\theta<1$, $E,Q,M,J>0$ and
$0<a,t<2$. There is no size assumption on $\eta,\zeta,e$ or a
positive lower bound on any opening.

\paragraph{Three short positive source identities.}
The original coquotient polynomials obey
\begin{align*}
 P_4&=TE+x(1-r)G,\\
 P_\alpha&=Rh(1+y)+\zeta ME,\\
 P_d&=p_0+\eta p_1,\qquad
 p_0=1+2xh+rx^2G+x\zeta,\quad
 p_1=x\{K\zeta+xh\theta\}.
 \tag{NC2}
\end{align*}
These are obtained by substituting $a\eta=1-r$ and $t\zeta=\theta$
in the unchanged seven native coquotients. In particular $P_d\ge1$.
The original nonnegative $D,\beta$ are precisely the inequalities
\[
 B_5:=K(1+y)(r+\eta)-(\eta T-1+r)P_d\ge0,
 \qquad B_\beta:=Kx(r+\eta)G-\theta P_d\ge0.
 \tag{NC3}
\]
The exact relative coefficients are
\begin{align*}
 \frac C d&=\frac{K(r+\eta)E}{2(1-r)P_d},&
 \frac A d&=\frac{\eta P_4}{2rP_d},\\
 \frac m d&=\frac{Kx(r+\eta)h(1+y)}{2M\zeta P_d},&
 \frac\alpha d&=\frac{x^2\zeta(r+\eta)P_\alpha}{2QP_d},\\
 \frac Dd&=\frac{B_5}{2J\eta P_d},&
 \frac\beta d&=\frac{B_\beta}{2\theta P_d}.
 \tag{NC4}
\end{align*}
The source inequalities make every denominator positive; their sizes
will enter the bounds below.

\begin{lemma}[A uniform native gap]
On every physical native C1 contact oriented by $0<k\le1$,
\[
                         E\ge\frac1{18}.             \tag{NC5}
\]
This includes the literal boundaries $D=0$ or $\beta=0$.
\end{lemma}
\begin{proof}
Define
\[
 \mathcal A=\theta T+K\zeta>0,\qquad
 \mathcal C=RG(1+rxh),
\]
\begin{align*}
 \mathcal B={}&h^2rx^2+hrx(2-h)+h^2x+hrx^2\zeta
             +hx\zeta+hx+h\\
             &+rx\zeta+x\zeta+2\zeta+1.
\end{align*}
Direct multiplication of (NC2)--(NC3), using $x\eta h=1-E$,
gives the polynomial identity
\[
                  hB_5=-\mathcal A E^2+\mathcal B E-\mathcal C.
 \tag{NC6}
\]
Every summand in $\mathcal B$ is nonnegative. Since
$0<x,r,h\le1$ and $h(2-h)\le1$, its six terms without $\zeta$
sum to at most $6$ and the coefficients of $\zeta$ sum to at most
$6$. Hence $\mathcal B\le6(1+\zeta)$. The native inequality
$\theta=t\zeta<2\zeta$ gives $1<h+2\zeta$, so
$1+\zeta<3G$. Also $\mathcal C\ge G>0$. From $B_5\ge0$ and
(NC6),
\[
 G\le\mathcal C\le\mathcal B E\le18GE.
\]
Cancel the positive pure source $G$. This proves (NC5) without
inverting an opening or assuming any of $\eta,\zeta$ bounded.
\end{proof}

\paragraph{An elementary division-free elimination.}
Keep $p_0,p_1$ from (NC2), and set
\begin{align*}
 a_2&=Kx\zeta-Tp_1,&
 a_1&=K(1+x\zeta r)-Tp_0+(1-r)p_1,&
 a_0&=Kr+(1-r)p_0,\\
 b_1&=KxG-\theta p_1,&
 b_0&=KxrG-\theta p_0.
\end{align*}
Then $B_5=a_2\eta^2+a_1\eta+a_0$ and
$B_\beta=b_1\eta+b_0$. Expanding these two expressions, or just
eliminating their quadratic and linear terms successively, gives
\begin{align*}
 &b_1^2B_5-a_2B_\beta^2-(b_1a_1-2a_2b_0)B_\beta\\
 &\qquad=b_1^2a_0-b_1a_1b_0+a_2b_0^2\\
 &\qquad=x^2h^2K^2RG^2
       \{1+2xh+rx^2h^2+(1-r)x\zeta\}.
 \tag{NC7}
\end{align*}
The last equality is a short polynomial identity in $x,r,h,\zeta$;
all five of its input coefficients are given above. No inverse of
$b_1$, no resultant premise, and no localization is used. For reference,
$b_1=xh\{K(1+\zeta)-x\theta^2\}>0$ as well.

\begin{theorem}[Actual unit/unit C1 contacts enter the five-event exit]
If both actual C1 half-angles $k,v_c$ are positive bounded units,
the original source has five consecutive events bounded below relative
to $d$. Consequently its entire original parent has the inherited
all-tests tensor-SOS certificate. All upper event scales and the
retained zero-coefficient boundaries are allowed.
\end{theorem}
\begin{proof}
After the exact Alice orientation, $x,y$ are bounded positive units.
Thus $\eta\zeta=y/x$ is a bounded positive unit. If $\eta$ were
unbounded, $\zeta$ would be infinitesimal. Then $\theta=t\zeta$
would be infinitesimal, contradicting $x\eta(1-\theta)<1$.
If $\zeta$ were unbounded, $\eta$ would be infinitesimal and
$r=1-a\eta$ would have residue one, contradicting $x\zeta r<1$.
Therefore $\eta,\zeta$ are both bounded positive units.

All polynomials in (NC2)--(NC3) now have bounded coefficients and
bounded arguments. If $h$ were infinitesimal, the exact specialization
\[
             B_\beta\big|_{\theta=1}
                    =-1-x\zeta(1-r)\le-1
\]
would contradict $B_\beta\ge0$. Thus $h$ is a positive bounded
unit. The uniform cap gives that $E$ is a unit. Equations (NC2) make
$P_4,P_\alpha,P_d$ positive bounded units: use $P_4\ge E$,
$P_\alpha\ge Rh(1+y)$, and $P_d\ge1$.

The first four equations of (NC4) now show that
$C,A,m,\alpha$ are all strong relative to $d$. Their denominators
are positive bounded quantities; they may be infinitesimal. Suppose
both $D/d$ and $\beta/d$ were infinitesimal, including a possible
literal zero. Their last two bounded denominator products in (NC4)
would make $B_5,B_\beta$ infinitesimal. Every coefficient on the
left of (NC7) is bounded, so that side would be infinitesimal. Its
right side is a positive unit, since $x,h,K,R,G$ are units and the
last bracket is at least one. This is a contradiction.

Thus $D$ or $\beta$ is strong. The original cycle
$(\lambda,D,C,A,\alpha,m,\beta,n)$ has respectively the five
consecutive strong events $(D,C,A,\alpha,m)$ or
$(C,A,\alpha,m,\beta)$. Apply the common five-consecutive-event
theorem, whose stated scope allows arbitrary upper weights and every
independent testing chart. No independent test is identified with
an actual source angle in this source-entry argument.
\end{proof}



% Source formulas: p51451360ff4a
\subsection{Consequences of the native source inequalities}

\label{component:p51451360ff4a}

\label{p51451360ff4a:sec:c1-residual-source-implications}

This section records source implications of (NC1)--(NC7), not further
all-tests certificates. All quantities are in the coefficient field.
Write $f\asymp g$ for two strictly positive quantities whose quotient
and reciprocal quotient are bounded. Set
\[
             L=x\zeta,\qquad S=1+L+y,\qquad \rho=r+\eta.
\]
The native signs and (NC5) imply the uniform estimates
\[
 S\le P_d\le6S,\qquad
 \frac{1+L}{54}\le P_4\le3(1+L),\qquad
 P_\alpha\asymp h(1+y)+\zeta M.                 \tag{NI1}
\]
Indeed $rx^2\zeta<x$ follows from $Q>0$, and
$\eta x^2h\theta=x\theta(1-E)\le1$, giving the $P_d$ bounds
from (NC2). For $P_4$, if $L\le2$ use $E\ge1/18$; if $L\ge2$,
then $r<1/L\le1/2$ and $P_4\ge L/2$. Its upper bound follows
immediately from (NC2). The last comparison uses $1\le R\le2$
and $1/18\le E<1$.
Consequently the unchanged original source satisfies
\begin{align*}
 C/d&\asymp\frac{\rho}{(1-r)S},&
 A/d&\asymp\frac{\eta(1+L)}{rS},\\
 m/d&\asymp\frac{x\rho h(1+y)}{M\zeta S},&
 \alpha/d&\asymp
 \frac{x^2\zeta\rho\{h(1+y)+\zeta M\}}{QS}.
 \tag{NI2}
\end{align*}
No opening has been bounded away from zero in these statements.
In particular, $E$ and $P_4$ do not supply additional hidden small
source denominators on C1.1, C1.4 or C1.5.

\paragraph{C1.1: two unavoidable large weights and its actual angle.}
Here $x,\zeta$ are infinitesimal and $\eta$ is unbounded.
The common original-$D$ cap gives $q=x\eta$ bounded with
residue at most $1/2$. Hence $y=q\zeta$ is infinitesimal:
the actual angle type is necessarily $(I,I)$, independently of $e$.
The exact formulas give
\[
 C/d\asymp\frac{\eta}{1-r},\qquad
 A/d\asymp\frac{\eta}{r},\qquad
 m/d\asymp\frac{e^2}{M},\qquad
 \alpha/d\asymp xq\zeta,\qquad
 \beta/d+\frac12\asymp\frac{e^2}{t}.             \tag{NI3}
\]
Thus $C,A$ are unbounded relative to $d$, and $\alpha$ is weak.
The common bounded polynomial formula for $B_5$ gives $D/d=O(1/\eta)$,
so $D$ is weak too, including its literal zero. In (NI3), $\beta\ge0$
forces $t=O(e^2)$. Therefore when $e$ is infinitesimal, $t$ is
infinitesimal, $M$ is a unit and $m$ is weak. When $e$ is a unit,
$m$ is strong; when $e$ is unbounded, both $m,\beta$ are unbounded.
These are restrictions on the original C1.1 row, not new certificates.
Its positive infinitesimal $\delta_{\rm tot}$ condition is automatic
because $C/d,A/d$ are unbounded.

\paragraph{C1.4: the natural scale is $L=x\zeta$.}
Now $x$ is infinitesimal and $\zeta$ is noninfinitesimal.
The following implications follow uniformly through all retained
coefficient boundaries.
\begin{enumerate}
\item If $L$ is unbounded, $Q>0$ forces $r$ infinitesimal, and
$1-r=a\eta<2\eta$ makes $\eta$ noninfinitesimal. Hence $y=L\eta$
is unbounded: the actual type is $(I,H)$. Here $\zeta$ is unbounded,
$t=\theta/\zeta$ is infinitesimal, and $M$ is a unit. Since
$h y/\zeta=x\eta h<1$, (NC2) gives $P_\alpha\asymp\zeta$.
Equations (NI1)--(NI2) then yield
\[
 C/d\asymp L^{-1},\qquad
 A/d\asymp r^{-1},\qquad
 \alpha/d\asymp L/Q,\qquad
 m/d\asymp (x\eta h)/\zeta\le O(\zeta^{-1}).
 \tag{NI4}
\]
In particular $C,m$ are weak, $A,\alpha$ unbounded, and
$(m/d)/(C/d)=O(x)$ is infinitesimal. The original total-scale
intersection is automatic.

\item If $L$ is a bounded unit, $\zeta$ is unbounded, $M$ is a
unit, $P_\alpha\asymp\zeta$, and $S\asymp1+\eta$. The identity
$1=r+a\eta$ gives $\rho/(1+\eta)$ a positive bounded unit.
Thus
\[
 C/d\asymp(1-r)^{-1},\qquad
 \alpha/d\asymp Q^{-1},\qquad
 A/d\asymp\frac{\eta}{r(1+\eta)},\qquad
 m/d=O(x).
 \tag{NI5}
\]
So $C,\alpha$ are strong and $m$ weak. Moreover $A$ is weak
exactly when $\eta$ is infinitesimal; that case has $r$ a unit and
$D/d$ unbounded, by $B_5\sim1$ in (NC3). The actual Bob angle is
infinitesimal, a unit, or unbounded exactly as $\eta$ is, because
$y=L\eta$. No restriction on $D,\beta$ is added in the other cases.

\item If $L$ is infinitesimal and $\eta$ is bounded, then $y$ is
infinitesimal. The $B_\beta\ge0$ equation bounds
$\theta\le Kx\rho G/P_d=O(x+L)$, so $h,M$ are units.
Now (NI2) gives $C$ strong, $m,\alpha$ weak, and $A$ weak exactly
when $\eta$ is infinitesimal. In that last case $D/d$ is unbounded,
as in the preceding item. The actual type is $(I,I)$.

\item If $L$ is infinitesimal and $\eta$ is unbounded, then
\[
 C/d\asymp\frac{\eta}{(1-r)(1+y)},\qquad
 A/d\asymp\frac{\eta}{r(1+y)}.                    \tag{NI6}
\]
Both are unbounded since $\eta/(1+y)=(\eta^{-1}+L)^{-1}$.
Also $Q$ is a unit and (NI2) bounds
\[
 \alpha/d\le O\!\left(x^2\zeta\eta h+
              \frac{x^2\zeta^2\eta}{1+y}\right)=O(L),
\]
so $\alpha$ is weak. Finally $J\sim1$ and
$P_d/(1+y)\sim1$ directly in (NC2); (NC4) gives
$D/d+1/2\sim1/2$, hence $D$ weak. This extends the
$C,A$-large/$D,\alpha$-weak source pattern of C1.1 into C1.4,
without claiming that its existing conductors apply. Here $y$ may
be infinitesimal, a unit, or unbounded. The total-scale intersection
is again automatic.
\end{enumerate}
These four alternatives exhaust C1.4, but they are necessary source
reductions only. In particular no entire C1.4 leaf is closed.

\paragraph{C1.5 after the actual unit/unit certificate.}
Here $k$ is a unit. The new theorem removes actual $(U,U)$; the
common both-$\eta,\zeta$-infinitesimal theorem is also subtracted.
Two useful remaining source types can be specified more tightly.

If $v_c$ is infinitesimal, $\eta\zeta$ is infinitesimal. Neither
$\eta$ nor $\zeta$ can be unbounded: $\zeta$ infinitesimal forces
$\theta$ infinitesimal and excludes unbounded $\eta$ by $E>0$;
$\eta$ infinitesimal gives $r\sim1$ and excludes unbounded $\zeta$
by $Q>0$. Thus exactly one of $\eta,\zeta$ is a unit in this residual
subcase. The exact stabilizer (AS5) exchanges these two alternatives.
It keeps $k$ a bounded unit, although a second orientation need not
keep $k\le1$. Choose the representative $\eta$ a unit and
$\zeta$ infinitesimal. Then
\[
 m/d,\beta/d\text{ are unbounded},\qquad
 C/d,A/d\text{ are strong},\qquad \alpha/d\text{ is weak}.
 \tag{NI7}
\]
Indeed $\theta=t\zeta$ is infinitesimal, $E$ is a unit,
$P_d,P_4,P_\alpha,r+\eta$ are units, and (NC4) gives
$m/d\asymp1/(M\zeta)$ and $\beta/d+1/2\asymp1/\theta$.
The reflected representative also has a unit $E$: its exact value
is $E'=h(1+y)/G$ in the original coordinates, with $h$ a unit by
the same $\theta=1$ exclusion used in the ordinary-angle proof.
There is no conclusion about the size of $D/d$ in (NI7).

If $v_c$ is unbounded, both $\eta,\zeta$ are noninfinitesimal and
at least one is unbounded. To see the first assertion, an infinitesimal
$\zeta$ would force $h\sim1$ and $\eta$ bounded by $E>0$;
an infinitesimal $\eta$ would force $r\sim1$ and $\zeta$ bounded
by $Q>0$. Two rate regimes are then useful.
\begin{itemize}
\item If $\eta$ is unbounded and $\zeta$ a unit, then $h$ is
infinitesimal. The exact alternative numerator identity
\[
 B_\beta=(x\eta h)\{K(1+\zeta)-x\theta^2\}
             +xrGT-\theta(1+2xh+x\zeta)             \tag{NI8}
\]
shows that $x\eta h$ is a unit: otherwise its residue at
$\theta=1$ would be $-1-x\zeta(1-r)<0$.
Thus $P_\alpha$ is a unit and $P_d\asymp\eta$.
Equations (NC4) give
\[
 C/d\asymp(1-r)^{-1},\quad A/d\asymp r^{-1},\quad
 m/d\asymp M^{-1},\quad\alpha/d\asymp Q^{-1},
 \quad D/d,\beta/d\text{ weak}.                    \tag{NI9}
\]
For the last assertion use $J\sim1$, $\theta\sim1$ and
$P_d/(Ky)\sim1$ in the original $D/d+1/2,\beta/d+1/2$
formulas. The total-scale condition $\delta_{\rm tot}$ infinitesimal
is therefore equivalent on this typed branch to at least one of
$r,1-r,M,Q$ being infinitesimal. If all four are units, the inherited
positive-total-scale theorem already applies.

\item If $\zeta$ is unbounded, $r$ is infinitesimal and $M$ a
unit. The noninfinitesimal $\eta$ obeys $\eta h\asymp1$.
For if $\theta$ is infinitesimal this follows from $E>0$ and
$h\sim1$; otherwise divide (NI8) by $\zeta$, use
$xrG/\zeta=O(r)$, and obtain that $x\eta h$ has a positive-unit
lower bound. Therefore (NC2) and (NC4) give
\[
 C/d\asymp\zeta^{-1},\quad m/d\asymp\zeta^{-1},\quad
 A/d\asymp r^{-1},\quad\alpha/d\asymp\zeta/Q.
 \tag{NI10}
\]
Thus $C,m$ are weak and $A,\alpha$ unbounded. No size restriction
on $D,\beta$ is inferred. This branch already has infinitesimal
$\delta_{\rm tot}$.
\end{itemize}
All implications retain the original physicality and exact source openings.



% Source formulas: pa133fae35e52
\subsection{The bound on $x\eta$}

\label{component:pa133fae35e52}

Throughout, $d$ is the original positive signed coefficient. The event
coefficient at the fourth cycle position is denoted $A_{\rm ev}$, and
$D_{\rm ev}=e_5$. These must not be confused with actual state coordinates.
The source has the canonical omissions $\lambda=n=0$ and is normalized
by $c=e_3=2$. We use the canonical actual-contact sign entry, including
its covered single-zero $D_{\rm ev}$ or $\beta$ boundaries. Explicitly,
the retained $m,c,A_{\rm ev},\alpha$ are positive, and $D_{\rm ev},\beta$
are positive or exactly one of them is zero. The noncoercive further-zero
faces are already separate exits. The sign and nonsingularity input is the native-entry theorem EN in Chapter~\ref{ch:entry}, proved for every finite noncommuting physical contact in this exact coefficient domain.

\paragraph{Exact transport from the actual contact.}
Write its positive-angle contact as
\[
 u_c=k,\quad v_c=\ell,\qquad
 \psi_c=(S,k(\theta-1),kSr,1),\quad
 S>0,\quad 0<r,\theta<1.
\]
The exact C1 reflection lets us orient $0<k\le1$. Define source-only
parameters
\[
 \eta=\frac{\ell}{kS},\quad \zeta=\frac{S\ell}{k},\quad
 f=\sqrt\eta,\quad g=\sqrt\zeta,\quad
 e=\frac{k}{S}=\frac{kf}{g},\quad
 a=\frac{1-r}{\eta},\quad t=\frac\theta\zeta.
                                                               \tag{CE1}
\]
Thus the classifier $\mathfrak T=k/(S\ell)$ is exactly $\zeta^{-1}$,
and its old state-ratio variable $B$ is exactly $\eta^{-1/2}$.
Put $\kappa=k^2$, $K=1+\kappa$. The original amplitude signs imply
\[
 R=1+\kappa r,\quad T=1+\kappa(1-\theta),\quad
 M=R-t>0,\quad J=T-a>0,
\]
\[
 E=1-\kappa\eta(1-\theta)>0,\qquad
 Q_c=1-\kappa\zeta r>0.                              \tag{CE2}
\]
Indeed the positive second amplitude is $k\zeta M$, the negative
fifth amplitude is $-\ell J$, the state gap is $SE$, and the remaining
state gap is $Q_c$. In particular $0<a,t<2$, without any shape-unit
hypothesis.

The seven native coquotients used by the uniform theorem are valid
under (CE1), not just asymptotically. Their constructed contact is
\[
 (W,X,Y,Z)=(\Delta k,-k(1-t\zeta),\Delta k^2(1-a\eta),1),
 \qquad \Delta=e^{-1}=S/k,
\]
which is exactly the original contact. The native contact identities
prove all six kernel and angular equations. To make the nonsingularity
premise explicit, put $Y_0=Sr$, $h_0=S+k\ell(\theta-1)$,
$R_2=\ell(S+k^2Y_0)-k\theta$, and
$R_5=k(S-Y_0)+k^2\ell(\theta-1)-\ell$. The common selected
six-column contact system, on $(d,m,A_{\rm ev},D_{\rm ev},\alpha,\beta)$
with $c=2$, has determinant
\[
 -8192\theta k^2\ell(1+k^2)^2(1+\ell^2)^2\mathcal D,
\]
\[
 \mathcal D=Y_0(1+k^2)(kY_0+\ell)(1-k\ell Y_0)h_0R_2(-R_5)>0.
\]
Each factor is strictly positive by the stated source sign theorem.
Thus the normalized coefficients equal the unique original source
coefficients in that common system. This
establishes full original-coefficient transport, including the exact
zero openings; it does not use an archived completion label.

\paragraph{A bounded polynomial consequence of the original fifth event.}
Assume now
\[
                 k\text{ infinitesimal},\qquad
                 \zeta=\mathfrak T^{-1}\text{ infinitesimal}.     \tag{CE3}
\]
Set $q=\kappa\eta$. Since $t<2$ and $E=1-q(1-t\zeta)>0$, $q$ is
bounded and $\bar q\le1$. The original fifth-event formula is
\[
 D_{\rm ev}=\frac{a}{Jc_3}(2c_5-J\eta c_d).
\]
Using $L_c=r$, $H_c=r+\eta$ and $V_c=q\zeta$, its sign is precisely
\[
 \mathcal B_5:=K(1+q\zeta)(r+\eta)-J\eta P_d\ge0.       \tag{CE4}
\]
Although $\eta$ need not be bounded, this whole expression has the
following bounded polynomial form. Define
\begin{align*}
 H={}&2+\kappa r+\zeta[K-(\kappa+2)t]
     +q\zeta(-\zeta t^2+t+1)
     +\kappa\zeta(1-r)(t-1),\\
 P_d={}&1+q\zeta+\kappa H.
\end{align*}
Then exact cancellation using $a\eta=1-r$ and $q=\kappa\eta$ gives
\[
 \mathcal B_5=Kr(1+q\zeta)+(1-r)P_d
       +q\{\zeta t(1+q\zeta)-(K-\kappa\zeta t)H\}.     \tag{CE5}
\]
Every variable on the right is bounded. Hence (CE4) gives
\[
       1-2\bar q\ge0,\qquad 0\le\bar q\le\tfrac12.     \tag{CE6}
\]
This is obtained before dividing by $1-q$ or by a coefficient opening.
In particular $E$ is a positive unit with residue at least $1/2$.

The fourth coquotient polynomial likewise has the bounded expression
\[
\begin{split}
 P_4={}&K(1-q)+\kappa(1-r)-\kappa\zeta t\\
 &+q\zeta(-\kappa\zeta t^2+2\kappa t+t)
   +\kappa\zeta(1-r)(1-t).
\end{split}                                                       \tag{CE7}
\]
Thus $\bar P_d=1$ and $\bar P_4=1-\bar q\ge1/2$.
These exact low-degree identities are the physicality input needed
for the entry, rather than an estimate on a truncated coefficient.



% Source formulas: p0acdb66bba2f
\section{Contact identities and complementary measurement charts}
The scalar contact polynomial supplies the first square roots. Polynomial division then expresses state couplings through these roots. The complementary measurement charts follow directly from the positive event terms and will be used throughout the remaining proof.

\subsection{Polynomial source coordinates and exact contact balances}

\label{component:p0acdb66bba2f}

\label{p0acdb66bba2f:app:c1-source-definitions}

The letter $d$ in the contact--conductor proof denotes the new parameter
$\delta$, not the Bell coefficient. Here the latter is denoted $d_{\rm B}$,
and the original event coefficient $e_5$ is denoted $D_{\rm e}$.
The testing factor is always $D=1+u^2$. The conductor monomial
$C=\delta b^2uw$ is distinct from the clearing multiplier
$\mathcal C=K^6Q^4T_n$.

\paragraph{Scope of the polynomial identities.} The seven coquotient formulas and the contact identities (SD9)--(SD11) are exact for arbitrary positive native parameters. The size assumptions below are needed only for the displayed bounded-Gram application.
\paragraph{Domain and seven native coquotients.}
The parameters $k,f,g>0$ belong to the source field; put $e=kf/g$ and $\Delta=e^{-1}$. The formulas are exact before a size condition is imposed. The shapes $a,t$ are positive bounded elements;
\[
 R=1+k^2(1-af^2),\qquad T=1+k^2-k^2g^2t,
 \qquad M=R-t>0,\quad J=T-a>0.
\]
All original event coefficients are assumed nonnegative. In particular,
the formulas below retain the requirements $D_{\rm e},\beta\ge0$.
No relation compares a testing rate with a coefficient rate.
Put
\[
 \kappa=k^2,\quad \eta=f^2,\quad \zeta=g^2,
 \quad V_c=\kappa\eta\zeta,\quad
 L_c=1-a\eta,\quad H_c=1+(1-a)\eta,\quad Q_c=1-\kappa\zeta+aV_c.
\]
Define
\begin{align*}
P_4={}&1+\kappa+\eta\kappa(a-\kappa-1)-\zeta\kappa t
 +V_c[-\zeta\kappa t^2+2\kappa t+a(1-t)+t],\\
P_\alpha={}&1+\kappa-a\eta\kappa
 +\zeta[(1+\kappa)-(2+\kappa)t]
 +V_c[-\zeta t^2+a(t-1)+t],\\
P_d={}&(1+\kappa)^2-a\eta\kappa^2
 +\zeta\kappa[(\kappa+1)-(\kappa+2)t]\\
 &+V_c[-\zeta\kappa t^2+\kappa a(t-1)+\kappa t+\kappa+1].
\end{align*}
The complete seven coquotients are
\begin{align*}
c_2&=2(1+\kappa)(1-t\zeta)L_c(1+V_c)H_cQ_c,\\
c_b&=2(1+\kappa)L_c[1+(1-t)\zeta]H_cQ_c,\\
c_3&=2(1+\kappa)L_cH_cQ_c(1-\kappa\eta+tV_c),\\
c_5&=2(1+\kappa)L_cH_cQ_c(1+V_c),\\
c_4&=2Q_cP_4,\qquad c_\alpha=2L_cH_cP_\alpha,
 \qquad c_d=4L_cQ_cP_d.
\end{align*}
At $\eta=\zeta=0$, the first six equal $2(1+k^2)$ and $c_d=4(1+k^2)^2$.
Here $c_b$ is the original $c_\beta$, and the subscript on $c_d$ is a
name rather than the conductor parameter.

\paragraph{Original source and full parent.}
The original event coefficients and multiplier are
\[
\begin{gathered}
\lambda=n=0,\qquad c=2,\qquad
 d_{\rm B}=\frac{2a\eta c_d}{c_3},\qquad
 m=\frac{2ae^2c_2}{Mc_3},\qquad A=\frac{2a\eta^2c_4}{c_3},\\
\alpha=\frac{2a\Delta^2k^6\eta^2c_\alpha}{c_3},\qquad
D_{\rm e}+\frac{d_{\rm B}}2=\frac{2ac_5}{Jc_3},\qquad
\beta+\frac{d_{\rm B}}2=\frac{2ae^2c_b}{tc_3},\qquad
\Lambda=\frac{RTc_3}{2a}.
\end{gathered}                                                    \tag{SD1}
\]
Let $F$ denote the original complete parent, including its signed core.
For a testing state $(W,X,Y,Z)$ define
\[
\begin{gathered}
 D=1+u^2,\quad D_v=1+v^2,\quad B=X+uZ,\quad Q_A=W+uY,\\
 r=uW-Y,\quad r_2=vQ_A-B,\quad r_5=r-v(Z-uX).
\end{gathered}
\]
With $H_2=4Tc_2$, $H_b=4Tc_b$, $H_3=4Rc_3$, $H_5=4Rc_5$,
the exact entire parent is
\begin{align*}
\Phi:=\Delta^2\Lambda F={}&H_2\frac RM r_2^2
 +H_b\frac RtD_v B^2+\Delta^2H_3\frac TaD_vr^2
 +\Delta^2H_5\frac TJr_5^2\\
 &+4RT\Delta^2\eta^2c_4DD_vY^2
 +4RT\Delta^4k^6\eta^2c_\alpha D(vY-Z)^2\\
 &+2RT\Delta^2\eta c_dD\{(vY-Z)^2-D_vZ^2\}.
                                                               \tag{SD2}
\end{align*}
Thus every original source event is accounted for.

\paragraph{Retained terms and state coordinates.}
In the remaining chart $Z=1$ with $u,v,W,X,Y$ infinitesimal, put
\[
 p=t/R,\quad h=a/T,\quad q_A=v(W+uY),\quad q_B=v(1-uX),
 \quad N=B-pq_A,\quad M_B=r-hq_B.
\]
Retain the four nonnegative terms
\[
 \Pi_1=H_2\frac tM N^2,\quad
 \Pi_2=H_b\frac MtD_vN^2,\quad
 \Pi_3=\Delta^2H_3\frac JaD_vM_B^2,\quad
 \Pi_4=\Delta^2H_5\frac aJM_B^2.                         \tag{SD3}
\]
Their subtraction is the exact identity
\begin{align*}
G=\Phi-\sum_j\Pi_j={}&H_2\{N^2-2q_AN+(1-p)q_A^2\}
 +H_bD_v\{N^2+2q_AN+pq_A^2\}\\
 &+\Delta^2H_3D_v\{M_B^2+2q_BM_B+hq_B^2\}
 +\Delta^2H_5\{M_B^2-2q_BM_B+(1-h)q_B^2\}
 +\mathcal W,\\
\mathcal W={}&4RT\Delta^2\eta^2c_4DD_vY^2
 +4RT\Delta^4k^6\eta^2c_\alpha D(vY-1)^2\\
 &+2RT\Delta^2\eta c_dD\{(vY-1)^2-D_v\}.
                                                               \tag{SD4}
\end{align*}
The testing coordinate $w=(W+uY)/D$ has the full inverse
\[
\begin{gathered}
 X=N+pDvw-u,\qquad q_B=v(D-uN-puDvw),\\
 W=w+\frac{u(M_B+hq_B)}D,\qquad
 Y=uw-\frac{M_B+hq_B}D.
\end{gathered}                                \tag{SD5}
\]
Its only testing denominator is the positive unit $D$.
On the negative-core branch, sign gauges give $u,v>0$, and $w>0$:
after eliminating the two positive core squares the remaining scalar is
\[
 \frac{uvDD_v}{\mathcal D}\{Duv w^2-2DD_vw+uv\},\qquad
 \mathcal D=u^2v^4+u^2v^2+u^2+v^4+2v^2+1>0,
\]
which is nonnegative if $w\le0$. The eliminated matrix is
$\left(\begin{smallmatrix}1+v^2+u^2v^2&uv\\uv&(D+v^2)/D\end{smallmatrix}\right)\succ0$.
The nonnegative-core branch and literal axes already give direct SOS.
Thus the radicals in the conductor proof separate between the two fields.

\paragraph{Short native proof of the exact contact identities.}
Write $V=1+V_c$ and $K=1+\kappa$ in this paragraph. The factored
coquotients above give immediately
\[
 c_b-c_2=\zeta c_3,\qquad c_3-c_5=-\kappa\eta c_b,
 \qquad c_2=(1-t\zeta)c_5.                            \tag{SD9}
\]
Hence $Vc_b-c_2=\zeta c_5$ and
$Vc_3-c_5=-\kappa\eta c_2$.
The remaining three balances are the following short polynomial
identities; substituting the displayed $P_4,P_\alpha,P_d$ proves them
by multiplication and collecting terms:
\begin{align*}
\mathfrak Y:={}&2\kappa c_2+2VL_cc_4
                 -2\kappa^2\zeta Q_cc_\alpha-Q_cc_d=0,\\
\mathfrak S:={}&\kappa L_cc_2+c_5+VL_c^2c_4+\kappa Q_c^2c_\alpha
                 +\tfrac12(\kappa\zeta L_c^2-2L_c-\eta)c_d=0,\\
\mathfrak V:={}&\kappa Rc_2+t\kappa V_cc_b+aV_cc_3+Tc_5
 +KV_cL_c^2c_4-\kappa^2\zeta KQ_cL_cc_\alpha
 -\tfrac K2(Q_cL_c+\eta)c_d=0.                        \tag{SD10}
\end{align*}
For example, after the common factor $4L_cQ_c$ is removed,
the first balance is just
\[
 \kappa K(1-t\zeta)VH_c+VP_4
       -\kappa^2\zeta H_cP_\alpha-Q_cP_d=0.
\]
The two other balances are equally direct identities in the same three
short polynomials; no expanded scalar or ideal calculation enters.

The exact contact, in the original coordinates of (SD2), is
\[
 u=k,\quad v=\Delta k^2\eta,\qquad
 (W,X,Y,Z)=(\Delta k,-k(1-t\zeta),\Delta k^2L_c,1).      \tag{SD11}
\]
At this point
\[
 r_2=Mk\zeta,\quad B=tk\zeta,\quad r=av,\quad
 r_5=-Jv,\quad vY-1=-Q_c,\quad D=K,\quad D_v=V.
\]
Set $\Psi=\Phi/(4RT)$ and let $E_s=\tfrac12\partial_s\Psi$,
with all other original testing coordinates held fixed.
Direct differentiation of the seven terms in (SD2), followed by (SD9),
gives at (SD11)
\begin{align*}
 \Psi&=K\zeta^2\mathfrak S,\\
 E_W&=kv\{\zeta c_2+\Delta^2(Vc_3-c_5)\}=0,\\
 E_X&=k\zeta\{Vc_b-c_2-\zeta c_5\}=0,\\
 E_Y&=\tfrac12K\Delta^3\kappa\eta^2\mathfrak Y,\\
 E_u&=k\zeta^2\{\mathfrak S-c_2+(1-t\zeta)c_5\},\\
 E_v&=v\Delta^2\mathfrak V.
\end{align*}
They all vanish by (SD9)--(SD10). Homogeneity in $(W,X,Y,Z)$ also gives
$E_Z=0$. This proves contact and full first-order stationarity directly
in the native source.



% Source formulas: p58894e12bcef
\subsection{A finite four-monomial circuit}

\label{component:p58894e12bcef}

Let
\[
a_1=\sqrt{wv},\quad a_2=\sqrt{\delta uv},\quad
a_3=b\sqrt{\delta kuw},\quad a_4=\delta k^{3/2}b.
\]
The radicals are separated coefficient/testing radicals. Put
\[
h_{13}=a_1^2-a_3^2,\quad h_{23}=a_2^2-a_3^2,\quad
h_{34}=a_3^2-a_4^2,
\qquad q_1=a_1a_2-a_3a_4,\quad q_2=a_1a_3-a_2a_4.
\]
Averaging the three-square arithmetic--geometric-mean identities for
pairings $12|34$ and $13|24$ gives exactly
\begin{equation}
\sum_{i=1}^4a_i^4-4a_1a_2a_3a_4
=h_{13}^2+h_{23}^2+h_{34}^2-h_{13}h_{23}+h_{23}h_{34}
 +q_1^2+q_2^2.\label{p58894e12bcef:eq:compact-circuit}
\end{equation}
The contrast block is
\[
G=\begin{pmatrix}1&-1/2&0\\-1/2&1&1/2\\0&1/2&1\end{pmatrix}\succ0;
\qquad \det G_{[1]}=1,\quad\det G_{[2]}=3/4,\quad\det G=1/2.
\]
Unlike a single pairing, the two pairings control all three independent
contrasts of the four squared monomials.

Write $K=1+k^2$, $D=1+u^2$ and
\[
g_0=\delta(v-b^2uw),\qquad g_4=uwv-\delta^2k^4b^2.
\]
The leading scalar used in the complete source is
\[
P=D(w^2+\delta^2)v^2+\delta^2Kb^4u^2w^2+
\delta^4Kk^6b^4-2\delta^2K^2b^2uwv
=g_0^2+g_4^2+\sum_i a_i^4-4\prod_i a_i.
\]
There is a second exact simplification:
\begin{equation}
g_4=u h_{13}+k h_{23}-ku g_0+k h_{34}.\label{p58894e12bcef:eq:g4-conductor}
\end{equation}
If $k,u$ are infinitesimal, the four conductor coefficients
are infinitesimal. Otherwise retain $g_4^2$ as a separate square.
The two-pairing contrast matrix displayed above is strict in either case.


% Source formulas: p466b31015532
\subsection{The restricted three-monomial contact lift}

\label{component:p466b31015532}

\label{p466b31015532:sec:three-circuit}

Work in the same tensor ring $T=K_c\otimes_{\mathbb R}K_t$ and bounded
subring $B$ as throughout this package. Every inverse and radical below
is of a pure-field or positive separated element. Let $k>0$ be a
pure-source infinitesimal and $s>0$ a positive pure-source bounded unit.
Let $u>0$ be a pure-testing infinitesimal and $t>0$ a positive
pure-testing bounded element (possibly infinitesimal). Define
\[
 A_1=u^2t,\qquad A_2=\frac{ku t^2}{s},\qquad A_3=k^2s,
 \qquad G=kut,\qquad d_i=A_i-G.
\]
Then $A_1A_2A_3=G^3$. Two exact conductor identities are
\begin{equation}\label{p466b31015532:T3:conductors}
 G(u-k)=k d_1,\qquad G(t-s)=s d_2.
\end{equation}
They permit a unit testing ratio $t$ without comparing a mixed tensor.

Put $Q=G\sum_i d_i$ and $L=\sum_i A_i^2-3G^2$.
For $a_i=A_i^{1/3}$, the elementary cubic identity gives
\begin{align}\label{p466b31015532:T3:sos}
 Q&=\frac12\sum_{1\le i<j\le3}\sum_{\ell=1}^3
       \{\sqrt{G a_\ell}(a_i-a_j)\}^2,\\
 L&=\sum_i d_i^2+2Q.\nonumber
\end{align}
Thus $L$ has a positive constant diagonal Gram on a finite list of
three difference roots and nine cubic-identity roots. No assertion
that every nonnegative tensor is SOS is used.

\paragraph{Exact finite lift.}
Let $\lambda\in B_c$ have positive real residue, and consider
\[
 F=\sum_{i=1}^3(\lambda+\epsilon_i)A_i^2
       -(3\lambda+\epsilon_g)G^2,
 \qquad h=\sum_i\epsilon_i-\epsilon_g.
\]
Write $\Delta_u=u-k$, $\Delta_t=t-s$. Assume explicit finite identities
\begin{align}
 \epsilon_i-\epsilon_{i0}&=K_{iu}\Delta_u+K_{it}\Delta_t,
       &\epsilon_i,\epsilon_{i0}&\in B_{\rm Inf},\label{p466b31015532:T3:coeff}\\
 h&=b_u\Delta_u+b_t\Delta_t+
          \sum_{j,\ell\in\{u,t\}}H_{j\ell}\Delta_j\Delta_\ell,
       &b_u,b_t&\in K_c,\label{p466b31015532:T3:taylor}\\
 2(\epsilon_{10}-\epsilon_{30})+k b_u&=0,
 &2(\epsilon_{20}-\epsilon_{30})+s b_t&=0.\label{p466b31015532:T3:jets}
\end{align}
Require $K_{iu},H_{uu},H_{ut},H_{tu}\in B$ and
$K_{it},H_{tt}\in B_{\rm Inf}$. The distinction is essential:
the second conductor has a unit coefficient, so boundedness alone of
the second-coordinate arrays does not suffice.

Then $F$ is tensor SOS. Indeed, expanding $A_i=G+d_i$ and applying
\eqref{p466b31015532:T3:conductors} and \eqref{p466b31015532:T3:jets} yields the exact identity
\begin{align}\label{p466b31015532:T3:error}
 F-\lambda L={}&\sum_i\epsilon_i d_i^2+2\epsilon_{30}Q
 +2\sum_i(kK_{iu}d_1d_i+sK_{it}d_2d_i)\\
 &+k^2H_{uu}d_1^2+ks(H_{ut}+H_{tu})d_1d_2
          +s^2H_{tt}d_2^2.\nonumber
\end{align}
In detail, the two bare-$G$ first-order terms are cancelled by
$G^2b_u\Delta_u=kb_uGd_1$ and
$G^2b_t\Delta_t=sb_tGd_2$. All remaining coefficients in
\eqref{p466b31015532:T3:error} are bounded infinitesimals. The strict diagonal
Gram for $\lambda L$ in \eqref{p466b31015532:T3:sos} therefore remains strictly
positive in real residue. The bounded-Gram lemma supplies a finite
tensor SOS. This argument uses no power series or mixed inverse.

\paragraph{Normal extension.}
Suppose a coupling to a normalized hard state normal has the exact form
$\ell=\sum_i c_iA_i$ with $c_i\in B_{\rm Inf}$,
$\sum_i c_{i0}=0$, and
\[
 \sum_i(c_i-c_{i0})=J_u\Delta_u+J_t\Delta_t,
 \qquad J_u\in B,\quad J_t\in B_{\rm Inf}.
\]
Then
\[
 \ell=\sum_i c_i d_i+kJ_u d_1+sJ_t d_2
\]
has infinitesimal coefficients on the same root list. It can be joined
to a hard normal block with strictly positive real residue. Extra
already-controlled scalar roots may be appended, provided their
coupling coefficients are likewise infinitesimal. The source-specific
arrays and hard-block bounds remain separate proof obligations.

\paragraph{Small-$L$ coordinates.}
For $x=k^2$, $s=\sqrt{\eta\zeta}$ and $t=v/u$, this gives
\[
 (A_1,A_2,A_3,G)=
 \left(uv,\sqrt{\frac{x}{\eta\zeta}}\frac{v^2}{u},
       \sqrt{\eta\zeta}\,x,\sqrt{x}\,v\right).
\]
The lemma applies only when the stated pure-field size assumptions and
all finite scalar and normal identities have been proved. A leading
three-term AM--GM form alone does not close a source domain.



% Source formulas: p89fd3edb3edd
\subsection{The complementary projective testing charts}

\label{component:p89fd3edb3edd}

\label{p89fd3edb3edd:sec:r6-c1-common-remote}
Normalize the original parent by its positive pure-source coefficient $d$.
Write $D_z=1+z^2$, and write the signed core as $2K$, where
\[
 K=(uW+uvX-Y-vZ)^2+D_u(vY-Z)^2+D_v(X+uZ)^2-D_uD_vZ^2.
 \tag{CE1}
\]
The event terms used below are
\[
4m(vW-X+uvY-uZ)^2,\quad4CD_v(uW-Y)^2,
\quad4AD_uD_vY^2,\quad4\alpha D_u(vY-Z)^2.
\]
Every omitted event and every nonnegative excess is retained as its exact
original square. Source and testing fields remain independent. ``Strong''
means bounded below by a positive real number; ``unbounded'' means larger
than every real number. Neither assertion constrains upper scales.

\begin{theorem}[Two adjacent unbounded events, with complete raw frames]
If $C,A$ are unbounded, the full original parent is tensor SOS at every
projective test for which $u$ is not infinitesimal, and every state.
No physicality assumption beyond nonnegative event weights is needed for
this particular assertion. By the exact cyclic relabeling below, if
$A,\alpha$ are unbounded, the conclusion holds whenever $v$ is not
projectively unbounded.
\end{theorem}
\begin{proof}
Select $C=A=\mu$, with a positive real $\mu$ chosen after the real testing
residues; unboundedness permits every such choice. The following are exact
full quadratic forms, before any restriction to a state ray or contact.
They therefore retain every first state normal and hard-normal coupling.

For $u$ a bounded nonzero unit and $v$ bounded, use the four raw roots
\[
 (W,Y,B,V),\qquad B=X+uZ,\quad V=vZ.
\]
The entire core equals
\[
 K=(uW-Y+uvB)^2+D_vB^2+D_uv^2Y^2+D_uu^2V^2
       -2D_uuV(W+vB).                         \tag{CE2}
\]
The two selected event forms remain $4\mu D_v(uW-Y)^2$ and
$4\mu D_uD_vY^2$. In their common real kernel $W=Y=0$, the core has
matrix on $(B,V)$
\[
 \begin{pmatrix}1+v^2D_u&-D_uuv\\-D_uuv&D_uu^2\end{pmatrix},
 \qquad \det=D_uu^2>0.                       \tag{CE3}
\]
Here and in residue matrices, bounded testing letters mean their real
residues. Thus a sufficiently large finite $\mu$ gives a strictly positive
residue on the entire four-root space, even when $v=0$ literally.

For $p=1/u$ infinitesimal and $v$ bounded, homogenize by $p^2$ and take
\[
 (W,Y,B,R),\qquad B=pX+Z,\quad R=vX.
\]
The exact full core is
\[
 K_p=B^2-2pvBW-2vBY+D_pR^2+2D_pRW
                  +(W-pY)^2+D_pv^2Y^2.       \tag{CE4}
\]
The selected terms are $4\mu D_v(W-pY)^2$ and
$4\mu D_pD_vY^2$. At $p=0$ their sum with $2K_p$ is
\[
 4\mu D_v(W^2+Y^2)+2(W+R)^2+2(B-vY)^2,       \tag{CE5}
\]
strict for every $\mu>0$. The formula allows literal $p=0$ and $v=0$;
the raw roots need not form independent state coordinates.

For projectively unbounded $v$, put $w=1/v$, homogenize by $w^2$, and
use $(W,X,Y,H)$ with $H=wZ$. The exact full core is
\begin{align*}
 K_w={}&u^2H^2-2u^2HY-2uHW+2uwHX+u^2w^2W^2\\
 &+2u^2wWX-2uw^2WY+(D_u+w^2)(X^2+Y^2)-2uwXY.
                                                        \tag{CE6}
\end{align*}
The selected terms are $4\mu D_w(uW-Y)^2$ and
$4\mu D_uD_wY^2$. If $u$ is a bounded nonzero unit, at $w=0$ their
common kernel has $W=Y=0$ and $K_w=D_uX^2+u^2H^2$ strictly positive.
A sufficiently large real $\mu$ again works. If also $p=1/u$ is
infinitesimal, homogenize (CE6) by $p^2$. The exact full core is
\begin{align*}
 K_{pw}={}&H^2-2pHW+2pwHX-2HY+w^2W^2+2wWX-2pw^2WY\\
 &+(D_p+p^2w^2)(X^2+Y^2)-2pwXY.              \tag{CE7}
\end{align*}
At $p=w=0$ the full selected residue is
\[
 4\mu(W^2+Y^2)+2X^2+2(H-Y)^2>0.             \tag{CE8}
\]
All these exact Gram coefficients are bounded. Strictness in (CE3) and
its counterpart following (CE6) implies a finite cap by the elementary
real Schur-complement argument: the event Gram is positive definite on
the complement of its kernel, and the core is positive definite on the
kernel. No source/testing rate comparison is involved. Apply the bounded
positive-residue Gram lemma, and restore every retained original square.

For the second assertion, compose the common Alice reflection with
the swap/reflection. The resulting one-step inverse cycle rotation is
\[
 (u',v')=(1/v,u),\quad
 (W',X',Y',Z')=(Y+vZ,-W-vX,Z-vY,vW-X),          \tag{CE9}
\]
and sends event index $i$ to $i-1$ in
$(\lambda,D,C,A,\alpha,m,\beta,n)$. In particular original $A,\alpha$
become new $C,A$. Its full-parent identity is
\[
 v^2F(\mathrm{rot}^{-1}c;1/v,u;\psi')=D_vF(c;u,v;\psi).
                                                        \tag{CE10}
\]
The state map has determinant of absolute value $D_v^2$, and its
homogeneous version at $v=\infty$ is invertible as well. Thus
$u'=1/v$ not infinitesimal is exactly $v$ not projectively unbounded.
Only pure testing factors are restored. Both signs, all literal endpoints,
and complex states (by real/imaginary splitting) are included.
\end{proof}

\begin{theorem}[Four consecutive strong events]
For an original physical parent with $C,A,\alpha,m$ strong, the full
parent is tensor SOS at every test except possibly $u$ a bounded nonzero
unit and $v$ projectively unbounded. All original upper coefficient
scales and all nonnegative unselected coefficients are allowed.
\end{theorem}
\begin{proof}
Cap these four events at any positive real $\mu$ smaller than all four
weights, and retain every excess. Noncommuting-residue tests use the
common physical-source theorem, applied to the original full parent.
For commuting-residue tests use the capped parent.

First suppose $v$ is bounded. On the three open commuting axes, the core
and the four caps have zero common kernel: at $u=0$, the core forces
$X=Y=0$, then $m$ and $\alpha$ force $W=Z=0$ when $v$ is a unit;
at $u=\infty$, $C,A$ force $W=Y=0$, then the core forces $X=Z=0$;
at $v=0$, $A,\alpha$ force $Y=Z=0$, then the core forces $W=X=0$
when $u$ is a unit. These are strictly positive bounded state Grams.

At $(u,v)=(0,0)$ the residue is
\[
 (2+4\mu)X^2+(2+8\mu)Y^2+4\mu Z^2.
\]
Only the $W$ ray remains. Normalize $W=1$ and take the five infinitesimal
roots $(u,v,X,Y,Z)$. The full polynomial has no constant or linear term;
its quadratic part is
\[
 2(u-Y)^2+2X^2+
 4\mu\{(v-X)^2+(u-Y)^2+Y^2+Z^2\},            \tag{CE11}
\]
which is strictly positive. At $(u,v)=(\infty,0)$ only the $X$ ray
remains. Normalize $X=1$ and use $(p,v,W,Y,Z)$. Its quadratic part is
\[
 2(W+v)^2+2(p+Z)^2+
 4\mu\{(p+Z)^2+W^2+Y^2+Z^2\},              \tag{CE12}
\]
also strictly positive. All higher monomials in each chart allocate to
products of two displayed roots with infinitesimal bounded coefficients.
This is an exact finite Gram, not a truncation certificate.

For $w=1/v$ infinitesimal use the full raw frame of (CE6). All four
selected terms, with every normal present, are exactly
\[
 4\mu\{(W-wX+uY-uH)^2+D_w(uW-Y)^2
              +D_uD_wY^2+D_u(Y-H)^2\}.        \tag{CE13}
\]
At $u=w=0$ their sum with $2K_w$ has residue
\[
 4\mu W^2+2X^2+(8\mu+2)Y^2+4\mu(Y-H)^2,    \tag{CE14}
\]
strict on all four roots. If $p=1/u$ is infinitesimal, homogenize again;
the exact terms become
\[
 4\mu\{(pW-pwX+Y-H)^2+D_w(W-pY)^2
              +D_pD_wY^2+D_p(Y-H)^2\}.        \tag{CE15}
\]
Together with (CE7), at $p=w=0$ the residue is
\[
 4\mu W^2+2X^2+4\mu Y^2+(8\mu+2)(Y-H)^2,    \tag{CE16}
\]
again strict. These frames require no state normalization, and no
division by either small testing coordinate. Literal $w=0$ is included.
The only chart not covered is ordinary $u$ and infinitesimal $w$.
Restore the original terms and all pure factors as above.
\end{proof}

\begin{lemma}[Matched weak outer events at the large-Bob corner]
Suppose $q>0$ is a pure-source infinitesimal and
\[
 C\ge c q^2,\quad m\ge m_*q^2,\quad
 A\ge a_*/q^2,\quad\alpha\ge b_*/q^2
                                                        \tag{CE17}
\]
for positive real $c,m_*,a_*,b_*$. At projectively unbounded $v$,
the complete parent is tensor SOS whenever $u$ is infinitesimal or
projectively unbounded. In the latter chart the lower bound on $m$
is unnecessary.
\end{lemma}
\begin{proof}
Keep $w=1/v$ and $H=wZ$ as in (CE6)--(CE7). If $u$ is infinitesimal,
use the five raw roots
\[
 E=qW,\quad X,\quad F=Y/q,\quad G=H/q,\quad V=uwW.
\]
The exact whole core (CE6) becomes
\begin{align*}
 K_w={}&V^2+(D_u+w^2)X^2
       +q^2\{u^2(G-F)^2+D_wF^2\}\\
 &-2uGE+2quwX(G-F)+2uVX-2uw^2EF.             \tag{CE18}
\end{align*}
Cap only $m,A,\alpha$ at the indicated weights. Their complete terms
are respectively
\[
 4m_*(E-qwX+uq^2(F-G))^2,\quad
 4a_*D_uD_wF^2,\quad4b_*D_u(F-G)^2.
\]
The selected residue is
\[
 2V^2+2X^2+4m_*E^2+4a_*F^2+4b_*(F-G)^2>0.   \tag{CE19}
\]
Every coefficient in (CE18) is bounded independently of all relative
source/testing rates.

If instead $p=1/u$ is infinitesimal, use
\[
 E=qW,\quad F=Y/q,\quad G=H/q,\quad
 R=X+wW,\quad S=pX.
\]
The exact whole core (CE7) is
\[
 K_{pw}=R^2+D_wS^2+q^2(G-F)^2+p^2D_wq^2F^2
       -2pE(G+w^2F)+2qwS(G-F).               \tag{CE20}
\]
Cap only $C,A,\alpha$. Their complete terms are
\[
 4cD_w(E-pq^2F)^2,\quad4a_*D_pD_wF^2,
 \quad4b_*D_p(F-G)^2.
\]
Together with $2K_{pw}$ their residue is
\[
 2R^2+2S^2+4cE^2+4a_*F^2+4b_*(F-G)^2>0.     \tag{CE21}
\]
Only the pure-source inverse $1/q$ occurs in a root. These exact raw
frames include every state normal and every literal endpoint. The
strict-residue lemma and restoration of excess terms prove the claim.
\end{proof}

\paragraph{Exact application to the original C1 source rows.}
The source premises here are precisely the already proved original
formulas (NC1)--(NC7), (NI3)--(NI10); no independent source family is
substituted. These common theorems give the following new quantified
test reductions, always intersected with prior certificates:
\begin{enumerate}
\item C1.1, every $e$ type: (NI3) gives $C,A$ unbounded;
only $u$ infinitesimal remains.
\item C1.4, $L$ infinitesimal and $\eta$ unbounded: (NI6)
again gives $C,A$ unbounded; only $u$ infinitesimal remains.
\item C1.4, $L$ unbounded: (NI4) and (CE20) leave only
$u$ bounded and $v$ unbounded.
\item C1.5 actual UH, $\eta$ unbounded and $\zeta$ unit:
(NI9) gives four strong events; only $u$ unit and $v$ unbounded remains.
\item C1.5 actual UH, $\zeta$ unbounded: (NI10) and (CE17)
leave the same sole strip, $u$ unit and $v$ unbounded.
\end{enumerate}
For the last line choose $q=\zeta^{-1/2}$. Formula (NI10) gives
$C,m\asymp q^2$, $\alpha\asymp1/(q^2Q)$ and $A\asymp1/r$;
the original $Q=1-x\zeta r>0$, with unit $x$, bounds $Aq^2$ below.
Thus (CE17) holds with positive real constants. The adjacent-event
theorem covers bounded $v$ and (CE18)--(CE21) remove both end types of
$u$ at unbounded $v$. Both actual UH source patterns now have the same
sole untreated testing strip, ordinary $u$ and unbounded $v$.

In the $L$-unbounded C1.4 line take $q=L^{-1/2}$.
Then (NI4) gives $C\asymp q^2$, $\alpha\asymp1/(q^2Q)$ and
$Aq^2\asymp1/(Lr)>1$. The $p$-chart (CE20), which needs no lower
bound on $m$, removes the projectively unbounded-$u$ tests there.

All bounds are relative to original $d$. The four-source-event result
does not require the additional NI9 total-scale condition that at least
one opening be infinitesimal. Every literal permitted $D=0$ or
$\beta=0$ is retained throughout, since neither coefficient was capped.
No claim of nonemptiness or full-row completion is made. In particular,
the $L$ unit and the $L$ infinitesimal/$\eta$ bounded parts of C1.4 are
unchanged, and these remote results do not supply the missing critical
contact conductors.



% Source formulas: p6cad7544d1b8
\subsection{Two exact complementary-current identities}

\label{component:p6cad7544d1b8}

\label{p6cad7544d1b8:sec:c1-uniform-all-tests}

Let $\mathcal H=2[(uW+uvX-Y-vZ)^2+D_u(vY-Z)^2+D_v(X+uZ)^2-D_uD_vZ^2]$ be the original testing core.

\begin{lemma}[Two strong currents]\label{lem:c1-two-current}
Suppose the original event coefficients are nonnegative, $d>0$, and
$d/c,d/D$ are infinitesimal, where $c=e_3$, $D=e_5$.
Then every testing chart with $1/v$ bounded is tensor-SOS, for every
other angle and every state.
\end{lemma}
\begin{proof}
Put $g=cD/(c+D)$, $h=D/(c+D)$, $R=uW-Y$,
$C=v(uX-Z)$, $B=X+uZ$, and $N=R+hC$.
Let $D_u=1+u^2$, $D_v=1+v^2$ and set
\[
 K_0=2\sqrt g\,C,\qquad
 B_0=\sqrt{2d}\sqrt{D_v/D_u}\,B,\quad
 c_u=D_u^{-1/2},\quad t_u=uD_u^{-1/2},\quad
 \eta_v=\sqrt{D_v}/v.
\]
The complete selected parent is exactly
\[
\begin{split}
4cD_vR^2+4D(R+C)^2+d\mathcal H
={}&4(c+D)N^2+4cv^2R^2+2d(R+C)^2\\
 &+2dD_u(vY-Z)^2+(K_0,B_0)Q(K_0,B_0)^T,
\end{split}
\]
where
\[
 Q=\begin{pmatrix}
 1-\dfrac{d}{2g}c_u^2\eta_v^2&
 \sqrt{\dfrac{d}{2g}}t_u\eta_v\\
 \sqrt{\dfrac{d}{2g}}t_u\eta_v&1
 \end{pmatrix}.
\]
To check the identity, use $Z=(uB-C/v)/D_u$ and complete the two current
squares. Since $d/g=d/c+d/D$ is infinitesimal and $c_u,t_u,\eta_v$
are bounded, $\bar Q=I_2$. Restore every unselected positive event.
All roots separate between fields. Clearing the original pure-testing
homogeneous angle factors gives the same identity at literal infinity:
$\eta_v^2=1+(1/v)^2$, and $c_u,t_u$ are the normalized homogeneous
Alice coordinates. Thus no infinite-angle endpoint is excluded.
\end{proof}

\begin{lemma}[The critical current coefficient]\label{lem:c1-critical-current}
Under the preceding strength assumptions, suppose
\[
 \gamma=\frac{A g}{d^2}\quad\hbox{is bounded},\qquad 4\bar\gamma>\bar c_u^2,
\]
where $A=e_4$ is the original event coefficient. Then every independent
test with $v$ infinitesimal satisfying this residue inequality is tensor-SOS, for every
state, including the homogeneous Alice-infinite endpoint.
\end{lemma}
\begin{proof}
Keep $R,C,B,N,g,h,c_u,t_u$ from the preceding proof, and put
$\epsilon=d/g$, $q=\sqrt\epsilon$ and
\[
 K_1=2\sqrt g\,C+\sqrt{d\epsilon}\,Y,\quad
 B_1=\sqrt d(B-uvY),\quad V_0=\sqrt d\,vY,
 \quad Y_0=\sqrt{d\epsilon}\sqrt{D_u}\,Y.
\]
Define
\[
 L_0=\frac{qc_u}{2}(K_1-c_uY_0)-t_uvB_1+c_uvV_0.
\]
Direct multiplication gives $L_0=\sqrt{dD_u}\,v(vY-Z)$ and the exact
full selected-parent identity
\[
\begin{split}
&4cD_vR^2+4D(R+C)^2+4AD_uD_vY^2+d\mathcal H\\
&\qquad=4(c+D)N^2+4cv^2R^2+2d(R+C)^2+\mathcal G,\\
\mathcal G={}&K_1^2+2D_v(B_1^2+V_0^2)
 +(4\gamma D_v-c_u^2)Y_0^2\\
&+2qv(K_1-c_uY_0)V_0-2L_0^2.
\end{split}
\]
The four-root Gram of $\mathcal G$ has residue
$\operatorname{diag}(1,2,2,4\bar\gamma-\bar c_u^2)$.
Its last entry is positive by the stated residue inequality. Every correction is polynomial in bounded
$q,v,c_u,t_u,\gamma$ and has zero residue. The strict-residue lemma and
restoration of the other events prove the claim. The displayed bounded
homogeneous factors also cover literal Alice infinity.
\end{proof}



\section{Bounded reciprocal contacts and a nonvanishing Alice scale}
For infinitesimal $\eta=f^2,\zeta=g^2$, one contact calculation covers bounded $k$ and bounded $\Delta=1/e$, including infinitesimal $k$. For unit $k$, a source reflection covers every positive ratio $g/f$. The remaining ordinary-Alice sources use the state completions derived below.

\subsection{One contact calculation for bounded $k$ and bounded reciprocal $e$}
\label{sec:c1-bounded-reciprocal}

Use the original C1 source in the normalization $C=2$, with the seven
native coquotients $c_2,c_b,c_3,c_5,c_4,c_\alpha,c_d$ and exact
source/parent SD1--SD2. Suppose $k>0$ is bounded, $f,g>0$ are
infinitesimal, and $\Delta=g/(kf)$ is bounded. Assume the original
positive shapes $a,t$ and positive gaps $M=R-t$, $J=T-a$.
They make $a,t$ bounded because $R,T\le1+k^2$. No bound $a,t<2$
is used; in particular the statement allows a bounded unit $k>1$.
Then the complete parent is tensor SOS on the testing chart with
$u$ bounded, $v$ infinitesimal, $Z$ a unit and all state coordinates
bounded. The proof allows bounded unit $u,W$ and both unit and
infinitesimal $k,\Delta$.

Normalize $Z=1$ in the testing field. Retain exactly the four positive
squares SD3 and use the full change of state coordinates SD5, with
$w=(W+uY)/(1+u^2)$. Put
\[
 K=1+k^2,\quad \delta=\Delta R/K,\quad b^2=\eta K/R,
 \quad Q=K+ak^2b^2,\quad T_n=K^2-t\delta^2k^4Qb^2.
\]
Here $\delta,k,a,t,u,w$ are bounded, $b$ and $v$ are infinitesimal,
and
\[
 \eta={Kb^2\over Q},\quad\Delta={\delta Q\over K},\quad
 \zeta={\delta^2k^2Qb^2\over K},\quad
 R={K^2\over Q},\quad T={T_n\over K}.\tag{SD6}
\]
Every denominator $K,Q,T_n,R,T$ is a positive pure-source unit;
its residue is respectively $K,K,K^2,K,K$.
The testing denominator $D=1+u^2$ is a positive bounded testing unit.
The positive source multiplier $\Delta^2\Lambda$ need not be a unit;
it is restored using pure-source square roots after the construction.

Define the polynomial pullback $C_i=K^2Q^3c_i$ by this substitution.
It is polynomial because a native monomial
$k^jf^{2r}g^{2s}a^\alpha t^\tau$ becomes
\[
 \delta^{2s}k^{j+2s}b^{2(r+s)}a^\alpha t^\tau
                         K^{2+r-s}Q^{3-r+s},\tag{SD7}
\]
and the seven native polynomials have $-2\le r-s\le3$.
Thus no shape inverse or negative exponent is introduced.
Let $\mathcal C=K^6Q^4T_n$ and define
\[
\begin{gathered}
 A=wv,\quad B=\delta v,\quad C_0=\delta b^2uw,
 \quad D_0=\delta^2k^3b^2,\quad D_v=1+v^2,\\
 V_n=K^2-tQuvw,\quad Y_n=KT_nuw-avV_n,\quad
 B_n=vY_n-KT_n,\quad A_c=C_5T_n+aK(C_3D_v-C_5).
\end{gathered}
\]
Put $G=\Delta^2\Lambda F-\sum_i\Pi_i$ exactly as in SD4.
The complete scalar of $\mathcal CG$ at the two state
normals $N=\delta M_B=0$ is the following five-term polynomial:
\[
\begin{split}
 F_0={}&4KQT_n^2[C_2(K^2-tQ)+C_bD_vtQ]D^2A^2
       +4Q^2A_cD^2V_n^2B^2\\
 &+4K^3C_4DD_v(KT_nC_0-ab^2BV_n)^2
       +4KQ^2C_\alpha DB_n^2D_0^2\\
 &+2K^2QC_dD\{b^2B^2(Y_n^2-K^2T_n^2+2aKT_nV_n)
                                     -2K^2T_n^2BC_0\}.
\end{split}
\]
This follows by the state substitution in SD4. In particular it lies
in $(A,B,C_0)^2+(D_0^2)$ before any contact division.
The complete two first-normal coefficients are
\[
\begin{split}
 S_B={}&-8aK^2b^4C_4D_vV_n
       +8\delta^2Q^2k^6b^4C_\alpha B_n+4KQb^2C_dB_n,\\
 S_C={}&8K^3T_nb^2C_4D_v,\\
 L_1={}&A[8K^3QT_n^2(C_bD_v-C_2)D]\\
       &+B[-8\delta K^2Q^2uDvV_nA_c+\delta K^3auvS_B]
                         +C_0[\delta K^3auvS_C],\\
 L_2={}&B[8K^2Q^2T_nDV_n(C_3D_v-C_5)-K^2T_nS_B]
                                      -C_0K^2T_nS_C.
\end{split}\tag{SD8}
\]
They include every first-order state term, and belong to
$(A,B,C_0)$ before contact division.

\paragraph{A single positive-weight estimate.}
Give $(\delta,k,b,a,t,u,w,v)$ weights $(0,0,1,0,0,0,0,2)$.
The zero-weight quantities are bounded; only $b,v$ require positive
weights. Each of $A,B,C_0,D_0$ has weight two. Directly from the
seven source polynomials,
\[
 C_i=2K^6+O_{\rm wt}(2)\ (i=2,b,3,5,4,\alpha),\quad
 C_d=4K^7+O_{\rm wt}(2),\quad Q=K+O_{\rm wt}(2),
\]
\[
 T_n=K^2+O_{\rm wt}(2),\quad V_n=K^2+O_{\rm wt}(2),
 \quad B_n=-K^3+O_{\rm wt}(2),\quad A_c=2K^8+O_{\rm wt}(2).
\]
The five scalar lines therefore have leading terms respectively
\[
 8K^{14}D^2A^2,\quad8K^{14}D^2B^2,\quad
 8K^{15}DC_0^2,\quad8K^{15}DD_0^2,
 \quad-16K^{16}DBC_0.
\]
These are identities modulo the coefficient ideal $(b^2,v)$ for
arbitrary bounded $k,\delta,u,w$, not identities only at contact.
The mixed term in the third line has a factor $b^2$; so does the
separate $B^2$ term in the fifth line. Consequently, with
\[
 P=D(A^2+B^2)+K(C_0^2+D_0^2)-2K^2BC_0,
 \qquad E_0=F_0-8K^2D\mathcal CP,
\]
one has $\operatorname{wt}(E_0)\ge6$.
For the complete normals use the exact differences
\[
 K(C_b-C_2)=\delta^2k^2Qb^2C_3,\qquad
 Q(C_3-C_5)=-k^2Kb^2C_b.
\]
Since $K,Q$ have positive unit constant terms in $b$, these give
weight at least two for both differences; also $S_B,S_C$ have
weight at least two. The displayed full normal expressions imply
$\operatorname{wt}(L_1),\operatorname{wt}(L_2)\ge4$.

The exact native contact is
$(u,w,v,N,M_B)=(k,\delta k,\delta k^2b^2,0,0)$.
The six native balances SD9--SD11 and the contact-zero retained
squares prove $E_0\in\mathfrak J^2$ and $L_i\in\mathfrak J$,
where $\mathfrak J=(u-k,w-\delta k,v-\delta k^2b^2)$.
Apply the restricted contact lift CC1--CC6.
For completeness, its coefficient filtration is particularly simple
here. The three contact substitutions preserve weight. Monic division
by $u-k$ and $w-\delta k$ loses no weight. A division by
$v-\delta k^2b^2$ loses two, but its quotient is always multiplied
by $b^2$ in the CC2 formula for $C_0S$.
Thus that formula preserves coefficient weight. The normal monomial
quotients have weight at least $4-2=2$; the scalar monomial quotients
have weight at least $6-4=2$. The successive CC reductions, including
the gradient subtraction $2k\rho_c(-2k,1,1,-1)$, preserve this bound.
Therefore every final scalar and normal correction coefficient has
weight at least two. It is a bounded tensor coefficient of zero
residue. No determinant inverse is a lifting coefficient.

\paragraph{The complete state-normal block.}
Before multiplication by $\mathcal C$, let
\[
 C_B=H_3D_vh_s+H_5(1-h_s),\quad h_s=a/T,
\]
\[
 B_Y=4RT\Delta^2\eta^2c_4DD_v
 +4RT\Delta^4k^6\eta^2c_\alpha Dv^2
 +2RT\Delta^2\eta c_dDv^2.
\]
The full matrix on $(N,\delta M_B)$ is defined equivalently by
coefficient extraction of SD4 after SD5. Its entries are
\[
\begin{split}
 A_{11}&=H_2+H_bD_v+\Delta^2C_Bu^2v^2
                                      +B_Yh_s^2u^2v^2/D^2,\\
 A_{12}&=-(\Delta^2/\delta)(H_3D_v-H_5)uv
                                      -(B_Y/\delta D^2)h_suv,\\
 A_{22}&=(\Delta^2/\delta^2)(H_3D_v+H_5)+B_Y/(\delta^2D^2).
\end{split}
\]
All negative powers cancel: $\Delta/\delta=Q/K$ is a source unit,
and $B_Y/\delta^2$ is an explicit bounded infinitesimal because each
summand contains $\eta^2$ or $\eta v^2$ after that cancellation.
Also $H_i\sim8K^2$, $C_B$ is bounded and
$H_3D_v-H_5=4R(c_3-c_5)+H_3v^2$ is infinitesimal.
Hence the block after multiplication by $\mathcal C$ has exact
residue $16\bar K^{14}I_2$, for arbitrary bounded $u,w$.

Finally use the six roots
\[
 x=B-C_0,\quad y=A-kC_0,\quad z=uB-kC_0,\quad
 h=kC_0-D_0,\quad q_1=a_1a_2-a_3a_4,\quad
 q_2=a_1a_3-a_2a_4,
\]
where $a_1=\sqrt{wv}$, $a_2=\sqrt{\delta uv}$,
$a_3=b\sqrt{\delta kuw}$ and $a_4=\delta k^{3/2}b$.
Writing $g_4=uwv-\delta^2k^4b^2$, the exact identities are
\[
\begin{split}
 P&=x^2+g_4^2+\sum_i a_i^4-4a_1a_2a_3a_4\\
  &=x^2+g_4^2+y^2+z^2+h^2-yz+zh+q_1^2+q_2^2.
\end{split}
\]
The $(y,z,h)$ matrix has successive principal minors
$1,3/4,1/2$; the other three coefficients are one.
Thus the six-root Gram is strictly positive and constant.
Retain $g_4^2$ separately; its coefficients
need not be infinitesimal when $k,u$ are units.
The full scalar and both normals are thus infinitesimal corrections
of this strict six-root Gram and the two strict state-normal roots.
The positive-residue Gram lemma applies. Restore $g_4^2$ and every
retained square, and undo the positive pure-source multiplier.
Literal axes and a nonnegative testing core have their direct original
event certificates. On the negative-core chart the exact SD core
calculation gives $u,v,w>0$, so all radicals above are separated
nonnegative radicals. This proves the stated local result for every
real state, and separately for both parts of a complex state.


\subsection{Every test in the bounded-reciprocal and ordinary-Alice sources}
\label{sec:c1-uniform-projective-join}
\label{component:p9f00f949fc8c}
\label{component:p64b054e815d0}
\label{component:p18a57c9c271d}
\label{component:p41b39e2f94e1}
\label{p41b39e2f94e1:thm:bounded-delta}
\label{p18a57c9c271d:thm:c1-ordinary-angle}

The preceding local calculation gives the complete projective result
when $k$ is positive bounded, $\eta=f^2,\zeta=g^2$ are infinitesimal
and $\Delta=g/(kf)$ is bounded. It also gives the entire source with
unit $k$ and arbitrary positive $g/f$, after the source reflection
below. We give the complementary testing and state cases explicitly.

\paragraph{Complementary testing angles.}
Let $g_c=CD/(C+D)$. The source coefficients SD1 give, with
$D_1=2c_5/c_3-J\eta c_d/c_3\sim2$,
\[
 {g_c\over a}={2D_1\over 2J+aD_1},\qquad
 2J+aD_1\sim2K,
\]
\[
 d/C,\ d/D\in\operatorname{Inf}_c,\qquad
 \gamma={Ag_c\over d^2}\sim{1\over4K^3}>0.
\]
Lemma~\ref{lem:c1-two-current} therefore applies when $1/v$ is bounded, including
Bob infinity, for every Alice angle. When $v$ is infinitesimal and
Alice is unbounded, Lemma~\ref{lem:c1-critical-current} has
residue diagonal
\[
 \operatorname{diag}(1,2,2,4\bar\gamma-\bar c_u^2),
 \qquad c_u=(1+u^2)^{-1/2}.
\]
Here $\bar c_u=0$ and $\bar\gamma>0$, so the matrix is strict.
Its homogeneous identity includes Alice infinity. Only bounded $u$
and infinitesimal $v$ remain. No upper bound on an event coefficient
was used; the complementary identities retain all original positive
excess squares.

\paragraph{The complementary state chart.}
Normalize a nonzero state by a coordinate of maximal absolute value.
Then all state coordinates are bounded. The testing signed core is a
pure-testing element. If it is nonnegative, it is a testing square
and every event term is already a separated square.

Suppose the core is negative. At $v=0$ it is
$2\{(X+uZ)^2+(uW-Y)^2\}$. Thus both displayed currents have zero
residue. If $Z$ is a unit, use
Section~\ref{sec:c1-bounded-reciprocal}; its coefficient bounds allow
bounded $u,w$ and every bounded $W,X,Y$ after $Z=1$ normalization.
If $Z$ is infinitesimal, then $X$ is infinitesimal and
$Y-uW$ infinitesimal. Some state coordinate is a unit, so $W$ is a
unit and $Q_A=W+uY$ is a unit with residue $(1+\bar u^2)\bar W$.

On this chart the exact five-root current identity uses
\[
 z=(r_2,B,\Delta r,\Delta r_5,\Delta\eta Y).
\]
Retain the four positive source-gap excesses, the whole alpha term
and the positive $v^2Y^2$ part of the core. The remaining diagonal is
\[
 H_2z_1^2+H_bD_vz_2^2+H_3D_vz_3^2+H_5z_4^2
                       +4RTc_4D_uD_vz_5^2,
\]
and its complete negative correction is
\[
 -4RTc_dD_u{\Delta Z\over Q_A}z_5(z_1+z_2)
 -2RTc_dD_u{\Delta^2\eta Z^2\over Q_A^2}(z_1+z_2)^2.
\]
The five positive coefficients are units and both correction
coefficients are infinitesimal because $\Delta$ is bounded and
$Z$ infinitesimal. This proves the state chart. The literal testing
axes have nonnegative core directly. The local and complementary
state arguments are therefore exhaustive.

\paragraph{Ordinary Alice with arbitrary state scale.}
Suppose $k$ is a positive bounded unit and $f,g$ are infinitesimal.
Only $\Delta$ unbounded remains. The exact complete-parent reflection
is
\[
 u'=1/u,\quad v'=v,\quad
 (W',X',Y',Z')=(vW-X,-W-vX,Z-vY,Y+vZ),\tag{AS3}
\]
\[
 u^2F(\operatorname{ref}(c);1/u,v;\psi')
                         =(1+v^2)F(c;u,v;\psi),\tag{AS4}
\]
with event permutation $(\lambda_{\rm ev}\ n)(m\ C)(A\ \alpha)
(D\ \beta)$. Put $H_f=1+(1-a)\eta$, $H_g=1+(1-t)\zeta$
and $V=1+k^2\eta\zeta$. At the actual source contact its native
parameters are
\[
\begin{gathered}
 k'=1/k,\quad
 \eta'=k^2\zeta H_f/H_g,\quad\zeta'=k^2\eta H_g/H_f,\\
 a'=M/(k^2H_f),\quad t'=J/(k^2H_g),\quad
 M'=aV/(k^2H_g),\quad J'=tV/(k^2H_f).
\end{gathered}\tag{AS5}
\]
These identities follow from $H_g-Q_c=\zeta M$ and
$H_f-E=\eta J$ and preserve all original signs and gaps.
Here $H_f,H_g$ are positive units, $k'$ is a unit,
$\eta',\zeta'$ are infinitesimal, and
\[
 \Delta'={\sqrt{\zeta'/\eta'}\over k'}
                         ={H_g\over\Delta H_f}
\]
is infinitesimal. Normalize the reflected coefficients by a positive
pure-source factor to $C'=2$ and apply the bounded-reciprocal result.
The homogeneous reflection restores the original projective endpoints
as well. This proves every ordinary-Alice source with small $f,g$,
without a separate comparable-ratio or shape-unit argument.


% Source formulas: p050c7ae6d9b8
\subsection{The positive Laurent form at the ordinary/small source}

\label{component:p050c7ae6d9b8}

\label{p050c7ae6d9b8:sec:c1-ui-leading}

This is the positive Laurent-model lemma used in the complete UI proof.
Write $x=k^2$, $r=1-a\eta$, $K=1+x$, $R=1+xr$ and
\[
 E_0=1-x\eta,\quad J=K-a,\quad
 P_4=KE_0+x(1-r),\quad P_d=1+2x+rx^2.
\]
At a real interior source point the leading coefficients of the
parent divided by $a$ are
\[
\begin{aligned}
 h&=\frac{2x\eta}{RE_0},& C&=\frac2a,&
 A&=\frac{2\eta^2P_4}{Kr(r+\eta)E_0},\\
 d&=\frac{4\eta P_d}{K(r+\eta)E_0},&
 D&=\frac2{JE_0}-\frac d2,&
 \alpha&=\frac{2x^2\eta R}{KE_0}.
\end{aligned}                                                     \tag{CL1}
\]
All are positive except $D\ge0$, including its literal zero. Here
$\alpha$ is the coefficient after extracting the weak $g^2$ scale.
For a test $(u,V)$ put $D_u=1+u^2$ and define
\begin{align*}
f={}&4hV^2(W+uY)^2+4C(uW-Y)^2+4AD_uY^2\\
 &+(4D+2d)(uW-Y-VD_uZ)^2+4\alpha D_uZ^2\\
 &-4dD_uVYZ-2dD_uV^2Z^2.                           \tag{CL2}
\end{align*}
Let $M_f$ be its symmetric three-state matrix and set $U=u^2$, $T=V^2$.
Direct collection gives
\[
 \det M_f=64D_u^2UT\,\mathcal L(U,T),                 \tag{CL3}
\]
where
\begin{align*}
\mathcal L={}&AC(D+d/2)U+
 \frac{h\alpha(A+C+D+d/2)}U+\frac{A\alpha(C+D+d/2)}T\\
 &+hD(A+C+d/2)\frac TU+hC(D+d/2)UT\\
 &+h(AD+2CD+Ad/2+Cd/2+Dd/2+d^2/4)T+L_0,\\
L_0={}&ACD-ADd/2+h\alpha(C+D+d/2)
                         -d^2(A+C+D)/4-d^3/8.       \tag{CL4}
\end{align*}
Every nonconstant coefficient is nonnegative. The coefficients of
$U,U^{-1},T,T^{-1}$ are strictly positive, also when $D=0$.
Thus in coordinates $(\log U,\log T)$ the Hessian is strictly
positive and the function is coercive. Substitution of (CL1) gives
\[
 \mathcal L(x,x\eta)=\partial_U\mathcal L(x,x\eta)
                   =\partial_T\mathcal L(x,x\eta)=0. \tag{CL5}
\]
Strict convexity proves that $\mathcal L$ is positive away from this
single contact. The $(W,Y)$ principal block of (CL2) is strictly
positive for $U,T>0$, because its $C$ and $A$ terms already have
trivial common kernel. Hence $f$ is positive semidefinite everywhere
and positive definite off the contact locus. Literal endpoints follow
by continuity. This proof uses no numerical determinant table.

There is a useful corrected formula at the large-$D$ boundary.
Impose its hard equation $uW-Y-VD_uZ=0$, eliminating $W$.
The resulting two-state determinant is $16D_u^2T\mathcal L_D$, where
\begin{align*}
\mathcal L_D={}&hT\{(C+A+d/2)/U+(2C+A+d/2)+CU\}\\
 &+h\alpha(1+1/U)+AC(U+1)+A\alpha/T-Ad/2-d^2/4.
                                                               \tag{CL6}
\end{align*}
In particular the coefficients of $T/U$ and $UT$ are
$h(C+A+d/2)$ and $hC$, respectively. Interchanging these coefficients
is incorrect. All three identities follow by polynomial collection in the displayed model.



% Source formulas: pa0b854a04f49
\subsection{The ordinary/small source: the complete state quadratic}

\label{component:pa0b854a04f49}

\label{pa0b854a04f49:sec:r6-c1-ui}
Use the unchanged native C1 source in the representative of (NI7):
$k,\eta$ are positive bounded units and $\zeta=g^2$ is positive
infinitesimal. All original openings and $D=0$ are retained.
The exact stabilizer recovers the other representative; no condition
$k\le1$ is imposed after that change. This section proves the missing
ordinary/small testing strip, and joins the remote theorem (CR1)--(CR13).

\paragraph{Exact bounded remainder.}
Normalize by the positive pure-source $a$. Write $x=k^2$,
$r=1-a\eta$, $K=1+x$, $R=1+xr$, $M=R-t$, and
\[
 J_s=K-xtg^2-a,\quad E=1-x\eta+tx\eta g^2,\quad
 Q_c=1-xg^2r,\quad y=x\eta g^2.
\]
Here $a,t,r,1-r,M,J_s$ are positive bounded openings; $E,Q_c,R,K,r+\eta$
are positive units. The exact polynomials $P_4,P_d,P_\alpha$ are (NC2).
The coefficients after division by $a$ are
\[
\begin{aligned}
 c&=2/a,& A_0&=\frac{2\eta^2P_4}{Kr(r+\eta)E},&
 d_0&=\frac{4\eta P_d}{K(r+\eta)E},\\
 D_0&=\frac{2(1+y)}{J_sE}-d_0/2,&
 \alpha_0&=\frac{2x^2\eta P_\alpha}{KQ_cE},&
 \alpha/a&=g^2\alpha_0.
\end{aligned}                                                     \tag{CU1}
\]
In particular $c,A_0$ are strong; $d_0,\alpha_0$ are units; and
$ca,A_0r,D_0J_s$ are bounded with continuous rational expressions
through every opening. Put
\[
 h_2=\frac{2x\eta(1-tg^2)(1+y)}E,\qquad
 h_b=\frac{2x\eta[1+(1-t)g^2]}E,
 \qquad h_b-h_2=2x\eta g^2.
\]
Both $h_2,h_b$ are positive units. The exact coefficients are
$m/a=h_2/(g^2M)$ and $\beta/a+d_0/2=h_b/(g^2t)$.

For the test put $D_u=1+u^2,D_v=1+v^2$ and
\[
 Q=W+uY,\quad B=X+uZ,\quad p=t/R,
 \quad N=(B-pvQ)/g,\quad L=vQ/g.
\]
The inverse $W=Q-uY$, $X=gN+pvQ-uZ$ is polynomial.
Retain the two positive squares
\[
 \frac{4h_2t}{MR}N^2,\qquad
 \frac{4h_bM}{tR}D_vN^2.                            \tag{CU2}
\]
Their remainder is exactly (CL7). Choose bounded caps
$c_*\le c$, $A_*\le A_0$, $D_*\le D_0$, with $c_*,A_*>0$,
$D_*\ge0$; bounded original weights can be kept unchanged.
Retain the three positive centered excesses
\[
\begin{aligned}
 &4(c-c_*)D_v(r_t-avZ)^2,\quad
 4(A_0-A_*)D_uD_v(Y-rvZ/\eta)^2,\\
 &4(D_0-D_*)(J_t+J_svZ)^2,
\end{aligned}                                                     \tag{CU3}
\]
where
\[
 r_t=uQ-D_uY,\quad J_t=r_t-D_uvZ+ugvN+upv^2Q.
\]
All five retained rows vanish at the actual contact
\[
 (u,v)=(k,k\sqrt\eta g),\quad
 (W,X,Y,Z)=(g/\sqrt\eta,k(tg^2-1),krg/\sqrt\eta,1).
                                                               \tag{CU4}
\]
Set $c_a=2-c_*a$, $A_r=A_0r-A_*r$, $D_J=D_0J_s-D_*J_s$.
After removing (CU2)--(CU3), the entire parent divided by $a$ is
\begin{align*}
G={}&\frac4R\{(h_2+h_bD_v)N^2+2(h_bD_v-h_2)LN
              +[h_2(1-p)+h_bpD_v]L^2\}\\
 &+4c_*D_vr_t^2+8c_aD_vvZr_t-4c_aaD_vv^2Z^2\\
 &+4A_*D_uD_vY^2+\frac{8A_r}{\eta}D_uD_vvZY
                   -\frac{4A_rr}{\eta^2}D_uD_vv^2Z^2\\
 &+4D_*J_t^2-8D_JvZJ_t-4D_JJ_sv^2Z^2\\
 &+4\alpha_0D_u(gvY-gZ)^2
       +2d_0\{J_t^2+D_u(v^2Y^2-2vYZ-v^2Z^2)\}.
                                                               \tag{CU5}
\end{align*}
Thus $G$ is an exact bounded Gram on $(N,Q,Y,vZ,gZ,L)$ whenever
$u$ is a bounded unit and $v$ infinitesimal. Every source denominator
in its coefficients is a unit. Formula (CU5), rather than a scalar
restriction, supplies all normal terms below.

\paragraph{Local contact lifting with all first normals.}
Put $\kappa=k\sqrt\eta$, $\varepsilon=\kappa g$,
$s=R/(k\eta)$, $\delta=u-k$, and define
\[
 E_Q=Q-svZ,\quad E_Y=Y-rvZ/\eta,\quad
 R_h=(vQ-s\varepsilon^2Z)/\varepsilon,
 \quad S=v\delta Z,\quad T=\varepsilon\delta Z.
\]
Let $c_L=4[h_2(1-p)+h_bpD_v]/R$, a positive bounded unit.
The sole potentially unbounded $L$ row can be replaced exactly:
\[
 L=\kappa R_h+kR gZ,\qquad
 L^2-\kappa^2R_h^2=2kR vQZ-k^2R^2g^2Z^2.           \tag{CU6}
\]
Consequently (CU5) equals $c_L\kappa^2R_h^2$
plus $8\kappa(h_bD_v-h_2)NR_h/R$ and a polynomial $G_0$.
Explicitly, in (CU5) replace its first line by
\[
 \frac4R(h_2+h_bD_v)N^2+8k(h_bD_v-h_2)gNZ
                 +c_L(2kRvQZ-k^2R^2g^2Z^2).        \tag{CU7}
\]
All coefficients of $G_0$ are bounded polynomials in $u,v,g$ with
pure-source unit denominators. The coefficient of $NR_h$ is
infinitesimal, since $h_bD_v-h_2=2x\eta g^2+h_bv^2$.

Here are exact prescriptions for every normal coefficient, avoiding
any implicit omission. Substitute $Q=E_Q+svZ$ and
$Y=E_Y+rvZ/\eta$ in $G_0$ and collect the quadratic polynomial
in $(N,E_Q,E_Y)$. The constant coefficient is $\phi(u,\xi)Z^2$,
where $\xi=v^2$ and $\phi\in(g^2,\xi)$. The linear coefficients
of $E_Q,E_Y$ are $vL_Q(u,\xi)Z,vL_Y(u,\xi)Z$; that of $N$ is
$gL_N(u,\xi)Z$, where $L_N\in(g^2,\xi)$. These parity and ideal
assertions follow term by term in (CU5),(CU7). The remaining quadratic
normal block is the full coefficient collection of the same expression.
For explicit pen-and-paper formulas, set
\[
 b=r/\eta,\quad F_0=us-D_ub,\quad
 G_0^{\rm sc}=F_0-D_u+up\xi s,\quad
 H_0=(8D_*+4d_0)G_0^{\rm sc}-8D_J.
\]
Then, writing $D_v=1+\xi$ in these polynomials,
\begin{align*}
 L_Q={}&2c_LkR+8D_vu(c_*F_0+c_a)+u(1+p\xi)H_0,\\
 L_Y={}&-8D_vD_u(c_*F_0+c_a)
       +8D_uD_v(A_*b+A_r/\eta)-D_uH_0\\
      &+8\alpha_0D_ug^2(b\xi-1)+4d_0D_u(b\xi-1),\\
 L_N={}&8k(h_bD_v-h_2)+u\xi H_0.
\end{align*}
The complete scalar coefficient is
\begin{align*}
 \phi={}&c_L(2kRs\xi-k^2R^2g^2)
   +4c_*D_vF_0^2\xi+8c_aD_vF_0\xi-4c_aaD_v\xi\\
 &+4A_*D_uD_vb^2\xi+8(A_r/\eta)D_uD_vb\xi
                    -4A_rrD_uD_v\xi/\eta^2\\
 &+4D_*(G_0^{\rm sc})^2\xi-8D_JG_0^{\rm sc}\xi-4D_JJ_s\xi\\
 &+4\alpha_0D_ug^2(b\xi-1)^2
   +2d_0\{(G_0^{\rm sc})^2\xi+D_u(b^2\xi^2-2b\xi-\xi)\}.
\end{align*}
These formulas display the two asserted ideals directly. They also
give a short explicit input for the monic divisions, without an
unprinted state-normal computation.

The source is an exact stationary zero and every removed square
vanishes there. Therefore
\[
 \phi=\phi_u=\phi_\xi=0,\qquad L_Q=L_Y=L_N=0
                       \quad\text{at }(u,\xi)=(k,\varepsilon^2).
\]
Successive monic polynomial division gives, with bounded coefficients,
\[
 \phi=\delta^2(\varepsilon^2B_0+\xi B_1)
       +\delta(\xi-\varepsilon^2)B_2
       +(\xi-\varepsilon^2)^2B_3,                  \tag{CU8}
\]
\[
 L_i=\delta L_{i0}+(\xi-\varepsilon^2)L_{i1} (i=Q,Y),\quad
 L_N=\delta(g^2L_{N0}+\xi L_{N1})
                         +(\xi-\varepsilon^2)L_{N2}.
\]
For clarity, each quotient is obtained by subtracting the value at
$u=k$ and dividing by the monic $u-k$, then by the monic
$\xi-\varepsilon^2$. The second $u$ quotient of $\phi$ stays in
$(g^2,\xi)$; its constant and first $u$ terms have the stated double
and simple zeros in $\xi$. The same argument gives the $L_N$ ideal.
No small source number is inverted in these divisions.

Put $R_s=(v^2-\varepsilon^2)Z/\varepsilon$. The exact root identities
\[
 \varepsilon R_s=(\varepsilon R_h-vE_Q)/s,\quad
 R_sT=(R_hT-SE_Q)/s,\quad
 S^2-T^2=\delta(R_hT-SE_Q)/s                       \tag{CU9}
\]
allocate (CU8) and every first normal on the six roots
$(N,E_Q,E_Y,R_h,S,T)$. For example $v\delta Z=S$,
$\delta(\xi-\varepsilon^2)Z^2=(R_hT-SE_Q)/s$, and
$(\xi-\varepsilon^2)^2Z^2=(\varepsilon R_h-vE_Q)^2/s^2$.
The $N$ normal has only infinitesimal couplings. All quadratic normal
coefficients are already bounded. This is the required exact full lift.

\paragraph{Strict matched Hessian and source boundary caps.}
First take finite interior source residues. The leading form is
(CL2). Put $V=k\sqrt\eta V_1$, $W=W_1/\sqrt\eta$,
$Y=kY_1/\sqrt\eta$, $Z=1$. At the contact
$(W_1,Y_1,V_1,u)=(1,r,1,k)$, the Schur complement of the two state
coordinates in half the Hessian, in coordinates $(V_1,u)$, is
\[
\begin{aligned}
 G_{11}&=\frac{32\eta x^2P_4}{KE_0},&
 G_{12}&=\frac{16\eta k^3(R^2-2KE_0)}{KE_0R},\\
 G_{22}&=\frac{32\eta x(P_4+xrP_d)}{R^2P_4},&
 \det G&=\frac{256\eta^2x^3(4K^2E_0-R^2)}{K^2E_0^2}.
\end{aligned}                                                     \tag{CU10}
\]
Here $E_0=1-x\eta$, $P_4=KE_0+x(1-r)$, $P_d=1+2x+rx^2$.
There are no opening denominators. The leading $D\ge0$ condition is
\[
 K(2+xr)E_0-R^2\ge0,
\]
so $4K^2E_0-R^2\ge K(4K-2-xr)E_0>0$.
The state block is uniformly bounded below since $c,A_0$ are strong.
Thus the matched Hessian is strict, uniformly through every permitted
source boundary after the appropriate hard directions are retained.

We give the boundary justification explicitly. Divergences can occur
only at $a=0$, $r=0$, or $J_s=0$, for $c,A_0,D_0$, respectively.
Since $a+J_s\sim K$, $c$ and $D_0$ cannot both diverge. If $a=0$,
then $r=1$, so $c$ is the only divergent weight. The only possible
divergent pair is $A_0,D_0$. Their hard-state rows are $E_Y$ and
$kE_Q-KE_Y$, with nonzero determinant $k$. All angular parts of
these centered rows are bounded. Solving these at most two real
linear equations therefore supplies a bounded competitor for the
minimizing state vector. The bounded finite part and (CU10) imply
that the state minimizer is bounded as well, by completing its
positive state square. The full matched Hessian hence has a uniform
positive lower bound on the Euclidean norm in a nearby finite-source
family. On a vector in the limiting hard kernel, adjust only the two
state coordinates to annihilate the nearby exact hard rows. This
adjustment tends to zero, and the divergent penalties vanish exactly.
The bounded finite part is consequently strictly positive on that
kernel. The elementary real finite-penalty lemma supplies finite caps.

Every allowed residue is approachable by such finite sources. For
fixed $x>0$, the real leading inequalities are
\[
 0\le r\le1,\qquad
 \eta_{\min}(r)=\frac{1-r}{K}\le\eta\le
 \eta_{\max}(r)=\frac{1-R^2/[K(2+xr)]}{x},
\]
with positive interval width $R/[xK(2+xr)]$. Perturb $r$ into
$(0,1)$ and its interval coordinate into $(0,1)$.
This makes $a,r,J$ and the $D$ numerator strictly positive while
converging to the given residue. Taking $t=R/2$ and then sufficiently
small real $g>0$ restores all original source signs, including
$\beta>0$. This construction, rather than substitution into a
singular leading denominator, justifies the boundary limit.

In the six-root residue Gram, $N$ is an isolated positive pivot.
The $(E_Q,E_Y)$ block is strict, $R_h$ has positive pivot
$c_L\kappa^2$, $S$ couples only to $E_Q,E_Y$, and $T$ only to $R_h$.
Setting $S=T$ is the matched Hessian under an invertible real change
of coordinates: $R_h/g$ has relative-angle derivative $2s\kappa\ne0$,
and $T/g=\kappa\delta$ supplies the ordinary-angle derivative.
After the hard Schur complements, the sum of the $S,T$ diagonals is
therefore strictly positive. Add a real multiple of the last zero
in (CU9) to redistribute this sum into two positive diagonals.
The extra coefficients contain $\delta$ and vanish in the residue.
The bounded strict-residue Gram lemma proves the full parent at
$u-k$ and $v$ infinitesimal, including every source opening.

\paragraph{Ordinary angles away from contact.}
Let now $\bar u^2\ne\bar k^2$, with $u$ a unit and $v$ infinitesimal.
Use the bounded frame $(N,r_t,Y,vZ,gZ,L)$ of (CU5), and the exact zero
\[
 L(gZ)-Q(vZ)=0,\qquad Q=(r_t+D_uY)/u.              \tag{CU11}
\]
For finite leading source weights, $N$ is an isolated positive pivot.
The remaining residue is
\[
 q^TPq+2Ub^Tq+c_2U^2+4hL^2+c_0T^2,
 \quad q=(r_t,Y)^T,\ U=vZ,\ T=gZ,
\]
where the source coefficients are (CL1) and
\[
\begin{gathered}
 s_D=4D+2d,\quad s_C=4C+s_D,\quad
 P=\operatorname{diag}(s_C,4AD_u),\\
 \ell=(1/u,D_u/u)^T,\quad b=(-s_DD_u,-2dD_u)^T,\\
 c_0=4\alpha D_u,\qquad c_2=s_DD_u^2-2dD_u.
\end{gathered}                                                     \tag{CU12}
\]
Apply the real rank-one allocation lemma (UA1)--(UA3), using (CU11).
We verify its hypotheses uniformly through all source boundaries.
Writing $d_a=\ell^TP^{-1}\ell$, $e_a=\ell^TP^{-1}b$,
$h_a=c_2-b^TP^{-1}b$, direct inversion yields
\[
\begin{aligned}
 d_a&=\frac{aJE_0}{8u^2(JE_0+a)}+\frac{D_u}{4Au^2},\\
 e_a&=-\frac{aD_u}{u(JE_0+a)}-\frac{dD_u}{2Au},\\
 h_a&=\frac{8D_u^2}{JE_0+a}-2dD_u-\frac{d^2D_u}{A},\\
 P^{-1}b&=(-aD_u/(JE_0+a),-d/(2A))^T.
\end{aligned}                                                     \tag{CU13}
\]
All are bounded continuous, since $JE_0+a$ is a positive unit and
$1/A$ has a continuous bounded formula proportional to $r$.
The same holds for $P^{-1}\ell$.

Normalize the Laurent polynomial (CL4) by $As_C$.
Its $U$ and $T^{-1}$ coefficients are, respectively,
$1/[2(JE_0+a)]$ and $\alpha/4$, positive units.
Its $T$ coefficient is
\[
 h\{(D+d/2)/s_C+2CD/(As_C)+d/(8A)\},
\]
also a positive unit: when $1/A=0$, $r=0$ forces $a$ a unit,
so $(D+d/2)/s_C=a/[4(JE_0+a)]$ is a unit; otherwise the last term
is positive. The $U^{-1}$ coefficient is
$h\alpha(1/s_C+1/(4A))$. If it vanishes, $r=0$ and $J=0$,
and the $T/U$ coefficient has the positive term $hD/s_C=h/4$.
All remaining normalized nonconstant coefficients are bounded and
nonnegative. The constant is bounded by the exact contact equation.
Thus strict convexity and coercivity in logarithmic coordinates
persist at every boundary. The normalized Laurent function retains
its unique stationary zero $(U,T)=(x,x\eta)$.

For fixed $u^2\ne x$, the rank-one determinant polynomial is
\[
 p(T)=c_0+(h_a+4hd_ac_0)T+4h(d_ah_a+e_a^2)T^2
      =\frac{16D_u}{As_C}T\mathcal L(u^2,T)>0.
\]
Its constant and leading coefficients are positive units by the
preceding coefficient bounds. The strict allocation interval
therefore survives at every source boundary. The two real Gram
blocks have a uniform positive lower bound; (CU13) controls their
state minimizers even if $P$ is unbounded. The centered divergent
rows have hard parts $r_t$, $Y$, $r_t$; the only simultaneous pair
is $Y,r_t$, whose determinant is one. Exact small state corrections
annihilate the large rows in the nearby finite family. The same
finite-penalty argument as above supplies caps in (CU3). The exact
bounded remainder (CU5), after (CU11), has strict residue.
This proves every ordinary/small test away from contact.

\paragraph{Exhaustive certificate.}
The two arguments cover the whole ordinary/small testing strip;
the exact gauges cover negative contact signs. The remote theorem
(CR1)--(CR13), including CR5a, covers its complement.
All transformations are homogeneous in the state; no state coordinate
was divided out in the tensor proof. Real and imaginary state parts
are treated separately. Hence the entire actual UI subdomain of C1.5,
and its exact stabilizer images, has an all-tests tensor-SOS proof.
The actual UH subdomain of C1.5 remains a separate obligation.



% Source formulas: p013c226a79b3
\subsection{The ordinary/small source: complementary testing charts}

\label{component:p013c226a79b3}

\label{p013c226a79b3:sec:c1-ui-remote}
The source has actual ordinary Alice angle and infinitesimal Bob
angle. Apply the exact native stabilizer (AS5) as in (NI7) and take
$\eta$ a positive bounded unit, $\zeta=g^2$ infinitesimal and $k$ a
bounded unit. The reflected orientation need not satisfy $k\le1$.
Its $E$ remains a positive unit by the exact source entry in (NI7).
Normalize by the positive pure-source $d$.
Equations (NC2)--(NC4) give
\[
 C/d,\ A/d,\quad g^2m/d,\quad g^2\beta/d,\quad
 \alpha/(dg^2)\quad\hbox{strong},\qquad D/d\ge0.       \tag{CR1}
\]
For example $P_d,P_\alpha,Q,r+\eta$ are units, whereas the
positive openings $M,t,r,1-r$ are bounded and may be infinitesimal.
The $-1/2$ in $\beta/d$ disappears after multiplication by $g^2$.
In this particular source subdomain $\beta=0$ is impossible; $D=0$
is allowed throughout. Choose positive real constants
$\mu,\nu,c,a,b$ below the five quantities in (CR1), respectively
in the order $g^2m/d,g^2\beta/d,C/d,A/d,\alpha/(dg^2)$.
Whenever needed we cap these coefficients at
\[
 m=\mu/g^2,\quad\beta=\nu/g^2,\quad C=c,\quad A=a,
 \quad\alpha=bg^2.                                  \tag{CR2}
\]
Every unused term, every positive excess, and the complete $D$
term is retained as a square. The signed core, divided by two, is
\[
 K=J^2+D_u(vY-Z)^2+D_v(X+uZ)^2-D_uD_vZ^2,
 \quad J=uW+uvX-Y-vZ.                              \tag{CR3}
\]
Only the core and the explicitly selected caps need a bounded Gram.
This observation keeps all upper coefficient scales out of the proof.

\begin{theorem}[C1 UI remote tests]
For every source satisfying (CR1), the complete original C1 parent
is tensor SOS at every independent testing pair except possibly
$u$ a bounded nonzero unit and $v$ infinitesimal. The conclusion
includes literal projective endpoints, both signs, all states,
and $D=0$.
\end{theorem}
\begin{proof}
All constants in residue forms below are real residues. Roots may be
unbounded, but all coefficients in their exact Gram are bounded.
First put $Q=W+uY$, $B=X+uZ$; then
\[
 M=vQ-B,\quad uW-Y=uQ-D_uY,\quad
 J=uQ-D_uY+uvB-vD_uZ.                               \tag{CR4}
\]
These identities account for every state normal.

If $u$ is a bounded unit and $v$ projectively infinite, put $w=1/v$
and homogenize by $w^2$. Use roots
$(E_Q,E_B,Y,Z)=(Q/g,B/g,Y,Z)$ and cap only $m,\beta$.
The exact homogenized core is
\[
 [w(ugE_Q-D_uY)+ugE_B-D_uZ]^2
 +D_u(Y-wZ)^2+D_wg^2E_B^2-D_uD_wZ^2.
\]
Its residue is $D_u(Y^2+u^2Z^2)$. Including the two caps gives
\[
 4\mu E_Q^2+4\nu E_B^2+2D_u(Y^2+u^2Z^2)>0.         \tag{CR5}
\]

If $u$ is infinitesimal and $v$ projectively infinite, the preceding
four-root frame loses the $Z$ pivot. Keep $w=1/v$ and instead use
$(E_Q,E_B,Y,U,T)$, where $U=uZ$ and $T=gwZ$. Put
$A_0=w(ugE_Q-D_uY)+ugE_B$. The exact homogenized core is
\[
 K_w=A_0^2+D_uY^2+D_wg^2E_B^2+D_uU^2
       -2D_ug(wE_Q+E_B)U+2wD_uuYU.
\]
All its coefficients are bounded, with residue $Y^2+U^2$.
Cap $m,\beta,\alpha$; their amplitudes after normalization are
$E_Q-wE_B$, $E_B$, and $gY-T$, with the positive factors
$D_w,D_u$ in the last two terms. The full selected residue is
\[
 4\mu E_Q^2+4\nu E_B^2+2Y^2+2U^2+4bT^2>0.       \tag{CR5a}
\]
This includes literal $u=0$ or $w=0$; no division by either was used.

If $u$ is infinitesimal and $v$ a bounded unit, use
$E_Q=Q/g,E_B=B/g,Y,U=uZ,Z_g=gZ$, and cap $m,\beta,\alpha$.
Substitute $W=gE_Q-uY$, $X=gE_B-uZ$ in (CR3).
All terms involving $Z$ collect into $uZ$; the resulting exact
core is a polynomial in $(u,v,g;E_Q,E_B,Y,U)$.
Its residue is $D_vY^2+v^2U^2$. The full selected residue is
\[
 4\mu(vE_Q-E_B)^2+4\nu D_vE_B^2
       +2D_vY^2+2v^2U^2+4bZ_g^2>0.                 \tag{CR6}
\]
The $\alpha$ amplitude is $gvY-Z_g$, so no inverse of $u$ was used
in constructing this Gram.

If both $u,v$ are infinitesimal, set
\[
 U=uQ,\quad V=vQ/g,\quad E_B=B/g,\quad Z_g=gZ,
 \quad S=uvZ,
\]
and include $Y$ as the sixth root. The exact core is
\begin{align*}
K={}&(U-D_uY+uvgE_B)^2+D_uv^2Y^2+D_vg^2E_B^2+D_uS^2\\
 &+2D_uuYS-2D_uuVZ_g-2D_uvgE_BS.                 \tag{CR7}
\end{align*}
In particular the apparently problematic term $UvZ$ is exactly
$uVZ_g$, with no mixed inverse. Cap all five weights in (CR2).
The entire selected residue is
\[
 4\mu(V-E_B)^2+4\nu E_B^2+(4c+2)(U-Y)^2
          +4aY^2+4bZ_g^2+2S^2>0.                  \tag{CR8}
\]

It remains to cover $u$ projectively infinite. Put $p=1/u$,
homogenize by $p^2$, and use $D_p=1+p^2$,
$Q=pW+Y$, $B=pX+Z$. Then
\[
 M=vQ-B,\quad r=W-pY,\quad
 J=D_p(W+vX)-p(Q+vB).                              \tag{CR9}
\]
For $v$ a bounded unit, use $(Q/g,B/g,W,X)$ and cap $m,\beta,C$.
The core after $Y=gE_Q-pW$, $Z=gE_B-pX$ is a bounded polynomial
with residue $(W+vX)^2$. The selected residue is
\[
 4\mu(vE_Q-E_B)^2+4\nu D_vE_B^2
           +4cD_vW^2+2(W+vX)^2>0.                 \tag{CR10}
\]

For infinitesimal $v$, use $(W,R,Y,V,E_B)$ with
$R=vX$, $V=vQ/g$, $E_B=B/g$. Cap $m,\beta,C,A$, retaining
$\alpha$ entirely. The exact core is
\begin{align*}
K_p={}&(W+D_pR-pY-pvgE_B)^2+D_vg^2E_B^2\\
 &+D_p\{v^2Y^2-2vgYE_B+2pYR-(vgE_B-pR)^2\}.
                                                               \tag{CR11}
\end{align*}
Every coefficient is bounded. The selected residue is
\[
 4\mu(V-E_B)^2+4\nu E_B^2+4cW^2+4aY^2
                    +2(W+R)^2>0.                 \tag{CR12}
\]

Finally, when both testing angles are infinite, also put $w=1/v$
and homogenize by $p^2w^2$. Use $(Q/g,B/g,W,X)$ and cap
$m,\beta,C$. Substitution $Y=gE_Q-pW$, $Z=gE_B-pX$ in the
homogenized core is polynomial and has residue $X^2$.
The selected residue is
\[
 4\mu E_Q^2+4\nu E_B^2+4cW^2+2X^2>0.               \tag{CR13}
\]
Each displayed positive form is strict on the complete root space.
The bounded positive-residue Gram lemma therefore supplies tensor
SOS. Restoring the retained terms proves the original parent.
The bounded unit/unit testing chart is already covered by the
common noncommuting-residue theorem. This leaves exactly the
ordinary/small testing strip stated in the theorem. The exact gauges
$(u,W,X)\mapsto(-u,-W,-X)$ and
$(v,W,Y)\mapsto(-v,-W,-Y)$ preserve every event square and the core.
All homogeneous formulas allow $p=0$ or $w=0$ literally.
Real and imaginary state components are treated separately.
\end{proof}

The ordinary/small testing strip is completed in Section~\ref{component:pa0b854a04f49}. Together with the remote charts proved here, it covers the entire actual UI source. The actual UH source is treated by the separate complete conductor in Section~\ref{component:pf7ffa921793c}.



% Source formulas: pf7ffa921793c
\subsection{The ordinary/unbounded source}

\label{component:pf7ffa921793c}

\label{pf7ffa921793c:sec:r7-c1-uh}
All coefficients below are divided by the original positive pure-source
$d$. We treat precisely (NI9) and (NI10), with no alteration of the native
source (NC1)--(NC7). The revision-6 common-event theorem already covers
all tests except positive unit $u$ and $w=1/v$ infinitesimal.
Put $x=k^2$, $K=1+x$, $\sigma=1/\eta$, $z=1/\zeta=q^2$,
$b=r/z$, $H=h/\sigma$, and retain $\sigma>0$ (it need not be a unit).
On (NI9), $\sigma$ is infinitesimal and $z$ is a unit; on (NI10),
$z$ is infinitesimal and $\sigma$ is bounded. Throughout, $k,H$ are
positive units, $b,z,\sigma$ are bounded, $Q=1-xb>0$, and
\[
 r=zb,\quad R=1+xzb,\quad M=R-z+\sigma zH,
 \quad J=1+\sigma(xH-1+zb),\quad\theta=1-\sigma H,
 \quad E_s=1-xH.
\]
The exact source polynomials, all with bounded arguments, are
\begin{align*}
 P={}&xK+\sigma x+\sigma z+\sigma zb x^2
 +x^2\sigma zH(1-\sigma H)+2x\sigma^2zH
 +\sigma^2z^2bx^2H,\\
 V={}&z(1+x\sigma H)E_s+x(1-zb)(1+\sigma zH),\\
 B_\alpha={}&R\sigma zH+RxH+ME_s.
\end{align*}
Writing $m=m_{\rm old}/z$, $c=C/z$, $a=Az$, $\ell=\alpha z$,
\[
 c={K(1+\sigma zb)E_s\over2(1-zb)P},\quad
 m={Kx(1+\sigma zb)H(\sigma z+x)\over2MP},\quad
 a={V\over2bP},\quad
 \ell={x^2(1+\sigma zb)B_\alpha\over2QP}.                 \tag{UH1}
\]
Define the polynomial quotients
\[
 N_D={K(\sigma z+x)(1+\sigma zb)-JP\over\sigma},\qquad
 N_\beta={Kx(1+\sigma zb)(1+\sigma zH)-\theta P\over\sigma}.
\]
Then $D=\sigma N_D/(2JP)$ and
$\beta=\sigma N_\beta/(2\theta P)$, exactly.

\paragraph{The bounded full parent.}
Let $D_u=1+u^2$, $D_w=1+w^2$. Its six event amplitudes are
\[
 M_w=W-wX+uY-uwZ,\ R_w=uW-Y,\ Y,\
 J_w=uwW+uX-wY-Z,\ Y-wZ,\ B_w=X+uZ.
\]
Thus the exact parent is
\[
 F=4mzM_w^2+4czD_wR_w^2+4(a/z)D_uD_wY^2
 +4DJ_w^2+4(\ell/z)D_u(Y-wZ)^2+4\beta D_wB_w^2+2K_w, \tag{UH2}
\]
where, without any restriction to a state ray,
\begin{align*}
 K_w={}&u^2w^2Z^2-2u^2wYZ-2uwWZ+2uw^2XZ
       +u^2w^2W^2+2u^2wWX-2uw^2WY\\
      &+(D_u+w^2)(X^2+Y^2)-2uwXY.
\end{align*}
For a potentially unbounded weight choose a finite cap, keeping bounded
weights unchanged. Subtract and retain the following nonnegative
centered excesses:
\begin{align*}
 &4(m-m_*)[qM_w-kM(wZ/q)]^2,\\
 &4(c-c_*)D_w[qR_w-x(1-r)(wZ/q)]^2,\\
 &4(a-a_*)D_uD_w[(Y-xbwZ)/q]^2,
 \quad4(\ell-\ell_*)D_u[(Y-xbwZ)/q]^2.          \tag{UH3}
\end{align*}
The first line uses $qM_w$ and the second $qR_w$; thus these are
exactly excesses of the scaled weights in (UH1). Their remainders
are bounded because $mM,c(1-r),ab,\ell Q$ are bounded.
When $\sigma$ is a unit, additionally retain
\[
 4(D-D_*)(J_w+JZ)^2,
 \qquad4(\beta-\beta_*)D_w(X+uhZ)^2.               \tag{UH4}
\]
Here $DJ$ and $\beta\theta$ are bounded. When $\sigma$ is
infinitesimal, no (UH4) cap is needed: $D,\beta=O(\sigma^2)$,
as proved below. Denote the complete remainder by $G$.
Every retained row vanishes at the actual contact
\[
 u=k,\quad w={q\sqrt\sigma\over k},\qquad
 (W,X,Y,Z)=({\sqrt\sigma\over q},-kh,kbq\sqrt\sigma,1).
                                                               \tag{UH5}
\]
No comparison between independent source and test rates occurs.

\paragraph{Exact homogeneous conductor, including every first normal.}
Put $h=\sigma H$, $\delta=u-k$, and use the six roots
\[
\begin{gathered}
 N_x=X+khZ,\quad N_y=(Y-kbq^2W)/q,\\
 U=wW-\sigma Z/k,\quad N=(wZ-q^2W/k)/q,\\
 S=q\delta W,\quad T=\sigma\delta Z.
\end{gathered}                                     \tag{UH6}
\]
Substitute $X=N_x-khZ$, $Y=qN_y+kbq^2W$ in the entire
bounded remainder $G$. Its complete quadratic normal block is the
coefficient collection in $(N_x,N_y)$. Its scalar and first-normal
coefficients have exactly the following supports:
\begin{align*}
 G_{00}&={A\over q^2}w^2Z^2+B w^2W^2+C\sigma^2 Z^2
                  +D_1 wWZ+E q^2W^2,\\
 G_{10}&=A_xwW+B_xZ+C_xw^2Z,\\
 G_{01}&=q(A_y+B_yw^2)W+C_ywZ/q.                    \tag{UH7}
\end{align*}
All eleven displayed coefficient families and the quadratic normal
coefficients are bounded polynomials in $u$, with bounded source
coefficients; $B_x$ is divisible by $\sigma$. In the unit-$\sigma$
case divisibility means a bounded quotient, not a new polynomial
identity. Formulas below specify every coefficient before caps, and
(UH3)--(UH4) specify their exact changes. No state normal is omitted.

Set $s=\sigma/k$, $a_0=q^2/k$ and
\[
 L=2A/k+2B\sigma/k+D_1,\quad
 J_0=C-B/k^2,\quad K_0=A/k^2+2B\sigma/k^2+D_1/k+E.
\]
The scalar is identically
\[
 A N^2+B U^2+qLWN+\sigma^2J_0Z^2+q^2K_0W^2.       \tag{UH8}
\]
Exact source contact and all stationary jets give
\[
 L(k)=J_0(k)=K_0(k)=0,\qquad K_0'(k)=-\sigma J_0'(k),
\]
\[
 B_x(k)+s[A_x(k)+a_0C_x(k)]=0,\qquad
 A_y(k)+C_y(k)/k+s a_0B_y(k)=0.                    \tag{UH9}
\]
These hold separately for every centered excess, hence for every cap.
Write the monic divisions
\[
 L=\delta L_1,\quad J_0=\delta j+\delta^2J_2,\quad
 K_0=-\sigma\delta j+\delta^2K_2,
\]
\[
 B_x+s(A_x+a_0C_x)=\sigma\delta H_x,\qquad
 A_y+C_y/k+s a_0B_y=\delta H_y.
\]
All quotients are bounded. The scalar and full first normals become
\begin{align*}
 G_{00}={}&A N^2+B U^2+(L_1+k\sigma j)SN-kjTU
                          +J_2T^2+K_2S^2,\\
 G_{10}={}&(A_x+a_0C_x)U+C_xqwN+H_xT,\\
 G_{01}={}&(C_y+s q^2B_y)N+qB_ywU+H_yS.            \tag{UH10}
\end{align*}
Together with the unchanged full quadratic normal block, this is an
exact bounded six-root Gram for $G$, not a scalar restriction.
The only root relation needed for strict completion is
\[
 T^2=\sigma S^2+k\sigma\delta SN-k\delta TU.        \tag{UH11}
\]
Its verification is one multiplication using (UH6). Adding any bounded
multiple of its zero changes no parent. No mixed inverse, root, order,
or source/test identification occurs in (UH6)--(UH11).

For the cap argument it is useful to print every centered row directly
in this same frame. In the order $m,c,(a,\ell),D,\beta$, they are
\begin{align*}
 &uzN_y-(kR+z\delta)N-qwN_x-(zQ/k)S,\\
 &S-zN_y-x(1-r)N,\qquad N_y-xbN,\\
 &(u-kbz)U+uN_x-qwN_y+(E_s/k)T,\qquad N_x+HT.
                                                               \tag{UH20}
\end{align*}
These are exact identities before taking any residue.

\paragraph{Explicit coefficient input.}
For (UH2) before caps put $j_u=u-kbz$, $p_u=1+ukbz$,
$t_u=kh-u$. The full scalar coefficients in (UH7) are
\begin{align*}
 A={}&4mz^2t_u^2+4\ell D_u+(4\beta+2)zt_u^2,\\
 B={}&(4cz+4D+2)j_u^2+4aD_uxb^2z,\\
 C={}&[4D(1+ukh)^2+4\beta(u-kh)^2+2D_uxh^2]/\sigma^2,\\
 D_1={}&8mzp_ut_u-8Dj_u(1+ukh)-8\ell D_ukb
          -4u^2k(h+bz)-4u+4uxhbz,\\
 E={}&4mp_u^2+4cj_u^2+4(a+\ell)D_uxb^2+2D_uxb^2z.
\end{align*}
The first normals are
\begin{align*}
 A_x={}&-8mzp_u+(8D+4)uj_u,\\
 B_x={}&-8Du(1+ukh)+8\beta(u-kh)-4D_ukh,\\
 C_x={}&-8mzt_u+(8\beta+4)(u-kh),\\
 A_y={}&8mzup_u-8czj_u+8(a+\ell)D_ukb+4D_ukbz,\\
 B_y={}&-8czj_u+8aD_ukb-(8D+4)j_u,\\
 C_y={}&8mz^2ut_u-8\ell D_u+8Dz(1+ukh)+4uzt_u.
\end{align*}
The coefficients of $N_x^2,N_y^2,N_xN_y$ are respectively
\[
 4mzw^2+4Du^2+4\beta D_w+2(D_u+w^2),
\]
\[
 4mz^2u^2+4cz^2D_w+4aD_uD_w+4\ell D_u+4Dzw^2
                            +2z(D_u+w^2),
 \qquad-8mzuqw-8Duqw-4uqw.                         \tag{UH12}
\]
Here products such as $Du$ mean $D\,u$, whereas $D_u=1+u^2$.

\paragraph{Infinitesimal $\sigma$: source collapse and strict residue.}
The exact polynomial $P_m=(N_D+xN_\beta)/\sigma$ is bounded.
At $\sigma=0$, $N_D=-xN_\beta$ and
\[
 N_\beta=x(K+z)H+bxz-x-z.
\]
Nonnegativity of $D,\beta$ forces
\[
 H=H_0+O(\sigma),\quad
 H_0={x+zQ\over x(K+z)},\qquad E_0=1-xH_0={R\over K+z}.
                                                               \tag{UH13}
\]
Indeed $JD+x\theta\beta=\sigma^2P_m/(2P)$ and $J,\theta$
are units. Thus $D,\beta=O(\sigma^2)$, including either literal zero.
Moreover
\[
 (D+x\beta)/\sigma^2\longrightarrow KH_0E_0/2>0.
\]
The other limiting coefficients are
\[
 c_0={R\over2x(1-zb)(K+z)},\quad m_0={xH_0\over2(1-zQ)},
 \quad a_0^{\rm wt}={H_0\over2b},\quad
 \ell_0={xR\over2Q(K+z)}.                           \tag{UH14}
\]
At a matched angle $u=k$, the residue Gram (UH10) splits into
$(N_y,N,S)$ and $(N_x,U,T)$, because their only off-block coefficients
contain $w$. Before caps, the $(N_y,N)$ block is positive definite;
its angular Schur complement is
\[
 \Gamma=8RH_0/K>0.                                \tag{UH15}
\]
The $(N_x,U)$ block is precisely
\[
 2K(N_x+xH_0U)^2+2KxH_0E_0U^2.                   \tag{UH16}
\]
The remaining $T$ couplings are bounded. As $\sigma,\delta$ have
zero residue, (UH11) permits an arbitrarily large positive $T^2$
pivot without changing any other residue coefficient. This proves a
strict full six-root residue Gram. For $u^2\ne x$, use the full raw
frame $(E_t,F_t,G_t,R_t,X,T_0)=(qW,Y/q,wZ/q,wW,X,\sigma Z)$.
Its leading first block is
\begin{align*}
 &4m_0[E_t+uz(F_t-G_t)]^2+4c_0(uE_t-zF_t)^2
       +4a_0^{\rm wt}D_uF_t^2+4\ell_0D_u(F_t-G_t)^2\\
 &\qquad+2z\{u^2(G_t-F_t)^2+F_t^2\}-4uE_tG_t.
\end{align*}
Put $v_0=u^2$ and $M_0=1-zQ$. In $(F_t,F_t-G_t)$ the hard block
is diagonal, with weights
\[
 B_h={2(x+zQ)[v_0(1-r)+R]\over bx(1-r)(K+z)},\qquad
 L_h={2R[v_0(x+zQ)+xM_0]\over Q(K+z)M_0}.
\]
The exact $E_t$ Schur complement is
\[
 {2(v_0-x)^2(1+v_0)RH_0\over
 [v_0(1-r)+R][v_0(x+zQ)+xM_0]}.                  \tag{UH17}
\]
It has a positive-unit limit through every opening when $v_0\ne x$.
The hard minimizers are
\[
 {F_t\over E_t}={ubK\over v_0(1-r)+R},\qquad
 {F_t-G_t\over E_t}={-uKQ\over v_0(x+zQ)+xM_0},
\]
also bounded through every opening. The other raw blocks are
$2\{u^2R_t^2+2u^2R_tX+D_uX^2\}$ and
$4(D+u^2\beta)T_0^2/\sigma^2$, both strictly positive.
This gives the full bounded away-angle lift, without a determinant-only
claim or any omitted first state normal.

\paragraph{Unit $\sigma$: angular Schur sign.}
Here (NI10) gives $z=0$ at residue. Put $h=\sigma H$,
$J=1+\sigma(xH-1)$, $\theta=1-\sigma H$. The source signs imply
\[
 0<\sigma<K,\qquad {1\over K+\sigma}\le H
       \le{x+\sigma\over x(K+\sigma)}.             \tag{UH18}
\]
These inequalities are stated on the interior; its closure includes
$\sigma=K$ when $J=\theta=0$.
The limits are
\[
 m={KxH\over2(K+\sigma)},\quad c={K(1-xH)\over2x(K+\sigma)},
 \quad a={1\over2b(K+\sigma)},\quad
 \ell={x\over2Q(K+\sigma)},
\]
\[
 D={K\over2J(K+\sigma)}-\tfrac12,
 \quad\beta={K\over2\theta(K+\sigma)}-\tfrac12.
\]
The two hard blocks in (UH10) are positive definite. Their angular
Schur coefficients $\gamma_S,\gamma_T$ have no $ST$ term and satisfy
\[
 \gamma_S+\sigma\gamma_T
 ={ -2p(H)\over\sigma(K+\sigma)}>0,\qquad
 p(H)=Kx(K+4\sigma)H^2-2K(x+2\sigma)H+x.           \tag{UH19}
\]
Convexity of $p$ and (UH18) prove strict negativity: its endpoint
values are respectively
\[
 {-\sigma(3\sigma x+4\sigma+4x+4)\over(K+\sigma)^2},\quad
 {-\sigma(4\sigma x+3\sigma+4x^2+4x)\over x(K+\sigma)^2}.
\]
Use (UH11) to distribute this positive total between $S^2,T^2$;
for example choose its multiple so both resulting Schur coefficients
are positive. This yields the required strict six-root Gram.

For completeness, angles $u^2\ne x$ are strictly positive as well.
Eliminating the two hard state variables from the leading form,
multiplying by the positive source factor $(K+\sigma)/2$, and
writing $U_0=u^2$, $P_0=(w/q)(Z/q)$, $T_0=(Z/q)^2$ gives a positive
Laurent function
\[
 A+B P_0^2+C P_0^2/T_0+D_0T_0-2E P_0,
\]
where
\[
 A=KxH+K(1-xH)U_0/x,\ B=x(1+U_0),\
 C={KU_0\over J+\theta U_0},
\]
\[
 D_0=(1+U_0)\left[{K(1+U_0)\over J+\theta U_0}-(K+\sigma)\right],
 \quad E={K\sqrt{U_0}(1+U_0)\over J+\theta U_0}.
\]
The strict criterion is $E<\sqrt{AB}+\sqrt{CD_0}$.
Squaring only to analyze these positive real residue quantities gives
\[
 (E^2-AB-CD_0)^2-4ABCD_0
 ={K^2(1+U_0)^2\over(J+\theta U_0)^2}(U_0-x)^2Q_2(U_0).
\]
The quadratic $Q_2$ has positive constant $H^2x^2J^2$, positive
leading coefficient $\theta^2(1-xH)^2$, and $Q_2(x)=xp(H)<0$.
Explicitly, with $E_s=1-xH$,
\[
 Q_2(U_0)=E_s^2\theta^2U_0^2
 +2E_s\{x\theta^2+E_s\theta J+(K+\sigma)\theta-2K\}U_0
 +x^2H^2J^2.
\]
The concave quadratic
$U_0(K+\sigma)-[x^2H+(1-xH)U_0](J+\theta U_0)$
is negative at zero and infinity and positive at $x$; at each of its
two zeros the preceding discriminant is negative. Consequently its
positive interval lies inside $Q_2<0$. This proves the criterion for
all $U_0\ne x$. Sign charts with $P_0\le0$ are immediate. These
square roots are an analysis of real residues only and are not tensor
roots or mixed localizations.

This inequality is directly a bounded Gram allocation, not a
comparison of mixed ratios. Before choosing projective coordinates,
put $E_t=qW$, $G_t=wZ/q$, $R_t=wW$. The exact relation
$E_tG_t=R_tZ$ permits distributing the cross term between the two
quadratic blocks $(E_t,G_t)$ and $(R_t,Z)$. The strict inequality
above is exactly the condition that both $2\times2$ blocks can be
made positive definite by a finite real allocation. It therefore
handles every independent rate of $q,w$.
At $J=0$ or $\theta=0$, the root-interval argument permits one endpoint
at zero or infinity. At the double face $J=\theta=0$, necessarily
$\sigma=K$, $H=1/K$. As $J+\theta U_0\to0$,
\[
 E-\sqrt{CD_0}\longrightarrow K\sqrt{U_0},\qquad
 AB-K^2U_0=(U_0-x)^2.
\]
Thus the constrained allocation remains strictly positive away from
$U_0=x$, including this double face without dividing by a vanishing
residue denominator.

\paragraph{Opening faces, tensor conclusion, and exact scope.}
The finite-cap argument is applied to the full Gram above. For
$\sigma$ infinitesimal the possible diverging rows have openings
$1-zb$, $1-zQ$, $b$, $Q$; at most two can diverge together. The pairs
$(c,m)$, $(m,a)$, $(c,\ell)$ have independent hard rows; other pairs
are excluded by $0<zb,xb<1$ and $1-zQ>0$. In the $(N_y,N)$ columns
of (UH20), their determinants have absolute values $kzK,k,x$,
respectively; each is a unit on its compatible face. The hard block
is uniformly coercive, since it contains positive-unit multiples
of $N_y^2$ and $(N_y-N)^2$. Its minimizing hard coordinates for
fixed angular $S$ are explicitly
\[
 {N_y\over S}={b(2bxz-x+1)\over K},\qquad
 {N\over S}={4bxz-x-2z+1\over xK}.
\]
They are bounded through every opening. Thus the limit of the
uncapped Schur complement is the Schur complement restricted to the
diverging centered rows' kernel. The quotients (UH15)--(UH16) stay
positive units there. The elementary finite-dimensional penalty
lemma now applies: a bounded symmetric form strictly positive on
the common kernel becomes strictly positive after sufficiently large
finite positive multiples of these rows' squares are added. This
proves the finite-cap claim, with the entire angular right-hand sides
in (UH20) retained. Away from $u^2=x$, (UH17) and its printed bounded
hard minimizers give the same constrained strictness directly.

For unit $\sigma$, only $b,Q,J,\theta$ can be openings.
The $b,Q$ rows coincide in the hard direction $N_y-xbN$ and cannot
both diverge. The $J,\theta$ rows on $(N_x,U)$ are independent even
if both vanish; their determinant is the unit $u$. Formula (UH19)
remains a positive unit through these faces. More explicitly, the
two hard blocks dominate positive-unit multiples of
$N_y^2+(N_y-N)^2$ and $N_x^2+(N_x+U)^2$. Their angular minimizers are
\[
 {N_y\over S}={b(1-HxK)\over K},\quad
 {N\over S}={1-HxK\over xK},
\]
\[
 {N_x\over T}={\sigma H^2xK-HxK-\sigma H(2x+1)+x\over\sigma K},
 \quad {U\over T}={HxK(K+\sigma)-xK-\sigma\over x\sigma K}.
\]
Their $D$ and $\beta$ row defects divided by $T$ are respectively
$k(HK-1)J/(\sigma K)$ and $-x\theta(HK-1)/(\sigma K)$.
Hence every relevant constrained limit exists with bounded
minimizers; the same penalty argument applies. On the angular
relation $T^2=\sigma S^2$, its strictly positive Schur value is
(UH19), and after taking finite caps the bounded zero (UH11)
distributes that value into two strict angular pivots. The away-angle
allocation and its single/double-face limits above give the same
conclusion there. All removed excesses remain exact positive
source/test squares. No positive lower bound on an original opening
has been imposed.

A strict bounded residue Gram lifts by the common tensor Gram
lemma. Every root in (UH6) is a finite sum of separated source and
testing factors. After restoring the retained squares this certifies
the complete original parent, including $D=0$ or $\beta=0$.
Together with revision-6 remote tests, this closes both (NI9) and
(NI10). Those are the entire actual UH portion of C1.5. The already
proved UU and UI portions and exact source stabilizers therefore close
the entire C1.5 row, subject to the exact checks and mathematical review
recorded alongside this section.



% Source formulas: p7d4e5164e84a
\section{Vanishing contact scales and the state normal directions}
When $k$ vanishes, some contact directions become singular. The following polynomial identities retain their exact factors before any rescaling. They control the scalar term and both state normal directions, including the complementary state charts.

\subsection{The endpoint polynomial source}

\label{component:p7d4e5164e84a}
\paragraph{Seven short source polynomials.}
The following definitions make the endpoint input reproducible without a stored expansion. Put
\[
K=k^2,\quad \eta=e^2b^2,\quad\zeta=k^2b^2,\quad V_c=K\eta\zeta.
\]
\[
L_c=1-a\eta,\qquad H_c=1+(1-a)\eta,\qquad Q_c=1-K\zeta+aV_c.
\]
Define
\begin{align*}
P_4={}&1+K+\eta K(a-K-1)-\zeta Kt
 +V_c[-\zeta Kt^2+2Kt+a(1-t)+t],\\
P_\alpha={}&1+K-a\eta K+\zeta[(1+K)-(2+K)t]
 +V_c[-\zeta t^2+a(t-1)+t],\\
P_d={}&(1+K)^2-a\eta K^2+\zeta K[(K+1)-(K+2)t]\\
&+V_c[-\zeta Kt^2+Ka(t-1)+Kt+K+1].
\end{align*}
The seven coquotients are
\begin{align*}
c_2&=2(1+K)(1-t\zeta)L_c(1+V_c)H_cQ_c,\\
c_\beta&=2(1+K)L_c[1+(1-t)\zeta]H_cQ_c,\\
c_3&=2(1+K)L_cH_cQ_c(1-K\eta+tV_c),\\
c_5&=2(1+K)L_cH_cQ_c(1+V_c),\\
c_4&=2Q_cP_4,\qquad c_\alpha=2L_cH_cP_\alpha,
\qquad c_d=4L_cQ_cP_d.
\end{align*}
Their residues are respectively $2,2,2,2,2,2,4$, so all are positive coefficient units on this domain. With $M=R_c-t$ and $J=T_0-a$, the unchanged source is
\begin{gather*}
\lambda=n=0,\quad c=2,\qquad
d=\frac{2a\eta c_d}{c_3},\quad m=\frac{2ae^2c_2}{Mc_3},\quad
A=\frac{2a\eta^2c_4}{c_3},\quad
\alpha=\frac{2aK^2\eta\zeta c_\alpha}{c_3},\\
D+\frac d2=\frac{2ac_5}{Jc_3},\qquad
\beta+\frac d2=\frac{2ae^2c_\beta}{tc_3},\qquad
\Lambda=\frac{R_cT_0c_3}{2a}>0.
\end{gather*}
The theorem retains the hypothesis $D,\beta\ge0$. Write
\[
r_2=v-X+uvY-uZ,\quad r_\beta=X+uZ,\quad r_3=u-Y,
\quad r_5=u+uvX-Y-vZ,\quad r_\alpha=vY-Z,
\]
and put $q_0=v_g$, $z_0=\eta u$. Then $G$ is exactly the sum of
\begin{align*}
G_2&=4T_0e^2c_2[(r_2+tq_0)^2-tR_cq_0^2],\\
G_\beta&=4T_0e^2c_\beta(1+v^2)[(r_\beta+Mq_0)^2-R_cMq_0^2],\\
G_3&=4R_cc_3(1+v^2)[(r_3+Jz_0)^2-T_0Jz_0^2],\\
G_5&=4R_cc_5[(r_5-az_0)^2-aT_0z_0^2],\\
G_4&=4R_cT_0\eta^2c_4(1+u^2)(1+v^2)Y^2,\\
G_\alpha&=4R_cT_0K^2\eta\zeta c_\alpha(1+u^2)r_\alpha^2,\\
G_d&=2R_cT_0\eta c_d(1+u^2)[r_\alpha^2-(1+v^2)Z^2].
\end{align*}
These seven factored expressions agree separately with the defining source. Their proof requires multiplication of short polynomials only.



% Source formulas: p935c9d6318ea
\subsection{Finite monic divisions at the endpoint}

\label{component:p935c9d6318ea}
\paragraph{Deriving the scalar and first-normal inputs by monic division.}
Here is a finite construction of every coefficient in the preceding table. It uses the seven factored terms, differentiation, and division by monic polynomials. Thus neither a stored expansion nor a rational-function ideal calculation is an input.

Let $\mathcal R=\Q[a,t,b,e,k]$ and use an indeterminate $U$ for $u^2$ in this paragraph only. Put $K=k^2$. Substitute $(X,Y,v)=(X_g,Y_g,v_g)$ in the peeled polynomial $E$. Directly from the seven terms, the resulting expression has the form
\[
\frac{E_g}{e^2b^4}=e^2\mathcal A(U)+ekZ\mathcal B(U)+k^2Z^2\mathcal C(U),
\qquad \mathcal A,\mathcal B,\mathcal C\in\mathcal R[U].
\]
The displayed division means coefficientwise polynomial extraction: the constant, linear and quadratic coefficients in $Z$ have the factors $e^2$, $ek$ and $k^2$, respectively. No evaluated tensor is divided. Substituting the graph into the factored terms gives the five short contact and axis identities
\begin{gather*}
(\mathcal A+\mathcal B+\mathcal C)(K)=0,\qquad
(\mathcal B+2\mathcal C)(K)=0,\qquad
(\mathcal A+\mathcal B+\mathcal C)'(K)=0,\\
\mathcal A(0)=\mathcal B(0)=0.
\end{gather*}
Consequently the following monic divisions have zero remainder:
\begin{align*}
\mathcal N&=\frac{\mathcal A+\mathcal B+\mathcal C}{(U-K)^2},&
\mathcal D&=\frac{\mathcal B+2\mathcal C}{U-K},\\
\mathcal R_0&=\frac{\mathcal C+K\mathcal D+K^2\mathcal N}{U},&
\mathcal J&=\frac{-\mathcal D-2K\mathcal N}{U}.
\end{align*}
For the last two divisions, evaluate the numerators at $U=0$ and use the first two division identities together with $\mathcal A(0)=\mathcal B(0)=0$. This is a polynomial-ring argument; the symbols $K$ and $U$ are not inverted after evaluation. Set
\[
A_{00}=\mathcal R_0(u^2),\qquad A_{01}=u\mathcal J(u^2),
\qquad A_{11}=\mathcal N(u^2).
\]
Comparing the three coefficients of $1,Z,Z^2$ proves exactly
\[
\frac{E_g}{e^2b^4}
=A_{00}\,[u(e-kZ)]^2+A_{01}\,u(e-kZ)(eu^2-k^3Z)
 +A_{11}(eu^2-k^3Z)^2.
\]
Multiplying by $e^2b^4$ gives the scalar identity used above.

For each $\xi\in\{X,Y,v\}$, let $R_\xi$ be the graph restriction of $\partial_\xi E$, divided coefficientwise by $e^2b^2,eb^2,e^2b^2$, respectively. The first two are linear in $Z$. The third has one quadratic coefficient, supplied solely by the signed term:
\[
[Z^2]R_v=-4R_cT_0c_d\,e k b^2u(1+u^2).
\]
Define
\[
C_X=C_Y=0,\qquad C_v=4R_cT_0c_d\,eb^2(1+u^2).
\]
Then $\widetilde R_\xi=R_\xi-C_\xi Z\,u(e-kZ)$ is linear in $Z$ and has the polynomial parity form
\[
\widetilde R_\xi=eu\,L_{\xi,0}(U)+kuZ\,L_{\xi,1}(U),
\qquad L_{\xi,i}\in\mathcal R[U].
\]
In particular, subtracting the quadratic term also leaves its corrected linear coefficient divisible by $ku$; this follows upon collecting the signed term with the two current terms. The contact identity is
$\big(L_{\xi,0}+L_{\xi,1}\big)(K)=0$. Form the monic quotient
\[
J_\xi=\frac{L_{\xi,0}+L_{\xi,1}}{U-K}.
\]
The promised full first-normal coefficients are now defined by the short formulas
\[
P_\xi=(L_{\xi,0}-UJ_\xi)(u^2)+C_\xi Z,
\qquad Q_\xi=uJ_\xi(u^2).
\]
Indeed direct multiplication gives
$R_\xi=P_\xi u(e-kZ)+Q_\xi(eu^2-k^3Z)$.
This proves the full first-normal identities, including the quadratic-in-$Z$ term; no first-normal term has been omitted.

For completeness, the weighted splits have an equally explicit finite rule. In a polynomial expanded in $(e,k,b,u)$, place a monomial with $k$-exponent at least $r$ in the first bucket and divide it by $k^r$; place every remaining monomial in the second bucket and divide it by $e^s b^s$. Use $(r,s)=(4,2)$ for $A_{00}$ and $(r,s)=(2,1)$ for
\[
A_{01},\quad P_X,\quad P_Y,\quad P_v-16kP_0.
\]
Substitution of the seven short terms into the monic formulas shows that every monomial left for the second bucket has the required $e^s b^s$ factor. Equivalently, the only polynomial remainders to check are
\[
A_{00}\equiv0\pmod{(k^4,e^2b^2)},\qquad
A_{01},P_X,P_Y,P_v-16kP_0\equiv0\pmod{(k^2,eb)}.
\]
These are ordinary finite polynomial truncations: discard the monomials divisible by the displayed generators. The two bucket quotients have respectively residues $(-8,-8)$ for $A_{00}$ and $(0,0)$ for each of the other four entries; $A_{11}$ has residue $8$, and all $Q_\xi$ have residue zero. To obtain these numbers, after the divisions set $e=k=u=Z=0$, keeping $a,t,b$ free. In particular the new term $C_vZ$ belongs to $eb\,\mathcal R[u,Z]$ and its second-bucket quotient contains the infinitesimal testing factor $Z$. No assumption $\bar b=0$ is used.

This is a reproducible coefficient construction: the source consists of seven short terms, the scalar construction of four monic quotients, and each first-normal construction of one monic quotient. The finite checks accompany these formulas as a convenience; the definitions and the multiplication identities, rather than expanded coefficient arrays, are the proof objects.



% Source formulas: p05180564ed05
\subsection{The endpoint product relation}

\label{component:p05180564ed05}
\label{p05180564ed05:sec:productrepair}
The new boundary mechanism is algebraic. A troublesome root is retained, and a second relation controls its product with an additional root. This supplies two positive pivots without comparing an infinitesimal in one field with an infinitesimal in the other.

\begin{lemma}[Product-relation completion]\label{p05180564ed05:lem:productrepair}
Let $z$ be a root vector over $T$, and $h,r\in T$. Suppose
\[
F=(z,h,r)^TQ(z,h,r),\qquad Q\in\Sym(B),\qquad
\bar Q=\operatorname{diag}(Q_0,-a,0),
\]
where $Q_0$ is a real positive-definite matrix and $a\in\R_{\ge0}$. Suppose $p\in B$ has nonzero residue, $A$ is a row with entries in $\mathfrak n$, and
\[
ph=Az+r,\qquad phr=(z,h,r)^TE(z,h,r),\quad E\in\Sym(\mathfrak n).
\]
Then $F\in\Sigma T^2$.
\end{lemma}
\begin{proof}
For a sufficiently large positive real $M$, add the exact zero
\[
M\big[p^2h^2+r^2-(Az)^2-2phr\big]=0.
\]
The residue Gram becomes $\operatorname{diag}(Q_0,M\bar p^2-a,M)\succ0$. The subtracted terms have zero residue by hypothesis. Lemma~\ref{p25af2c58da7e:lem:gram} applies. In particular $p$ is never inverted.
\end{proof}

\subsubsection*{Exact contact scope and full-polynomial input}
This application uses the canonical C1 contact model, with $\lambda=n=0$. Its source parameters are positive $e,k,b,a,t$ satisfying
\[
e,k\in\mathfrak m_c,\quad b\in B_c,\qquad
R_c=1+k^2(1-ae^2b^2),\quad T_0=1+k^2-tk^4b^2,
\]
\[
0<a<T_0,\qquad0<t<R_c.
\]
The signed weight satisfies $d>0$, and all eight event coefficients are nonnegative original physical C1 contact coefficients. This is a specified contact chart, not every coefficient vector with the same omissions. The independent testing chart is $W=1$ with $u,v,X,Y,Z$ infinitesimal. In this section there is no restriction on $k^3/e$ or $k/e$.

The complete seven-term contact polynomial is denoted $G$. Its short factored definition and unchanged original coefficient vector are given below. 

Define
\begin{gather*}
v_g=ekb^2u,\quad X_g=-uZ+tv_g,\quad Y_g=(1-ae^2b^2)u,\\
N_1=X-X_g,\quad N_2=Y_g-Y,\quad \delta v=v-v_g,\quad D_t=Z-uX,\\
U_0=1+u^2,\quad c_0=1-tk^2b^2,\quad P_0=T_0+u^2=U_0+k^2c_0,\\
S=u+uvX-Y-vZ+(T_0-a)e^2b^2u,\\
s_0=eb^2u(e-kZ),\qquad q=eb^2(eu^2-k^3Z).
\end{gather*}
\par\noindent\textit{Polynomial source. }The source is defined (Section~\ref{component:p7d4e5164e84a}).\par
Direct substitution gives the exact nonlinear relation
\[
L=P_0s_0-S+N_2-D_t\delta v+uv_gN_1-uc_0q=0. \tag{E1}
\]
Peel the two complete current squares:
\[
E=G-H_3(1+v^2)N_2^2-H_5S^2,\qquad
H_3=4R_cc_3,\quad H_5=4R_cc_5,\quad \bar H_3=\bar H_5=8.
\]
Each peeled term divided by $e^2$ is polynomial in bounded variables. The full graph scalar has the identity
\[
E_g=A_{00}s_0^2+A_{01}s_0q+A_{11}q^2,
\]
and the entire first normals, with $g_0=s_0/(eb^2)$ and $g_2=q/(eb^2)$ denoting their displayed polynomial expressions, satisfy
\[
E_X=e^2b^2(P_Xg_0+Q_Xg_2),\quad
E_Y=eb^2(P_Yg_0+Q_Yg_2),\quad
E_v=e^2b^2(P_vg_0+Q_vg_2).
\]
The exact polynomial splits and residues needed are
\begin{center}\small
\begin{tabular}{ccc}\toprule
Quantity&Polynomial split&Quotient residues\\\midrule
$A_{00}$&$k^4a_0+e^2b^2a_1$&$-8,-8$\\
$A_{01}$&$k^2c_1+ebc_2$&$0,0$\\
$P_X$&$k^2p_X+ebq_X$&$0,0$\\
$P_Y$&$k^2p_Y+ebq_Y$&$0,0$\\
$P_v-16kP_0$&$k^2p_v+ebq_v$&$0,0$\\\bottomrule
\end{tabular}
\end{center}
Also $\bar A_{11}=8$ and $\bar Q_X=\bar Q_Y=\bar Q_v=0$. These equalities remain valid with $b$ symbolic; boundedness, rather than infinitesimality, suffices for $b$.

\par\noindent\textit{Coefficient construction. }The finite coefficient construction is given (Section~\ref{component:p935c9d6318ea}).\par

Exact finite Taylor expansion in $(N_1,-N_2,\delta v)$ accounts for every higher-normal term. Because $E/e^2$ is polynomial, every term of normal degree at least two allocates to $x=eN_1$, $y=e\delta v$, $N_2$ with bounded coefficients. Terms of degree at least three have an extra infinitesimal normal. Terms involving $N_2$ retain an $e$ factor. The only additional nonzero residue is
\[
B_0=\begin{pmatrix}16&-8\\-8&8\end{pmatrix}
\quad\text{on }(x,y).
\]
It is visible without expansion of the full Hessian: at $u=Y=Z=0$, the quadratic residue of $E/e^2$ is $8(X-v)^2+8X^2$, with $b$ still free.

\subsubsection*{The new identity and its nine-root proof}
If the signed testing core is nonnegative, its pure testing square root and the original event terms already give SOS. Otherwise the sign gauges $u,v\ge0$ preserve every term, and the core Schur complement implies $u,v>0$ and $D_t>0$. To check the latter, after $Z=D_t+uX$ the Schur complement on $(W,D_t)$ is
\[
\frac{2(1+v^2)}{N_0}
\begin{pmatrix}
u^2v^2(1+u^2)&-uv(1+u^2)(1+v^2)\\
-uv(1+u^2)(1+v^2)&u^2v^2
\end{pmatrix},
\]
where $N_0=u^2v^4+u^2v^2+u^2+v^4+2v^2+1>0$. The eliminated block is positive definite. For $W=1,D_t\le0$, the remaining form is nonnegative, proving the claim.

Consequently $w=\sqrt{ek}\sqrt{D_t}\,\delta v$ is a legal separated root. Set
\begin{gather*}
h=k^2s_0-uq=e^2b^2u(k^2-u^2),\qquad z=ebs_0,\qquad r=k^2D_t\delta v,\\
A_0=k^2(S-N_2)-k^3b^2u^2x-uU_0q.
\end{gather*}
Multiplication of (E1) by $k^2$ gives $P_0h=A_0+r$. The additional exact identity is
\begin{align*}
P_0hr={}&ku^2(P_0-k^2)y(S-N_2)
-ku(P_0+c_0k^2u^2)yq\\
&+u^2(P_0-k^2)w^2-k^4b^2u^2T_0xy. \tag{E2}
\end{align*}
Every displayed coefficient has zero residue and is a polynomial in separately bounded factors. To derive (E2), substitute the two soft roots to obtain
\[
u(P_0-k^2)s_0-(1+c_0u^2)q
=ekb^2(k^2-u^2)(D_t+uN_1).
\]
Then use $D_t=U_0Z-uN_1-tu v_g$, $T_0=1+k^2c_0$, and (E1) to replace $S-N_2$ in (E2). The terms in $D_t(\delta v)^2$ cancel; the remaining terms give $P_0hr$. This is a whole-polynomial calculation.

Let $J_0$ denote the right-hand side of (E2). Thus
\[
Z_0=P_0^2h^2+r^2-A_0^2-2J_0=0.
\]
Add to the entire $G$ the exact zero expression
\[
-16kyL+16(ebL)^2+16Z_0.
\]
The first term cancels the exceptional $16kP_0$ first normal and supplies $16w^2$. Use $k^2s_0=h+uq$ in every other scalar and first-normal term. The scalar becomes
\[
a_0(h+uq)^2+a_1z^2+[c_1(h+uq)+c_2z]q+A_{11}q^2.
\]
Its only negative residue diagonals are $-8$ on $h,z$. The safe relation
\[
ebL=P_0z-ebS+ebN_2-bD_ty+ekb^3u^2x-ebu c_0q
\]
has only infinitesimal coefficients apart from $P_0$; hence $16(ebL)^2$ changes the $z$ diagonal to $8$. Every coefficient of $A_0$ is infinitesimal, as is every coefficient in (E2). Thus $16Z_0$ changes the $h$ diagonal to $8$ and supplies a new $r$ diagonal $16$, with no other residue change.

The full Gram on
\[
\boxed{(x,y,N_2,S,w,q,h,z,r)}
\]
therefore has residue
\[
\boxed{B_0\oplus\operatorname{diag}(8,8,16,8,8,8,16)\succ0.}
\]
Its first two LDL pivots are $16,4$; the other seven pivots are given above. The strict-residue lemma proves SOS of the complete seven-term polynomial for every coefficient ratio $k^3/e$.

\subsubsection*{Restoration and coverage}
The complete original-parent restoration is
\begin{align*}
\Lambda F={}&G
+4T_0e^2c_2\frac{t}{M}(r_2-Mq_0)^2
+4T_0e^2c_\beta\frac{M}{t}(1+v^2)(r_\beta-tq_0)^2\\
&+4R_cc_3\frac{J}{a}(1+v^2)(r_3-az_0)^2
+4R_cc_5\frac{a}{J}(r_5+Jz_0)^2.
\end{align*}
Every additional term is a nonnegative coefficient factor times a pure-testing factor times a square. This identity is four instances of
\[
S^2+2(A+B)qS+(A+B)Aq^2+\frac BA S^2
=\frac{A+B}{A}(S+Aq)^2.
\]
Use respectively $(A,B,q)=(M,t,q_0),(t,M,q_0),(a,J,z_0),(J,a,-z_0)$ and $S=r_2-Mq_0,r_\beta-tq_0,r_3-az_0,r_5+Jz_0$. Comparing the seven resulting coefficients with the displayed source proves the identity. Since $\Lambda$ is a positive pure coefficient factor, the pure-factor restoration in \hyperref[component:p9b00698ff36a]{RFA13} (p.~\pageref{component:p9b00698ff36a}) recovers the original parent. Normalizing the original state to $W=1$ and undoing that normalization use only testing factors.

The result applies to every positive bounded $b$, with no
condition on $k^3/e$. The unbounded-$b$ case is proved separately in
Section~\ref{component:pabbcd0905dcd}; the full source and state partition
is given in Section~\ref{component:p845b3cb8b114}.


% Source formulas: pa4f6d1e57f34
\subsection{The uniform small-$e$ source on the bounded $Z$ chart}

\label{component:pa4f6d1e57f34}

\label{pa4f6d1e57f34:sec:uniform-c1-z-new}

\paragraph{Domain and conclusion.}
Use the original physical uniform C1 source with $d=1$,
$\lambda=n=0$, $k,f,g,e=kf/g$ positive infinitesimals, and
all original nonnegative event weights. Retain
\[
 r=1-af^2,\quad R=1+k^2r,\quad
 T=1+k^2-tk^2g^2,\quad M=R-t>0,\quad J=T-a>0,
 \qquad a,t>0.                                                   \tag{EZ1}
\]
In particular $a,t$ are bounded, $r,R,T$ are positive units,
and $M,J$ may approach zero at arbitrary rates.
For every independent testing $u,v,X,Y$ infinitesimal and
$Z=1$, $W$ arbitrary bounded of either sign, the entire original
parent is a separated tensor SOS. This includes both bounded
and unbounded $b=g/k$, and every permitted $\beta=0$ face.
It does not claim the complementary $W=1,Z$-infinitesimal chart.

\subsubsection*{The exact source and a common set of bounded coefficients}

Use the seven short native coquotients
$c_2,c_\beta,c_3,c_5,c_4,c_\alpha,c_d$ in the retained
Section~\ref{component:p7d4e5164e84a}, now with $\eta=f^2,\zeta=g^2$.
Their residues are $2,2,2,2,2,2,4$.
In the present $d=1$ normalization the exact weights are
\[
\begin{gathered}
 C=\frac{c_3}{af^2c_d},\quad A=\frac{f^2c_4}{c_d},\quad
 m=\frac{k^2c_2}{g^2Mc_d},\quad
 \alpha=\frac{k^4g^2c_\alpha}{c_d},\\
 D+\tfrac12=\frac{c_5}{Jf^2c_d},\qquad
 \beta+\tfrac12=\frac{k^2c_\beta}{g^2tc_d}.
\end{gathered}                                                   \tag{EZ2}
\]
Every use below is of these actual weights. In particular
$\beta\geq0$ is retained, not replaced by a formal source condition.

The following two pure-source coordinate choices are exhaustive:
\[
\begin{array}{c|ccccc}
 &s&\omega&P&G&\kappa=P/G\\ \hline
 b\text{ bounded}&e&b^2&k^2&k^2&1\\
 b\text{ unbounded}&f&1&k^2&g^2&(k/g)^2
\end{array}                                                     \tag{EZ3}
\]
In both rows $f^2=s^2\omega$, $g^2=\omega G$,
$s,\sqrt P,\sqrt G$ are infinitesimal, and $\omega,\kappa$
are positive bounded. In the second row $\kappa$ is infinitesimal.
Here $\beta\geq0$ gives exactly
$t\leq2\kappa c_\beta/c_d$, so $t=\kappa\tau$ with
$\tau>0$ bounded; consequently $M$ is a unit in this row.

Define bounded source coefficients
\[
 c_*=\frac{c_3}{Tc_d},\quad d_*=\frac{c_5}{Tc_d},\quad
 a_*=\frac{c_4}{c_d},\quad \gamma_*=\frac{c_\alpha}{c_d},
 \qquad B_*=\omega\beta t
              =\frac{\kappa c_\beta}{c_d}-\frac{\omega t}{2}.
                                                               \tag{EZ4}
\]
The first four have residue $1/2$; $B_*\geq0$ may vanish.
Use the further table
\[
\begin{array}{c|ccc}
 &\mu&\chi&\psi\\ \hline
 b\text{ bounded}&c_2/(Rc_d)&t&tM\\
 b\text{ unbounded}&\kappa c_2/(Mc_d)&0&0 .
\end{array}                                                     \tag{EZ5}
\]
All entries are bounded. In the first row $\mu$ has residue
$1/2$; in the second $\mu$ is infinitesimal.

\subsubsection*{Exact removal of the source poles}

The original quarter-parent is
\[
\begin{split}
 F={}&mr_2^2+CD_vr_3^2+AD_uD_vY^2
 +(D+\tfrac12)r_5^2+(\alpha+\tfrac12)D_u(vY-1)^2\\
 &+(\beta+\tfrac12)D_v(X+u)^2-\tfrac12D_uD_v,
\end{split}                                                     \tag{EZ6}
\]
where $r_2=vW-X+uvY-u$, $r_3=uW-Y$,
$r_5=uW+uvX-Y-v$, $D_u=1+u^2$, $D_v=1+v^2$.
Put
\[
 z=sW,\qquad q_0=f^2uW^2=\omega uz^2,\qquad
 z_0=f^2uW=s\omega uz.
\]
Retain the following complete nonnegative terms of the original
polynomial:
\[
 \frac{CJ}{T}D_v(r_3-az_0)^2,\qquad
 \frac{(D+1/2)a}{T}(r_5+Jz_0)^2,\qquad
 \beta D_v(X+u-tq_0)^2.                                        \tag{EZ7}
\]
For bounded $b$, also retain
\[
                       \frac{mt}{R}(r_2-Mq_0)^2.              \tag{EZ8}
\]
Each vanishes at the actual source contact. They are literal
positive squares, even when $\beta=0$ or a source opening is small.
The first two removals use $a+J=T$ and the exact identity
\[
 \frac aT[(L+Jq)^2-TJq^2]+\frac JT(L-aq)^2=L^2.
\]
The last removal uses
$r_2^2-(t/R)(r_2-Mq_0)^2
=(M/R)(r_2^2+2tq_0r_2-tMq_0^2)$.
There is no mixed denominator.

Use all three unrestricted normal coordinates
\[
 N_x=X+u-tq_0,\qquad
 N_y=(Y-ruW)/s,\qquad N_v=(v-f^2uW)/s.
                                                               \tag{EZ9}
\]
Only $N_x$ is necessarily bounded. The exact bounds needed later are
$N_x,sN_y,sN_v\in\mathrm I$; they follow from the stated testing
chart and bounded $W$. No boundedness of $N_y,N_v$ is assumed.
Define
\[
 X_*=-u+tq_0+N_x,\qquad
 y_*=ruz+s^2N_y=sY,\qquad
 v_*=\omega uz+N_v=v/s,
\]
\[
 L_2=zv_*+uv_*y_*-X_*-u,\quad
 L_3=T\omega uz-N_y,\quad
 L_5=-N_y+v_*(uX_*-1).
\]
After multiplying the original quarter-parent by the positive
pure-source $\omega$ and removing precisely the scaled terms
EZ7--EZ8, the \emph{entire} remaining polynomial is
\[
\begin{split}
\mathcal F={}&
 \mu(L_2^2+2\chi q_0L_2-\psi q_0^2)
 +c_*D_v(L_3^2-TJ\omega^2u^2z^2)\\
 &+\omega^2a_*D_uD_vy_*^2
 +d_*(L_5^2-aT\omega^2u^2z^2)\\
 &+(\omega^2P^2G\gamma_*+\omega/2)
                         D_u(v_*y_*-1)^2
 +\frac{\omega}{2}D_v(X_*+u)^2\\
 &+B_*D_v\{2q_0(X_*+u)-tq_0^2\}
                         -\frac{\omega}{2}D_uD_v,
 \qquad D_v=1+s^2v_*^2 .
\end{split}                                                     \tag{EZ10}
\]
This is a short full-parent definition of every scalar and normal
coefficient used below. In particular no normal was set to zero
in obtaining it. All its coefficient arrays are bounded through
$a,J,M,t$; the unbounded original weights occur only in the
retained complete squares.

\subsubsection*{The entire scalar is a cubic circuit with finite corrections}

Set $U=u^2,V=z^2$. The polynomial obtained from EZ10 at
$N_x=N_y=N_v=0$ is exactly
\[
 \mathcal F_0=\omega^2\Phi(U,V),\qquad
 \Phi=\gamma_*P^2G(1+U)+UV\,H(U,V),                            \tag{EZ11}
\]
where $H$ has bounded finite polynomial coefficient arrays.
The apparent divisions by $\omega^2$ and $UV$ in this definition
are coefficientwise polynomial extractions: substitution in EZ10
shows those factors before any evaluation. No testing coordinate
is inverted.

At the actual source contact $u^2=P,z^2=G$ all three normal
coordinates and all retained terms vanish. The original contact
value and stationary equations give
\[
 \Phi(P,G)=\Phi_U(P,G)=\Phi_V(P,G)=0.                            \tag{EZ12}
\]
The exact native coquotients in EZ2 independently verify these
identities. Substituting EZ11, and cancelling nonzero pure-source
factors in the source field, gives
\[
 H(P,G)=-\gamma_*P(1+P),\quad
 H_U(P,G)=\gamma_*,\quad
 H_V(P,G)=\gamma_*\kappa(1+P).
\]
Write $\Delta=U-P$, $\Gamma=V-G$.
Finite monic division, first in $U-P$ and then in $V-G$,
therefore gives the complete scalar identity
\[
\begin{split}
 \Phi={}&\gamma_*\{\mathcal M+P\kappa U\Gamma^2\}\\
 &+UV\{\Delta^2H_{20}+\Delta\Gamma H_{11}
                                      +\Gamma^2H_{02}\},\\
 \mathcal M={}&U^2V+\kappa UV^2+P^2G-3PUV .
\end{split}                                                     \tag{EZ13}
\]
All quotients are finite bounded coefficient polynomials; no
inverse of $P,G$ is introduced by their monic construction.
The extra term $\gamma_*P\kappa U\Gamma^2$ is retained as a
nonnegative square.

In the inverse-$b$ row one further coefficient factor is needed:
\[
                         H_{02}=\kappa\widetilde H_{02},
 \qquad\widetilde H_{02}\text{ has bounded arrays}.              \tag{EZ14}
\]
Here is a finite source-division certificate, rather than an
estimate involving independent testing rates. Write
$P=\kappa G$, $t=\kappa\tau$ before substituting the native
coquotients into EZ10. Every denominator is among the pure-source
units $M,T,c_d$ and their factors. At formal $\kappa=0$,
\[
 P=t=\mu=B_*=0,\quad
 c_*=d_*=\frac{1+(1-a)s^2}{2},\quad
 a_*=\frac1{2(1-as^2)} .
\]
The formula gives $H(0,V)=0$ identically in $V$.
After clearing unit denominators, ordinary polynomial division
by $\kappa$ therefore has zero remainder. Apply the two monic
$V-G$ divisions to this quotient. This proves EZ14 with bounded
arrays. In the bounded-$b$ row $\kappa=1$, so it needs no such
factor.

\subsubsection*{Six explicit strict roots and their conductors}

The following identities apply to both rows of EZ3. Put
\[
 \mathsf h=G^{1/6},\quad\mathsf c=\kappa^{1/6},\quad
 \mathsf r=\sqrt[3]u,\quad\mathsf s=\sqrt[3]z,
\]
\[
 (L_1,L_2,L_3)=
 (\mathsf r^4\mathsf s^2,\
   \mathsf c^2\mathsf r^2\mathsf s^4,\
   \mathsf h^6\mathsf c^4),\qquad
 (f_1,f_2,f_3)=
 (\mathsf r^2\mathsf s,\
   \mathsf c\mathsf r\mathsf s^2,\
   \mathsf h^3\mathsf c^2).
\]
Define $R_{2i-1}=f_i(L_1-L_2)$ and
$R_{2i}=f_i(L_1-L_3)$ for $i=1,2,3$.
Then, exactly,
\[
 \mathcal M=\sum_{i=1}^3
       (R_{2i-1}^2-R_{2i-1}R_{2i}+R_{2i}^2).                  \tag{EZ15}
\]
Each pair has the fixed strictly positive Gram
$\left(\begin{smallmatrix}1&-1/2\\-1/2&1\end{smallmatrix}\right)$.
Real cube roots of signed testing quantities are permitted:
$z=sW$ has the separated cube root
$s^{1/3}\sqrt[3]W$. The $L_i=f_i^2$ identity is valid for
either sign of $u,z$. No mixed radical occurs.

For explicit finite allocation define
\[
 J_1=\mathsf r^3R_1+\mathsf c^2\mathsf r\mathsf s^2R_2,\qquad
 J_2=-\mathsf h^6R_1+\mathsf s^6R_2
                              -\mathsf h^3\mathsf s^3R_5.
\]
Direct multiplication gives
\[
\begin{aligned}
 uz\Delta&=J_1,&u^2z\Gamma&=J_2,\\
 \kappa uz\Gamma
   &=\mathsf c^2\mathsf r\mathsf s^2(R_2-R_1)
                                  -\mathsf c^3\mathsf s^3R_3,\\
 UV\Delta^2&=J_1^2,\\
 UV\Delta\Gamma
   &=\mathsf s^6R_2^2-\mathsf h^6R_1^2
             -\mathsf h^3\mathsf s^3R_1(R_5+R_6),\\
 \kappa UV\Gamma^2
   &=\mathsf s^6(2R_1^2-3R_1R_2+R_2^2+R_3^2)
                      +\mathsf h^3\mathsf s^3R_1(R_5-R_6).
\end{aligned}                                                   \tag{EZ16}
\]
All coefficients on the right are infinitesimal: $\mathsf r,
\mathsf s,\mathsf h$ are infinitesimal and $\mathsf c$ is
bounded. It is harmless that $\mathsf c=1$ in the first row.
Equations EZ13--EZ16 therefore put the entire scalar, apart
from its retained positive square, on the six rows $\omega R_i$
with a bounded Gram whose real residue is $\tfrac12$ times the
strict matrix in EZ15.

\subsubsection*{Every complete first normal and all higher terms}

Define $L_x,L_y,L_v$ as half the three derivatives of the
whole polynomial EZ10 at zero normals. Expanding its short
factored form gives exactly
\[
 L_x=\omega uz^2T_x(U,V),\qquad
 L_y=\omega uzT_y(U,V),\qquad
 L_v=\omega uzT_v(U,V),                                       \tag{EZ17}
\]
where each $T_i$ is a polynomial with bounded coefficient arrays.
This supplies a finite construction of every coefficient:
differentiate EZ10, set the three normals to zero, and extract
the displayed monomial. No derivative is dropped.
The exact actual-source contact equations give $T_i(P,G)=0$.
Consequently the monic divisions are
\[
 T_i=(U-P)A_i(U,V)+(V-G)B_i(V).                               \tag{EZ18}
\]
For inverse $b$, each $B_i=\kappa\widetilde B_i$ with
bounded arrays. Indeed the same source substitution used in
EZ14 gives $T_i(0,V)=0$ identically at $\kappa=0$.
Divide the cleared source polynomial by $\kappa$ and then by
the monic $V-G$. This is an exact coefficient construction;
it uses neither $1/u$, $1/z$, nor an evaluated mixed quotient.
For bounded $b$, $\kappa=1$ again suffices.

Combining EZ16--EZ18 proves that both complete state normals
and the complete testing-angle normal have infinitesimal
coefficients on the strict six-row list $\omega R_i$.
The additional factor $z$ in $L_x$ is bounded and infinitesimal.
This argument uses the full first-normal polynomial, not just
its value or linear approximation at contact.

The full normal expansion of EZ10 has a bounded quadratic Gram
whose real residue, in order $(N_x,N_y,N_v)$, is
\[
 \begin{pmatrix}
  \bar\mu+\bar\omega/2&0&0\\
 0&1&1/2\\
 0&1/2&1/2
 \end{pmatrix}\succ0.                                         \tag{EZ19}
\]
For bounded $b$, $\bar\mu=1/2$; for inverse $b$,
$\omega=1$ and $\bar\mu=0$. Both give a strictly positive
first pivot. The other block has pivots $1,1/4$.

Here is also a finite justification for all higher terms, since
the scaled normals $N_y,N_v$ may be unbounded. Expanding EZ10
in the three normals gives degree at most four. A monomial
of normal degree at least three containing $N_x$ has total
$N_y,N_v$ degree at most two. Assign those factors, and any
needed extra $N_x$ factor, to two Gram rows; the remaining
factor contains $N_x\in\mathrm I$. For a higher monomial
without $N_x$, each of the $N_y,N_v$ exponents is at most two,
and its coefficient has a factor $s^2$. If its degree is
three, the unassigned coefficient is a bounded multiple of
$s(sN_y)$ or $s(sN_v)$. If its degree is four, it is a
bounded multiple of $(sN_y)(sN_v)$ (or of one of their
squares). All are infinitesimal by EZ9. These support and
divisibility statements follow directly from the factored
quadratics in EZ10 and are checked on the entire polynomial.
Thus no unbounded scaled normal is hidden in a Gram coefficient.

The scalar, all first normals, and the complete hard block now
give a finite bounded Gram with strictly positive real residue.
The separated strict-Gram lemma supplies tensor SOS.
Restore the whole retained terms EZ7--EZ8 and the positive
square in EZ13, and divide by the pure-source positive $\omega$.
This proves the stated original-parent theorem on both source
rows, without a comparison of source and testing rates.

\paragraph{State and boundary scope.}
The signs of $W$ and of the testing angles were not restricted
in EZ10--EZ19. The polynomial and cube-root identities include
literal zero testing coordinates. If a bounded normalized
$Z=1$ state has $X$ or $Y$ a unit, the original pure-testing
core has positive residue at the double-small angle corner;
its positive source multiple is already a separated square,
and all original event terms are nonnegative squares.
Thus this proof supplies the entire $Z$-unit state chart.
Complex states split into their real and imaginary quadratic
evaluations. Section~\ref{component:p845b3cb8b114} combines this
certificate with the two strong-current outer testing identities
and the complementary $W$-unit chart to cover all tests.



% Source formulas: pabbcd0905dcd
\subsection{The uniform small-$e$ source at unbounded $b$}

\label{component:pabbcd0905dcd}

\label{pabbcd0905dcd:sec:c1-inverse-b-local-new}

\paragraph{Statement and certificate ring.}
Let $K_c,K_t$ be independent real closed extensions of $\mathbb R$,
with bounded rings $B_c,B_t$ and infinitesimal ideals
$\mathfrak m_c,\mathfrak m_t$. The certificate ring is
$\mathcal T=K_c\otimes_{\mathbb R}K_t$, with bounded coefficient ring
$\mathcal B=B_c\otimes_{\mathbb R}B_t$ and zero-residue ideal
$\mathfrak n=\mathfrak m_c\otimes B_t+B_c\otimes\mathfrak m_t$.
We use the finite strict-residue Gram lemma in this ring.
All radicals and inverses below are pure source, pure testing, or
products of separately specified such factors.

Use the actual uniform native C1 stationary source, with
$\lambda=n=0$, $d>0$, and the original nonnegative event coefficients.
Normalize $d=1$ by a positive pure-source scale. Suppose
\[
 k,f,g>0\text{ are infinitesimal},\qquad k/g\in\mathfrak m_c,
 \qquad r=1-af^2,\quad h=1-tg^2,
\]
\[
 R=1+k^2r,\quad T_0=1+k^2h,\quad
 0<a<T_0,\quad 0<t<R,\quad \beta\geq0.                 \tag{IW1}
\]
All original strict gaps remain imposed. Then the complete original
parent is tensor-SOS whenever the independent positive testing angles
$u,v$ are infinitesimal, for every real or complex testing state.
The proof below works in the rotated chart $w=1,z>0$ with $z$ arbitrary;
in particular it includes the previously missing original
$W$-unit, $Z$-infinitesimal chart. It uses no bound on $g/f$ and no
comparison between any testing and source infinitesimals.
The exact sign, zero-axis and state reductions are given at the end.

\subsubsection*{The actual C1 source and bounded source arrays}

For clarity, the six weights used here come from the following
unchanged native formulas. Put $x_n=k^2$, $\eta=f^2$, $\zeta=g^2$,
$y_n=x_n\eta\zeta$, $K=1+x_n$, $G=h+\zeta$, and
\[
 M=R-t,\quad J=T_0-a,\quad E_n=1-x_n\eta h,
 \quad Q_n=1-x_n\zeta r,\quad \theta=t\zeta.
\]
Define
\[
 P_d=1+2x_nh+rx_n^2G+x_n\zeta
       +\eta x_n(K\zeta+x_nh\theta),
\]
\[
 P_4=T_0E_n+x_n(1-r)G,\qquad
 P_\alpha=Rh(1+y_n)+\zeta M E_n.
\]
The actual relative coefficients, now with $d=1$, are
\[
\begin{array}{ll}
 m=\dfrac{Kx_n(r+\eta)h(1+y_n)}{2M\zeta P_d},&
 C=\dfrac{K(r+\eta)E_n}{2(1-r)P_d},\\[4pt]
 A=\dfrac{\eta P_4}{2rP_d},&
 D+\tfrac12=\dfrac{K(1+y_n)(r+\eta)}{2J\eta P_d},\\[4pt]
 \beta+\tfrac12=\dfrac{Kx_n(r+\eta)G}{2\theta P_d},&
 \alpha=\dfrac{x_n^2\zeta(r+\eta)P_\alpha}{2Q_nP_d}.
\end{array}                                                     \tag{IW2}
\]
No auxiliary C2 stationary source is used.
Since $a,t$ are bounded and $k,f,g$ are infinitesimal,
$r,h,P_d,R,T_0$ have residue one. The inequality
$\beta+\tfrac12\geq\tfrac12$ in (IW2) gives
$t=O(k^2/g^2)$. Thus $t$ and $p=t/R$ are infinitesimal, and
$M=R(1-p)$ is a positive unit. The two independent openings
$a$ and $J=T_0-a$ are not capped.

Introduce the exact compact coordinates
\[
 j=\frac{T_0}{g^2R},\quad B=\frac{k^2g^2R}{T_0},\quad
 T=\frac{f^2T_0}{R},\quad p=\frac tR,\quad H=\frac a{T_0},
 \qquad x=jB,\quad y=TB,\quad q_z=jT.                    \tag{IW3}
\]
Here $j$ is unbounded, $0<H<1$, and
\[
 T,x,B,p,\sigma:=1/j,\delta:=j^2B\in\mathfrak m_c,
 \qquad \delta\sim k^2/g^2,\qquad p=O(\delta).
                                                               \tag{IW4}
\]
The quantity $q_z$ can have any positive size. The exact inverse is
\[
\begin{gathered}
 R=\frac{1+x}{1+HTx},\quad T_0=\frac{1+x}{1+Bp},\quad
 r=\frac{1-HT}{1+HTx},\quad h=\frac{j-p}{j(1+Bp)},\\
 f^2=\frac{T(1+Bp)}{1+HTx},\quad
 g^2=\frac{1+HTx}{j(1+Bp)},\quad a=HT_0,\quad t=pR.
\end{gathered}                                                   \tag{IW5}
\]
In particular both strict gaps are retained exactly.

Here is a short rational dictionary sufficient to check every source
estimate. Write $D_+=D+\tfrac12$, $b_+=\beta+\tfrac12$,
$\mu=m(1-p)$, and set
\[
\begin{split}
 P={}&1+B(j+1-p)+BT(1+Hj)\\
 &+B^2T(Hj+jp-p^2+p)+B^2HT^2j(1+Bp),\\
 L={}&1+T(1-H+Bp),\qquad E=1+B(HT+p-1),\\
 F_\alpha={}&j+1-2p+BHTj(1-p)+BTp(j-p).
\end{split}
\]
Substitution of (IW5) in (IW2) gives exactly
\[
\begin{split}
 \mu&=\frac{Bj(j-p)(1+y)(1+Bp)L}{2(1+x)P},\\
 C&=\frac{(1+Bp)L[1+BTp-BTj(1-H)]}{2HT(1+x)P},\\
 A&=\frac{T(1+Bp)[1+BT(2Hj-Hp+H-j+p)+B^2H^2T^2j]}
              {2(1-HT)P},\\
 D_+&=\frac{(1+y)(1+Bp)(1+HTx)L}{2T(1-H)(1+x)P},\\
 b_+&=\frac{Bj(1+Bp)L(j+1-p+HTx)}{2p(1+x)P},\\
 \alpha&=\frac{B^2LF_\alpha}{2EP}.
\end{split}                                                     \tag{IW6}
\]
The denominators $P,E,1+x,1-HT$ are positive bounded units.
After $j=1/\sigma$, $B=\delta\sigma^2$, each of the following
normalized arrays is a rational function with polynomial numerator
and denominator in the bounded variables $\sigma,\delta,T,p,H$;
its denominator has a nonzero real constant at
$\sigma=\delta=T=p=0$, independently of $H$:
\[
 \frac{2\mu}{\delta},\quad\frac{2b_+p}{\delta},\quad
 2TCH,\quad2TD_+(1-H),\quad\frac{2A}{T},\quad
 \frac{2\alpha}{B^2j}\quad\hbox{all have residue }1.
                                                               \tag{IW7}
\]
This finite polynomial criterion proves boundedness and the stated
residues at every relative rate and at both $H$ openings. It also
shows $\beta p\leq b_+p=O(\delta)$, including $\beta=0$.
A useful exact cancellation is
\[
                  K_d:=CH-D_+(1-H)=-\frac{b_+p}{j}.    \tag{IW8}
\]

\subsubsection*{The entire parent and both state normals}

Put $D_u=1+u^2$, $D_v=1+v^2$. The quarter-parent is
\[
\begin{split}
 \mathcal P={}&m(vW-X+uvY-uZ)^2+CD_v(uW-Y)^2+AD_uD_vY^2\\
 &+D_+(uW+uvX-Y-vZ)^2
 +(\alpha+\tfrac12)D_u(vY-Z)^2\\
 &+b_+D_v(X+uZ)^2-\tfrac12D_uD_vZ^2.
\end{split}                                                     \tag{IW9}
\]
On the rotated chart
$w=(W+uY)/D_u=1$, $z=(Z-uX)/D_u>0$, use the two unrestricted normals
\[
 N=\frac{X+uZ-pD_uv}{D_u},\qquad
 M_0=\frac{uW-Y-HD_uvz}{D_u}.
\]
Their exact inverse is
\[
 X=N+pv-uz,\quad Z=z+uN+puv,\quad
 W=1+u(M_0+Hvz),\quad Y=u-M_0-Hvz.                     \tag{IW10}
\]
Only the pure-testing unit $D_u$ is inverted. Retain the positive
term $mpD_u^2N^2$, and put $G=(\mathcal P-mpD_u^2N^2)/D_u$.
Direct substitution, before any state specialization, gives
\[
\begin{split}
 G={}&\mu D_u[N^2-2vN+(1-p)v^2]
      +CD_uD_v(M_0+Hvz)^2+A D_v(u-M_0-Hvz)^2\\
 &+D_+D_u(M_0-(1-H)vz)^2\\
 &+(\alpha+\tfrac12)
       [u(1-p)v-(1+Hv^2)z-uN-vM_0]^2\\
 &+b_+D_uD_v(N+pv)^2-\tfrac12D_v(z+uN+puv)^2.
\end{split}                                                     \tag{IW11}
\]
In particular $G$ is exactly quadratic in the two normals:
\[
 G=F+2L_NN+2L_MM_0+H_NN^2+H_MM_0^2
                         +2(\alpha+\tfrac12)uvNM_0,
                                                               \tag{IW12}
\]
where $F=G|_{N=M_0=0}$ and
\[
\begin{split}
 h_0&=\mu+b_+\geq\tfrac12,\\
 H_N&=h_0D_u+b_+v^2+\beta u^2v^2+\alpha u^2,\\
 H_M&=D_u(CD_v+D_+)+AD_v+(\alpha+\tfrac12)v^2.
\end{split}                                                     \tag{IW13}
\]
The potentially unbounded $\beta$ is fully retained in (IW11)--(IW13).
There is no omitted higher-normal polynomial.

\subsubsection*{An exact separated scalar circuit}

Define the rows
\[
 U=u/\sqrt x,\quad V=v/\sqrt y,\quad Z_*=z/\sqrt{q_z},
 \quad Q_*=U-VZ_*,\quad B_*=V-UZ_*,\quad R_*=Z_*-UV,
\]
\[
 \Phi=U^2+V^2+Z_*^2+U^2V^2Z_*^2-4UVZ_*.
\]
None of $U,V,Z_*$ is assumed bounded. On the positive chart,
\[
 \Phi=(U-V)^2+Z_*^2(1-UV)^2+2UV(1-Z_*)^2             \tag{IW14}
\]
is a legal separated SOS; for example
$\sqrt{UV}=(xy)^{-1/4}\sqrt{uv}$.

For an explicit finite definition of the actual masses, put
$K_0=CH^2+D_+(1-H)^2$ and
\[
\begin{array}{ll}
 M_a=Ax,&M_b=[\mu(1-p)+b_+p^2]y,\\
 M_c=\alpha q_z,&
 M_{bc}=[K_0+AH^2+2H\alpha+H-\tfrac12]yq_z,\\
 M_{b^2}=b_+p^2y^2,&
 M_{b^2c}=H^2(A+C+\alpha+\tfrac12)y^2q_z,\\
 M_{ab^2c}=CH^2xy^2q_z,&M_{ab^2}=\beta p^2xy^2.
\end{array}
\]
Set
\[
\begin{gathered}
 f_c=M_b,\quad h_c=M_{ab^2},\quad
 k_c=(M_a-M_b+M_{ab^2c}-M_{bc}+h_c)/2,\\
 r_c=M_c-M_b-k_c,\quad q_c=M_a-f_c,\quad
 g_c=M_{ab^2c},\quad d_c=M_{b^2c}-k_c,
 \quad e_c=b_+p^2y^2-g_c.
\end{gathered}
\]
The complete scalar from (IW11) satisfies the exact identity
\[
\begin{split}
 F={}&q_cQ_*^2+r_cR_*^2+d_cV^2Q_*^2+e_cV^2B_*^2\\
 &+f_c\Phi+g_cV^2\Phi+k_cZ_*^2(V^2-1)^2+h_cV^2R_*^2.
\end{split}                                                     \tag{IW15}
\]
The last term is the extra nonnegative square supplied by the actual
$\beta$ event. 

The same finite bounded-array criterion as (IW7), with the displayed
mass definitions, gives
\[
 \frac{2f_c}{Tx^2}\sim1,\quad
 \frac{2q_c}{Tx}\sim1,\quad
 \frac{2r_c}{TjB^2}\sim1,\quad
 \frac{k_c}{T^2B^2j}\in\mathfrak m_c,\quad
 \frac{g_c}{f_c}\in\mathfrak m_c.                     \tag{IW16}
\]
Here $\sim1$ means residue one. In particular the three principal
masses $f_c,q_c,r_c$ are positive, and
$|k_c|/r_c$, $|k_c|/(q_cy)$ are infinitesimal.
All normalized mass denominators remain units after
$j=1/\sigma$, $B=\delta\sigma^2$; the companion finite-array
checks test their residues independently of $H$.

For the signed terms use the exact identities
\[
\begin{split}
 (1+U^2)\Phi&=R_*^2+B_*^2+U^2Q_*^2+U^2(1-UVZ_*)^2,\\
 (1+V^2)\Phi&=Q_*^2+R_*^2+V^2B_*^2+V^2(1-UVZ_*)^2,\\
 Z_*(V^2-1)&=-VQ_*-R_*.
\end{split}                                                     \tag{IW17}
\]
The $e_c,g_c$ combination in (IW15) is exactly
\[
 b_+p^2y^2V^2B_*^2-g_c\Phi
       +g_c[Q_*^2+R_*^2+V^2(1-UVZ_*)^2].
\]
A $k_c$ of either sign costs at most
$3|k_c|V^2Q_*^2+2|k_c|R_*^2$, including its occurrence in $d_c$.
The corresponding relative costs are
$3|k_c|v^2/(q_cy)$ and $2|k_c|/r_c$, both in $\mathfrak n$.
Every inequality here abbreviates a finite SOS identity, using
$2A^2+2B^2-(A+B)^2=(A-B)^2$ and nonnegative pure-source factors.
Consequently the scalar supplies strict fixed real fractions of
\[
                       q_cQ_*^2,\qquad r_cR_*^2,
                       \qquad f_c\Phi.                \tag{IW18}
\]
All other nonnegative terms remain available as whole squares.

\subsubsection*{Finite budgets for the complete normal rows}

Call a row free if any fixed positive real multiple of its square can
be subtracted at infinitesimal cost from (IW18), by a finite separated
SOS identity. This definition allows a bounded normal coupling to be
absorbed by a sufficiently large fixed real helper pivot.
Put
\[
 q=vz-Tu,\quad b=v-uz/j,\quad R_z=z-uv/B,\quad
 R_u=v(u^2-x),\quad R_v=u(v^2-y),
\]
\[
 J_u=u(u^2-x),\qquad J_v=v(v^2-y),\qquad
 \widehat q=q/\sqrt T.
\]
Since $Q_*=-q/(T\sqrt x)$, the first mass in (IW18) has a positive
unit coefficient on $\widehat q^2$.

First, $R_u=\sqrt y\,x(-B_*-UR_*)$. Applying the first identity
of (IW17) to $B_*^2$, and the actual $r_c$ pivot to $U^2R_*^2$,
shows that $R_u$ is free. The three relative costs are bounded by
constant multiples of
\[
                         B,\qquad u^2/j,\qquad u^2,
                                                               \tag{IW19}
\]
all infinitesimal. Second,
$R_v=\sqrt x\,y(-Q_*-VR_*)$, so $\sqrt T R_v$ is free;
its two costs are $O(y^2)$ and $O(Tv^2)$.

The same-angle rows have exact decompositions
\[
\begin{split}
 J_v&=\sqrt y\,v(v+\sqrt{y/x}\,u)(V-U)+(y/x)R_u,\\
 J_u&=\sqrt x\,u(u+\sqrt{x/y}\,v)(U-V)+(x/y)R_v.
\end{split}                                                     \tag{IW20}
\]
If $\rho=O(\delta)$ is a nonnegative source factor, then
$\rho J_v$ is free. Indeed $\rho\sqrt{y/f_c}=O(\sqrt\delta)$,
$\sqrt{y/x}=\sqrt{T/j}$ is bounded, and
$\rho y/x$ is infinitesimal. Thus every coefficient after division
by the strict $\sqrt{f_c}(V-U)$ pivot is infinitesimal.
Also $\sqrt T\,TxJ_u$ is free: the two coefficients of its
first summand relative to $\sqrt{f_c}(U-V)$ are bounded multiples of
\[
                  T\sqrt x\,u^2,\qquad\sqrt{T\delta}\,uv,
\]
and its second summand is exactly
$\delta(\sqrt T R_v)$, because $Tx^2/y=j^2B=\delta$.
This last use of $\delta\in\mathfrak m_c$ is the additional
source fact specific to the inverse-$b$ sector.

Finally,
\[
 j\alpha b=-(\alpha/B)R_u-\alpha uR_z.                \tag{IW21}
\]
Here $\alpha/B=O(x)$, and the square of $\alpha uR_z$ costs
$O(x^2u^2)$ relative to $r_cR_*^2$. Hence the entire left side
of (IW21) is free. These identities involve no bound on $z$.

\paragraph{The unrestricted first normal.}
Define the bounded source quantities
\[
 K_z=H(\alpha+\tfrac12)-\tfrac12,\qquad
 K_R=-\mu-\alpha(1-p)-\tfrac12+b_+p+TK_z,
\]
\[
 \rho=b_+p+\beta px\geq0.
\]
Direct differentiation of the entire (IW11), followed by the actual
source contact identity, gives
\[
 L_N=-j\alpha b+(K_R+\beta pv^2)R_u
                   +\rho J_v+K_zuvq.                 \tag{IW22}
\]
For completeness, before imposing the source, the remainder is
$v S_N$, with
\[
 S_N=j\alpha-\mu+b_+p+K_Rx+b_+py+\beta pxy=0
\]
identically under (IW6). Thus no constant or first-normal term is
silently dropped. Equations (IW7) imply
$K_R\sim-\tfrac12$, $\rho=O(\delta)$ and
$\beta p=O(\delta)$. The first three terms in (IW22) are free by
(IW19)--(IW21); the last has coefficient
$K_z\sqrt T\,uv\in\mathfrak n$ on $\widehat q$.
Dividing by $\sqrt{h_0}$ preserves all these bounds because
$h_0\geq\tfrac12$.

\paragraph{The unrestricted second normal.}
Use the exact $K_d$ from (IW8), and put
\[
\begin{split}
 C_q={}&K_dD_u+AH+\alpha+\tfrac12
        +v^2[CHD_u+AH+(\alpha+\tfrac12)H],\\
 c_2={}&T[CH+AH+(\alpha+\tfrac12)H]
                      -A-(\alpha+\tfrac12)(1-p).
\end{split}
\]
The complete second normal has the exact division
\[
 L_M=C_qq+(c_2+TCHu^2)R_v+T(K_d+CHy)J_u.             \tag{IW23}
\]
Before substituting the actual source its only remainder is $uS_M$,
where
\[
\begin{split}
 c_0&=T[K_d+AH+\alpha+\tfrac12]-A,\\
 S_M&=c_0+TK_dx+yc_2+TCHxy=0.
\end{split}
\]
The last equality is an exact identity from (IW6). Furthermore
$c_2$ is bounded (in fact infinitesimal), $TCH$ is a unit, and
$(K_d+CHy)/x$ is bounded with residue $-\tfrac12$.
It follows from the two free rows in (IW20) that
\[
 \sqrt{TH(1-H)}[(c_2+TCHu^2)R_v+T(K_d+CHy)J_u]
\]
is free with bounded coefficients. For the other term,
\[
 \sqrt{TH(1-H)}C_qq
     =T\sqrt{H(1-H)}C_q\widehat q.
\]
Its coefficient is in $\mathfrak n$: multiplying $C_q$ by $T$
clears the only $CH$ pole, and every term then contains either
$T$ or $v^2$. This statement is made on the full expression in
(IW23), before any relative source/testing rate is considered.

\subsubsection*{Strict full Gram and state restoration}

Normalize the hard rows by pure-source factors,
\[
 \widehat N=\sqrt{h_0}\,N,\qquad
 M_0=\sqrt{TH(1-H)}\,\widehat M.
\]
Their full hard matrix is bounded: $b_+/h_0\leq1$,
$\beta/h_0\leq1$, $\alpha/h_0=O(\alpha)$, and the two
source poles of $H_M$ are exactly cleared by $TH(1-H)$.
Equations (IW7), (IW13) give its real residue
\[
                           \operatorname{diag}(1,\tfrac12).
                                                               \tag{IW24}
\]
Its off-diagonal coefficient has the factor
$uv\sqrt{TH(1-H)}/\sqrt{h_0}$ and is infinitesimal.
The $\widehat q$ couplings in (IW22)--(IW23) are infinitesimal.
All remaining complete-normal couplings are bounded coefficients
of free rows. Buy sufficiently large fixed real pivots on the
finitely many free rows, use the elementary real Young square
completion, and apply the strict-residue Gram lemma to the
remaining finite matrix. The strict scalar pivots (IW18) lose only
infinitesimal amounts. This proves $G\in\Sigma\mathcal T^2$.
No state normal has been set to zero in this conclusion.
Restoring $mpD_u^2N^2$, the pure-testing factor $D_u$, and the
positive source scale $d$ proves the assertion for (IW9).

The proof has used $u,v>0$ and $w=1,z>0$, with arbitrary $z$.
The exact sign/state lemma SJ1--SJ10 applies to the same original
core: sign gauges make $u,v$ nonnegative while preserving every
term; a negative core at positive angles forces the rotated
$w,z$ to have the same nonzero sign. Dividing all state coordinates
by the nonzero pure-testing $w$ gives precisely this chart.
When the core is nonnegative it is itself a pure-testing square,
and every event term is already positive. Zero angles have the
direct core SOS exits. The zero state is immediate, and a complex
state splits into its real and imaginary quadratic evaluations.
Thus (IW1) is proved for all states at the double-infinitesimal
angles. The existing two-current outer-angle identities give the
separate all-tests join if needed; they are not required as a premise
of the local conductor proved here.



% Source formulas: p845b3cb8b114
\subsection{The uniform small-$e$ testing and state partition}

\label{component:p845b3cb8b114}

\label{p845b3cb8b114:sec:uniform-c1-complete-current-join}

\paragraph{Theorem and unchanged source scope.}
Use the original physical canonical C1 stationary source with
$\lambda=n=0$, $d>0$, every original nonnegative event coefficient,
and every original strict source gap. Suppose
\[
 k,f,g,e=kf/g>0\text{ are infinitesimal},\qquad
 r=1-af^2,\quad R=1+k^2r,\quad T_0=1+k^2-tk^2g^2,
\]
\[
                         a,t>0,\qquad J=T_0-a>0,\quad M=R-t>0.
                                                               \tag{UJ1}
\]
Then the entire original parent is a finite tensor sum of squares
for every independent projective testing pair and every real or
complex testing state. There is no restriction on $b=g/k$, $g/f$,
or any relative infinitesimal rate. In particular the original
$\beta=0$ face is retained.

This is a mathematical join of the actual bounded-$b$ endpoint
product-repair proof, the whole $Z$-chart proof EZ1--EZ19,
the inverse-$b$ local proof IW1--IW24, and the two full-parent
outer-current identities below. No unavailable EZ or IW
coefficient witness is used as a premise.

\subsubsection*{Exact original-source bounds for the two outer identities}

Normalize $d=1$ by a positive pure-source scale. The retained
Section~\ref{component:p7d4e5164e84a} prints seven original coquotients.
Their respective residues on (UJ1) are
\[
 (c_2,c_\beta,c_3,c_5,c_4,c_\alpha,c_d)
                    \sim(2,2,2,2,2,2,4).
\]
Their arguments are bounded, so they are positive source units.
The actual weights satisfy
\[
 C=\frac{c_3}{af^2c_d},\qquad
 D=\frac{B_c}{2Jf^2c_d},\qquad A=\frac{f^2c_4}{c_d},
 \qquad B_c=2c_5-Jf^2c_d\sim4.
\]
Consequently
\[
 C^{-1},D^{-1}\in\mathfrak m_c,\qquad
 g_B=\frac{CD}{C+D}>0,\qquad
 \gamma=A g_B=
 \frac{c_4c_3B_c}{c_d^2(aB_c+2Jc_3)}\sim\frac14.
                                                               \tag{UJ2}
\]
This is uniform in both openings $a,J$: the last denominator has
residue $16\cdot4(\bar a+\bar J)=64$, since $a+J=T_0\sim1$.
No positive lower bound on either opening is assumed.
These bounds use the actual original $C,D,A$, not selected
substitute weights. Any other positive event is restored whole.

\paragraph{Complementary testing angles.}
The actual source estimates UJ2 verify the hypotheses of
Lemmas~\ref{lem:c1-two-current} and~\ref{lem:c1-critical-current}.
For $v$ infinitesimal and $1/u$ bounded, $\bar\gamma=1/4$
and $\bar c_u^2<1$, so the latter real residue condition is strict.
Their whole homogeneous identities include both projective endpoints.
Only the two infinitesimal testing angles remain.
\subsubsection*{Every remaining source and state chart}

At $u,v$ infinitesimal normalize a nonzero testing state by a
coordinate of largest absolute value. This is a pure-testing
operation; all four coordinates are bounded and at least one
is a unit. If $X$ or $Y$ is a unit, the original pure-testing
core has positive real residue $2(X^2+Y^2)$. It is a positive
testing square, and every original event remains a positive
separated square. We can therefore assume $X,Y$ infinitesimal.

If $Z$ is a unit, normalize $Z=1$. The remaining $W$ is arbitrary
bounded and may have either sign. EZ1--EZ19 applies in each of
its two actual-source coordinate rows:
\[
\begin{array}{c|c|c}
 \text{source}&\text{exact coordinates}&\text{actual proof}\\ \hline
 b=g/k\text{ bounded}& f=eb,\quad g=kb&
                     \mathrm{EZ1\!-\!EZ19,\ bounded\ }b\\
 b=g/k\text{ unbounded}&
        (k,f,g),\quad k/g\in\mathfrak m_c&
                     \mathrm{EZ1\!-\!EZ19,\ inverse\ }b.
\end{array}                                                     \tag{UJ5}
\]
The extra inverse-$b$ source inequality used by EZ is derived from
the actual $\beta\geq0$: $tb^2\leq2c_\beta/c_d$. Thus it adds no
source hypothesis to (UJ1). Its full proof retains both state
normals and the testing-angle normal, including unbounded scaled
normal coordinates through their exact $sN_y,sN_v$ coefficients.

If $Z$ is infinitesimal, $W$ must be a unit. Normalize $W=1$;
then $u,v,X,Y,Z$ are all infinitesimal.
For bounded $b$, use the actual retained
Section~\ref{component:p05180564ed05} proof. Its source variables are exactly
$(e,k,b,a,t)$ with $f=eb$, $g=kb$; its gaps
$R_c=1+k^2(1-ae^2b^2)$ and
$T_0=1+k^2-tk^4b^2$ coincide with (UJ1).
It allows every bounded positive $b$, including infinitesimal $b$,
both strict-gap openings, and all original coefficient faces.
Its whole seven-term polynomial, complete first normals,
higher-normal support and nine-root product identity are printed
with Section~\ref{component:p7d4e5164e84a} and
Section~\ref{component:p935c9d6318ea}. Its final original-parent
restoration is an exact sum of complete positive terms.
There is no remaining condition on $k^3/e$ or $k/e$ in that proof.

For unbounded $b$, use IW1--IW24. Its exact source hypotheses
$k,f,g,k/g\in\mathfrak m_c$, physical $\beta\geq0$, and the
original gaps follow directly from (UJ1) and this ratio class.
It proves the complete original parent with both unrestricted
state normals for every positive rotated state ratio, so it
includes the present $W$-unit chart. The retained $\beta$ term,
including its zero face, and both source openings are part of
that proof. The state and sign reduction is the actual
SJ1--SJ10 identity, not a chart label.

These cases exhaust every source ratio and every state.
The zero state is immediate. A complex state splits into its
real and imaginary evaluations because the parent matrix is real.
Undo all pure-testing normalizations and the positive source
scale $d$. This proves the theorem on (UJ1) at all independent
projective tests.



% Source formulas: pf5d2f33f694c
\subsection{Every positive $e$ in the uniform source}

\label{component:pf5d2f33f694c}

\label{pf5d2f33f694c:sec:uniform-c1-all-e-rederived}

Use the original canonical C1 source and the exact native transport
EN13, with every original coefficient sign and strict gap retained.
Suppose
\[
 k,f,g>0\text{ are infinitesimal},\qquad e=kf/g>0 .
 \tag{UE1}
\]
There is no additional restriction on the size of \(e\), either
source ratio \(g/k,g/f\), or either separately labelled testing field.
Then the entire original parent is tensor SOS at every independent
projective testing pair and every real or complex testing state.

Here are the actual proof inputs and the complete join.  EN13 gives
\(0<a,t<2\), \(r=1-af^2\sim1\), and the original strict
\(M=R-t>0,J=T-a>0\).  Split the positive pure-source \(e\) into
\[
 e\text{ infinitesimal}\quad\hbox{or}\quad \Delta=1/e\text{ bounded}.
 \tag{UE2}
\]
In the first case UJ1--UJ5 is the full-parent theorem.  It
includes both source ratios \(g/k\) bounded and unbounded, both
state charts, both complete state normals and all angular/projective
joins; the actual bounded-ratio product-repair, EZ and IW proof bodies are its inputs.

In the second case use the actual retained
Section~\ref{component:p41b39e2f94e1}, Theorem \emph{Uniform C1, bounded
reciprocal \(e\)}.  Its hypothesis is exactly UE1 with
\(\Delta\) bounded, and its proof is complete in that body.
For clarity, the scope extension used there is substantive.
After its pure-source substitution
\[
 \delta=\Delta R/(1+k^2),\quad b^2=f^2(1+k^2)/R,
\]
both \(\delta\) and the pure-testing rotated state coordinate \(w\)
are allowed to be bounded units.  In its polynomial ring the weights
\[
 (\delta,k,b,a,t,u,w,v)=(0,2,6,0,0,2,0,16)
 \tag{UE3}
\]
make the full scalar error and both complete first normals have weights
at least \(42,32,26\).  The displayed unlocalized monic conductor,
restricted to \((A_m,B_m,C_m)^2+(D_0^2)\), gives a scalar Gram
whose every correction coefficient has weight at least \(2\);
the first-normal coefficient bounds are \(12,6\).
The exact hard block on its two normalized state normals has
residue \(16I_2\), including every quadratic term.  Its explicit
eleven-soft-root identity supplies a strict positive diagonal.
Thus the entire parent, with all four retained terms restored,
has a bounded strict-residue Gram.  This is a finite polynomial
division proof; neither \(\delta^{-1}\) nor a mixed element is a
lifting coefficient.

The same body proves the complementary \(W=1,Z\) infinitesimal
chart with five exact current roots.  Its two full correction
coefficients contain respectively \(\Delta Z\) and
\(\Delta^2 f^2Z^2\), so vanish in residue even when \(\Delta\)
is a unit.  The original source signs and both strict gap openings
remain intact.

For the outer angle joins, no conclusion from the old infinitesimal-\(e\)
assertion is imported.  Use the two exact whole-parent identities in
Lemmas~\ref{lem:c1-two-current} and~\ref{lem:c1-critical-current}.
Their source hypotheses hold on all of UE1.  Indeed the seven native
coquotients have residues \(2,2,2,2,2,2,4\), regardless of \(e\).
Writing \(B_c=2c_5-Jf^2c_d\sim4\), the original weights after
normalizing \(d=1\) satisfy
\[
 C^{-1},D^{-1}\text{ infinitesimal},\qquad
 A\frac{CD}{C+D}
 =\frac{c_4c_3B_c}{c_d^2(aB_c+2Jc_3)}\sim\frac14.
 \tag{UE4}
\]
Here \(a+J=T\sim1\), so the denominator has residue \(64\)
without an individual lower bound on either opening.
Those identities therefore cover every testing chart outside
\(u,v\) infinitesimal, including literal projective infinities.
Literal axes, the zero state, signs and complex states are covered
by their exact full-parent normalization identities.

The two alternatives UE2 exhaust every positive source \(e\).
All scalings are within one labelled field, and every subtracted
term is restored.  This proves UE1 at all tests, with literal
\(\beta=0\) and every allowed original coefficient boundary.
The two real closed extensions are arbitrarily chosen labelled
factors; no algebraic-disjointness assumption is imposed.



% Source formulas: p9afbe73899af
\subsection{The source with infinitesimal $k,\zeta$ and unbounded $\eta$}

\label{component:p9afbe73899af}

\label{p9afbe73899af:sec:r7-c1-small}
This section concerns the original C1.1 source, not an enlarged family.
All coefficients are divided by $d>0$. Put
\[
 b=\eta^{-1/2},\quad g=\sqrt\zeta,\quad q=k^2\eta,
 \qquad x=b^2q,\quad a=b^2(1-r).
 \tag{SC1}
\]
Here $b,g$ are positive infinitesimals, $q>0$ is bounded, and
$\overline q\le1/2$. The new bounded two-normal reduction below is
asserted for $q$ a positive unit. The exact polynomial identities
remain valid more generally, but their indicated boundedness argument
would not control $g/\sqrt q$ when $q$ is infinitesimal.
No all-tests closure is inferred from the reduction.

Keep $K,R,T,M,J,E,Q,P_4,P_d,P_\alpha$ from (NC1)--(NC2), and set
\[
\begin{aligned}
 c_b&=b^2C/d=\frac{K(1+b^2r)E}{2(1-r)P_d},&
 a_b&=b^2A/d=\frac{P_4}{2rP_d},\\
 D_b&=D/(db^2),&\alpha_b&=\alpha/(db^2g^2),\\
 h_2&=\frac{Kq(1+b^2r)(1-tg^2)(1+qg^2)}{2P_d},&
 h_b&=\frac{Kq(1+b^2r)[1+(1-t)g^2]}{2P_d}.
\end{aligned}
\]
Thus $m/d=h_2/(g^2M)$ and $\beta/d+1/2=h_b/(g^2t)$.
The exact relation
\[
 h_b-h_2=\frac{Kq(1+b^2r)E g^2}{2P_d}
 \tag{SC2}
\]
removes the apparent opening dependence from the first normal of the
paired $m,\beta$ currents. The quantities $D_b,\alpha_b,h_2,h_b$ are
bounded; $\alpha_b,h_2,h_b$ are units in the stated $q$-unit domain.
Their residues are respectively $(1-2q)/2,q^2/2,q/2,q/2$, with
bounded source letters on the right understood at residue.

\paragraph{Exact full-parent reduction.}
For bounded affine tests define
\[
 Q_b=(W+uY)/b,\quad Y_b=Y/b^2,\quad p=t/R,\quad
 N=\frac{(X+uZ)/b-pvQ_b}{g},
\]
\[
 H=uQ_b/b,\quad L=vQ_b/g,\quad V=uvZ/b,\quad
 B_b=gN+pvQ_b,\quad D_u=1+u^2,\quad D_v=1+v^2.
\]
All roots use only pure-source inverses. Retain the two nonnegative
original terms
\[
 \frac{4h_2t}{MR}N^2,
 \qquad \frac{4h_bM}{tR}D_vN^2.
\]
The complete original parent divided by $db^2$, minus these terms,
is exactly
\begin{align*}
 G={}&\frac4R\{(h_2+h_bD_v)N^2+2(h_bD_v-h_2)LN
              +[h_2(1-p)+h_bpD_v]L^2\}\\
 &+4D_v\{c_b(H-D_uY_b)^2+a_bD_uY_b^2\}
   +4D_bJ_b^2+4\alpha_b g^2D_u(vb^2Y_b-Z)^2+2K_b,
                                                        \tag{SC3}\\
 J_b={}&b^2(H-D_uY_b)+ubvB_b-D_uvZ,\\
 K_b={}&b^2(H-D_uY_b)^2+2buv(H-D_uY_b)B_b+u^2v^2B_b^2\\
 &+D_uv^2b^2Y_b^2+D_uV^2
       -2D_uV(Q_b+vB_b)+2D_uu^2vZY_b.
\end{align*}
These are entire quadratic polynomials in the state; no restriction
of a state normal has been made. In particular, (SC3) is a direct
input for collecting both complete first normals and their full hard
quadratic, rather than an identity solely on a contact graph.

\paragraph{A harmonic normal removes both $r$ openings.}
Put
\[
 s=c_b+a_b,\quad \rho=c_b/s,\quad h=c_ba_b/s,
 \quad \rho_0=Kr/R,\qquad
 N_y=\sqrt{s}\,(D_uY_b-\rho_0H).
\]
Then the complete $C,A$ line in (SC3), divided by $4D_v$, equals
\[
 \left[N_y+\sqrt{s}(\rho_0-\rho)H\right]^2+hH^2
 -\frac{a_bu^2}{D_u}
           \left(\frac{N_y}{\sqrt{s}}+\rho_0H\right)^2.
 \tag{SC4}
\]
The inverse $1/D_u$ is pure testing, and every coefficient in (SC4)
is bounded uniformly through either $r\to0$ or $r\to1$.
Indeed $s$ is bounded below, $h$ is a positive unit, and
$a_b/s$, $a_b\rho_0/\sqrt{s}$ and $a_b\rho_0^2$ are bounded.
There is no mixed harmonic denominator.

Here is the exact source cancellation that also centers the normal
without an opening loss. Define
\[
 S_0=Kr(1+b^2r)E+(1-r)P_4,
\]
\[
 T_0=q^2(1-tg^2)^2+qrg^2+2qtg^2-qg^2-2q+r.
\]
Then
\[
 \rho=\frac{Kr(1+b^2r)E}{S_0},\qquad
 Kr-\rho R=-\frac{b^2r(1-r)K T_0}{S_0},\qquad
 h=\frac{K(1+b^2r)E P_4}{2P_dS_0}.             \tag{SC5}
\]
Both $S_0$ and the other denominators are positive units. Also
$s\asymp[r(1-r)]^{-1}$, so $\sqrt{s}(\rho_0-\rho)$ is
infinitesimal without a lower bound on either opening.
At the actual contact,
\[
 (u,v)=(b\sqrt q,\sqrt q g),\quad
 Q_b=RgZ,\quad Y_b=r\sqrt q gZ,
\]
and both $N,N_y$ vanish exactly. Thus (SC3)--(SC5) remove every
unbounded source weight while retaining the actual contact and every
state normal.

\paragraph{The leading angular circuit.}
For $u,v$ infinitesimal, the leading scalar after the harmonic state
normal has been removed is twice
\[
 (1-q)H^2+qL^2+V^2+(1-2q)(vZ)^2+q^2(gZ)^2-2VQ_b.
\]
Here the leading relation $H(vZ)=VQ_b$ is exact for the displayed
roots. It splits as
\[
 (1-2q)(H-vZ)^2+
 qH^2+qL^2+V^2+q^2(gZ)^2-4qVQ_b.              \tag{SC6}
\]
The final four positive terms form one AM--GM circuit, since
\[
 (\sqrt qH)(\sqrt qL)V(qgZ)=(qVQ_b)^2.
\]
On the positive negative-core chart its radical factors separate
between the source and testing fields. This proves positivity of the
leading scalar, including $q=1/2$, where its separate binomial drops
out. The leading hard normal coefficients are $4q$ on $N$ and $4$
on $N_y$. This organizing form alone is not yet a certificate for
all the higher terms in (SC3).

A useful exact conductor can be stated without comparing any source
and testing rate. Normalize $Z=1$ on its pure-testing state chart,
put $c=\sqrt q$, $w=(W+uY)/R$, and now use the contact-exact roots
\[
 H_0=uw/b^2,\ L_0=vw/(bg),\ V_0=uv/b,\ D_0=c^2g,
 \quad A_0=cH_0,\ B_0=cL_0,\ C_0=V_0.
\]
Their contact is $(u,v,w)=(bc,cg,bg)$. Set
\[
 E_0=H_0-v,\quad
 S_v=A_0-D_0-cE_0,\quad
 S_u=C_0-D_0-S_v,\quad S_w=B_0-D_0-S_v.
\]
Direct multiplication gives
\[
\begin{aligned}
 H_0(u-bc)&=(u-bc)E_0+bS_u,\\
 H_0(v-cg)&=(v-cg)E_0+(v/c)S_v,\\
 H_0(w-bg)&=(w-bg)E_0+(bg/c)S_w.
\end{aligned}                                                     \tag{SC7}
\]
Only the positive-unit inverse $1/c$ occurs. These identities lift
contact-vanishing coefficients times $H_0$ to the displayed root
module with bounded coefficients on the small pure-testing chart.
They do not alone prove a full scalar square lift: a first-jet
syzygy among the redundant angular roots still has to be allocated,
particularly when the coefficient of $E_0^2$ vanishes at $q=1/2$.

\paragraph{Exact small-cost angular pivots.}
The four-term circuit has the exact symmetric decomposition
\[
 \sum_{i=1}^4A_i^2-4\sqrt{A_1A_2A_3A_4}
 =\frac13\sum_{i<j}(A_i-A_j)^2
  +\frac23\sum_{\{ij\}\mid\{kl\}}
       (\sqrt{A_iA_j}-\sqrt{A_kA_l})^2,
\]
where the second sum is over the three pair partitions. Use
$(A_1,A_2,A_3,A_4)=(A_0,B_0,C_0,D_0)$ above and assume the positive
chart $u,v,w>0$. Then the following identities are exact:
\[
\begin{aligned}
 c\sqrt{uw}\,E_0
 &=\sqrt{uw}(A_0-D_0)
      +b\sqrt{cg}(\sqrt{A_0D_0}-\sqrt{B_0C_0}),\\
 \sqrt{vw}(u-bc)
 &=\frac{b^{3/2}}{\sqrt c}
             (\sqrt{A_0C_0}-\sqrt{B_0D_0}),\\
 \sqrt{uw}(v-cg)
 &=b\sqrt{g/c}
             (\sqrt{B_0C_0}-\sqrt{A_0D_0}),\\
 \sqrt{uv}(w-bg)
 &=\frac{b^{3/2}\sqrt g}{c}
             (\sqrt{A_0B_0}-\sqrt{C_0D_0}).
\end{aligned}                                                     \tag{SC9}
\]
The coefficient factors on the right are infinitesimal; all radicals
separate between the source and testing fields. Consequently any fixed
positive real multiples of the squares of the four left-hand roots can
be inserted as additional positive Gram pivots at an infinitesimal cost
in the displayed strict circuit Gram. To see this without a limiting
assertion, write the four identities as $P=TZ$, with $Z$ the displayed
circuit roots and every entry of $T$ in the bounded maximal ideal. Then
$Z^TGZ=\tau P^TP+Z^T(G-\tau T^TT)Z$, and the second matrix retains
the strictly positive real residue of $G$ for every fixed real $\tau>0$.
This is a genuine full-rates Gram mechanism. It still requires showing
that the entire first-normal and scalar remainders of (SC3) allocate on
these weighted roots; the identities do not make that last membership
automatic.

\paragraph{A quantified remote improvement when $e$ is unbounded.}
The source identity $e^2=q/g^2$ and (NI3) imply that $m,\beta$ are
both unbounded when $e$ is unbounded. Three inverse-cycle rotations
send this adjacent pair to $C,A$ and transform tests as
$(u,v)\mapsto(v,1/u)$. Therefore the adjacent-event theorem covers
all tests with $v$ not infinitesimal. Combined with the previous
$C,A$ certificate, only $u,v$ infinitesimal remain in this $e$ type.
This includes every $q$-unit source considered above.

There is also a complete state-chart certificate there. Normalize the
pure-testing state by a coordinate of maximal absolute value. If
$W$ is a unit, cap $C,A,m,\beta$ at a sufficiently large positive
real $\mu$, retain all excesses, and leave every other event as its
original nonnegative square. At $u=v=X=Y=0$, the quadratic part in
$(u,v,X,Y)$ is
\[
 2(uW-Y)^2+2(X+uZ)^2-4uvWZ+
 4\mu\{(uW-Y)^2+Y^2+(vW-X-uZ)^2+(X+uZ)^2\}.
 \tag{SC8}
\]
The event part is strictly positive for every real $W\ne0$ and bounded
real $Z$: its zero equations give successively $Y=u=X=v=0$. A sufficiently
large finite $\mu$ therefore absorbs the entire signed quadratic core.
The exact remainder has degree at least three in these four small
roots, with bounded coefficients in $W,Z$, so it allocates with zero
residue. This proves the entire original parent on this chart.
If the core has negative sign and $W$ is infinitesimal, then $Z$ is
a unit, since otherwise the core residue forces $X=Y=0$, contrary
to maximal-coordinate normalization. Thus the remaining state chart
has $Z=1$, $W$ infinitesimal, and both testing half-angles infinitesimal.
This is a full quantified narrowing, not closure of C1.1.

\paragraph{Completion and exact outstanding obligation.}
The following section supplies a finite scalar square lift and both
complete first normals for the $q$-unit C1.1 subdomain, including
$D=0$. It uses a contact-centered $D$ term followed by the geometric
mean conductor; an unweighted $E_0^2$ perturbation by itself is not
controlled by the circuit uniformly at arbitrary relative rates.
For $q$ infinitesimal, the uniform normal and protected contact lift
in Sections~\ref{pc376284fc7ae:sec:r7-c1-q-inf-local}--\ref{pd2ecc6ac1891:sec:r7-c1-q-inf-remote}
resolve the additional $g/\sqrt q$ issue and all remote tests.
No assumption of bounded $e$ or $1/e$ is imposed. 



% Source formulas: p95350f69a353
\subsection{Positive-unit $q=x\eta$}

\label{component:p95350f69a353}

\label{p95350f69a353:sec:r7-c1-q-unit}
\begin{theorem}
On the original C1.1 source, suppose $q=k^2\eta$ is a positive bounded
unit. The entire original parent is tensor SOS at every independent
projective test and every real or complex state. All positive opening
rates and the permitted literal boundary $D=0$ are included.
\end{theorem}
The proof uses the exact full quadratic (SC3), not its leading scalar.
The $q$-infinitesimal part of C1.1 is not included in this theorem.

\paragraph{One more centered term.}
Use the variables of (SC1)--(SC5), and put $s_D=J/R$. Retain the
additional nonnegative square
\[
                  4D_b(J_b+s_DH)^2.                 \tag{SC10}
\]
It vanishes at actual contact, where $H=RvZ$ and $J_b=-JvZ$.
Its removal replaces $4D_bJ_b^2$ in (SC3) exactly by
\[
                  -8D_bs_DHJ_b-4D_bs_D^2H^2.
\]
In particular it removes the entire separate $(vZ)^2$ term.
This step allows $D_b=0$ literally and does not divide by $D_b$.

Set $Z=1$ on the remaining state chart. Put $c=\sqrt q$ and
$w=(W+uY)/R$, and use the four separated positive monomials
\[
 A=\frac{cuw}{b^2},\quad B=\frac{cvw}{bg},\quad
 C=\frac{uv}{b},\quad D=c^2g,\quad
 \Gamma=(ABCD)^{1/4}.
 \tag{SC11}
\]
Their exact contact is $(u,v,w)=(k,cg,bg)$, with $k=bc$, and
$A=B=C=D$ there. Notice that $H=RA/c$, $L=RB/c$, $V=C$,
$g=D/q$, and $Hv=R\Gamma^2/q$.

Here are explicit formulas for the entire scalar at $N=N_y=0$,
after (SC10). Write $a_t=1-\rho_0+pv^2$. Define
\begin{align*}
 P_H={}&4D_v\{c_b(1-\rho_0)^2+a_b\rho_0^2/D_u\}
       +2b^2\{a_t^2+v^2\rho_0^2/D_u\}\\
 &+4\alpha_bg^2b^4v^2\rho_0^2/D_u
       -8D_bs_Db^2a_t-4D_bs_D^2,\\
 P_L={}&4[h_2(1-p)+h_bpD_v]/R,\\
 P_{Hv}={}&-4[D_u(1+pv^2)-u^2\rho_0]
             -8\alpha_bg^2b^2\rho_0+8D_bs_DD_u.
\end{align*}
The exact scalar is
\[
 f=f_AA^2+f_BB^2+f_CC^2+f_DD^2-f_\Gamma\Gamma^2,
 \tag{SC12}
\]
where
\[
 f_A=R^2P_H/q,\quad f_B=R^2P_L/q,\quad
 f_C=2D_u,\quad f_D=4\alpha_bD_u/q^2,
 \quad f_\Gamma=-RP_{Hv}/q.
\]
These are rational functions of $u,v$ over bounded source coefficients,
with only positive-unit denominators. For example,
\[
 c_b(1-\rho_0)^2=
   \frac{K(1+b^2r)E(1-r)}{2P_dR^2},\qquad
 a_b\rho_0^2=\frac{P_4K^2r}{2P_dR^2}.
\]
Thus neither $r$ opening survives in their denominators.
For $u,v,b,g$ infinitesimal their residues are
\[
               \overline f_A=\overline f_B=
 \overline f_C=\overline f_D=2,\qquad
               \overline f_\Gamma=8.               \tag{SC13}
\]
This uses the exact source residues in (SC2), including
$\overline D_b=(1-2\overline q)/2$. It does not assume a positive
lower bound on $D_b$.

Here is an explicit finite denominator certificate used below.
Every function just displayed is $P(u,v)/(1+u^2)$ for a polynomial
$P$ with bounded pure-source coefficient array. For a monomial use
\[
 x^m-x_0^m=(x-x_0)\sum_{j=0}^{m-1}x^{m-1-j}x_0^j.
\]
Applying this identity to each monomial supplies finite first and
second difference arrays with bounded source coefficients. For the
denominator use only the exact identity
\[
 \frac1{1+u^2}-\frac1{1+k^2}
 =-\frac{(u-k)(u+k)}{(1+u^2)(1+k^2)}.
\]
Thus the first and second finite Taylor arrays about $(k,cg)$
are bounded and use only pure-testing $1+u^2$ and pure-source
$1+k^2$ denominators, both units. The same construction after
$u=U^2,v=V^2$ gives the half-variable arrays below. This is a
finite polynomial identity construction, not an analytic series.

\paragraph{Every first normal and the complete hard block.}
Let $\delta_\rho=\rho_0-\rho$. The full first-normal coefficients
of the remaining quadratic, with no normal suppressed, are
\begin{align*}
 \mathcal L_N={}&\frac8R(h_bD_v-h_2)L
       +4buvg(a_t-2D_bs_D)H-4D_uvgV,\\
 \mathcal L_y={}&8D_v\sqrt{s}\,\delta_\rho H\\
 &+\frac H{\sqrt s}\left[
 -\frac{8D_va_bu^2\rho_0}{D_u}+8D_bs_Db^2
 +\frac{8\alpha_bg^2b^4v^2\rho_0}{D_u}
 -4b^2a_t+\frac{4b^2v^2\rho_0}{D_u}\right]\\
 &+\frac{4bu}{\sqrt{s}}V
       -\frac{8\alpha_bgb^2v}{\sqrt{s}}(gZ).
                                                        \tag{SC14}
\end{align*}
Each is a bounded infinitesimal-coefficient linear combination of
$A,B,C,D$. At actual contact it is zero exactly, since the source
is a stationary zero and every removed term vanishes there.
The coefficients of $N^2,N_y^2,NN_y$ are respectively
\[
\begin{aligned}
 H_{NN}&=4(h_2+h_bD_v)/R+2g^2u^2v^2,\\
 H_{yy}&=4D_v\left(1-\frac{a_bu^2}{D_us}\right)
       +\frac{4\alpha_bg^2b^4v^2}{D_us}
       +\frac{2b^2}{s}\left(1+\frac{v^2}{D_u}\right),\\
 H_{Ny}&=-4buvg/\sqrt{s}.
\end{aligned}                                                     \tag{SC15}
\]
Thus the hard Gram has residue $\operatorname{diag}(4q,4)$.
In (SC15), $H_{Ny}$ denotes the polynomial coefficient of $NN_y$,
not the off-diagonal matrix entry.

\paragraph{A finite geometric-mean square lift.}
Here is a detailed lift for the particular raw module (SC12).
It is not a claim about arbitrary contact-vanishing polynomials.
Define $\alpha_i=A_i-\Gamma$, with
$(A_1,A_2,A_3,A_4)=(A,B,C,D)$, and put
\[
 T_4=\sum_iA_i^2-4\Gamma^2.
\]
A fixed positive Gram on the following finite roots represents $T_4$:
\[
 \alpha_i,\qquad
 Q_{ij|kl}=\sqrt{A_iA_j}-\sqrt{A_kA_l},\qquad
 \sqrt\Gamma(\sqrt{A_i}-\sqrt{A_j}),\qquad
 \sqrt\Gamma((A_iA_j)^{1/4}-(A_kA_l)^{1/4}).          \tag{SC16}
\]
Use all four $\alpha_i$, all three pair partitions, and all six
unordered differences in the third group. To verify strictness, average
positive multiples of the symmetric circuit decomposition in (SC9) and
\[
 T_4=\sum_i\alpha_i^2+2\Gamma\left(\sum_iA_i-4\Gamma\right).
\]

over the three decompositions
\[
 \sum_iA_i-4\Gamma=(\sqrt{A_i}-\sqrt{A_j})^2
 +(\sqrt{A_k}-\sqrt{A_l})^2
 +2((A_iA_j)^{1/4}-(A_kA_l)^{1/4})^2.
\]
The additional differences $A_i-A_j$ in the first decomposition are
$\alpha_i-\alpha_j$ and give nonnegative corrections to the already
strict $\alpha$ block.

Write $\epsilon_i=f_i-2$, $\epsilon_\Gamma=f_\Gamma-8$, and
$\sigma=\sum_i\epsilon_i-\epsilon_\Gamma$. Then
\[
 f-2T_4=\sum_i\epsilon_i\alpha_i^2
       +2\Gamma\sum_i\epsilon_i\alpha_i+\Gamma^2\sigma.
 \tag{SC17}
\]
All $\epsilon$ have zero residue. The exact contact value implies
$\sigma(k,cg)=0$. Clear the pure-testing denominator $D_u$ if needed;
monic division gives
$\sigma=(u-k)\sigma_u+(v-cg)\sigma_v$ with bounded coefficients.
For the monomials (SC11), the exact conductors are
\[
\begin{aligned}
 \Gamma(u-k)&=\sqrt{uk}\,Q_{13|24},\\
 \Gamma(v-cg)&=\sqrt{vcg}\,Q_{23|14},\\
 \Gamma(w-bg)&=\sqrt{wbg}\,Q_{12|34}.
\end{aligned}                                                     \tag{SC18}
\]
Thus the last term in (SC17) becomes $\Gamma$ times a linear
combination of the $Q$ roots with infinitesimal bounded coefficients.
The exact identity
\[
 Q_{ij|kl}-\tfrac12(A_i+A_j-A_k-A_l)
 =-\tfrac12(\sqrt{A_i}-\sqrt{A_j})^2
  +\tfrac12(\sqrt{A_k}-\sqrt{A_l})^2               \tag{SC19}
\]
changes these terms into infinitesimal quadratic corrections on (SC16)
and $\Gamma\sum_i\mu_i\alpha_i$, where each $\mu_i$ is infinitesimal
and is rational in $\sqrt u,\sqrt v$ with bounded coefficient arrays
and positive-unit denominators.

The exact three contact jets imply that the four source values
$\mu_i(k,cg)$ coincide. Indeed the differential vectors of
$\alpha_i/\Gamma$ at contact are the four rows
\[
 (1/2,-1/2,1/2),\quad(-1/2,1/2,1/2),\quad
 (1/2,1/2,-1/2),\quad(-1/2,-1/2,-1/2)
\]
in logarithmic coordinates $(u/k,v/(cg),w/(bg))$. Their sole linear
relation is the sum of all four rows. Cancelling the nonzero source
factors in this equality establishes coincidence; no inverse of a
small number enters a Gram coefficient.
Call the common infinitesimal value $\mu_*$. The common term is
\[
 \mu_*\Gamma\sum_i\alpha_i
   =\mu_*\Gamma\left(\sum_iA_i-4\Gamma\right),
\]
an infinitesimal correction on the explicitly displayed variance roots.
For each remaining $\mu_i-\mu_*$, monic division in $\sqrt u,\sqrt v$
has bounded coquotients. Apply the half-variable conductors
\[
\begin{aligned}
 \Gamma(\sqrt u-\sqrt k)
   &=(uk)^{1/4}\sqrt\Gamma
          ((AC)^{1/4}-(BD)^{1/4}),\\
 \Gamma(\sqrt v-\sqrt{cg})
   &=(vcg)^{1/4}\sqrt\Gamma
          ((BC)^{1/4}-(AD)^{1/4}).
\end{aligned}                                                     \tag{SC20}
\]
Their coefficient factors are infinitesimal. Multiplying by the
remaining $\alpha_i$ proves that the entire scalar correction in
(SC17), not merely its low-order jet, is a bounded zero-residue
quadratic form on (SC16). This uses only two finite monic divisions.
The boundedness of their coefficient arrays follows from the explicit
rational functions (SC12), whose test denominators are powers of
$D_u$, and from the displayed square-root substitutions; bounded
function values alone are not being used as a substitute for this fact.

The same argument for a first normal is shorter. If
$\mathcal L=\sum_i a_iA_i$, write
$\mathcal L=\sum_i a_i\alpha_i+\Gamma\sum_i a_i$.
Its exact contact zero gives $\sum_i a_i(k,cg)=0$; one monic division
and (SC18) allocate it on (SC16), with infinitesimal coefficients.
Apply this to both complete expressions (SC14).
Together with (SC15), this yields a bounded strictly positive-residue
Gram for the entire original remainder. Restore (SC10), the paired
terms, and all positive pure-source and pure-testing normalizations.

\paragraph{Exhaustion and legal radicals.}
For $q$ a unit, $e^2=q/g^2$ is unbounded. The exact adjacent-event
rotations and the $W$-unit state argument (SC8) therefore reduce every
uncovered negative-core test to $u,v,W$ infinitesimal and $Z$ a unit.
Literal axes are direct original-core exits. Use exact sign gauges and
normalize $Z=1$. The common pure-testing Schur identity forces
$W+uY>0$ on the remaining negative-core chart: with
$w_t=(W+uY)/D_u$, the eliminated scalar is
\[
 \frac{uvD_uD_v}{\mathcal D}
       (D_uuv w_t^2-2D_uD_vw_t+uv),\qquad
 \mathcal D=u^2v^4+u^2v^2+u^2+v^4+2v^2+1>0,
\]
which is nonnegative when $w_t\le0$; its eliminated two-square
matrix is strictly positive. Hence $w=(W+uY)/R>0$.
Every monomial in (SC11), and every product under a radical in
(SC16)--(SC20), factors into a positive pure-source quantity and a
positive pure-testing quantity. The occurrence of $R$ in $w$ is a
positive pure-source unit, not a mixed denominator. All test
coefficients remain bounded because $u,v,w$ are infinitesimal.
If the general toric contact lemma is used, its third coordinate
is literally the pure-testing $W+uY$, with pure-source contact
$Rbg$; this direction is inactive in every coefficient function.
Complex state evaluations split into real and imaginary parts.
This supplies every independent test and proves the theorem.



% Source formulas: p7d43f3fccb7d
\subsection{A uniform normal at infinitesimal $q$}

\label{component:p7d43f3fccb7d}

This intermediate lemma by itself is not an certificate of the full
$q$-infinitesimal C1.1 source. The following two completion sections
supply the scalar allocation and remote joins. Keep (SC1)--(SC3), but do
not discard the whole second paired term. Its addition to the
$N^2$ coefficient in (SC3) gives exactly
\[
 \frac4R(h_2+h_bD_v)+\frac{4h_bM}{tR}D_v
       =H_0+\frac{4h_b}{t}v^2,
 \qquad H_0=\frac{4h_2}{R}+\frac{4h_b}{t}.
\]
Retain the last summand as the positive square
$4h_b(vN)^2/t$, and use the pure-source normalization
\[
                         \widehat N=\sqrt{H_0}\,N.
\]
The exact original inequality $\beta\ge0$ gives
$h_b/t\ge g^2/2$. The explicit formula for $h_2$ gives
$h_2/q$ a positive bounded unit, uniformly through both openings.
Consequently $H_0\ge c_0(q+g^2)$ for some positive real $c_0$.
This argument also permits $\beta=0$ literally. It does not require
$e=\sqrt q/g$ or its reciprocal to be bounded.

The coefficient of $\widehat N^2$ from this source part is exactly
one. Its remaining contribution $2g^2u^2v^2/H_0$ is bounded and
infinitesimal on the small-angle testing chart. Moreover
\[
 \frac{g}{\sqrt{H_0}},\quad
 \frac{\sqrt q}{\sqrt{H_0}}
\]
are bounded. Since $h_b-h_2=O(qg^2)$, the complete first normal
(SC14), after this normalization, is an infinitesimal-coefficient
linear combination of $B,H,V$ for $u,v$ infinitesimal. The first
coefficient is bounded because
\[
 \frac{h_bD_v-h_2}{\sqrt{qH_0}}
            =O(g^2+v^2),
\]
and the other two coefficients are bounded by constant multiples of
$buv$ and $v$. The hard cross coefficient with $N_y$ is likewise
infinitesimal.

The issue left by this lemma is scalar and angular: $H$ is unweighted in this
statement, whereas the four-term circuit controls $\sqrt q H$.
One must retain a strict $D$-based contact normal and allocate every
unweighted $H$ coefficient on that normal and the circuit. The
naive full extraction (SC10) fails in this branch: its raw $P_H$
satisfies the exact formal specialization
\[
       P_H\big|_{q=0,\,v=0}=-\frac{2ru^2}{1+u^2},
\]
so $P_H/q$ need not be bounded. Thus the corrected hard-normal
normalization removes the former $g/\sqrt q$ obstruction. The later
protected-square lift establishes the missing scalar allocation.



% Source formulas: pc376284fc7ae
\subsection{Infinitesimal $q$ at the two small testing angles}

\label{component:pc376284fc7ae}

\label{pc376284fc7ae:sec:r7-c1-q-inf-local}
\begin{theorem}
On the original C1.1 source with $q=k^2\eta$ infinitesimal, the entire
original parent is tensor SOS whenever both independent testing
half-angles $u,v$ are infinitesimal. All original opening rates and
$\beta=0$ are retained. The boundary $D=0$ is empty here because
$D_b$ has residue $1/2$.
\end{theorem}
This statement includes every state; it does not yet assert the
ordinary or unbounded $v$ testing charts.

Use the preceding source notation and put $c=\sqrt q$, $e=c/g$,
$h=H/R$, and $E_*=h-v$. Work first with $Z=1$ and $u,v,h>0$.
Define
\[
 A=ch,\quad B=cL/R,\quad C=uv/b,\quad D=qg,\quad
 \Gamma=(ABCD)^{1/4},\quad d_i=A_i-\Gamma.
\]
Then
\[
 \Gamma^2=qhv,\qquad BC=e v^2h,\qquad
 u h=uE_*+bC.
 \tag{QI1}
\]
Here $h=u(W+uY)/(Rb^2)$ is separated into a positive pure-source
factor and a positive pure-testing factor. Every radical below is
therefore legal; no mixed inverse or order comparison is used.
The exact original inequality $\beta\ge0$ gives $t=O(e^2)$, while
$t$ is bounded. Hence $t/e$ and $\sqrt{t/e}$ are bounded.

\paragraph{Exact scalar, with the $D$ term restored.}
Write $U=u^2,V=v^2$, $D_u=1+U,D_v=1+V$, and
$a_t=1-\rho_0+(t/R)V$. At both state normals zero the entire original
scalar is
\[
 a h^2+b_*hv+c_*v^2+P_L L^2+2D_uC^2+4\alpha_b g^2D_u,
 \tag{QI2}
\]
where
\begin{align*}
 a={}&4D_vR^2\{c_b(1-\rho_0)^2+a_b\rho_0^2/D_u\}
 +2b^2R^2\{a_t^2+V\rho_0^2/D_u\}\\
 &+4\alpha_b g^2b^4VR^2\rho_0^2/D_u
 +4D_bb^4R^2a_t^2,\\
 b_*={}&-4R\{D_u(1+tV/R)-U\rho_0\}
 -8\alpha_b g^2b^2R\rho_0-8D_bb^2Ra_tD_u,\\
 c_*={}&4D_bD_u^2.
\end{align*}
All these rational functions have bounded source coefficient arrays;
only $D_u$ and positive pure-source units occur in their denominators.
The $r$ openings cancel exactly as in (SC12).
In particular
\[
 D_b=\frac{N_D}{2JP_d},\qquad
 N_D=\frac{K(1+qg^2)(1+b^2r)-JP_d}{b^2}
\]
has a polynomial numerator: the numerator before division vanishes
identically at $b^2=0$, where $P_d=1+qg^2$. This is a finite monic
division, so it introduces no unbounded source coefficient array.
Also $\alpha_b/q^2$, $h_2/q$, and $h_b/q$ have explicit bounded
arrays from their source formulas. These observations justify the
$q$ coquotients below uniformly through both source openings.

Use formal specialization at $q=0$ only to construct finite polynomial
coquotients; it is not a new physical source. Put $z=b^2$,
$J_0=1-z(1-r)$, and $S=1-r+tV$. The three specialized coefficients are
\begin{align*}
 a_0={}&2D_v\{(1+zr)(1-r)+r/D_u\}
       +2z\{S^2+Vr^2/D_u\}+2z^2S^2/J_0,\\
 b_0={}&-4\{D_u(1+tV)-Ur\}-4zSD_u/J_0,\\
 c_0={}&2D_u^2/J_0.
\end{align*}
Define $\widetilde b_0=b_0+4tD_uV$, and the bounded finite
coquotients
\[
 X=a_0+\widetilde b_0+c_0=U X_U+V X_V,
 \qquad Y=-\widetilde b_0-2c_0=U Y_U+V Y_V.
 \tag{QI3}
\]
For definiteness, first subtract the value at $U=0$, divide by $U$,
and then divide that value by $V$. Their two $V$ coquotients are
\[
 X_V=2+2zr(1-2t)+\frac{2z(1+zr)t^2V}{J_0},
 \qquad Y_V=\frac{4zt}{J_0}.
 \tag{QI4}
\]
Thus $X_V=2+\varepsilon_V$, $\varepsilon_V\in B_{\rm Inf}$;
$X_U,Y_U$ have bounded arrays. Likewise write
$a=a_0+qa_1$, $b_*=b_0+qb_1$, $c_*=c_0+qc_1$ by finite numerator
division. All $a_1,b_1,c_1$ have bounded arrays. Direct source
specialization gives
\[
 \overline{a_1-c_1}=2,\qquad
 \overline{b_1+2c_1}=-8,\qquad
 \overline{R^2P_L/q}=2,\qquad
 \overline{4\alpha_bD_u/q^2}=2.
 \tag{QI5}
\]

The following is the exact scalar identity to be allocated:
\begin{align*}
 (\mathrm{QI2})={}&(c_0+qc_1)E_*^2+2D_uC^2
 +X_U(uE_*+bC)^2+Y_Uu(uE_*+bC)E_*\\
 &+X_V(vh)^2+Y_Vv(vh)E_*
 -4(t v/e)D_uBC\\
 &+(a_1-c_1)A^2+(b_1+2c_1)\Gamma^2
 +(R^2P_L/q)B^2+(4\alpha_bD_u/q^2)D^2.
 \tag{QI6}
\end{align*}
It uses only (QI1) and $qv^2=2\Gamma^2-A^2+qE_*^2$.
In particular the apparently dangerous negative $t$ term is already
an infinitesimal-coefficient $BC$ term, since $v$ is infinitesimal.

\paragraph{Two extra contact squares and their finite first jets.}
Set
\[
 F_*=vh-gD,\quad
 R_F=\frac{\sqrt{g\Gamma}}{c}(\Gamma-D),\quad
 V_0=cv,\quad R_E=\sqrt\Gamma(\sqrt A-\sqrt{V_0}).
\]
All four roots vanish at actual contact. Exact multiplication gives
\begin{align*}
 (vh)^2&=F_*^2+2R_F^2+2gF_*d_4+g^2D^2-4g^2\Gamma d_4,\\
 c\Gamma E_*&=2\Gamma d_1-R_E^2,\\
 \Gamma F_*&=-2g\Gamma d_4+R_F^2/g.
 \tag{QI7}
\end{align*}
The root $R_E$ receives a strict real square pivot from $E_*^2$:
\[
 qE_*^2=(A-V_0)^2
       =4R_E^2+(\sqrt A-\sqrt{V_0})^4.
 \tag{QI8}
\]
Reserve this fraction $q$ of the strict $E_*$ pivot; the remaining
coefficient still has positive real residue. The leading $2(vh)^2$
in (QI6), by (QI7), gives strict $F_*$ and $R_F$ pivots, with only
infinitesimal quadratic corrections and the printed first-jet term.
The remaining leading angular form is $2(\sum A_i^2-4\Gamma^2)$,
with the strict finite toric square list from (SC16).

For precision, the only additional bare-$\Gamma$ normal couplings
left by (QI6), apart from ordinary toric first jets, have the form
$\Gamma E_*a_*+\Gamma F_*b_*'$ with
\[
 a_*=bu(2X_U+Y_U)+Y_Vvg,
 \qquad b_*'=2\varepsilon_Vg.
\]
At contact $u=bc,v=cg$ these satisfy
\[
 a_{*0}/c=b^2(2X_{U0}+Y_{U0})+Y_{V0}g^2\in B_{\rm Inf},
 \qquad b_{*0}'/g=2\varepsilon_{V0}\in B_{\rm Inf}.
\]
Their differences from contact have finite bounded monic Taylor
arrays. The conductors (SC18) allocate those differences against
$E_*,F_*$ with infinitesimal coefficients. Apply (QI7) to their
constant terms. What remains is a bounded infinitesimal toric
quadratic plus $\Gamma\sum\mu_i d_i+\Gamma^2\sigma$.
The exact scalar value and three original first jets imply its
contact value and first-jet balances. The same two finite divisions
as in (SC17)--(SC20), or the Hadamard contact-lift identity, allocate
it on the strict toric root list. No infinite series or unsupported
ideal-membership assertion is used: (QI3), the $q$ coquotients,
and the two displayed contact coefficient ratios are the finite
certificates needed for the additional roots.
Every coefficient function in these Taylor steps depends only on
$u,v$. Thus only the first two conductors in (SC18)--(SC20) are
used; the third state-coordinate jet supplies the inactive Hadamard
balance. No bound on $(W+uY)/R$ after $Z$ normalization is assumed.

\paragraph{Both complete state normals.}
Use $\widehat N=\sqrt{H_0}N$, $H_0=4h_2/R+4h_b/t$, as in the
preceding uniform-normal lemma. Add back the centered $D$ term
from (SC10), including its normals. Its scalar root is exactly
\[
 J_*=(T+b^2tV)h-D_uv
 =E_*+b^2c(1-tg^2)A-buC+b^2tvF_*+b^2tvgD.
 \tag{QI9}
\]
Its full root is $J_*+ubvgN-b^2N_y/\sqrt s$.
All additional hard coefficients are infinitesimal after the
$\widehat N$ normalization.

For the first $N$ normal in (SC14), its $H$ term is a bounded
infinitesimal multiple of $u h=uE_*+bC$; its other terms are
infinitesimal multiples of $B,C$. For the harmonic normal, multiply
its $H$ coefficient by $\sqrt s$. After the exact opening
cancellations, its value at $q=U=V=0$ is zero. Finite numerator
coquotients therefore give
\[
 \sqrt s\,[H\text{-coefficient in }\mathcal L_y]
       =qK_q+UK_U+VK_V
 \tag{QI10}
\]
with bounded arrays. Use
\[
 qH=RcA,\quad u^2H=Ru(uE_*+bC),\quad
 v^2H=Rv(F_*+gD).
\]
The two remaining terms of $\mathcal L_y$ have infinitesimal
coefficients on $C,D$ since $\alpha_b/q^2$ is bounded. Thus both
complete first normals, including (QI9), are infinitesimal
combinations of $E_*,F_*,A,B,C,D$. Their exact contact zeros make
the coefficient sum of $A,B,C,D$ zero at contact. One finite monic
division and (SC18) allocate these on the same strict square list.
The hard block has residue $\operatorname{diag}(1,4)$ and all its
other entries have zero residue. This proves a bounded positive
residue Gram for the entire original parent, after restoring all
retained positive terms.

\paragraph{States and signs.}
The restriction $Z=0$ is positive by the original pure-testing state
restriction. For $Z\ne0$, divide the state by the pure-testing $Z$;
its size is immaterial, since every coefficient array above depends
only on $u,v$. On the negative-core chart, exact sign gauges give
$u,v>0$ and the common pure-testing Schur identity forces
$W+uY>0$. Hence $h>0$, as required. The other core-sign chart and
literal testing axes are direct original exits. Complex states
split into real and imaginary parts. This completes the asserted
small-$u$, small-$v$ theorem.



% Source formulas: pd2ecc6ac1891
\subsection{Infinitesimal $q$ at the other Bob testing angles}

\label{component:pd2ecc6ac1891}

\label{pd2ecc6ac1891:sec:r7-c1-q-inf-remote}
\begin{theorem}
The original C1.1 source with $q=k^2\eta$ infinitesimal is tensor SOS
at every independent test. Together with the $q$-unit theorem this
closes the entire original C1.1 source.
\end{theorem}
The common $C,A$-large certificate already handles $u$ not
infinitesimal. The preceding local theorem handles $u,v$
infinitesimal. Here are the two remaining $v$ charts. Every positive
opening rate and the literal permitted coefficient boundary is kept.

\paragraph{Ordinary $v$: the full normal residue is strict.}
Let $v$ be a positive bounded unit, use $Z=1$ on the negative-core
chart, and retain the roots of (QI1). Normalize $N$ by $H_0$ as
before, but now restore the positive term $4h_b(vN)^2/t$.
Write
\[
 \tau=\operatorname{res}t,\qquad
 \rho=\operatorname{res}(t/e^2),\qquad
 a^2=\frac{\tau}{2(1+\tau)},\quad
 \gamma^2=\frac{\rho}{2(1+\tau)}.
\]
The original $\beta\ge0$ equation implies $0\le\rho\le1$;
$0\le\tau\le1$ follows from $M=R-t>0$ and $R\sim1$.
The entire residue on $(\widehat N,B,C)$ has diagonal entries
\[
 1+\frac{v^2}{1+\tau},\qquad 2(1+\tau v^2),\qquad 2,
\]
and polynomial cross coefficients
\[
 4av^2,\qquad -4\gamma v,\qquad -4\sqrt{\tau\rho}\,v.
\]
Its exact square decomposition is
\begin{align*}
 &2(C-\gamma v\widehat N-\sqrt{\tau\rho}\,vB)^2
       +\widehat N^2+2B^2\\
 &\qquad +(1-\rho)v^2
    \left(\frac{\widehat N}{\sqrt{1+\tau}}+
                         \sqrt{2\tau}\,B\right)^2.
 \tag{QR1}
\end{align*}
Thus it is strictly positive definite, also at $\rho=1$, the
possible $\beta=0$ endpoint, and at $\tau=0$ or $1$.
The other residue blocks are exactly
\[
                       2E_*^2+2(vh)^2+4D_vN_y^2.
\]
There are no nonzero leading cross entries with these blocks.

For completeness, the boundedness check uses the entire raw scalar
(QI2) and both full normal formulas (SC14), with (QI9) restored.
Terms containing $u h$ use $u h=uE_*+bC$; all their coefficients
are infinitesimal. The $N$--$B$ and $N$--$C$ coefficients with
nonzero residue are precisely those retained in (QR1). The harmonic
normal has infinitesimal coefficient on $h$ for bounded $v$:
after the opening cancellation its coefficient lies in
$(b^2,q,u^2)$ with bounded arrays on this chart. Its other
coefficients are infinitesimal as well.
All remaining source-constant terms and weak monomials enter the
same strict Gram through the exact bounded pure-testing identities
\[
                   h=(vh)/v,\qquad
                   1=((vh)-vE_*)/v^2.
\]
For example $A=ch$ and $D=qg$ are infinitesimal-coefficient
combinations of these two roots. Consequently the actual full Gram
has the positive residue just displayed. Restore all retained
terms. The same sign, $Z=0$, and complex-state joins as in the
local theorem apply; no restriction on the size of the normalized
state was used.

\paragraph{Unbounded $v$: a five-root endpoint frame.}
Put $w=1/v$, so $w$ is infinitesimal, and multiply the original
parent by $w^2$. The exact source formulas give positive real
constants $\mu,d_0$ with
\[
 C/d\ge\mu/b^2,\qquad A/d\ge\mu/b^2,\qquad
 D/d\ge d_0 b^2.
\]
For instance sufficiently small fixed constants below $1/4$ work,
since $b^2C/d,b^2A/d$ are uniformly bounded below and
$D/(db^2)\sim1/2$. Retain all excesses and the entire $m,\alpha,
\beta$ terms. No assumption on $e$ is needed.
Use the five roots
\[
              P=uW/b,\quad Q=Y/b,\quad X,\quad J_b'=bZ,
                     \quad V_b=uwZ.
\]
The homogenized signed core is exactly
\begin{align*}
 K={}&V_b^2-2buV_bQ-2wPJ_b'+2wV_bX+b^2w^2P^2
       +2buwPX-2b^2w^2PQ\\
 &+(D_u+w^2)(X^2+b^2Q^2)-2buwXQ.
 \tag{QR2}
\end{align*}
The three selected event terms are
\[
 4\mu D_w(P-Q)^2+4\mu D_uD_wQ^2,
 \qquad
 4d_0(b^2wP+buX-b^2wQ-J_b')^2.
 \tag{QR3}
\]
Combining these with $2K$, the bounded Gram has strictly positive
residue
\[
 4\mu\{(P-Q)^2+Q^2\}+4d_0(J_b')^2+2X^2+2V_b^2.
\]
All other displayed coefficients have zero residue. This is a full
state-quadratic certificate, not a restriction to a contact graph.
It includes the homogeneous endpoint $w=0$ directly and uses no
inverse involving independent source and testing variables.

The infinitesimal, unit, and unbounded $v$ charts exhaust all tests
remaining after the common $C,A$ certificate. Hence the entire
$q$-infinitesimal C1.1 source, and consequently C1.1 itself, is closed.



% Source formulas: p69995d829bb2
\subsection{Unit $\eta,e,r$}

\label{component:p69995d829bb2}

\begin{theorem}
Use the unchanged physical native C1 source, $c=2$, $\lambda=n=0$,
with all original event signs and positive openings $a,t,M,J$.
Suppose $k,\zeta$ are positive infinitesimals, and
\[
 \eta,\quad r=1-a\eta,\quad e=k\sqrt{\eta/\zeta}
\]
are positive bounded units. Then the complete original parent is tensor
SOS at every independent projective test and every real or complex state.
The openings may be arbitrarily small, and $\beta=0$ is included.
\end{theorem}

\paragraph{The double-infinity corner.}
With $p=1/u$, $q=1/v$, the exact homogeneous parent is
\begin{align*}
F_{\infty\infty}={}&4m(pW-pqX+Y-qZ)^2
 +4c(1+q^2)(W-pY)^2+4A(1+p^2)(1+q^2)Y^2\\
&+(4D+2d)(qW+X-pqY-pZ)^2
 +(4\alpha+2d)(1+p^2)(Y-qZ)^2\\
&+(4\beta+2d)(1+q^2)(pX+Z)^2
 -2d(1+p^2)(1+q^2)Z^2.
\end{align*}
At $Z=1$, after retaining $\alpha,\beta$, its complete quadratic part is
\[
 4cW^2+2dX^2+4D(X-p)^2+4AY^2+(4m+2d)(Y-q)^2.
 \tag{BE1}
\]
In particular the $(Y,q)$ determinant is $4A(4m+2d)>0$.
This is an exact full-parent expansion.

An even shorter full-parent certificate retains every event except
$c,A$. The remaining polynomial belongs to $(W,X,Y,qZ)^2$, with bounded
coefficients at $p,q=0$. Its residue Gram is the quadratic form
\[
 4cW^2+2dX^2+4AY^2+2d(Y-qZ)^2.                 \tag{BE2}
\]
It is strictly positive when $c/d,A/d$ are bounded below by positive
units. All cross terms of the complete polynomial are allocated on
these four roots, so (BE2) handles arbitrary states and both literal
projective endpoints, without a normal restriction.

\paragraph{Uniform lower bounds and outer tests.}
Write $H=r+\eta$. The short coquotients have residues $2rH$ for
$c_2,c_b,c_3,c_5,c_\alpha$, $2$ for $c_4$, and $4r$ for $c_d$.
The source identities imply
\[
 \frac md=\frac{e^2c_2}{M\eta c_d},\quad
 \frac cd=\frac{c_3}{a\eta c_d},\quad
 \frac Ad=\frac{\eta c_4}{c_d},\quad
 \frac Dd=\frac{2c_5-J\eta c_d}{2J\eta c_d}.
 \tag{BE3}
\]
Every ratio is bounded below by a positive real number. For the last,
$J+a=T\sim1$ and
\[
 \overline{2c_5-J\eta c_d}=4\bar r(\bar H-\bar J\bar\eta)
 =4\bar r>0.
\]
No reciprocal of an opening is required to be bounded. Choose positive
real lower bounds for the four ratios, cap these four coefficients at
the corresponding multiples of $d$, and retain their nonnegative
excess event squares. All selected-parent Gram coefficients divided
by $d$ are then bounded.

The physical-source theorem covers every noncommuting residue test.
For $u$ infinitesimal and $v$ a bounded unit, the selected residue form is
\[
 4m(vW-X)^2+2d(1+v^2)(X^2+Y^2)
 +4(c+A)(1+v^2)Y^2+4D(Y+vZ)^2\succ0.
\]
For $u=\infty$ in residue and $v$ a bounded unit, retain the signed
core at that commuting residue and use the four positive squares
$cW^2$, $AY^2$, $m(vY-Z)^2$, $D(W+vX)^2$; their common kernel is zero.
For $v=\infty$ in residue and $u$ a bounded unit, use the core
$2d(1+u^2)(X^2+Y^2)$ and the $c(uW-Y)^2,D(uX-Z)^2$ terms.
The $(0,\infty)$ corner has the strictly positive state residue
\[
 4mW^2+2dX^2+\{4(c+A)+2d\}Y^2+4DZ^2.
\]
The $(\infty,\infty)$ corner is (BE2).
At $(\infty,0)$ retain all events except $c,A$ and use the exact
four-root allocation on
\[
 (pX+Z,\,W,\,Y,\,vX),\qquad p=1/u.
\]
Its residue is
\[
 2d(pX+Z)^2+2d(W+vX)^2+4cW^2+4AY^2\succ0.
\]
The full selected polynomial lies in the square of this displayed
root ideal; the pure-testing change $Z=(pX+Z)-pX$ makes this immediate.
The other coefficients have zero residue. These homogeneous
identities include literal infinities. Literal zero axes have a
nonnegative pure-testing core and are immediate original-event SOS.
Only bounded $u$ and infinitesimal $v$ remain.

\paragraph{State entry and the $W$ endpoint.}
Normalize the state by a maximal coordinate. Set
$D_u=1+u^2$ and $w=(W+uY)/D_u$. On the negative pure-core branch,
if $w$ is infinitesimal then $Z$ is a unit: otherwise the residue core
$2\{(X+uZ)^2+(uW-Y)^2\}$ forces every normalized state coordinate
to have zero residue. Normalize $Z=1$; the exact pure-testing Schur
identity in (SD5) gives $w>0$ after the standard angle gauges.

If $w$ is a unit and $u$ is a unit, a zero of that core residue has
$Y=uw\ne0$. The positive lower bound for $A/d$ gives a strict state
residue after clipping the positive events. Thus only $u,v$
infinitesimal and $W\sim w$ a unit need an additional state check.
When $Z$ is infinitesimal, retain all events except clipped $m,A$ and
use the roots $(X,Y,uW,vW)$. Since $Z/W$ is a pure-testing
infinitesimal, the complete selected polynomial has residue
\[
 2d\{X^2+(uW-Y)^2\}+4m(vW-X)^2+4AY^2\succ0.
 \tag{BE4}
\]
Every remaining coefficient contains $u$, $v$ or $Z/W$; there is no
constant or first-normal term outside this frame. This includes
literal $Z=0$. If $Z$ is a unit, normalize $Z=1$ and use the next
raw scalar frame with $w$ a unit.

\paragraph{The exact remainder and its complete normals.}
Use exactly (SD3)--(SD8), their four positive terms, and the positive
clearing factor $\mathcal C$. Put
\[
 \delta=\Delta R/K,\quad E=b^2=\eta K/R,\quad
 A_0=wv,\ B_0=\delta v,\ C_0=\delta Euw,
 \ D_0=\delta^2k^3E,\quad x=B_0-C_0.
\]
Here $\delta,E$ are positive bounded units, and $a,t$ are bounded.
The factored five-term scalar is in
$(A_0,B_0,C_0)^2+(D_0^2)$, and the two complete first normals are
in $(A_0,B_0,C_0)$. These are the full formulas, without restricting
normals or replacing native coefficients by leading surrogates.

At $k=v=0$ and at the relevant small-$w$ or small-$u$ endpoint, put
$r_0=1-aE$ and $H_0=r_0+E$. Direct collection of the short formulas
gives the raw scalar residue
\[
 8r_0H_0D_u^2 A_0^2+
 8D_u\{x^2+u^2r_0H_0B_0^2\}+8r_0H_0D_uD_0^2.
 \tag{BE5}
\]
The hard block on $(N,\delta M_B)$ is
\[
 \operatorname{diag}(16r_0H_0,16r_0H_0+8E^2/D_u),
\]
and the complete first normals have residues $0,16Ex$.
The coupled $(\delta M_B,x)$ determinant is
$128D_ur_0H_0>0$.

If $u$ is a unit and $w$ infinitesimal, (BE5) is a strict Gram on
$(A_0,x,B_0,D_0,N,\delta M_B)$. If $u$ is infinitesimal and $w$ a
unit, use the exact relations
\[
 B_0=(\delta/w)A_0,\qquad C_0=(\delta/w)A_0-x.
\]
The factor $\delta/w$ is a product of a pure-coefficient unit and a
pure-testing unit; it is not a mixed inverse. The resulting residue
is the strict form
\[
 8r_0H_0A_0^2+8x^2+8r_0H_0D_0^2
 +16r_0H_0N^2+(16r_0H_0+8E^2)(\delta M_B)^2
 +16E(\delta M_B)x.                                  \tag{BE6}
\]
This supplies precisely the unit-$w$ chart for which the older
$\Delta$-small retained-current argument is not applicable.

For transparency, the only nontrivial raw coefficient collection in
(BE5) is
\[
 r_0H_0D_u+a^2E^2+r_0E(2a-1)=1+u^2r_0H_0,
 \qquad aE+r_0=1.
\]
The cross coefficient is $-16D_u(aE+r_0)=-16D_u$.
In the full normal formula $S_B$ has residue $-16E$ and $S_C$ has
residue $16E$. These prove the claimed full normal residues directly.

\paragraph{The sole contact chart.}
For $u,w,v$ infinitesimal and $Z=1$, use the unchanged monic
contact--conductor identities and the two-pairing circuit. Use weights
\[ (\delta,k,b,a,t,u,w,v):\quad(0,1,0,0,0,1,1,2). \]
The exact source collection is
\[
 (E_0)_{\le6}=8\rho T_4-8(\rho+E)k^2x^2,
 \quad (L_1)_{\le3}=0,\quad (L_2)_{\le3}=16Ex,
 \qquad \rho=r_0H_0-1.
\]
The source/contact substitutions preserve this grading. After the
displayed terms are removed, the scalar and normal lift coefficients
have positive weights, hence zero residues. The term $k^2x^2$ also
has a zero-residue coefficient. Thus the soft leading form is
$8x^2+8r_0H_0T_4$, with the same strict two-pairing circuit frame.
The hard block and its only first-normal coupling are those of (BE6),
and their determinant is $128r_0H_0>0$. Retain the old positive
$g_4^2$ square. The full eight-root remainder Gram is strict.
This invokes only the polynomial monic conductor already proved;
no opening or contact difference is inverted.

Finally restore the four complete positive terms, the retained
positive square, and all pure-coefficient and pure-testing factors.
The cases above exhaust both projective angle lines and every
projective state. Complex states split into their real and imaginary
evaluations. Neither this theorem nor the correction of (BE1) asserts
coverage when $r$ is infinitesimal, $\eta$ unbounded, or $e$
infinitesimal; those are separate source domains.



% Source formulas: p3763434080c2
\subsection{The centered-$A$ source with unit $e$}

\label{component:p3763434080c2}

\begin{theorem}[The centered-$A$ unit-$e$ slice]
Use the physical native C1 source with its unchanged original events,
$c=2$, $\lambda=n=0$, all positive openings $a,t,M,J$, and all original
nonnegative signs, including literal $\beta=0$.
Suppose $k,\zeta,r=1-a\eta$ are positive infinitesimals, while
$\eta$ and $e=k\sqrt{\eta/\zeta}$ are positive bounded units.
Then the full parent is tensor SOS at every independent projective test
and every real or complex state. No bound on $k^2/r$ is imposed.
\end{theorem}

The theorem is restricted to unit $e$. Its proof uses $\delta^{-1}$
as a bounded pure coefficient number; it does not assert the cases of
infinitesimal or unbounded $e$.

\paragraph{Centering the sole large event.}
The exact native contact has $Y_c=r u_cW_c$. Consequently
\[
 L_A=Y-r uW
\]
vanishes at contact with no asymptotic replacement. The native formulas
give $a$ a positive unit, $A$ unbounded and $Ar$ a positive unit of
residue $2$. Put $A_*=1/a$, a positive bounded pure coefficient unit.
The complete term
\[
             4(A-A_*)D_uD_v L_A^2                         \tag{CA1}
\]
is positive and has zero value and full first derivative at the exact
native contact. Retain it as a square.

It is essential to center this term before taking residues.
The remaining $A$ channel is exactly
\[
 4A_*D_uD_vL_A^2+
 4(Ar)D_uD_v\{2uWY-r u^2W^2\}.                         \tag{CA2}
\]
Only $Ar$ and $r$ occur in its coefficients. In particular this is a
bounded expression for every positive rate of $r$ relative to $k$.

\paragraph{A source chart with no hidden $k^2/r$ quotient.}
Keep the exact contact variables $\delta,k,E=b^2$ from (SD6), and put
\[
 K=1+k^2,\quad \ell=K-aE=Qr>0,\quad
 Q=K^2-\ell k^2,\quad a=(K-\ell)/E.
\]
Here $\delta,E,a$ are positive bounded units and $k,\ell$ are positive
infinitesimals. The exact source substitution is
\[
 \eta=KE/Q,\quad r=\ell/Q,\quad
 \Delta=\delta Q/K,\quad \zeta=\delta^2k^2QE/K,
 \quad R=K^2/Q,\quad T=K-k^2t\zeta.                    \tag{CA3}
\]
All denominators in (CA3) are positive pure coefficient units.
For $i=2,b,3,5,\alpha,d$ let $\widetilde c_i=c_i/\ell$,
defined by cancelling the explicit $L_c=\ell/Q$ factor in the short
native formulas. This cancellation gives rational expressions with
unit denominators; it does not leave an inverse of $\ell$.
For example, writing $H_c=r+\eta$ and $V_c=k^2\eta\zeta$,
\[
 \widetilde c_3=\frac{2KH_cQ_c(1-k^2\eta+tV_c)}Q,
 \qquad \widetilde c_d=\frac{4Q_cP_d}Q.
\]
Their residues are $2E$ for $i=2,b,3,5,\alpha$, and $4$ for $i=d$;
$c_4$ has residue $2$.

Define the following bounded pure coefficient quantities:
\[
 \begin{array}{ll}
 B_2=\dfrac{8a\widetilde c_2}{R\widetilde c_3\Delta^2},&
 B_b=\dfrac{8a\widetilde c_b}{R\widetilde c_3\Delta^2},\\[4pt]
 B_3=8a/T,&B_5=\dfrac{8a\widetilde c_5}{T\widetilde c_3},\\[4pt]
 L=Ar=\dfrac{2a\eta^2c_4}{Q\widetilde c_3},&
 d=\dfrac{2a\eta\widetilde c_d}{\widetilde c_3},\\[4pt]
 \alpha=\dfrac{2a\Delta^2k^6\eta^2\widetilde c_\alpha}
                    {\widetilde c_3}.&
 \end{array}                                                     \tag{CA4}
\]
These formulas explicitly account for the cancellation needed below.

\paragraph{The exact full remainder, before any collection.}
Retain the four original (SD3) terms, divided by the positive source
factor $\Delta^2\Lambda$, and retain (CA1). Use the unchanged pure
testing coordinate $w=(W+uY)/D_u$ and the hard normals $N,M_B$.
For $Z=1$ put $p=t/R$, $h=a/T$ and
\[
 q_A=vD_uw,\quad q_B=v(D_u-uN-puD_uv w),
\]
\[
 W=w+u(M_B+hq_B)/D_u,\qquad
 Y=uw-(M_B+hq_B)/D_u.
\]
The complete remainder is exactly
\begin{align*}
\mathscr G={}&B_2\{N^2-2q_AN+(1-p)q_A^2\}
 +B_bD_v\{N^2+2q_AN+pq_A^2\}\\
&+B_3D_v\{M_B^2+2q_BM_B+hq_B^2\}
 +B_5\{M_B^2-2q_BM_B+(1-h)q_B^2\}\\
&+4A_*D_uD_v(Y-r uW)^2
 +4LD_uD_v(2uWY-r u^2W^2)\\
&+4\alpha D_u(vY-1)^2
 +2dD_u\{(vY-1)^2-D_v\}.                         \tag{CA5}
\end{align*}
There are no normal terms above degree two. Every coefficient of
(CA5) is bounded. The only testing denominator is a power of the
positive unit $D_u$; every source denominator in (CA3)--(CA5) is a
positive unit. The possibly unbounded opening ratios remain exclusively
in the four retained positive terms.

\paragraph{The exact low-weight collection.}
The outer/state lemma in Section~\ref{component:p364e14238053} reduces the
independent tests to $Z=1$ and $u,w,v$ infinitesimal, with $w>0$ on
the negative pure-core branch. Use the unchanged circuit monomials
\[
 A_0=wv,\quad B_0=\delta v,\quad C_0=\delta Euw,
 \quad D_0=\delta^2k^3E,
\]
\[
 x=B_0-C_0,\quad y=A_0-kC_0,\quad
 z=uB_0-kC_0,\quad h_0=kC_0-D_0.
\]
Let $q_1,q_2$ be the same two separated radical pairings as in the
contact--conductor proof, and let
\[
 T_4=y^2+z^2+h_0^2-2yz+2q_1^2.
\]
Give $(\delta,k,E,\ell,t,u,w,v)$ weights $(0,1,0,1,0,1,1,2)$.
Then $x$ has weight two and $y,z,h_0,q_1,q_2$ have weight three.
Denote the scalar of (CA5) by $S$, and its two complete first normals
by $L_N,L_M$. Direct multiplication of (CA3)--(CA5) through the
indicated weights gives
\[
 \begin{aligned}
 S_{\le6}&=\frac{4P}{E^4\delta^2}x^2
 -\frac{16(k+u)}{E^2\delta^2}xz
 +\frac8{E\delta^2}T_4,\\
 (L_N)_{\le3}&=0,\\
 (L_M)_{\le3}&=
 \frac{8\{E^2+2E-2\ell(E+1)\}}{E^2\delta}\,x,
 \end{aligned}                                                   \tag{CA6}
\]
where the short coefficient polynomial is
\[
 \begin{split}
 P={}&E^2(E+4)-E(E^2+6E+4)\ell+2(E+1)(E+2)\ell^2\\
 &+E^2(8-3E)k^2+4E^2ku+E^2(E+4)u^2.
 \end{split}
\]
The full hard-normal block, from the quadratic part of (CA5), has
residue
\[
       \operatorname{diag}\left(\frac{16}{E\delta^2},
                                \frac{4(E^2+4)}E\right).          \tag{CA7}
\]
Here and in the residue Grams, $E,\delta$ denote their positive
residues. Formula (CA6), rather than a large expansion, is the finite
coefficient identity used by the proof. Its 37 raw scalar monomials
collect to the displayed three terms. 

\paragraph{Why all higher terms have harmless coefficients.}
The scalar of (CA5), after clearing the stated units, belongs to
$(A_0,B_0,C_0)^2+(D_0^2)$. This follows directly from $q_A=wvD_u$,
$q_B|_{N=M_B=0}\in(v)$, and $Y,W$ in (CA5): the centered scalar
$Y-r uW$ lies in $(v,uw)$; its square is in the same raw monomial
square ideal since $\delta,E$ are units. The sole constant state term
comes from $\alpha$, which has the factor $k^6$, and is a bounded
unit multiple of $D_0^2$. Both full first normals belong to
$(A_0,B_0,C_0)$ by the same displayed expression.

All retained terms vanish to first order at the exact contact
\[ (u,w,v,N,M_B)=(k,\delta k,\delta k^2E,0,0). \]
Thus the scalar belongs to the square of its contact ideal and the
first normals belong to that ideal. The unchanged monic
contact--conductor identities therefore apply, over the coefficient
ring in which only the displayed positive units have been inverted.
They never divide by $k$, $\ell$, $M$, $J$, $t$, or a contact difference;
the displayed $1/a$ is the inverse of a known positive unit.
The contact substitutions preserve the new weights. After (CA6) is
subtracted, scalar weights are at least seven and normal weights at
least four. Their lifted coefficients consequently have positive
weights, hence zero residues. Clearing the finitely many unit
denominators before this polynomial lift and restoring them afterward
changes no residue sign and introduces no mixed inverse.

\paragraph{Strictness, full normals, and restoration.}
The leading scalar is
\[
 \frac{4(E+4)}{E^2\delta^2}x^2+\frac8{E\delta^2}T_4.
\]
The standard two-pairing circuit supplies a strict Gram on
$(y,z,h_0,q_1,q_2)$. The only nonzero first-normal residue is
$8(E+2)x/(E\delta)$, the coefficient of the full normal $M_B$.
Accordingly the coupled $(M_B,x)$ Gram is
\[
 \begin{pmatrix}
 4(E^2+4)/E&4(E+2)/(E\delta)\\
 4(E+2)/(E\delta)&4(E+4)/(E^2\delta^2)
 \end{pmatrix},\qquad
 \det=\frac{256}{E^3\delta^2}>0.                      \tag{CA8}
\]
Together with the independent $N$ pivot in (CA7), this gives a
bounded strict eight-root Gram for the entire remainder (CA5).
Restore (CA1), all four original positive terms, and all separated
normalization factors. The outer/state lemma exhausts the remaining
tests; real and imaginary state evaluations give the complex case.

This proves precisely the stated unit-$e$, unit-$\eta$,
infinitesimal-$r$ slice. Combined with the bounded-unit-$r$ theorem,
it covers every positive $r$ when $\eta,e$ are bounded units and
$k,\zeta$ infinitesimal. It makes no assertion about unbounded
$\eta$ or either nonunit alternative for $e$.



% Source formulas: pccaa793fd0bf
\subsection{The centered-$A$ source with bounded reciprocal $e$}

\label{component:pccaa793fd0bf}

\begin{theorem}
Use the unchanged physical native C1 source with $k,\zeta,r=1-a\eta$
positive infinitesimals and $\eta$ a positive bounded unit. Suppose
$e=k\sqrt{\eta/\zeta}$ is bounded below by a positive real number;
it may be unbounded. Then the entire original parent is tensor SOS at
every independent projective test and every real or complex state.
Every positive opening rate, literal $\beta=0$, and every rate of
$k^2/r$ is included.
\end{theorem}

Use all exact definitions and the centered positive term (CA1) in
Section~\ref{component:p3763434080c2}. The sole change is that
$\delta=\Delta R/K$ is now any positive bounded coefficient, so it
must not be inverted in a bounded Gram. The source chart (CA3),
$\ell=Qr=K-aE$, and the exact cancelled coquotients remain unchanged.
Their only denominators are positive units, independent of the rate
of $\ell$ relative to $k$.

\paragraph{Scale the entire remainder and its full normal.}
Let $\mathscr G$ be the exact full remainder (CA5). Put
\[
 \widehat G=\delta^2\mathscr G,\qquad
 \widehat M=\delta M_B,\quad \widehat Y=\delta Y,
 \quad\widehat W=\delta W.
\]
The original testing variables have not been rescaled against a
coefficient rate; these displayed quantities are explicit separated
tensor roots. Their complete formulas are
\[
 q_A=vD_uw,\quad q_B=v(D_u-uN-puD_uv w),
\]
\[
 \widehat W=\delta w+u(\widehat M+\delta hq_B)/D_u,
 \quad\widehat Y=\delta uw-(\widehat M+\delta hq_B)/D_u.
\]
Set
\[
 \widehat B_2=\frac{8a\widetilde c_2K^2}{R\widetilde c_3Q^2},
 \qquad
 \widehat B_b=\frac{8a\widetilde c_bK^2}{R\widetilde c_3Q^2}.
\]
The entire exact scaled remainder is
\begin{align*}
\widehat G={}&\widehat B_2\{N^2-2q_AN+(1-p)q_A^2\}
 +\widehat B_bD_v\{N^2+2q_AN+pq_A^2\}\\
&+B_3D_v\{\widehat M^2+2\delta q_B\widehat M+\delta^2hq_B^2\}
 +B_5\{\widehat M^2-2\delta q_B\widehat M+\delta^2(1-h)q_B^2\}\\
&+4A_*D_uD_v(\widehat Y-r u\widehat W)^2
 +4LD_uD_v(2u\widehat W\widehat Y-r u^2\widehat W^2)\\
&+4\alpha D_u(v\widehat Y-\delta)^2
 +2dD_u\{(v\widehat Y-\delta)^2-\delta^2D_v\}.
 \tag{CS1}
\end{align*}
All coefficients are bounded; every occurrence of $\delta$ has
nonnegative exponent. There is no inverse of $\delta$, $\ell$, $k$,
$M$, $J$ or $t$. The inverse $1/a$ in $A_*$ is a positive unit.
Every first-normal and hard-normal term is retained in (CS1).

\paragraph{Uniform contact grading.}
Keep the exact roots $A_0,B_0,C_0,D_0,x,y,z,h_0,q_1,q_2$ and polynomial
$P$ in (CA6). Use the same weights
$\operatorname{wt}(\delta,k,E,\ell,t,u,w,v)=(0,1,0,1,0,1,1,2)$.
Weight zero means bounded here, not bounded away from zero.
The exact finite coefficient collection from (CS1) is
\[
 \begin{aligned}
 (\widehat G|_{N=\widehat M=0})_{\le6}
 &=\frac{4P}{E^4}x^2-\frac{16(k+u)}{E^2}xz+\frac8E T_4,\\
 (\widehat L_N)_{\le3}&=0,\\
 (\widehat L_M)_{\le3}
 &=\frac{8\{E^2+2E-2\ell(E+1)\}}{E^2}x.
 \end{aligned}                                                    \tag{CS2}
\]
The full hard block on $(N,\widehat M)$ has residue
\[
        \operatorname{diag}(16/E,4(E^2+4)/E).          \tag{CS3}
\]
These are precisely the old identities after the indicated scaling,
but their validity for an infinitesimal $\delta$ requires the
nonnegative-exponent assertion in (CS1), not cancellation of a
purported unit $\delta$.

The scalar is in
$(A_0,B_0,C_0)^2+(D_0^2)$ over the ring with $E$ and the displayed
positive units inverted. Both complete first normals are in
$(A_0,B_0,C_0)$. Direct monomial allocation in (CS1) proves these
memberships with nonnegative $\delta$ exponents. At zero normals,
$q_A=A_0D_u$, $\delta q_B=B_0D_u(1-puv w)$, and both
$\widehat Y,u\widehat W$ belong to $(B_0,C_0)$ with bounded-unit
coefficients. The sole constant from the $\alpha$ term satisfies
$\alpha\delta^2/D_0^2=2a\widetilde c_\alpha/\widetilde c_3$ exactly.
Thus these are memberships of the entire rational expression, not merely
its finite collection. The unchanged monic
contact conductor lifts them without a $\delta$ inverse. Its contact
substitutions preserve the weights, even if $\delta$ is infinitesimal.
After (CS2), all remaining scalar weights are at least seven and all
remaining normal weights at least four, hence the lifted correction
coefficients have zero residue. 

The leading soft form is
\[
            \frac{4(E+4)}{E^2}x^2+\frac8E T_4.
\]
Use the strict two-pairing circuit on $T_4$. Its only hard coupling
is the full first normal $8(E+2)x/E$ on $\widehat M$, and
\[
 \det\begin{pmatrix}4(E^2+4)/E&4(E+2)/E\\
                    4(E+2)/E&4(E+4)/E^2\end{pmatrix}
                   =256/E^3>0.                       \tag{CS4}
\]
Together with $16/E$ on $N$, this gives the strict eight-root Gram.

\paragraph{Exhaustion and restoration.}
The independent outer/state lemma Section~\ref{component:p364e14238053}
uses the unit-$e$ hypothesis only to infer a positive lower bound on
$m/d$ in (CAO2). When $e$ is unbounded, $m/d$ is instead unbounded,
so that same lower bound and every clipping argument hold. Its
$c/d,A/d,D/d$ bounds, exact centered root identity, ordinary-$u$
certificate, unit-$w$ certificate, and literal projective endpoint
arguments do not contain $e$ or $\delta$. Thus it reduces every
remaining independent test to exactly the contact chart settled above.
Restore $\delta^2$ times every retained original term, including
(CA1), and undo the positive pure coefficient factor $\delta^2$.
All factors and radicals remain separated between the two fields.
Complex states split into real and imaginary evaluations.

This closes the formerly missing infinitesimal-$r$, unbounded-$e$
slice, and incorporates the previous unit-$e$ theorem. Together with
the common unit-$r$ results, it covers all positive native $r$
when $\eta$ is a unit and $e$ is bounded below. It does not yet
claim the infinitesimal-$e$ slice.



% Source formulas: p364e14238053
\subsection{Complementary centered-$A$ testing and state charts}

\label{component:p364e14238053}

\label{p364e14238053:sec:centered-a-outer}

\begin{lemma}
Use the unchanged physical native C1 source, normalized by $c=2$ and
$\lambda=n=0$, with all original event signs and positive openings
$a,t,M,J$. Suppose $k,\zeta,r=1-a\eta$ are positive infinitesimals and
$\eta,e=k\sqrt{\eta/\zeta}$ are positive bounded units.
Then the outer-test and state-entry arguments reduce the all-tests
problem to precisely
\[
 u,v,w>0\text{ infinitesimal},\qquad Z=1,
 \qquad w=(W+uY)/(1+u^2).
 \tag{CAO1}
\]
All other independent projective tests and states admit a full-parent
tensor SOS. No comparison of $k^2$ with $r$ is made, and arbitrarily
small positive openings and literal $\beta=0$ are included.
\end{lemma}

\paragraph{Source ratios with the common factor cancelled.}
Here $r$ always denotes the native source parameter, not the event root
$uW-Y$. Put
\[
 K=1+k^2,\quad H_c=r+\eta,\quad V_c=k^2\eta\zeta,\quad
 S_3=1-k^2\eta+tV_c,
\]
The two required short native polynomials are
\begin{align*}
P_4={}&1+k^2+\eta k^2(a-k^2-1)-\zeta k^2t\\
 &+V_c[-\zeta k^2t^2+2k^2t+a(1-t)+t],\\
P_d={}&(1+k^2)^2-a\eta k^4
 +\zeta k^2[(k^2+1)-(k^2+2)t]\\
 &+V_c[-\zeta k^2t^2+k^2a(t-1)+k^2t+k^2+1].
\end{align*}
They are positive bounded units with residue one. The exact source
identities, after cancelling the common factor $r$ in the coquotients,
are
\begin{align*}
 d&=\frac{4a\eta P_d}{KH_cS_3},&
 Ar&=\frac{2a\eta^2P_4}{KH_cS_3},\\
 \frac md&=\frac{e^2K(1-t\zeta)(1+V_c)H_c}{2M\eta P_d},&
 \frac cd&=\frac{KH_cS_3}{2a\eta P_d},\\
 \frac Ad&=\frac{\eta P_4}{2rP_d},&
 \frac Dd&=\frac{KH_c(1+V_c)-J\eta P_d}{2J\eta P_d}.
 \tag{CAO2}
\end{align*}
The positive openings obey
$M+t=1+k^2r$ and $J+a=1+k^2-k^2\zeta t$; hence $a,t,M,J$ are
bounded. Since $a\eta=1-r$, one has
\[
 \bar\eta\ge1,\qquad \bar d=4/\bar\eta,\qquad Ar\sim2,
 \qquad A\text{ positive unbounded}.
 \tag{CAO3}
\]
The numerator of $D/d$ in (CAO2) has residue
$\bar H_c-\bar J\bar\eta=1$. Therefore $m/d,c/d,D/d,A/d$ all have
positive real lower bounds. This conclusion uses no inverse bound on
an opening. Moreover,
\[
 \begin{cases}
 \displaystyle \bar D=\frac{2}{\bar\eta(\bar\eta-1)},
       &\bar\eta>1,\\[4pt]
 D\text{ positive unbounded},&\bar\eta=1.
 \end{cases}
 \tag{CAO4}
\]

\paragraph{The previously proved outer routing still applies.}
Clip $m,c,A,D$ at positive real multiples of $d$, retaining all their
nonnegative excess original event squares. The resulting coefficients
divided by $d$ are bounded. The outer-test argument in
Section~\ref{component:p69995d829bb2}, through (BE2)--(BE4), uses only
these lower bounds and the independent physical-source certificate at
noncommuting residue tests. Its exact homogeneous identities and
four-root ideal memberships are unchanged. It therefore includes all
projective infinity charts, the $(0,\infty)$, $(\infty,0)$ and
$(\infty,\infty)$ corners, and the unit-$W$, infinitesimal-$Z$ state
chart. The original-event SOS on literal zero axes is unchanged as
well. This invokes none of the bounded-$r$ scalar or normal conclusions
(BE5)--(BE6). Only bounded $u$ and infinitesimal $v$ remain.

\paragraph{An exact centered root identity.}
Write $D_u=1+u^2$, $D_v=1+v^2$ and introduce
\[
 h=A^{-1/2},\quad L=Ar,\qquad
 H=\sqrt A(Y-ruW),\quad B=X+uZ,\quad U=uW,\quad V=vZ.
 \tag{CAO5}
\]
Thus $Y=rU+hH$, while $h,r$ are infinitesimal and $L\sim2$.
All coefficient radicals in (CAO5) lie in the coefficient field;
$B,U,V$ lie in the testing field. The root $H$ is their finite tensor
combination, with no mixed inverse or mixed radical.

For clarity, the complete original finite parent used here is
\begin{align*}
F={}&4m(vW+uvY-X-uZ)^2+8D_v(uW-Y)^2
       +4AD_uD_vY^2\\
 &+(4D+2d)(uW+uvX-Y-vZ)^2
   +(4\alpha+2d)D_u(vY-Z)^2\\
 &+(4\beta+2d)D_v(X+uZ)^2-2dD_uD_vZ^2.
\end{align*}

Choose $0<D_*\le D$ bounded, and let $m_*=0$ or choose a positive
real $m_*\le m$, as specified below. Retain the four nonnegative
original event contributions with weights $m-m_*,D-D_*,\alpha,\beta$.
For $C=4D_*+2d$, the exact remaining parent is
\begin{align*}
F_*={}&8D_v[(1-r)U-hH]^2
 +4D_uD_vH^2+4LrD_uD_vU^2+8LhD_uD_vUH\\
&+C[(1-r)U-hH+uvB-D_uV]^2+2dD_vB^2\\
&+2dD_u\{v^2(rU+hH)^2-2(rU+hH)V-V^2\}
 +4m_*R_m^2,                                      \tag{CAO6}\\
R_m={}&vW+uv(rU+hH)-B.
\end{align*}
In particular, the centered $A$ contribution is the exact identity
\[
 4AD_uD_vY^2
 =4D_uD_vH^2+4LrD_uD_vU^2+8LhD_uD_vUH.
\]
This retains its complete hard root. The last two coefficients have
zero residue. Formula (CAO6) follows directly from the original parent
and includes the complete signed core; it is not an eliminated scalar
or a restriction to state normals. Every displayed coefficient is
bounded. The term involving $Z$ in that core has already cancelled
to the displayed expression in $V=vZ$.

\paragraph{Every state when $u$ is a bounded unit.}
Set $m_*=0$, so the entire original $m$ event is retained. Formula
(CAO6) is an exact Gram allocation on $(H,B,U,V)$. At infinitesimal
$v$ its residue is
\[
 4D_uH^2+2dB^2+8U^2+C(U-D_uV)^2-2dD_uV^2.
 \tag{CAO7}
\]
Residues are understood in this and the following determinant
calculations. The $U,V$ determinant is
\[
 D_u\{8CD_u-2d(8+C)\}.
\]
If $\bar\eta>1$, take $D_*=D$, which is bounded by (CAO4).
Then $C=8/(\bar\eta-1)$ in residue and this determinant is
\[
 \frac{64u^2(1+u^2)}{\bar\eta-1}>0.
\]
If $\bar\eta=1$, choose a sufficiently large positive real $D_*$;
this is permitted because $D$ is unbounded. The determinant is then
\[
 32(1+u^2)(D_*u^2+2u^2-2)>0.
\]
The first diagonal is $8+C>0$, so the complete four-root residue
Gram is positive definite in both cases. The cap is a real constant
chosen from the nonzero testing residue of $u$; no division of tensor
elements is introduced. The bounded full Gram is therefore tensor
SOS, and restoring the retained events proves the claim for every
state in this chart.

\paragraph{The remaining unit-$w$ state chart.}
Suppose $u,v$ are infinitesimal and $w$ is a unit. Then $W$ is a unit
on the negative-core branch. When $Z$ is infinitesimal the inherited
(BE4) argument already applies. Otherwise normalize the state so that
$Z$ is a unit and put $q=W/Z$, a pure-testing unit. Choose a positive
real lower bound $m_*\le m$. Now
\[
 R_m=qV+uv(rU+hH)-B,
\]
so (CAO6) is again an exact bounded Gram on $(H,B,U,V)$, with residue
\[
 4H^2+2dB^2+4m_*(qV-B)^2
       +8U^2+C(U-V)^2-2dV^2.                    \tag{CAO8}
\]
The $B$ pivot is positive. Its Schur complement adds
\[
 \mu V^2,\qquad \mu=\frac{4dm_*q^2}{d+2m_*}>0
\]
to the $U,V$ block. This is just a residue determinant computation;
no mixed inverse is used to define a root or a full-parent coefficient.
If $\bar\eta>1$, take $D_*=D$ as before. The determinant becomes
\[
 \frac{8\bar\eta\mu}{\bar\eta-1}>0.
\]
If $\bar\eta=1$, choose a sufficiently large positive real cap
$D_*$ so that $\mu(C+8)-64>0$. The unboundedness of $D$ permits this
choice. Thus the complete residue Gram in (CAO8) is strict, including
all hard-soft terms, and restoring the retained events gives the
full-parent tensor SOS.

\paragraph{Completion of state entry.}
Normalize the state by a maximal coordinate and first remove the
nonnegative pure-testing-core branch, which has its direct original
SOS. On the negative-core branch with $v$ infinitesimal, the residue
core is
\[
 2\{(X+uZ)^2+(uW-Y)^2\}.
\]
Consequently these two roots have zero residue. If $w$ and $Z$ were
both infinitesimal, the exact inverse relations
\[
 W=w+\frac{u(uW-Y)}{D_u},\quad
 Y=uw-\frac{uW-Y}{D_u},\quad X=(X+uZ)-uZ
\]
would make every normalized state coordinate infinitesimal. Hence
small $w$ forces $Z$ to be a unit. The bounded-unit $u$ chart was
settled by (CAO7), and all unit-$w$ states by (CAO8) and (BE4).
Normalize the only remaining states by $Z=1$. The same exact
pure-testing Schur identity (SD5) and testing sign gauges give
$u,v,w>0$. This is exactly (CAO1). Literal axes and literal state
zeros have already been included. Complex states split into their
real and imaginary evaluations because the parent is real.



% Source formulas: p890eea5f538b
\subsection{Unit $\eta,r$ and unbounded $e$}

\label{component:p890eea5f538b}

\begin{theorem}[Unit $\eta$ and unbounded $e$]
Use the unchanged physical native C1 source, its original event signs,
and the positive shapes and gaps $a,t,M,J$. Suppose $k,\zeta$ are
positive infinitesimals, $\eta$ is a positive bounded unit,
$r=1-a\eta$ is a positive bounded unit, and $e=k\sqrt{\eta/\zeta}$
is unbounded. Then the full original parent is tensor SOS at every
independent projective test and every real or complex state.
No positive lower bound on $a,t,M,J$ is imposed.
\end{theorem}

Write $H=r+\eta$ and $\Delta=1/e$. The seven native coquotients
have residues $2rH$ for $c_2,c_b,c_3,c_5,c_\alpha$, $2$ for
$c_4$, and $4r$ for $c_d$. All are positive units. In particular the
four retained terms of the contact--conductor proof remain positive
and the remainder has bounded coefficients. We reuse their exact
definitions, not asymptotic substitutes for the original source.

\paragraph{Outer tests.}
The exact formulas give
\[
 \frac md=\frac{e^2c_2}{M\eta c_d},\qquad
 \frac\beta d=\frac{e^2c_b}{t\eta c_d}-\frac12.
\]
Both ratios are unbounded. Apply the exact full-parent Alice reflection
(AS3)--(AS4), which exchanges $m\leftrightarrow c$ and
$\beta\leftrightarrow D$ while leaving $v$ unchanged. The already
proved two-strong-current identity then covers every test with $1/v$
bounded, for every Alice angle and every state. Its homogeneous form
includes both literal Alice endpoints, so no inverse-angle point is lost.

For $v$ infinitesimal and Alice infinity, use the critical-current
identity with $g_c=cD/(c+D)$, $\epsilon=d/g_c$ and
$\gamma=Ag_c/d^2$. Its actual sufficient conditions on this chart
are bounded $\epsilon$ and positive-unit $\gamma$: in its exact
four-root identity the correction $L_0$ has factors $c_u$ or $v$,
where $c_u=(1+u^2)^{-1/2}$ is infinitesimal. Thus $\sqrt\epsilon$
need only be bounded. The residue Gram is
$\operatorname{diag}(1,2,2,4\bar\gamma)$.
Use the exact relation $J+a=T$ with $T-1$ infinitesimal, and the
cancelled original coefficient formula
\[
 \frac Da=\frac{2c_5/c_3-J\eta c_d/c_3}{J}.
\]
Its numerator has residue $2/\bar H$. The unchanged native formulas give
\[
 \frac d a\sim\frac{4\eta}{H},\quad
 \frac A a\sim\frac{2\eta^2}{rH},\quad
 \frac D a\sim\frac2{JH},\qquad
 \epsilon\sim\frac{2\eta(JH+a)}H,\quad
 \gamma\sim\frac{H}{4r(JH+a)}.
\]
The identity
$\overline{JH+a}=1+\bar\eta(1-\bar a)^2$ shows that the two required
conditions hold even when $a$ or $J$ is infinitesimal. The calculation
uses the cancelled harmonic coefficient $g_c/a$, so no lower bound
on an opening is inserted. Homogeneous angle normalization includes
literal Alice infinity. It remains to treat bounded $u$ and
infinitesimal $v$.

\paragraph{State entry.}
Normalize the state by a coordinate of maximal absolute value and put
$D_u=1+u^2$ and $w=(W+uY)/D_u$. If $w$ is a unit, the exact five-root
current identity (UW2)--(UW3) applies with $Q=W+uY=D_uw$ a unit.
Its first four diagonal residues are $8rH$ and its fifth is
$8D_u$; both negative corrections have factors $\Delta$ or
$\Delta^2\eta$. Thus this chart is strict.

If $w$ is infinitesimal, the nonnegative pure-testing core is a direct
exit. On the negative-core branch $Z$ is a unit. Indeed its residue
at $v=0$ is $2\{(X+uZ)^2+(uW-Y)^2\}$; if $Z,w$ and both currents
had zero residue, every normalized state coordinate would have zero
residue. Normalize $Z=1$. The exact core Schur identity gives $w>0$
after the usual positive-angle sign gauges. All transformations here
belong solely to the testing field.

\paragraph{The ordinary-$u$ chart needs no contact division.}
Use the exact source change
\[
 \delta=\Delta R/K,\quad b^2=\eta K/R,\quad K=1+k^2,
 \quad Q=K+ak^2b^2,\quad T_n=K^2-t\delta^2k^4Qb^2.
\]
All source denominators are positive units, $\delta,k$ are
infinitesimal and $b^2\sim\eta$. Let
\[
 A_0=wv,\quad B_0=\delta v,\quad C_0=\delta b^2uw,
 \quad D_0=\delta^2k^3b^2,\qquad x=B_0-C_0.
\]
The five-term exact scalar belongs to
$(A_0,B_0,C_0)^2+(D_0^2)$ and both complete first normals belong
to $(A_0,B_0,C_0)$. Collecting these displayed factors before taking
residues is legitimate over the bounded coefficient ring. At
$\delta=k=w=v=0$, keeping bounded $u$ free, the scalar Gram is
\[
 8rH D_u^2A_0^2+
 8D_u\{x^2+u^2rH B_0^2\}+8rH D_uD_0^2.
 \tag{UE1}
\]
The hard block on $(N,\delta M_B)$ is
$\operatorname{diag}(16rH,16rH+8\eta^2/D_u)$, and the complete
first normals have residues $0$ and $16\eta x$. The latter is the
coefficient of $\delta M_B$ in the full polynomial, so the symmetric
Gram off-diagonal entry between $\delta M_B$ and $x$ is $8\eta$.
The coupled $(\delta M_B,x)$ block has determinant
\[
 (16rH+8\eta^2/D_u)\,8D_u-(8\eta)^2=128D_urH>0.
 \tag{UE2}
\]
If $u$ is a unit, the remaining $B_0$ pivot in (UE1) is strict,
as are the $A_0,D_0,N$ pivots. Every omitted coefficient has zero
residue. This proves the chart by a bounded strict Gram on these
six roots, without any mixed inverse. Equation (UE1) follows directly
from $a\eta+r=1$: its $(B_0,C_0)$ part before collection is
$8D_u\{(1+u^2rH)B_0^2-2B_0C_0+C_0^2\}$.

\paragraph{The infinitesimal-$u$ chart uses the old conductor.}
Now $\delta,k,u,w,v$ are infinitesimal, while $b,a,t$ are bounded.
Use the unchanged roots
\[
 y=A_0-kC_0,\quad z=uB_0-kC_0,\quad h=kC_0-D_0,
 \quad q_1=a_1a_2-a_3a_4,
\]
and the second pairing $q_2=a_1a_3-a_2a_4$ of the four separated
radicals. Set
\[
 T_4=\sum_{i=1}^4a_i^4-4\prod_{i=1}^4a_i
     =y^2+z^2+h^2-2yz+2q_1^2,
 \qquad \rho=(1-ab^2)(1+(1-a)b^2)-1.
\]
Give $(\delta,k,b,a,t,u,w,v)$ weights $(1,1,0,0,0,1,2,3)$.
The root $x$ has weight four and the other soft roots have weight
five. Substitution at the exact contact preserves this grading.
Writing $E_0$ for the complete scalar correction relative to the
old leading polynomial, the five-term factorization and the two
complete normal formulas give
\[
 \begin{aligned}
 E_0&=8\rho T_4-8(\rho+b^2)k^2x^2+O_{\rm wt}(12),\\
 L_1&=O_{\rm wt}(6),\qquad
 L_2=16b^2x+O_{\rm wt}(6).
 \end{aligned}\tag{UE3}
\]
These identities require only multiplication of the short source
polynomials. In particular the residue correction is the single
factor $\rho$ multiplying the same circuit, rather than a new
polynomial positivity problem.

The unlocalized monic contact--conductor lift applies unchanged.
After the terms explicitly displayed in (UE3) are removed, scalar
coefficient weights are at least two and normal coefficient weights
are at least one. Therefore they all have zero residue. There is no
division by $k,b,\delta,a,t$ or a contact difference. The leading
strict soft Gram is now
\[
                  8x^2+8rH T_4.
\]
The two-pairing circuit identity supplies a strict Gram on
$(y,z,h,q_1,q_2)$, and $rH$ is a positive unit. The complete hard
block is $\operatorname{diag}(16rH,16rH+8\eta^2)$; its only
nonzero first-normal residue is $16\eta x$. Its coupling to the
$8x^2$ pivot has determinant $128rH>0$. Retain the old positive
$g_4^2$ square. The resulting eight-root remainder Gram is strict.

Finally restore all four positive terms, the positive square just
retained, and the positive pure-coefficient multiplier. The state and
angle cases above exhaust the independent tests. Complex states split
into their real and imaginary quadratic evaluations. The result
closes the entire stated source domain, rather than only an optimizing
or matching testing chart.



% Source formulas: pa7e744614f49
\subsection{Unit $\eta$ and infinitesimal $e$: local polynomial}

\label{component:pa7e744614f49}

\begin{proposition}[Local part of the small-$e$ theorem]
Use the unchanged physical native C1 source, with $k,\zeta,e$ positive
infinitesimals and $\eta$ a positive bounded unit. Every native positive
$r=1-a\eta$, and every positive opening rate, is allowed.
At every independent test with $u,v,w$ infinitesimal, $Z=1$, and $w>0$,
the entire original parent is tensor SOS. Literal $\beta=0$ is included.
\end{proposition}
This proposition is the contact part only. The independent outer/state
join is supplied separately; no assertion here identifies a testing rate
with a source rate.

\paragraph{The source entry forced by the original sign.}
Keep $K,Q,E,\ell,\delta$ from (CA3), with $E=b^2$ and
$\ell=Qr=K-aE$. Here $E$ is a positive unit, $0<\ell<K$, and no lower
bound on $a$ or $\ell$ is assumed. Put
\[
 \varepsilon=1/\delta,\qquad j=\delta k,\qquad t=\varepsilon^2\tau.
\]
Both $\varepsilon,j$ are infinitesimal: $e=\varepsilon K/Q$ and
$\zeta=j^2QE/K$. The original condition $\beta\ge0$ gives
\[
 \frac{t}{e^2}\le\frac{2c_b}{\eta c_d},
\]
whose right side is bounded after cancelling the common native $r$
factor. Thus $\tau$ is positive bounded. The exact chart is
\[
 k=\varepsilon j,\quad K=1+\varepsilon^2j^2,\quad
 Q=K^2-\ell\varepsilon^2j^2,\quad a=(K-\ell)/E,
\]
\begin{align*}
 \eta&=KE/Q,\quad r=\ell/Q,\quad H=r+\eta,\quad
 e^2=\varepsilon^2K^2/Q^2,\\
 \zeta&=j^2QE/K,\quad R=K^2/Q,\quad T=K-k^2t\zeta. \tag{SE1}
\end{align*}
All denominators displayed here are positive units. The cancelled
coquotients $\widetilde c_i=c_i/\ell$ are the same short exact formulas
as in (CA3); the cancellation introduces no inverse of $\ell$.
Their residues are $2(E+\ell)$ for $i=2,b,3,5,\alpha$, and $4$ for
$i=d$; $c_4$ has residue $2$. These formulas hold uniformly as
$\ell$ tends to zero or to one.

\paragraph{One bounded centered remainder, including both normals.}
Set
\[
 A_* =\frac{ArH}{1+rH},\qquad
 A-A_* =\frac{A}{1+rH}>0.                              \tag{SE2}
\]
Retain the entire square $4(A-A_*)D_uD_v(Y-r uW)^2$.
It and its full first derivative vanish at exact native contact.
Divide the remaining parent by the positive pure coefficient $a$;
every factor $a$ cancels before coefficients are used. Write
\[
 \begin{array}{ll}
 b_2=8\widetilde c_2e^2/(R\widetilde c_3),&
 b_b=8\widetilde c_be^2/(R\widetilde c_3),\\
 b_3=8/T,&b_5=8\widetilde c_5/(T\widetilde c_3),\\
 L=Ar/a=2\eta^2c_4/(Q\widetilde c_3),&
 d_0=d/a=2\eta\widetilde c_d/\widetilde c_3,\\
 A_0=A_*/a=LH/(1+rH),&
 \alpha_0=\alpha/a=2\varepsilon^4j^6Q^2\eta^2
       \widetilde c_\alpha/(K^2\widetilde c_3).
 \end{array}                                                     \tag{SE3}
\]
Every quantity in (SE3) is bounded with unit denominators, and
$b_2,b_b$ contain $\varepsilon^2$.
Use the original normals $N,M_B$, put $p=t/R$, $h=a/T$, and write
\begin{align*}
 q_A&=vD_uw,\qquad q_B=v(D_u-uN-puD_uv w),\\
 W&=w+u(M_B+hq_B)/D_u,\qquad Y=uw-(M_B+hq_B)/D_u.
\end{align*}
The entire bounded remainder is
\begin{align*}
G_s={}&b_2\{N^2-2q_AN+(1-p)q_A^2\}
 +b_bD_v\{N^2+2q_AN+pq_A^2\}\\
 &+b_3D_v\{M_B^2+2q_BM_B+hq_B^2\}
 +b_5\{M_B^2-2q_BM_B+(1-h)q_B^2\}\\
 &+4A_0D_uD_v(Y-r uW)^2
 +4LD_uD_v(2uWY-r u^2W^2)\\
 &+4\alpha_0D_u(vY-1)^2
 +2d_0D_u\{(vY-1)^2-D_v\}
 +\nu D_vN^2,\qquad \nu=d_0(R-t)/R.                 \tag{SE4}
\end{align*}
This is (SD4), after the indicated exact centered subtraction and
normalization, except that a positive part of the original $\Pi_2$ is
kept inside the remainder. Indeed the original coefficient of
$D_vN^2$ in $\Pi_2/a$ is
\[
 \frac{4\beta+2d}{a}\frac{R-t}{R}\ \ge\ 2\nu.
\]
Thus subtracting only $\Pi_2/a-\nu D_vN^2$ remains legitimate, even
when $\beta=0$. The other three original terms are retained unchanged.
All full first normals and every hard quadratic term are present in
(SE4); there are no terms of higher normal degree.

\paragraph{A conductor that never divides by $\varepsilon$.}
Define the separated monomials and differences
\[
 A=\varepsilon wv,\quad B=v,\quad C=Euw,\quad
 D_0=\varepsilon^2j^3E,\quad k=\varepsilon j,
\]
\[
 x=B-C,\quad y=A-kC,\quad z=uB-kC,\quad h_0=kC-D_0.
\]
The two standard radical pairings are obtained by dividing the original
four squared circuit monomials by $\delta$: they are
$A,uB,kC,D_0$, and their product is $(k^2BC)^2$.
Their strict circuit is
$T_4=y^2+z^2+h_0^2-2yz+2q_1^2$.
The exact contact is $(u,w,v)=(k,j,\varepsilon j^2E)$.
Besides the three basic conductor identities,
\[
 C(u-k)=z-ux,\quad
 \varepsilon C(w-j)=y-\varepsilon wx,\quad
 C(v-\varepsilon j^2E)=E(wz+jy-kwx),                \tag{SE5}
\]
use the additional, division-free identity
\[
 C v(w-j)=Ew(w-j)z+Ejw y-Ekw^2x.                    \tag{SE6}
\]
For a monic contact division
$S=(u-k)T_u+(w-j)T_w+(v-\varepsilon j^2E)T_v$,
whenever $T_w=\varepsilon U+vV$, the bounded conductor is
\[
 \begin{aligned}
 CS={}&x[-uT_u-\varepsilon wU-Ekw^2V-EkwT_v]\\
 &+y[U+EjwV+EjT_v]
 +z[T_u+Ew(w-j)V+EwT_v].
 \end{aligned}                                      \tag{SE7}
\]
No division by $\varepsilon,j,k,r,a,t$, or a contact difference occurs.

Here is the full rational membership argument required for applying
(SE7), not only a finite truncation. At zero normals let
$V_0=1-puv w$. Then
\[
 q_A=D_u wv,\quad q_B=D_uvV_0,\quad
 uW=C/E+u^2hvV_0,\quad Y=C/E-hvV_0.
\]
The scalar of (SE4) therefore has an exact raw allocation
\[
 f_AA^2+f_{BB}B^2+f_{BC}BC+f_{CC}C^2+\rho D_0^2.     \tag{SE8}
\]
Expand only the nonconstant part of $(vY-1)^2$ into $B,C$; its constant
has $\alpha_0/D_0^2=2\widetilde c_\alpha/\widetilde c_3$.
The coefficients in (SE8) have unit denominators independent of $w$.
Every dependence on $w$ in them comes through $V_0$, and hence has a
factor $\varepsilon^2$; modulo $\varepsilon$ they are independent of
$w$ and even in $u,v$. Both complete first normals are in $(A,B,C)$
with bounded coefficients by the same exact formulas.

Clear the finitely many unit denominators, chosen even in $u,v$ and
independent of $w$, and use the unchanged scalar square-lift algorithm.
For its first scalar conductor, the only $w$-dependent coefficient
$T_w$ is divisible by $\varepsilon$ (in fact by $\varepsilon^2$).
At $\varepsilon=u=0$, its residual polynomial $S$ is even in $v$,
and contact makes its constant zero; thus $S\in(v^2)$ there.
More explicitly, modulo $\varepsilon$ the first divided differences
satisfy $T_u\in(u)$, $T_w=0$, and $T_v\in(v)$ after the preceding
monic substitutions. It follows
that after this first conductor, the sole possible $w$ dependence
modulo $\varepsilon,u$ in the four second-stage remainders is of the
form $w v$ times a polynomial in $v$. Hence their $T_w$ all belong to
$(\varepsilon,v)$, exactly the hypothesis of (SE7). The complete
normal raw coefficients are independent of $w$ modulo $\varepsilon$,
so their one conductor has $T_w\in(\varepsilon)$. The retained contact
squares and the original native balances supply the exact contact
value and first-jet identities throughout. This proves the full
scalar-square and complete first-normal lifts over the bounded ring.

\paragraph{Finite residues of that full lift.}
Give $(\varepsilon,j,u,w,v,N,M_B)$ weights $(1,1,2,1,3,5,3)$;
$E,\ell,\tau$ have weight zero. Every positive-weight variable is
infinitesimal, while all weight-zero coefficients are bounded.
The contact substitutions preserve the grading; $x$ has weight three
and the other circuit roots weight five. Both (SE5) and (SE6) preserve
weights, so after the scalar collection through weight ten and full
first normals through weight five, every lifted remainder coefficient
has positive weight. Restore the unit denominators afterward.

Set $H_0=E+\ell$ and $V_*=1+\ell H_0$. Direct exact multiplication of
(SE1)--(SE4), followed by (SE7), gives scalar residue
\[
       \frac{8(E+2)}{H_0V_*}x^2+8T_4,                         \tag{SE9}
\]
and complete first-normal residues
\[
       L_N=0,\qquad L_M=\frac{16E(H_0+1)}{H_0V_*}\,x.
\]
The hard block has residue
\[
 \operatorname{diag}\left(\frac{4E}{H_0},
           \frac{8(E^2+2E\ell+2\ell^2+2)}{V_*}\right).       \tag{SE10}
\]
These statements include coefficient corrections of positive weight;
they are not asserted as literal equalities of the full rational scalar.

The only hard-soft coupled determinant is
\[
 \det G_{M,x}
 =\frac{64(E^2+2E\ell+4E+4\ell)}{H_0^2V_*}>0.              \tag{SE11}
\]
Together with $4E/H_0>0$ and the strict two-pairing circuit this gives
the full bounded strict eight-root Gram. All original positive terms,
the centered square, and the positive pure-coefficient factor $a$ are
then restored. All inverses and radicals remain separated between the
two independent real closed fields.



% Source formulas: p05a86be0046c
\subsection{Unit $\eta$ and infinitesimal $e$: complementary tests}

\label{component:p05a86be0046c}

\label{p05a86be0046c:sec:small-e-outer}

\begin{lemma}
Use the unchanged physical native C1 source with $c=2$,
$\lambda=n=0$, positive original openings, and nonnegative original
event coefficients. Suppose $k,\zeta,e=k\sqrt{\eta/\zeta}$ are positive
infinitesimals, $\eta$ is a positive bounded unit, and
$r=1-a\eta>0$. Then every independent projective test and state is
tensor SOS except possibly the chart
\[
 u,v,w>0\text{ infinitesimal},\qquad Z=1,
 \qquad w=(W+uY)/(1+u^2).
 \tag{SEO1}
\]
There is no restriction on the positive relative scales of $a,r,J$
or on $k^2/r$. This lemma supplies a join, not a certificate for the
remaining chart (SEO1).
\end{lemma}

\paragraph{The small weak-event scale.}
Put $\kappa=k^2$, $K=1+\kappa$, $H=r+\eta$,
$V_c=\kappa\eta\zeta$, $S=1-\kappa\eta+tV_c$,
$R=1+\kappa r$, $T=K-\kappa\zeta t$,
$M=R-t$, and $J=T-a$. Let $P_4,P_d$ be the short native
polynomials printed in the centered-A outer lemma. All of
$K,H,S,P_4,P_d$ are positive bounded units. Set
\[
 L=KHS,\quad N=KH(1+V_c)-J\eta P_d,\quad
 d_0=4\eta P_d,\quad A_0=2\eta^2P_4,\quad D_0=2N.
\]
After exact cancellation of the native $r$ factors, the source is
\[
 d=\frac{ad_0}{L},\qquad A=\frac{aA_0}{rL},\qquad
 D=\frac{aD_0}{JL},\qquad
 \frac{m}{ae^2}=
 \frac{2(1-t\zeta)(1+V_c)}{MS}.                 \tag{SEO2}
\]
Here $\bar N=1$. Native $\beta\ge0$ gives the exact bound
\[
 \frac{t}{e^2}\le\frac{2c_b}{\eta c_d}
 =\frac{K[1+(1-t)\zeta]H}{\eta P_d}.
\]
Thus $t$ is infinitesimal, $M\sim1$, and $m/(ae^2)\sim2$.
In particular
\[
 \frac{k^2}{e^2}=\frac\zeta\eta
 \quad\hbox{is infinitesimal}.                       \tag{SEO3}
\]
The ratios $d/a,A/a,D/a$ have positive real lower bounds, while
$d/a$ is a bounded unit. These statements allow arbitrary positive
opening rates and literal $\beta=0$.

\paragraph{Choose the critical parameter without a small denominator.}
There are two source cases. If $a$ is a bounded unit, keep $A,D,d$
and replace $c=2$ by the critical number
\[
 c_* =\frac{d(d+2A)(d+2D)}{2(4AD-d^2)}.
 \tag{SEO4}
\]
This is positive and bounded, and $c_*-2=O(k^2)$.
Here is an exact short verification of the potentially delicate bound.
Define
\[
 Q=4A_0D_0-rJd_0^2,
\]
so $\bar Q=16\bar\eta^2\bar a\bar H>0$. Then
\[
 c_*-2=
 \frac{ad_0(d_0r+2A_0)(d_0J+2D_0)-4QL}{2QL}.
 \tag{SEO5}
\]
After $r=1-a\eta$ is inserted, its polynomial numerator vanishes
identically at $\kappa=0$, with $\zeta$ still free. It is therefore
$\kappa$ times a polynomial in bounded native variables; its
denominator is a positive unit. No inverse of $r$ or $J$ occurs.

If $a$ is infinitesimal, then $r,J\sim1$. Divide the whole parent
by $a$ and replace the original $A$ by
\[
 A_* =\frac{d^2(4+d+2D)}{2(8D-2dD-d^2)}.       \tag{SEO6}
\]
The denominator is positive. In this case
$A_*/a,D/a,d/a$ are positive bounded units, $c/a=2/a$ is
unbounded, and $(A_*-A)/a=O(k^2)$. To verify this without a
rate comparison, put
\[
 B_0=8D_0L-a(2d_0D_0+d_0^2J).
\]
Its residue is $16\bar H>0$, and
\[
 \frac{A_*-A}{a}=
 \frac{r d_0^2\{4JL+a(d_0J+2D_0)\}-2A_0B_0}
 {2rLB_0}.                                           \tag{SEO7}
\]
Again the numerator vanishes identically at $\kappa=0$, so it has
an exact polynomial factor $\kappa$, with a unit denominator.

Use $\sigma=1$ in the first case and $\sigma=a$ in the second,
and henceforth normalize all parent coefficients by $\sigma$.
Let $\epsilon$ be the positive part of the critical increase in the
chosen normalized coefficient. From (SEO2)--(SEO7),
\[
 \epsilon\text{ is infinitesimal},\qquad
 \frac{\epsilon}{m/\sigma}=O(\zeta)\text{ is infinitesimal}.
 \tag{SEO8}
\]
If the chosen critical value is smaller than the actual coefficient,
retain the positive excess term instead. In either source case
the chosen critical parent obeys
\[
 2c(4AD-d^2)=d(d+2A)(d+2D),                  \tag{SEO9}
\]
with the letters in (SEO9) now denoting the normalized critical
coefficients. In the first case $c,d$ are units and $A,D$ are
bounded below by positive units; in the second $A,D,d$ are units
and $c$ may be unbounded.

\paragraph{A critical-wall frame permitting the large coefficients.}
Write $D_u=1+u^2$, $D_v=1+v^2$ and
\[
 B=X+uZ,\quad U=uW,\quad V=vZ,\quad J_t=U-Y,
 \quad C=4D+2d,
\]
and define the pure-coefficient numbers
\[
 y=\frac d{2A},\quad
 b=\frac{4AD-d^2}{AC},\quad s=1-b,\quad q=y+b.
\]
The critical identity gives the useful exact relations
\[
 b=\frac C{4c+C},\quad Cs=4cb,\quad
 cb=\frac{d(2A+d)}{4A}.                         \tag{SEO10}
\]
In both source cases, $y,b,s,q$ are bounded and $b,s>0$.
The entire signed core plus the three critical terms is exactly
\begin{align*}
 F_*={}&4cD_vJ_t^2+4AD_uD_vY^2
       +C(J_t+uvB-D_uV)^2+2dD_vB^2\\
       &+2dD_u\{v^2Y^2-2YV-V^2\}.              \tag{SEO11}
\end{align*}

In the first source case use
\[
 H_A=\sqrt A(Y-yV),\qquad
 H_D=\sqrt C(J_t+uvB-(b+u^2)V).
 \tag{SEO12}
\]
The inverse is $Y=yV+H_A/\sqrt A$ and
$J_t=(b+u^2)V-uvB+H_D/\sqrt C$.
All coefficients in (SEO11) after this substitution are bounded.
Indeed the large $A,C$ occur only in the normalized squares and in
the bounded combinations $Ay^2$, $\sqrt A\,y$, $Cs^2$,
$\sqrt C\,s$; use $Cs=4cb$ and bounded $c$ for the last two.
At $u=v=0$, the hard block on $(B,H_A,H_D)$ is
\[
 \operatorname{diag}(2d,4,1+4c/C)\succ0.
\]
On its exact zero section the scalar coefficient of $V^2$ is
\[
 g u^2+\gamma v^2+4cu^4+g(y+2)u^2v^2+4cu^4v^2,
 \quad g=\frac{d(2A+d)}A,\quad
 \gamma=\frac{Dd(2A+d)^2}{A^2(2D+d)}.             \tag{SEO13}
\]
Both $g,\gamma$ are positive bounded units. Every hard-$V$ entry
has angular degree at least two, with bounded coefficients; this is
an exact full-parent polynomial statement. Allocate it on $uV,vV$,
leaving a factor $u$ or $v$ in its coefficient. The resulting exact
five-root Gram on $(B,H_A,H_D,uV,vV)$ therefore has strict residue.

In the second source case set $L_c=4c+C$ and instead use
\[
 H_A=\sqrt A(Y-yV),\qquad
 H_c=\sqrt{L_c}(J_t-bV).                         \tag{SEO14}
\]
Its coefficients are bounded because $A,D,d$ are units and
$c/L_c$, $cb/\sqrt{L_c}$ and $cb^2$ are bounded by (SEO10).
Its hard block is $\operatorname{diag}(2d,4,1)$.
The exact zero-section scalar is
\[
 g u^2+\gamma v^2+C u^4+gy\,u^2v^2.
 \tag{SEO15}
\]
The same complete hard-$V$ allocation gives a strict five-root
Gram on $(B,H_A,H_c,uV,vV)$. Thus this is a full critical-wall
certificate in both source cases, even with the stated unbounded
coefficients. Every radical and denominator used in the roots is
pure coefficient; no mixed inverse is used.

\paragraph{All unit-$W$ states at the small-angle corner.}
Normalize $W=1$, with $Z,Y,X$ bounded, and take $u,v$ infinitesimal.
Retain the complete original $\alpha,\beta$ terms. Put
\[
 R_m=v-X+uvY-uZ=v(1+uY)-B,\qquad
 V=\frac{Z}{1+uY}(R_m+B).                        \tag{SEO16}
\]
The displayed denominator is a positive pure-testing unit.
In the first source case the only negative correction to the
critical parent is $-4\epsilon D_vJ_t^2$; (SEO12) and (SEO16)
express $\sqrt\epsilon J_t$ as a bounded combination of the five
critical roots and $\sqrt{m/\sigma}\,R_m$, with every coefficient
having zero residue. In the second case the only negative
correction is $-4\epsilon D_uD_vY^2$; use
$Y=yV+H_A/\sqrt A$ and (SEO16) for the identical conclusion.
The only possibly delicate coefficient is a bounded testing
multiple of $\sqrt{\epsilon/(m/\sigma)}$, and it is infinitesimal
by (SEO8). The extra original $m$ square contributes the strict
diagonal $4$. Hence the complete perturbed Gram remains strict.
This includes unit or infinitesimal $Z$, literal $Z=0$, and both
signs of the bounded testing state. It uses the entire original
weak $m$ event, without a false positive-unit lower bound on $m$.

\paragraph{Finite outer angles without the weak event.}
The same critical parent also gives the two finite outer charts.
In (SEO11) regard $B,Y,J_t,V$ as independent roots. Exactly,
\[
 F_*(0,v)=F_*(0,0)
       +v^2\{2dB^2+(4A+2d)Y^2+4cJ_t^2\}.       \tag{SEO17}
\]
The zero-angle form has three positive hard pivots and null ray
$(B,Y,J_t,V)=(0,y,b,1)V$. On this null ray the additional
coefficient is $\gamma v^2>0$. The normalized hard frames
(SEO12) or (SEO14) keep every coefficient bounded. Thus when
$u$ is infinitesimal and $v$ a bounded unit the full four-root
residue is strict. The infinitesimal critical-parameter deficit
is absorbed directly on these four roots; retain the entire
$m,\alpha,\beta$ terms.

For bounded-unit $u$ and infinitesimal $v$, set $V'=D_uV$.
The other exact identity is
\[
 F_*(u,0)=F_*(0,0;V')
        +u^2\{4AY^2+2d(V')^2/D_u\}.             \tag{SEO18}
\]
Use the same normalized critical roots with $V'$ in place of $V$.
In the first source case include $uvB$ in the $H_D$ root so that
the unbounded $D$ event remains isolated for nonzero small $v$.
The extra null mass in (SEO18) is positive in residue because
$d$ and $u$ are units. The bounded full four-root Gram is strict,
and the deficit is again a zero-residue correction. These two
identities include every state, not merely the normal restriction.

For completeness the same argument covers both finite angles of
nonzero residue, without a separate physical-source certificate.
The exact determinant of (SEO11) on $(B,Y,J_t,V)$ is
\[
 64cdD_u^2D_v^2
 \{A(2D+d)u^2+D(2A+d)v^2\}.                    \tag{SEO20}
\]
The determinant in the normalized frame (SEO12) is this number
divided by $AC$, namely
\[
 32cdD_u^2D_v^2
 \left\{u^2+\frac{D(2A+d)}{A(2D+d)}v^2\right\}.
\]
It is a positive unit whenever at least one finite angle has
nonzero residue. The hard restriction $V=0$ is explicitly
\[
 \left(4D_uD_v+\frac{2dD_uv^2}{A}\right)H_A^2
 +H_D^2+4cD_v(H_D/\sqrt C-uvB)^2+2dD_vB^2,
\]
which has a strict bounded residue. Hence its full residue Gram
is positive definite by (SEO20). In frame (SEO14), the determinant
is (SEO20) divided by $A(4c+C)$. Here $A,D,d$ are units and
$c/(4c+C)\sim1/4$, so it too is a positive unit. Its hard
restriction is
\[
 \left(4D_uD_v+\frac{2dD_uv^2}{A}\right)H_A^2
 +\frac{4cD_v}{4c+C}H_c^2
 +C(H_c/\sqrt{4c+C}+uvB)^2+2dD_vB^2,
\]
again strict. All entries of both full transformed Grams are
bounded as already proved. The infinitesimal deficit is a
zero-residue correction on these four roots. This settles every
finite angle pair outside the double-infinitesimal corner.

\paragraph{Projective infinity charts.}
For these charts normalize by $a$, clip the positive coefficients
$c/a,A/a,D/a$ at suitable positive real lower bounds, and retain
their positive original excess squares. Retain the whole original
$m,\alpha,\beta$ terms. The normalized $d/a$ is a unit.
All the infinity-infinity and infinity-zero four-root identities
in (BE2) and the subsequent outer routing apply without $m$.
For $u=\infty$ in residue and $v$ a bounded unit, the signed-core
residue itself retains
$2d\{(W+vX)^2+(vY-Z)^2\}$; together with $cW^2,AY^2$ it is
strict, so no $m(vY-Z)^2$ lower bound is needed.
For $v=\infty$ and $u$ a bounded unit the old signed-core and
$c,D$ root argument is unchanged.
At the remaining $(u,1/v)=(0,0)$ corner, replace the obsolete
state-$W$ mass argument by the exact root ideal
$(X,Y,U=uW,Z)^2$ after retaining $m$. In $q^2F(u,1/q)$ its
complete selected residue is
\[
 2dX^2+4c(U-Y)^2+(4A+2d)Y^2+4DZ^2\succ0.
 \tag{SEO19}
\]
Every remaining occurrence of $W$ carries $u$, so this is a full
bounded Gram including all states. Literal infinity is included
by these homogeneous identities. Literal zero axes have their
unchanged original-event SOS. Finite tests with both noncommuting
residues were included directly in (SEO20).

\paragraph{The residual state chart.}
Only $u,v$ infinitesimal remain. Normalize the state by a maximal
coordinate. On a negative pure-testing-core branch, $X,Y$ have
zero residue. The unit-$W$ chart was just proved for every bounded
$Z$. Otherwise $Z$ is a unit, so normalize $Z=1$ and $W$ is
infinitesimal. Then $w=(W+uY)/D_u$ is infinitesimal. The inherited
exact testing Schur identity and sign gauges give (SEO1).
Complex states split into their real and imaginary evaluations.



% Source formulas: pd4b92bf7de71
\section{Nonvanishing \texorpdfstring{$L=x\zeta$}{L=x zeta}}
For positive-unit or unbounded $L$, the reciprocal $z=1/L$ is bounded. When $\eta$ is unbounded, the uniform completion below combines two nonnegative source weights through their sum and covers every relative size of $h$ and the infinitesimal $\sigma=1/\eta$. The other source and testing domains retain their separate completions in complementary coordinates.

\subsection{Unbounded $L$ and unit $\eta$}

\label{component:pd4b92bf7de71}

\label{pd4b92bf7de71:sec:c14-large-l-unit}
This section proves this full original-parent subdomain, retaining every
state normal and every independent source/testing rate. It makes no claim
for infinitesimal $\sigma$. The common common remote theorem leaves
$u$ bounded and $w=1/v$ infinitesimal; both $u$ infinitesimal and $u$ a
bounded unit are covered below. All calculations use the exact NC source,
with $d=1$, $n=\lambda=0$, and its original nonnegative $D,\beta$ faces.

\paragraph{Source chart and bounded products.}
Put $x=k^2$, $z=q^2=L^{-1}$, $r=bz$, $Q=1-b$, $\theta=1-h$,
$K=1+x$, $R=1+xbz$, $M=R-xz\theta$, and
\[
 J=1+xh-\sigma(1-bz),\qquad E=1-xh/\sigma.
\]
Here $x,z$ are positive infinitesimals, $0<b<1$, and $\sigma$ is a positive
bounded unit. The NC source gives exactly
\begin{align*}
 P&=\sigma+K+\sigma z(1+bx+2xh)+\sigma bx^2hz^2+x^2h\theta z,\\
 V_0&=z(1+xh)E+(1-bz)(1+xhz),\\
 C&=cz,&c&=\frac{K(1+\sigma bz)E}{2(1-bz)P},\\
 m&=\mu x^2z,&\mu&=\frac{K(1+\sigma bz)h(1+\sigma z)}{2\sigma MP},\\
 A&=a/z,&a&=\frac{V_0}{2bP},\\
 \alpha&=\ell/z,&\ell&=\frac{(1+\sigma bz)[(x/\sigma)Rh(1+\sigma z)+ME]}{2QP},\\
 D&=\frac{K(1+\sigma z)(1+\sigma bz)}{2JP}-\frac12,&
 \beta&=\frac{K(1+\sigma bz)(1+xhz)}{2\theta P}-\frac12.
\end{align*}
In particular $c,\mu,ab,\ell Q,DJ,\beta\theta$ are bounded,
$c,\mu,ab,\ell Q,DJ$ are positive units, and $ab-\ell Q\in z\mathcal O$.
At $x=z=0$ their residues are
\[
\begin{gathered}
c_0=\frac1{2(1+\sigma)},\quad
\mu_0=\frac h{2\sigma(1+\sigma)},\quad
(ab)_0=(\ell Q)_0=\frac1{2(1+\sigma)},\\
(DJ)_0=\frac{\sigma^2}{2(1+\sigma)},\quad
(\beta\theta)_0=\frac{(1+\sigma)h-\sigma}{2(1+\sigma)}.
\end{gathered}
\]
Source nonnegativity gives $0<\sigma\leq1$ and
$h\geq\sigma/(1+\sigma)$ at residue. Thus $h$ is a positive unit.
The openings $b,Q,J,\theta$ remain independent; $\beta=0$ is retained.
The face $D=0$ cannot occur here since $DJ$ is a positive unit.

Write $D_u=1+u^2$, $D_w=1+w^2$ and
\begin{gather*}
 M_w=W-wX+uY-uwZ,\quad R_w=uW-Y,\quad
 J_w=uwW+uX-wY-Z,\quad B_w=X+uZ,\\
 K_w=J_w^2+D_u(Y-wZ)^2+D_wB_w^2-D_uD_wZ^2.
\end{gather*}
The complete reciprocal parent is
\[
 F_w=4mM_w^2+4czD_wR_w^2+4(a/z)D_uD_wY^2
 +4DJ_w^2+4(\ell/z)D_u(Y-wZ)^2+4\beta D_wB_w^2+2K_w.
\]
Choose bounded positive-unit caps $a_*\leq a$, $\ell_*\leq\ell$,
$D_*\leq D$, and a bounded nonnegative cap $\beta_*\leq\beta$.
Keep as nonnegative squares the four centered excesses
\begin{gather*}
4(a-a_*)D_uD_wN_y^2,\quad4(\ell-\ell_*)D_uN_y^2,\\
4(D-D_*)(J_w+JZ)^2,\quad4(\beta-\beta_*)D_wN_x^2,\\
N_x=X+uhZ,\qquad N_y=(Y-bwZ)/q.
\end{gather*}
These roots vanish at the exact source contact
\[
u=k,\quad w=q\sqrt\sigma,\quad
(W,X,Y,Z)=(\sqrt\sigma/(kq),-kh,bq\sqrt\sigma,1).
\]
Set $d_J=(D-D_*)J$ and $b_\theta=\beta\theta$.
Their coefficients and the centered excesses preserve the original parent.

\paragraph{Bounded full-normal remainder.}
Define $G=wZ/q$, $R_h=G-uqW$ and
\[
c_N=8D_u[abD_w-\ell Q],\qquad c_G=4D_u[ab^2D_w+\ell Q^2].
\]
The $A,\alpha$ remainder is
$4D_u(a_*D_w+\ell_*)N_y^2+c_NN_yG+c_GG^2$.
Here $c_N\in(z,w^2)\mathcal O$ and $c_G$ is a positive unit.
Keep $c_GR_h^2+c_NN_yR_h$ in the bounded Gram and use
$G^2-R_h^2=2uwWZ-u^2zW^2$. The rest is the bounded polynomial
\begin{align*}
G_0={}&4D_u(a_*D_w+\ell_*)N_y^2+c_NuqWN_y+c_G(2uwWZ-u^2zW^2)\\
&+4mM_w^2+4czD_wR_w^2+4D_*J_w^2-8d_JZJ_w-4d_JJZ^2\\
&+4\beta_*D_wB_w^2+8(b_\theta-\beta_*\theta)uD_wZB_w
 -4(b_\theta-\beta_*\theta)u^2\theta D_wZ^2+2K_w.
\end{align*}
All its coefficients are bounded. Put $N_z=Z-uwW/\sigma$ and substitute
$Z=N_z+uwW/\sigma$, $X=N_x-uhZ$, $Y=qN_y+bwZ$.
For $t=u^2$, $\xi=w^2$, define
\[
 T_0=1+th+b\xi,\quad M_0=1+t(h-Q)\xi/\sigma,\quad
 R_f=1-b\xi/\sigma,\quad J_f=1-T_0/\sigma.
\]
In the formulas below $c_N,c_G$ mean their even-polynomial expressions.
The exact expansion is
\[
G_0=\phi W^2+wL_xWN_x+uqL_yWN_y+uwL_zWN_z+H(N_x,N_y,N_z),
\]
where
\begin{align*}
\phi={}&4\mu x^2zM_0^2+4cz(1+\xi)tR_f^2+c_Gt(2\xi/\sigma-z)\\
 &+4D_*t\xi J_f^2-8d_Jt\xi J_f/\sigma-4d_JJt\xi/\sigma^2
 +4b_\theta\theta(1+\xi)t^2\xi/\sigma^2\\
 &+2[t\xi J_f^2+(1+t)tQ^2\xi^2/\sigma^2
 +(1+\xi)t^2\theta^2\xi/\sigma^2-(1+t)(1+\xi)t\xi/\sigma^2],\\
L_x={}&-8\mu x^2zM_0+t[(8D_*+4)J_f-8d_J/\sigma
 +(8b_\theta+4\theta)(1+\xi)/\sigma],\\
L_y={}&8\mu x^2zM_0-8cz(1+\xi)R_f+c_N-(8D_*+4)\xi J_f
 +8d_J\xi/\sigma-4(1+t)Q\xi/\sigma,\\
L_z={}&8\mu x^2z(h-Q)M_0-8cz(1+\xi)bR_f+2c_G
 -(8D_*+4)T_0J_f-8d_JJ_f+8d_J(T_0-J)/\sigma\\
 &+8b_\theta\theta(1+\xi)t/\sigma+4(1+t)Q^2\xi/\sigma
 +4(1+\xi)t\theta^2/\sigma-4(1+t)(1+\xi)/\sigma.
\end{align*}
For a compact explicit specification of the entire quadratic block, put
\[
\begin{gathered}
m_N=-wN_x+uqN_y+uw(h-Q)N_z,\quad
r_N=-qN_y-bwN_z,\\
j_N=uN_x-wqN_y-T_0N_z,\quad
b_N=N_x+u\theta N_z.
\end{gathered}
\]
Then every quadratic normal coefficient is retained by the finite formula
\begin{align*}
H={}&4D_u(a_*D_w+\ell_*)N_y^2+4mm_N^2+4czD_wr_N^2
 +(4D_*+2)j_N^2-8d_JN_zj_N-4d_JJN_z^2\\
&+(4\beta_*+2)D_wb_N^2+8(b_\theta-\beta_*\theta)uD_wN_zb_N
 -4(b_\theta-\beta_*\theta)u^2\theta D_wN_z^2\\
&+2D_u(qN_y-QwN_z)^2-2D_uD_wN_z^2.
\end{align*}
The exact source substitutions give, at $(t,\xi)=(x,\sigma z)$,
\[
\phi=\phi_t=\phi_\xi=L_x=L_y=L_z=0.
\]
Moreover $\phi(0,\xi)=4\mu x^2z$,
$L_x(0,\xi)=-8\mu x^2z$, and $\phi(t,\sigma z)\in z\mathcal O[t]$.

\paragraph{Exact six-root lift in the critical strip $u$ infinitesimal.}
Monic division, first by $\xi-\sigma z$ and then by $t-x$, gives bounded
polynomial coefficients satisfying
\begin{align*}
\phi&=z(t-x)^2 A+t(t-x)(\xi-\sigma z)B+t(\xi-\sigma z)^2C,\\
L_x&=(t-x)(tA_x+zB_x)+t(\xi-\sigma z)C_x,\\
L_y&=(t-x)A_y+(\xi-\sigma z)B_y,\\
L_z&=(t-x)A_z+(\xi-\sigma z)B_z.
\end{align*}
For example, the constant coefficient of
$L_x(t,\sigma z)/(t-x)$ equals $8\mu xz$, giving its $(t,z)$ support.
No small or mixed inverse enters these monic quotients. Also $A(0)=4\mu$.
Take the six finite separated roots
\[
 N_x,\quad N_y,\quad N_z,\quad R_h,\quad
 A_\theta=q(t-x)W,\quad V=uw(t-x)W.
\]
The complete scalar lift follows from
\begin{align*}
 z(t-x)^2W^2&=A_\theta^2,\\
 t(t-x)(\xi-\sigma z)W^2&=u\sigma A_\theta R_h-\sigma VN_z,\\
 t(\xi-\sigma z)^2W^2&=\sigma^2(qR_h-wN_z)^2.
\end{align*}
All first normals lift by the following finite list, which includes their
whole coefficient support:
\begin{align*}
wt(t-x)W&=uV,&wz(t-x)W&=qwA_\theta,\\
wt(\xi-\sigma z)W&=u\sigma qwR_h-u\sigma w^2N_z,\\
uq(t-x)W&=uA_\theta,&uq(\xi-\sigma z)W&=\sigma zR_h-\sigma qwN_z,\\
uw(t-x)W&=V,&uw(\xi-\sigma z)W&=\sigma qwR_h-\sigma w^2N_z.
\end{align*}
Together with the displayed $H$ these give an exact bounded six-root Gram,
not merely a restriction to the contact line. At residue its first five
diagonal pivots are respectively
\[
2+4\beta_*,\quad4(a_*+\ell_*),\quad
4[D_*+d_J(1+\sigma)],\quad \frac2{1+\sigma},\quad4\mu_0.
\]
They are positive units through all allowed openings. The only potentially
nonzero remaining cross entry is $N_zV$. The exact relation
\[
 V^2=\sigma tA_\theta^2+u\sigma(t-x)A_\theta R_h
       -\sigma(t-x)VN_z
\]
introduces an arbitrarily large finite positive $V^2$ pivot with only
infinitesimal changes to all other entries. Its real Schur complement
therefore makes the residue strictly positive definite. The exact bounded
Gram is positive definite by the standard finite-dimensional local
criterion. All excess squares are nonnegative with separated coefficients.

\paragraph{Away from the angular contact: $u$ a bounded unit.}
The same source chart has a positive-unit angular gap at every such real
residue. A convenient raw finite frame is
$E_t=qW,F_t=Y/q,G_t=wZ/q,R_t=wW,X,Z$; products obey $E_tG_t=R_tZ$.
The residue parent divided by $2/(1+\sigma)$ is exactly
\[
u^2E_t^2+D_u\{F_t^2/b+(F_t-G_t)^2/Q\}
 +\frac{(uR_t+uX-Z)^2}{J_0}
 +\frac{(X+uZ)^2}{\theta}-(1+\sigma)D_uZ^2,
\quad J_0=1-\sigma.
\]
Eliminate the positive $F_t$ and $X$ directions and allocate
the common product between the remaining two quadratic blocks. Up to a
common positive factor the two blocks have diagonal products $AB,CD_0$
and total off-diagonal magnitude $E_0$, with $U=u^2>0$ and
\[
\begin{gathered}
\delta=1-\sigma+\theta U,\qquad
A=U,\quad B=1+U,\quad C=U/\delta,\\
D_0=(1+U)[(1+U)/\delta-(1+\sigma)],\quad
E_0=\sqrt U(1+U)/\delta.
\end{gathered}
\]
Both diagonal products are strictly positive, since
$1+U-(1+\sigma)\delta=\sigma^2+U[(1+\sigma)h-\sigma]>0$.
All square roots in this sentence are real-residue quantities used to
choose real allocation constants, not mixed tensor radicals.
The required strict inequality is $E_0<\sqrt{AB}+\sqrt{CD_0}$.
Set $N=U(2\sigma-\theta U)$. Directly,
\[
E_0^2-AB-CD_0=(1+U)N/\delta,\quad
N^2-4U^2[1+U-(1+\sigma)\delta]
 =U^3(\theta^2U-4h).
\]
If $N\leq0$ strictness is immediate. If $N>0$, then
$\theta^2U<2\sigma\theta\leq2h<4h$, again giving strictness.
The single-opening faces follow by their constrained limits. At the only
double opening $J=\theta=0$, so $\sigma=h=1$, the exact limit is
$E_0-\sqrt{CD_0}\longrightarrow\sqrt U$, whereas $AB-U=U^2>0$.
Here is a full-frame square completion which also justifies those cap
limits. For clarity first take $u>0$ and select a finite real
$\rho$ in the nonempty interval
\[
\max\{-\sqrt{AB},E_0-\sqrt{CD_0}\}<\rho
 <\min\{\sqrt{AB},E_0+\sqrt{CD_0}\}.
\]
Its nonemptiness follows from the strict allocation inequality above.
Put
\[
r_\rho=D_u/u-\rho\delta/U,\qquad
s_\rho=D_u[2\rho/u-(1+\sigma)]-\rho^2\delta/U>0.
\]
The preceding complete residue form equals
\begin{align*}
&UE_t^2+D_uG_t^2-2\rho E_tG_t
 +\frac{D_u}{bQ}(F_t-bG_t)^2\\
&\quad+\frac{\delta}{J_0\theta}
 \left[X+\frac{U\theta R_t+u(J_0-\theta)Z}{\delta}\right]^2
 +\frac U\delta(R_t-r_\rho Z)^2+s_\rho Z^2.
\end{align*}
The equality uses only $E_tG_t=R_tZ$ and is checked exactly.
For $u<0$ take the correspondingly signed $\rho$; all sign choices are
real constants in this allocation. All triangular minimizer coefficients
are bounded through every opening because $U$ is a unit and
$\delta=J_0+U\theta$. Their positive hard weights stay bounded below.
At the double face choose $|u|<|\rho|<\sqrt{U(1+U)}$, with matching
sign; then $s_\rho$ has a positive-unit limit. This proves uniform
coercivity of the constrained residue, not just determinant positivity.

For completeness the exact centered cap rows in the raw frame are
\[
 F_t-bG_t,\quad
 uR_t+uX-qwF_t+(J-1)Z,\quad X+uhZ.
\]
Only one of $b,Q$ can tend to zero. The other two rows have hard parts
on $(R_t,X)$ with determinant $u$, a unit; each single row has a unit
hard entry. Their simultaneous kernel is solved exactly by
\[
F_t=bG_t,\qquad X=-uhZ,\qquad
R_t=\{(u^2h-J+1)Z+qwbG_t\}/u.
\]
More generally, prescribed changes $(e_F,e_D,e_\beta)$ in these three
rows are produced by
$\Delta F_t=e_F$, $\Delta X=e_\beta$,
$\Delta R_t=(e_D-u e_\beta+qw e_F)/u$.
This is a bounded separated right inverse: its only inversion is the
pure-testing unit $u$. For the finite-cap existence argument fix $u$
at its ordinary real residue, so the same inverse is an ordinary real
finite-dimensional right inverse. No order on the tensor product is
used. It lifts the limiting constrained kernel to the exact centered
kernel while keeping every opening independent. The
bounded residual form restricted to that kernel is the strictly positive
constrained form just computed. An elementary block Schur complement
then provides finite sufficiently large caps on every divergent row;
bounded rows retain their full coefficients. This is the finite-cap
lemma used in the UH proof. It yields a strictly positive bounded
full-frame Gram. Exact coefficients differ from their residues
infinitesimally, completing this strip.
All divisions by $\delta$ and radicals in this paragraph belong only
to the ordinary real residue calculation and its face limits. The exact
tensor Gram uses the polynomial raw frame, real allocation constants,
and finite source caps; it introduces no mixed inverse or radical.

Consequently this section closes the $L$ unbounded, $\sigma$ bounded-unit
subdomain of C1.4 under the original parent, including all state normals,
all independent opening rates, and $\beta=0$. It does not close the
$\sigma$ infinitesimal subdomain or claim a whole-row count change.



% Source formulas: p080aa62de16d
\subsection{Unit $L$ and infinitesimal $\eta$: C1 substitution in LII}

\label{component:p080aa62de16d}

\label{p080aa62de16d:sec:r7-c1-L-unit-eta-small}

\begin{theorem}
The entire original C1 source subdomain
\[
 x=k^2\in\Inf_{>0},\qquad L=x\zeta\text{ a positive bounded unit},
 \qquad \eta\in\Inf_{>0}
\]
has a full-parent tensor-SOS certificate for every independent
projective testing pair.  All original openings and the retained
boundaries $D=0$ or $\beta=0$ remain allowed whenever they occur
in this source subdomain.
\end{theorem}

The proof is an exact adapter to the strong-II local and remote
certificates.  Its additional event is retained with a centered
square; it is not deleted from the parent.

\paragraph{The native source in the shared contact geometry.}
Keep the native notation of (NC1)--(NC4), but in this section denote
its $r,h,a$ by $r_n,h_n,a_n$ to avoid collisions.  Normalize $d=1$.
The exact native contact with $Z_c=1$ is
\[
 (u_c,v_c;W_c,X_c,Y_c,Z_c)
 =\left(\sqrt x,\sqrt{x\eta\zeta};
 \sqrt{\zeta/\eta},-\sqrt x\,h_n,
 \sqrt x\,r_n\sqrt{\zeta/\eta},1\right).
\]
It therefore has the same contact geometry used in (LII2), with
\[
 z=x,\quad y=L\eta,\quad B=L,\quad H=Lr_n=B-a_ny,
 \quad r=h_n=1-\theta,\quad t=a_n,\quad
 j=1+zr-t=J_n.                                           \tag{LC1}
\]
In particular $z,y$ are positive infinitesimals.  Here $r_n\sim1$,
$0<Lr_n<1$, and hence $b_0:=\operatorname{res}B$ satisfies
$0<b_0\le1$.  We use $b_0$ only for this source residue.

Set, as in the local strong-II proof,
\[
 c=Cyt,\quad d_p=Dyj,\quad
 \mu=m/z^2,\quad a_A=A/y,\quad
 \sigma=\alpha(1-H),\quad \tau=\beta(1-r)=\beta\theta.
\]
These are bounded source quantities, and $c,d_p,\mu,a_A,\sigma$
are positive units.  Writing $r_0=\operatorname{res}r$, exact
substitution in (NC2)--(NC4), with opening factors canceled first,
gives
\[
\begin{gathered}
 c,d_p\sim\frac{b_0}{2(1+b_0)},\qquad
 a_A\sim\frac1{2b_0(1+b_0)},\\
 \mu\sim\frac{r_0}{2b_0(1+b_0)},\qquad
 \sigma\sim\frac{b_0^2}{2(1+b_0)},\\
 \operatorname{res}\tau=
 \frac{b_0}{2(1+b_0)}-\frac{1-r_0}{2}.                    
\end{gathered}
\tag{LC2}
\]
Since $\tau\ge0$, the last identity implies
$r_0\ge1/(1+b_0)>0$, proving the asserted unit lower bound for
$\mu$.  Also $t+j=1+zr\sim1$, and $c,d_p$ being units proves
$C,D$ unbounded uniformly through either opening $t=0$ or $j=0$
in residue.  Thus $D=0$ is in fact absent from this specific
subdomain by the exact source implications; it has not been excluded
by an additional hypothesis.  Literal $\beta=0$ is included.

The source-harmonic quantity $yCD/(C+D)$ and the cap behavior of
$\alpha$ have precisely the residues (RR1)--(RR2), with
$h=b_0$.  The remote proof (RR3)--(RR13) therefore applies unchanged,
retaining the extra positive term
$4\beta D_v(X+uZ)^2$ throughout.  This covers every test outside
the local region where $u,v$ are both infinitesimal, including the
original projective endpoints.

\paragraph{The additional event and its exact local remainder.}
For the local region use exactly the graph, roots, harmonic pair,
and caps (LII2)--(LII4), with the source quantities (LC1)--(LC2).
Put
\[
 R=v(U-z)W,\qquad
 E_\beta=\xi+u\zeta+R/B
          =(X+uZ)-\frac{z(1-r)v}{B}W.                    \tag{LC3}
\]
This root vanishes at the actual contact.  Retain the entire
positive term $4\beta D_vE_\beta^2$, in addition to the first,
second and fourth terms of (LII5); there is no $n$ term since
$n=0$.  The exact remainder from the $\beta$ event is
\[
 \begin{split}
 \Delta_\beta={}&
 \frac{8\tau z}{B}D_vvW(\xi+u\zeta)\\
 &+\frac{4\tau z}{B^2}D_vV\{2U-z(1+r)\}W^2.
 \end{split}                                             \tag{LC4}
\]
All its coefficients are bounded even if $\beta$ is unbounded:
only $\tau=\beta(1-r)$ occurs.  The exact full remainder is (LII6)
with its $n$ line deleted, plus (LC4).  Thus no state normal,
boundary coefficient, or first contact jet has been omitted.

The $R_hW$ gauge (LII7) is unchanged, because (LC4) contains no
$R_h$.  The coefficient representation (LII8), its exact contact
and first-normal vanishings (LII9)--(LII10), and its bounded monic
divisions (LII12)--(LII13) all remain valid.  Explicitly, (LC4)
adds $8\tau z(1+V)/B$ to $P_\xi$ and $P_\zeta$, adds zero to
$P_N$, and adds its printed scalar line to $\phi$.  Hence the
essential support $P_\xi\in(U,z)$ and $\deg_U\phi\le2$ remain
literal polynomial statements.

The only changed data in (LII11) are
\[
\begin{split}
 d_s&=\frac{4z^2}{B^2}
   \left\{\mu(B+zr)^2+\frac{r^2}{2}-\tau(1+r)\right\},\\
 e_s&=\frac{4z^2}{B^2}
   \left\{\frac{r^2}{2}-\tau(1+r)\right\}.
\end{split}                                             \tag{LC5}
\]
Both quotients by $z^2$ are bounded; $e_s$ need not be positive.
Its contribution to the final Gram has an infinitesimal coefficient
because it is multiplied by $V$.  The decisive exact residue
cancellation, using (LC2), is
\[
 \begin{split}
 \operatorname{res}\left\{
 \mu B^2+\frac{r^2}{2}-\tau(1+r)\right\}
 &=\frac{b_0r_0}{2(1+b_0)}+\frac{r_0^2}{2}
 -\left\{\frac{b_0}{2(1+b_0)}-\frac{1-r_0}{2}\right\}(1+r_0)\\
 &=\frac1{2(1+b_0)}.
 \end{split}                                             \tag{LC6}
\]
Consequently
\[
 \operatorname{res}(d_s/z^2)=\frac2{b_0^2(1+b_0)},\qquad
 \operatorname{res}A_\xi=\frac4{b_0(1+b_0)}.              \tag{LC7}
\]
The second equality is unchanged from (LII18): the new $P_\xi$
term is a bounded multiple of $z$ independent of $U$, so it has
zero $U$-quotient residue.  The four hard-normal residues (LII17)
are unchanged as well.

Apply the exact allocations (LII15) to all the coefficients,
including (LC4).  They give the same bounded six-root Gram
$(N,\xi,\zeta,R_h,R,T)$ as in (LII16).  Its finite $(\xi,R)$
block is exactly (LII19), with $h=b_0$, and has strictly positive
Schur complement $2/[b_0(1+b_0)^2]$.  The same exact conductor
(LII20) supplies an arbitrarily large real $T^2$ pivot while
changing every other residue entry by zero.  The complete bounded
Gram is strict.  Restoring the four retained terms, including
$4\beta D_vE_\beta^2$, proves the local tensor SOS for all
independent rates and literal affine axes.  The remote certificate
above completes the theorem.

\paragraph{Source inequalities.}
  Stationarity
itself is the unchanged native source identity, not an added
restriction.  The result removes the whole $L$-unit, $\eta$-infinitesimal
subdomain from C1.4.  It does not assert the other $L$-unit source
regimes or the entire C1.4 row.



% Source formulas: p686209b87229
\subsection{Unit $L,\eta$: C1 substitution in LII}

\label{component:p686209b87229}

\label{p686209b87229:sec:r7-c1-L-eta-unit}

\begin{theorem}
For the original C1 source with $x\in\Inf_{>0}$ and both
$L=x\zeta$ and $\eta$ positive bounded units, the entire original
parent is tensor-SOS for every independent projective testing pair.
All original positive openings and literal $\beta=0$ are included.
\end{theorem}

The exact source formulas force $D$ strong on this subdomain, so
$D=0$ is absent by source implication.  No opening is removed by
an extra hypothesis, and the proof does not invoke a table or atlas.

\paragraph{The source and the only testing strip left by four events.}
Use the exact shared geometry (LC1), and write $\rho=r_n$ for the
native shape parameter and $r=1-\theta$ for the geometric one.
Thus
\[
 B=L,\quad y=B\eta,\quad H=B\rho=B-yt,\quad
 t=(1-\rho)/\eta,\quad j=1+zr-t,
 \quad z=x\in\Inf_{>0}.
\]
Here $B,y,\eta$ are units.  For the boundary coefficient formulas
below set
\[
 P=1+B+y,\qquad P_4=1+B(1-\rho),\qquad S=H+y=P-P_4.
                                                               \tag{LU1}
\]
All three quantities are positive units, uniformly through
$\rho\to0$, $1-\rho\to0$, $j\to0$, and $1-H\to0$.
The native formulas imply that $D,C,A,\alpha$ are strong.
One forward cyclic relabeling sends these four consecutive events
to $C,A,\alpha,m$ and sends testing coordinates to
$(u',v')=(v,1/u)$.  The full-parent four-consecutive theorem
(CE11)--(CE16) therefore covers every test except
\[
                    u\in\Inf,\qquad v\text{ a bounded nonzero unit}.
                                                               \tag{LU2}
\]
The cycle is the inverse of (CE9)--(CE10), so its literal endpoint
scope is the original homogeneous parent's scope.

It remains to prove (LU2).  All roots and coefficients used below
are valid for both testing signs, including literal $u=0$.

\paragraph{A bounded full-parent remainder with every opening retained.}
Use the graph, harmonic roots and gauge (LII2)--(LII8), with
$k=\sqrt y$ now a source unit.  The bounded source data are
\[
 c=Cyt,\quad d_p=Dyj,\quad
 \mu=m/z^2,\quad \sigma=\alpha(1-H),\quad
 \tau=\beta(1-r),\quad A_H=AH.
\]
They are bounded; $c,d_p,\mu,\sigma,A_H$ are units.  In residue,
with the boundary symbols of (LU1), their exact values are
\[
 c=\frac S{2P},\quad d_p=\frac{BP_4}{2P},\quad
 \mu=\frac{Sr(1+y)}{2B^2P},\quad
 \sigma=\frac{BS}{2P},\quad A_H=\frac{yP_4}{2P},\quad
 \tau=\frac S{2P}-\frac{1-r}{2}.                         \tag{LU3}
\]
These are residue identities, not substitutions of source and
testing fields.  They are obtained after multiplying away every
possibly small opening in the exact native formulas.
Since $\tau\ge0$, $r\ge P_4/P$ in residue and is a positive unit.
Also $t+j\sim1$, so $\ell=y(C+D)$ has a positive lower bound,
$q=D/(C+D)$ is bounded in $[0,1]$, and
$g=yCD/(C+D)$ is a positive unit.

Choose positive bounded caps $\ell_*\le\ell$,
$A_*\le A$, and $\alpha_*\le\alpha$.  Choose a bounded cap
$0\le\beta_*\le\beta$ as specified below.  Retain the five
contact-centered positive terms
\[
\begin{gathered}
 4(\ell-\ell_*)N^2,\qquad 4CV\eta_t^2,\qquad
 4(A-A_*)D_uD_v\eta_t^2,\qquad
 4(\alpha-\alpha_*)D_u(v\eta_t-\zeta)^2,\\
 4(\beta-\beta_*)D_v(\xi+u\zeta+R/B)^2,
 \qquad \eta_t=k(N+qR_h),\quad R=v(U-z)W.
\end{gathered}
\tag{LU4}
\]
The exact remainder is (LII6), with its $n$ line deleted, its $A$
term replaced by
\[
 4A_*D_uD_v\eta_t^2
 +\frac{8A_H}B D_uD_vuW\eta_t
 +\frac{4A_HH}{B^2}D_uD_vUW^2,
                                                               \tag{LU5}
\]
and with (LC4) plus $4\beta_*D_v(\xi+u\zeta+R/B)^2$ added.
This specifies every coefficient of the full remainder.  Only
bounded source products occur, even at an unbounded original weight.

The same dependent-root gauge and exact contact equations give
(LII9)--(LII13).  The new constant scalar coefficients are
\[
 \frac{d_s+e_sV}{z^2}=
 \frac4{B^2}\left[
 \mu(B+zr)^2+
 \left\{\frac{r^2}{2}-\tau(1+r)+\beta_*\right\}(1+V)
 \right].                                                   \tag{LU6}
\]
The support $P_\xi\in(U,z)$ and scalar degree at most two in $U$
remain exact.  Apply all of (LII15).  As $v$ is a unit, eliminate
$T$ by the exact pure-factor identity
\[
                        T=(uk/v)R.                         \tag{LU7}
\]
Its coefficient is a bounded infinitesimal tensor.  Thus the entire
remainder has a bounded Gram on
$(N,\zeta,R_h,\xi,R)$.  No source/testing ratio or mixed inverse
has been introduced.  We now verify its complete residue.

\paragraph{A three-square identity supplies all three state/angle normals.}
The following boundary computation is an exact rational identity
before taking any opening limit.  Define
\[
\begin{aligned}
 C_0&=\frac{\rho+\eta}{2(1-\rho)P},&
 A_0&=\frac{\eta P_4}{2\rho P},\\
 D_0&=\frac{P_4}{2(\eta-1+\rho)P},&
 \alpha_0&=\frac{B^2(\rho+\eta)}{2(1-B\rho)P}.
\end{aligned}
\]
These may have unbounded openings; only the final capped identity
below is used as a bounded Gram.  For an indeterminate $p$, put
\[
 e=Y-\rho p,\qquad f=Z-vp/B,\qquad T_v=(V-y)p.
\]
The entire leading three-state form is exactly
\[
\begin{split}
 &4C_0D_v(p-Y)^2+4A_0D_vY^2
 +(4D_0+2)(p-Y-vZ)^2+(4\alpha_0+2)(vY-Z)^2-2D_vZ^2\\
 &\quad=\frac{4D_0}{B^2}(T_v+Be+Bvf)^2
       +4\alpha_0(f-ve)^2+4(C_0+A_0+\tfrac12)D_ve^2.
\end{split}                                                    \tag{LU8}
\]
This identity holds for every $v$, not only near the actual Bob
angle.  In the contact allocation at $z=u=0$,
$e=k(N+qR_h)$ and $T_v=-B(kR_h+v\zeta)$.
Subtract the residue contributions of the retained terms (LU4).
The harmonic identity then gives the complete three-normal block
\[
 \boxed{4\ell_*N^2+4gR_h^2
 +4(A_*+\tfrac12)D_vy(N+qR_h)^2
 +4\alpha_*\{\zeta-vk(N+qR_h)\}^2.}                     \tag{LU9}
\]
It has strictly positive residue on $(N,\zeta,R_h)$ because
$\ell_*,g,\alpha_*$ are positive units.  No normal has been
suppressed.  Every rational cancellation used to obtain (LU9)
has bounded source products on both sides, so the identity extends
continuously through every opening listed after (LU1).

\paragraph{The remaining two-root block and a finite cap.}
There is no finite residue coupling between these three roots and
$(\xi,R)$.  Such terms carry $u$, $z$, or the coefficient in (LU7).
For completeness, the exact first-normal support gives this block
without dropping a term.  To first order in $(U,z)$ write
$P_\xi=UA(V)+zZ(V)$; both coefficient functions are affine in $V$,
and
\[
 Z(V)=\frac{D_v}{B}(-4r+8\tau-8\beta_*).
\]
The exact contact equation gives $A(y)+Z(y)=0$ in residue.
After the first allocation of (LII15), the coefficient of $\xi R$
is therefore $A(V)-(V-y)(A'+Z')=-Z(V)$.
Using (LU3) and (LU6), the complete block is
\[
 M_{\xi R}=
 \begin{pmatrix}
 2D_v(1+2\beta_*)&\dfrac{2D_v}{B}(P_4/P+2\beta_*)\\[1mm]
 \dfrac{2D_v}{B}(P_4/P+2\beta_*)&
 \dfrac2{B^2}\{(P_4/P+2\beta_*)D_v-(S/P)r(V-y)\}
 \end{pmatrix}.                                           \tag{LU10}
\]
Here all letters denote real residues when positivity is assessed.

If $\beta$ is bounded, take $\beta_*=\beta$.  The native relation
$(1-r)(2\beta+1)=S/P$ follows from (LU3), and (LU10) has determinant
\[
 \frac{4S D_v}{B^2P}(1+2\beta)r(1+y)>0.                 \tag{LU11}
\]
Its first diagonal entry is positive, including $\beta=0$.
If $\beta$ is unbounded, take $\beta_*$ to be a sufficiently large
positive real cap.  Increasing it adds the positive square
$4\beta_*D_v(\xi+R/B)^2$.  On its kernel $\xi=-R/B$,
the uncapped block equals
\[
          \frac{2S}{B^2P}\{1+ry+(1-r)V\}R^2>0.           \tag{LU12}
\]
The elementary real two-by-two Schur argument therefore supplies
a finite cap making (LU10) strictly positive.  Unboundedness of the
original $\beta$ permits that cap.

Together (LU9)--(LU12) prove strict positive residue of the full
five-root Gram.  The bounded Gram lemma gives its tensor SOS;
restoring (LU4) proves (LU2).  The cyclic four-event complement
completes every testing chart and original projective endpoint.

\paragraph{Verification scope.}
  The exact source contact and
monic contact divisions are the unchanged full-parent argument of
(LII7)--(LII15), with the bounded modifications (LU4)--(LU6).
This theorem closes the $L$-unit, $\eta$-unit subdomain of C1.4;
the $\eta$-unbounded subdomain is not included.



% Source formulas: p10ee3d96a3f1
\subsection{Large $\eta$ with bounded reciprocal $x\zeta$: one proof for every $h/\sigma$}
\label{sec:c1-unification-bounded-z-all-H}

Let $K_c,K_t$ be separately labelled real closed extensions of $\mathbb R$,
$T=K_c\otimes_{\mathbb R}K_t$, and $\Sigma T^2$ the finite sums of
squares in $T$. Bounded coefficients belong to $B_c\otimes B_t$ and
their real residue is taken in both factors. We use the bounded-Gram
fact: a bounded symmetric matrix with positive-definite real residue
is a Gram matrix of a finite sum of squares in $T$. Roots in the Gram
need not be bounded. A positive source or testing factor can be absorbed
into a square using the real closedness of its own field.

\begin{theorem}\label{thm:c1-unification-bounded-z-all-H}
Use the stationary C1 source with $d=1$, $\lambda_{\rm ev}=n=0$,
$m,C,A,\alpha>0$, $D,\beta\ge0$, and the native strict gaps. Assume
explicitly that
\[
 x=k^2\in\operatorname{Inf}_{c,>0},\qquad
 \sigma=1/\eta\in\operatorname{Inf}_{c,>0},\qquad
 z=1/(x\zeta)>0\text{ is bounded}.
\]
There is no size restriction on $H=h/\sigma$. Then the complete
reciprocal parent is in $\Sigma T^2$ for every independent pair
$u,w\in\operatorname{Inf}_t$ and every real or complex state.
There is no state-ratio or testing-ratio restriction. The literal
faces $D=0$ and $\beta=0$ are included.

Equivalently, the theorem applies to the rational source below under
$0<b<1$, $0<r=bz<1$, $0<h<1$, $E,M,J>0$, $D,\beta\ge0$,
and the stated assumptions on $x,\sigma,z$.
\end{theorem}

\paragraph{The source and its uniform estimates.}
The rational identities and complete parent expansions below are
algebraic for every positive $z$ with the native gaps. Only their
boundedness estimates and the theorem require $z$ bounded.
Put
\[
 r=bz,\quad Q=1-b,\quad 0<b,r<1,\quad \theta=1-h,\quad
 e=xh/\sigma=xH,\quad E=1-e,\quad K=1+x,
\]
\[
 R_s=1+xr,\quad M=R_s-xz\theta>0,\quad
 J=1+xh-\sigma(1-r)>0,
\]
\[
 P=\sigma+K+\sigma z(1+bx+2xh)
              +\sigma bx^2hz^2+x^2h\theta z.
\]
Define the opening-cleared products
\[
\begin{split}
 p&=\frac{z(1+xh)E+(1-r)(1+xhz)}{2P},\\
 l&=\frac{(1+\sigma r)\{eR_s(1+\sigma z)+ME\}}{2P},\qquad
 \kappa=\frac{K(1+\sigma r)(1+\sigma z)}{2MP}.
\end{split}
\]
The event coefficients are
\[
\begin{gathered}
 m=\kappa x^2zH,\quad
 C=\frac{Kz(1+\sigma r)E}{2(1-r)P},\quad
 A=p/r,\quad \alpha=l/(zQ),\\
 D=\frac{K(1+\sigma z)(1+\sigma r)-JP}{2JP},\qquad
 \beta=\frac{K(1+\sigma r)(1+xhz)-\theta P}{2\theta P}.
\end{gathered}\tag{P1}\label{eq:c1-eta-large-source}
\]
Only the original nonnegative $D,\beta$ sources are considered.  Set
$a=Az=p/b$, $\ell=\alpha z=l/Q$, and
\[
 H_0=1+zQ,\quad B_s=P+xz(J-\sigma),
\]
\[
 \mathcal H=R_sH_0+z\{(1-r)xh(2+xr)
       -e[2+x(1+h+r\theta)]+Ke^2\}.
\]

Direct multiplication gives the two exact source identities
\begin{align}
 2\theta P\beta&=hB_s-\sigma H_0,\tag{BE3}\label{eq:c1-unification-source}\\
 2P(JD+x\theta\beta)&=\sigma^2\mathcal H.\tag{BE4}\label{eq:c1-unification-source-sphere}
\end{align}
They contain the entire scale estimate required below. Initially
$0<e<1$ follows from $E>0$, so $\mathcal H$ is bounded; also
$P,J,M,B_s\sim1$. Eliminating $\beta$ between the identities gives
\begin{equation}\label{eq:c1-unification-e-bound}
 eB_s=xH_0+\sigma\mathcal H-2PJD/\sigma
       \le xH_0+\sigma\mathcal H.
\end{equation}
Thus $e=O(x+\sigma)$ is infinitesimal. Consequently
\[
 \mathcal H-H_0\in\operatorname{Inf}_c,\quad
 D=O(\sigma^2),\quad \beta\theta=O(h+\sigma),\quad
 p-H_0/2,\ l-1/2,\ \kappa-1/2\in\operatorname{Inf}_c,
 \qquad H_0=1+zQ\ge1.
\]
The quantities $p,l,\kappa$ are bounded and positive, and
$m=\kappa xze$ is infinitesimal. The individual positive coefficients
$C,a=p/b,\ell=l/Q$ may be unbounded; they will stay in whole squares.

The following estimates do not assume that $H$ is bounded. Set
\[
 \nu=(2D-2\beta\theta+h)/\sigma,\quad
 \lambda=(8Dh+4\beta\theta^2+2h^2)/\sigma-4bp\sigma z,\quad
 K_A=4\kappa H-\lambda,
\]
and define
\[
 \lambda_0=\lambda-2H,\quad d_G=4D/\sigma^2,\quad
 \rho=\lambda x/\sigma,\quad L_0=\lambda\sigma.
\]
The source identities imply
\begin{align}
 \nu&=2D/\sigma+[H_0-ez(J-\sigma)]/P,\label{eq:c1-unification-nu}\\
 \lambda&=2H+2e\theta z(J-\sigma)/P
            -2\theta H_0/P+8Dh/\sigma-4bp\sigma z.
            \label{eq:c1-unification-lambda}
\end{align}
Hence $\nu,d_G,\lambda_0$ are bounded,
$\lambda_0+2\theta H_0$ is infinitesimal, and $L_0-2h$ is infinitesimal.
For $K_A=4\kappa H-\lambda$ one also has
\begin{equation}\label{eq:c1-unification-KA}
 K_A-\{2\theta H_0-2h(1-r)\}\in\operatorname{Inf}_c.
\end{equation}
Here is a finite polynomial derivation of the only possible pole in
this assertion. Put
\[
 N_\kappa=K(1+\sigma r)(1+\sigma z)-MP
          =xF_x+\sigma F_\sigma,
\]
where $F_x=(N_\kappa-N_\kappa|_{x=0})/x$ and
$F_\sigma=N_\kappa|_{x=0}/\sigma$ are polynomials in the bounded
source arguments. Direct substitution gives
$F_\sigma=r-1+\sigma rz$. Therefore
$H(2\kappa-1)=(eF_x+hF_\sigma)/(MP)$ has residue $-h(1-r)$,
which proves \eqref{eq:c1-unification-KA}. In particular $K_A$ is bounded.

For later use write $R_\gamma=R_0-8pu^2$, where
\[
 R_0=-4\{2(C+D)(r-1)+2p+r-1\}.
\]
There is the uniform estimate
\begin{equation}\label{eq:c1-unification-R0}
 R_0=O(x+\sigma+e),\qquad
 \lambda>0\ \Longrightarrow\ R_0\sqrt\lambda\in\operatorname{Inf}_c.
\end{equation}
To verify the first estimate without any source-pole division, temporarily
treat $e$ as an independent variable. The product $N=JPR_0$ is the
polynomial
\[
 -4\{-JKz(1+\sigma r)(1-e)
 +(r-1)K(1+\sigma z)(1+\sigma r)
 +J[z(1+xh)(1-e)+(1-r)(1+xhz)]\}.
\]
It vanishes identically at $x=\sigma=e=0$. Consecutive differences
in $x,\sigma,e$ express it as $xN_x+\sigma N_\sigma+eN_e$, with
bounded polynomial quotients. For the second estimate use
$\lambda=O(1+H)$ from \eqref{eq:c1-unification-lambda} and
$e=O(x+\sigma)$ from \eqref{eq:c1-unification-e-bound}. The bound
\[
 R_0^2\lambda=O((x+\sigma)^2(1+H)),\qquad
 (x+\sigma)^2H=xe+2\sigma e+\sigma h
\]
has an infinitesimal right side. This proves the square-root estimate
for every possible $H$.

The same source formulas also give the identities used by the
complementary testing charts. With $B=B_s$, put
\[
 d_0=\sqrt{2PJD/(\sigma^2\mathcal H)},\qquad
 b_0=\sqrt{2Px\theta\beta/(\sigma^2\mathcal H)}.
\]
The source identity gives
\[
 D=O(\sigma^2),\quad x\theta\beta=O(\sigma^2),\quad
 JD+x\theta\beta\asymp\sigma^2.\tag{BE5}
\]
Then exactly
\[
 d_0,b_0\ge0,\quad d_0^2+b_0^2=1,\quad
 D={\sigma^2\mathcal H\over2PJ}d_0^2,\quad
 \beta={\sigma^2\mathcal H\over2Px\theta}b_0^2.\tag{BE6}
\]
\[
 H={H_0\over B}+{\sigma\over x}{\mathcal H\over B}b_0^2.\tag{BE7}
\]
\[
 m=m_0+m_1b_0^2,\quad
 m_0=\kappa x^2zH_0/B\asymp x^2z,\quad
 m_1=\kappa x\sigma z\mathcal H/B\asymp x\sigma z.\tag{BE8}
\]
\[
 {C(1-r)\over z}\sim\tfrac12,\qquad
 Ar\sim H_0/2,\qquad \alpha Qz\sim\tfrac12.\tag{BE9}
\]
No inverse of a sphere coordinate is used. Another useful uniform
source identity is $p-l\in zB_c$: with $e$ temporarily independent,
$2P(p-l)$ is a polynomial in bounded $x,\sigma,z,b,h,e$, and
at $z=0$ its two source numerators are both one
($e+(1-e)=1$ for the second). Its finite quotient by $z$ is
therefore polynomial and bounded. Restoring $e=xh/\sigma$ proves
the assertion through both radial openings.

\paragraph{The complete state polynomial.}
Write $t=u^2$, $\xi=w^2$, $D_u=1+t$, $D_w=1+\xi$, $q=\sqrt z$,
$g=h-Q$.  The original reciprocal parent is
\[
\begin{split}
 F_w={}&4m(W-wX+uY-uwZ)^2+4CD_w(uW-Y)^2+4AD_uD_wY^2\\
 &+(4D+2)(uwW+uX-wY-Z)^2+(4\alpha+2)D_u(Y-wZ)^2\\
 &+(4\beta+2)D_w(X+uZ)^2-2D_uD_wZ^2.
\end{split}\tag{P4}\label{eq:c1-reciprocal-parent}
\]
Use the homogeneous roots
\[
\begin{gathered}
 N_x=X+uhZ,\quad N_y=(Y-bwZ)/q,\quad
 \Gamma=\sigma Z-uwW,\quad R=wZ/q-uqW,\\
 A_\theta=q(t-x)W,\quad P_\beta=uZ-(x/\sigma)wW,
 \quad X_e=X+ewW,\quad V=uw(t-x)W,\\
 J_C=N_y+bR,\qquad J_\ell=N_y-QR.
\end{gathered}
\]
The following definitions collect every remaining coefficient:
\[
\begin{gathered}
 \nu=(2D-2\beta\theta+h)/\sigma,\qquad
 \lambda=(8Dh+4\beta\theta^2+2h^2)/\sigma-4bp\sigma z,\\
 \lambda_0=\lambda-2H=4D(1+h)/\sigma-2\theta\nu-4bp\sigma z,\\
 K_A=4\kappa H-\lambda,\qquad B_x=8\kappa e,\\
 A_y=8(p-l)+4z\{(1+2m)g+2Dh+2p\sigma\},\\
 R_x=-4\{2Db+(1+2m)g-2\beta\theta\},\\
 R_\gamma=-4\{2(C+D)(r-1)+2p+r-1+2pt\},\\
 \chi=8(pb+lQ)-4h+4eh(2D+1)-8Dh\\
 \hspace{20mm}+4z\{4Dbh+4pb\sigma+(1+2m)g^2+2\beta\theta^2\},\\
 k_1=2(1+2m)g^2+8Dhb+8pb\sigma+4\beta\theta^2,
\end{gathered}\tag{SP3}\label{eq:c1-contact-coefficients}
\]
\[
\begin{split}
 E_{yy}&=2z\{(1+2m)t+(1+2D)\xi\},\\
 E_{RR}&=z\{8bD+8bp\sigma+4b\sigma(2C+2D+1)(r-1)
                                  +k_1t+2b^2(2D+1)\xi\},\\
 E_{yR}&=4z\{t[2Dh+2p\sigma+(1+2m)g]+b(2D+1)\xi
       +\sigma(2C+1)(r-1)+2D[1+\sigma(r-1)]+2p\sigma\}.
\end{split}
\]
Direct expansion of the original parent, using only the displayed source
identities, gives the full four-state identity
\begin{align}
 F_w={}&\mathcal E_++[2(1+\xi)+4m\xi]N_x^2
       +[4a+2z(2C+1)]J_C^2+4\ell J_\ell^2\nonumber\\
 &+E_{yy}N_y^2+E_{RR}R^2+E_{yR}N_yR
        +(4D/\sigma^2)\Gamma^2\nonumber\\
 &-4\nu N_x(u\Gamma)
       +R_\gamma(N_y+bR)(qw\Gamma)\nonumber\\
 &-4uqw(2m+2D+1)N_xN_y+uqwR_xN_xR
           +qwB_xN_xA_\theta+uA_yN_yA_\theta\nonumber\\
 &+K_AA_\theta^2+\chi uA_\theta R+\mathscr S,
 \label{eq:compression-full-parent}\tag{SP6}\\
 \mathscr S={}&\lambda\sigma P_\beta^2
       +(\lambda x/\sigma)(t-x)(\xi-\sigma z)W^2,\nonumber
\end{align}
where all terms in the following remainder are retained tensor squares:
\[
\begin{split}
 \mathcal E_+={}&4\beta D_wN_x^2+2(2D+1)(uX_e)^2
              +4bpzt\Gamma^2\\
 &+4a(D_uD_w-1)J_C^2+4\ell tJ_\ell^2+4Cz\xi J_C^2.
\end{split}
\]
This identity includes all four state coordinates and both complete Gamma normal terms.  


All $E$ coefficients are infinitesimal. The state-normal cross
coefficients outside $\mathscr S$, namely
$\nu,R_\gamma,R_x,B_x,A_y,\chi$, are bounded; so are $K_A,d_G$.
The possibly unbounded scalar coefficient $\lambda$ is handled below.
Indeed the only $C$ poles outside positive squares occur through
$C(r-1)$; $\beta$ occurs there through the bounded product
$\beta\theta$. Every non-Gamma cross coefficient has an explicit
$u$ or $w$ factor and is infinitesimal. These conclusions use the
uniform estimates just proved, not a bound on $H$.

Since $4a,4\ell>1$, clip these two whole positive radial squares at
one and retain their source excesses. The resulting radial form is
\[
 J_C^2+J_\ell^2=(N_y+bR)^2+(N_y-QR)^2,
\]
whose real residue is positive definite because $b+Q=1$.
The $N_x$ coefficient has residue two. The exact testing sign gauges
$(u,Y,Z)\mapsto(-u,-Y,-Z)$ and $(w,X,Z)\mapsto(-w,-X,-Z)$
allow $u,w\ge0$; literal zero values are allowed.

\paragraph{A single construction for $\lambda>0$.}
In this paragraph $\rho>0$. Direct elimination of $\beta$ in
\eqref{eq:c1-unification-source} gives the exact identity
\begin{equation}\label{eq:c1-unification-S}
\begin{split}
 S:=d_G+\rho={}&2\mathcal H/P
 +d_G(1-J+2hx)+2eH_0/P\\
 &-2e^2z(J-\sigma)/P-4bpxz.
\end{split}
\end{equation}
Thus $S-2H_0$ is infinitesimal. In particular $S$ is a positive
bounded source unit and $d_G,\rho$ are both bounded. Also
$q_0:=\lambda k$ is bounded, since $q_0^2=L_0\rho$, and
$\lambda x=\rho\sigma$ is infinitesimal.

Define the separated monomials, keeping the signs of $W,Z$,
\[
 B_1=\sqrt{\lambda\sigma}\,uZ,\quad
 B_2=\sqrt\rho\,uwW,\quad
 B_3=\sqrt\lambda\,x\sqrt zW,\quad
 B_4=\sqrt{\lambda x}\,wZ/\sqrt z.
\]
Exactly
\[
 \mathscr C:=\mathscr S+\lambda xR^2
        =\sum_i B_i^2-4\lambda xuwWZ.
\]
If $WZ>0$, change the overall state sign so $W,Z>0$, put
$G=(B_1B_2B_3B_4)^{1/4}$ and $d_i=B_i-G$. Then
\[
 \mathscr C=\sum_i d_i^2+\mathscr C_+,
\]
\[
 \mathscr C_+=2G\{(\sqrt{B_1}-\sqrt{B_2})^2
 +(\sqrt{B_3}-\sqrt{B_4})^2
 +2[(B_1B_2)^{1/4}-(B_3B_4)^{1/4}]^2\}\in\Sigma T^2.
\]
All radicals separate into factors in $K_c$ and $K_t$. If $WZ\le0$,
put $d_i=B_i$ and $\mathscr C_+=-4\lambda xuwWZ\in\Sigma T^2$.
In both cases the differences $d_i-d_j$ equal $B_i-B_j$.

Put $T_\theta=k\sqrt z(u-k)W$ and $\Gamma_*=\sqrt{d_G}\Gamma$.
The following exact conductors require no inverse of $D$ or sphere
coordinate:
\begin{align}
 \sqrt\rho\,u\Gamma={}&k(d_1-d_2)-(u-k)(d_2-d_3)
                      -\sqrt{\rho\sigma}\,T_\theta,\label{eq:c1-unification-conductor1}\\
 \sqrt\rho\sqrt z w\Gamma={}&\sqrt\rho\sigma zR
 -(\sqrt z w+z\sqrt\sigma)(d_2-d_3)\nonumber\\
 &+\sqrt\rho(\sigma z/k+\sqrt\sigma\sqrt z w/k)T_\theta.
 \label{eq:c1-unification-conductor2}
\end{align}
Use $1=(d_G+\rho)/S$ to write
\[
 u\Gamma={\sqrt{d_G}u\over S}\Gamma_*
              +{\sqrt\rho\over S}(\sqrt\rho\,u\Gamma).
\]
After \eqref{eq:c1-unification-conductor1}, every coefficient on
$(\Gamma_*,T_\theta,d_i)$ is infinitesimal. For the other complete
normal use
\[
 \sqrt z wR_\gamma\Gamma=R_0\sqrt z w\Gamma
                              -8p\sqrt z wu(u\Gamma),
\]
and apply the same identity $1=(d_G+\rho)/S$ only to the first term.
In \eqref{eq:c1-unification-conductor2} its only potentially singular
$T_\theta$ coefficients reduce exactly to
\[
 {R_0\rho\sigma z\over kS}={R_0q_0z\over S},\qquad
 {R_0\rho\sqrt\sigma\sqrt z w\over kS}
       ={(R_0\sqrt\lambda)\sqrt\rho\sqrt z w\over S}.
\]
They are infinitesimal by \eqref{eq:c1-unification-R0}; every other
coefficient visibly is too. Thus both entire Gamma normal terms
have infinitesimal coefficients on a single frame, including $D=0$.

It remains to handle a possibly signed $K_A$. The same construction
will work for every value of this source coefficient. Put
\[
 c=\max(0,1-K_A),\qquad a_0=k/\sigma,\qquad
 J_A=\sqrt z(u-k)W,\quad F_A=u(Z-a_0wW),
\]
\[
 E_u=\sqrt uJ_A,\quad E_k=\sqrt kJ_A,\qquad
 U_* =\sqrt{q_0}E_u,\quad K_* =\sqrt{q_0}E_k.
\]
One direct multiplication gives, for every state and all source sizes,
\begin{equation}\label{eq:c1-unification-US}
 \mathscr S=L_0F_A^2+2q_0E_u^2+q_0E_k^2
                                  +2q_0\sqrt uE_uR.
\end{equation}
Average this with $\mathscr S=\mathscr C-\lambda xR^2$:
\[
 \mathscr S=\tfrac12L_0F_A^2+U_*^2+\tfrac12K_*^2
 +\sqrt{q_0u}\,U_*R+\tfrac12\mathscr C-\tfrac12\lambda xR^2.
\]
The two cross/radial coefficients are infinitesimal. The term
$L_0F_A^2/2$ is retained as a positive square.

If $c>0$, \eqref{eq:c1-unification-KA} implies that $h$ is a positive
unit: in fact its residue is at least $1/4$. The inequality obtained
from \eqref{eq:c1-unification-e-bound},
\[
 a_0^2(hB_s-\sigma H_0)\le\mathcal H,
\]
then proves $a_0$ bounded. Also $L_0\sim2h$ is a positive unit,
so $\lambda$ is unbounded. The exact angular certificate proved
below is
\begin{equation}\label{eq:c1-unification-angular-budget}
 \Delta_A:=\varepsilon_C\mathscr C+u^2R^2
                         +\varepsilon_G\Gamma^2-A_\theta^2\in\Sigma T^2,
\end{equation}
where
\[
 \varepsilon_C=2/\lambda+u+k,\qquad
 \varepsilon_G=2a_0u+(u+k)/L_0.
\]
Both coefficients are infinitesimal whenever $c>0$. Subtract
$c(d_G/S)\Delta_A$ from the parent, keeping that nonnegative
quantity separately. The angular part of the remainder is
\[
 (K_A+c)A_\theta^2-c(\rho/S)A_\theta^2.
\]
The first coefficient is a positive unit. The second square is an
infinitesimal Gram correction on $(U_*,K_*)$, by the exact relation
\begin{equation}\label{eq:c1-unification-weighted-angular}
 \rho A_\theta^2=
       (\sqrt{a_0u}\,U_*+\sqrt{a_0k}\,K_*)^2.
\end{equation}
The other subtractions have coefficients
$c d_G\varepsilon_C/S$ on $\mathscr C$,
$c d_Gu^2/S$ on $R^2$, and $c\varepsilon_G/S$ on $\Gamma_*^2$;
all are infinitesimal. If $c=0$, none of these terms is used.
This verifies the same construction in both situations without a
new choice of roots or a source/test rate comparison.

Finally split the positive angular coefficient between $A_\theta$
and $T_\theta$, retaining the exact nonnegative difference
\[
 A_\theta^2-T_\theta^2=z(u-k)^2W^2(u^2+2uk)\in\Sigma T^2.
\]
The complete bounded Gram on
\[
 (N_x,N_y,R,\Gamma_*,A_\theta,T_\theta,d_1,\ldots,d_4,U_*,K_*)
\]
has positive-definite residue: its radial block is $J_C^2+J_\ell^2$,
the $N_x$ pivot is two, the Gamma pivot is one, the angular pivots
are positive, and the circuit and last two pivots are respectively
$1/2,\ldots,1/2,1,1/2$. All other entries have zero residue.
The remaining circuit coefficient
$1/2-cd_G\varepsilon_C/S$ has positive real residue $1/2$.
The one-dimensional bounded-Gram fact makes it a tensor SOS,
so its product with the retained $\mathscr C_+$ is also a tensor SOS.
Restore the retained positive squares and apply the bounded-Gram fact.

\paragraph{Proof of the angular certificate, including state signs.}
For $U,V,Z_*>0$ and positive separated $a_*$ set
\[
 \Psi=1+U^2V^2+U^2Z_*^2+V^2Z_*^2-4UVZ_*,\quad T_*=U^2-1,
\]
\[
 Q_*=VZ_*-U,\quad R_*=Z_*-UV,\quad G_*=U(V-Z_*).
\]
The polynomial identity
\begin{align*}
 &[2+a_*(U+1)]\Psi+U^2Q_*^2
       +[2U+(U+1)/a_*]R_*^2-T_*^2\\
 &=(T_*+UQ_*)^2+2Q_*^2
 +(U+1)(\sqrt{a_*}G_*-R_*/\sqrt{a_*})^2\\
 &\quad+a_*(U+1)(1-VZ_*)^2
       +2a_*(U+1)VZ_*(U-1)^2
\end{align*}
is verified by multiplication. Expand each $U+1$ factor before
taking square roots. The right side is a finite separated SOS.
For $W\ne0$ substitute
\[
 U=u/k,\quad V=w/(\sqrt\sigma\sqrt z),\quad
 Z_*={\sqrt\sigma\over k\sqrt z}{Z\over W},\quad
 a_*=\lambda k,\quad M_*=k^2\sqrt zW.
\]
Then $\mathscr C=\lambda M_*^2\Psi$,
$A_\theta=M_*T_*$, $R=k\sqrt zWQ_*$ and
$\Gamma=k\sqrt\sigma\sqrt zWR_*$. Multiplication by $M_*^2$
gives \eqref{eq:c1-unification-angular-budget}. These are legitimate
separated factors even when an intermediate ratio is unbounded.
If $Z_*<0$, apply the displayed identity first to $|Z_*|$.
Replacing $|Z_*|$ by $Z_*$ increases $\Psi$ by $8UV|Z_*|$ and
each of $Q_*^2,R_*^2$ by $4UV|Z_*|$. The three resulting additional
terms have positive separated coefficients and are SOS. Thus the
same certificate holds for either state sign. At $Z_*=0$ the original
identity applies directly. If $W=0$, $A_\theta=0$, so
\eqref{eq:c1-unification-angular-budget} follows immediately from
$\mathscr C\in\Sigma T^2$. No state coordinate is inverted there.

\paragraph{The case $\lambda\le0$.}
The definition of $\lambda$ gives
$H^2\le2bpz$ and $\beta\theta^2/\sigma^2\le bpz$.
Thus $H$ is bounded, $h$ is infinitesimal, $\theta\sim1$,
$M_\beta=\beta\theta/\sigma^2$ is bounded and
$H-H_0,d_G-2H_0$ are infinitesimal. In particular
\[
 L=\lambda/\sigma
   =8(D/\sigma^2)h+4M_\beta\theta+2H^2-4bpz
\]
is bounded. Direct multiplication gives
\[
 u^2\Gamma+V=\sigma uP_\beta,
\]
\[
 A_\theta^2+u\Gamma P_\beta+uA_\theta R
      =\sigma P_\beta^2+(x/\sigma)(u^2-x)(w^2-\sigma z)W^2,
\]
where $V=uw(u^2-x)W$. Consequently
\[
 K_AA_\theta^2+\mathscr S
 =4\kappa H A_\theta^2+Lu^2\Gamma^2+L\Gamma V
                                      +\lambda uA_\theta R.
\]
Add $1+L^2$ times the exact zero
\[
 V^2-\sigma u^2A_\theta^2-u\sigma(u^2-x)A_\theta R
                                      +(u^2-x)V\Gamma=0.
\]
The $(\Gamma,V)$ real residue block is
\[
 \begin{pmatrix}2\bar H_0&\bar L/2\\
                 \bar L/2&1+\bar L^2\end{pmatrix},\qquad
 \det=2\bar H_0+(2\bar H_0-1/4)\bar L^2>0.
\]
The $N_x$, radial and angular blocks are strict; every remaining
cross coefficient, including both complete Gamma normal terms,
is infinitesimal. The bounded-Gram fact proves this case, including
$\lambda=0$.

All constructions are literal on the testing axes. No inverse of
$D$ or $\beta$ occurs, so both permitted source faces are included.
The parent is real quadratic in the state; applying the proof to
the real and imaginary parts proves the complex-state assertion.
This completes the theorem.
\subsection{The six-root identity and unit-Alice angular relation}

\label{component:pb1104217f67b}

\label{pb1104217f67b:sec:r13-joint-radial-beta-zero}

\paragraph{Quantified domain and conclusion.}
Retain the original physical C1.4 source and all its residual-complement
conditions. Suppose $x=k^2$, $z=L^{-1}$ and $\sigma=\eta^{-1}$ are
positive infinitesimals, $0<b<1$, and the original coefficient $\beta=0$.
There are no further source hypotheses. In particular $h/\sigma\sim1$
and $D/\sigma^2\sim1/2$ follow from BE3--BE7, rather than being imposed.
The conclusion below covers every bounded independent real test $u$,
every infinitesimal or zero independent reciprocal test $w$, all their
signs, and real or complex states. The previously proved remote
complement supplies the other tests on this large-$L$ source chart.
Thus this proves the literal $\beta=0$ part of C14.1. The adjoining
Section~\ref{p909aa5765721:sec:r13-beta-positive-extension} extends the actual certificate
to all sources with $\beta\theta/\sigma^2$ bounded; the present result
is its beta-zero corollary and supplies the residue calculation used
there. C14.2 is not included.

The calculation joins the entire radial scalar quotient with every
radial first normal \emph{before} deciding boundedness. Their exact
cancellations remove the BR3 pole, including its separate
$8(ab)b/\sigma$ summand through $Q=1-b$ infinitesimal.

\paragraph{Notation and the original contact identities.}
Use the complete reciprocal parent \eqref{eq:c1-reciprocal-parent} and the large-$L$ source
notation. Thus $m=\mu x^2z$, $C=cz$, $A=a/z$, and
$\alpha=\ell/z$. In expressions below $p=ab$ and $l=\ell Q$ are
the two opening-cleared products; $p$ is not the source parameter $b$.
Set $t=u^2$, $\xi=w^2$, $\xi_0=\sigma z$, $\delta=\xi-\xi_0$,
$g=h-Q$, $D_u=1+t$, $D_w=1+\xi$, and $q=\sqrt z$.
The source identities established in the handoff give
\[
 \phi=\phi_t=\phi_\xi=L_x=L_y=L_z=0
 \quad\hbox{at }(t,\xi)=(x,\xi_0),
\]
as well as $\phi(0,\xi)=4m$ and $L_x(0,\xi)=-8m$.
Here $\phi,L_x,L_y,L_z$ are the \emph{complete} EN2 polynomials,
with $D_*=D$, $\beta_*=0$ and $\beta=0$.

Choose the fixed caps $a_* = \ell_* =1/4$. They are smaller than
the actual $a,\ell$, since $p,l\sim1/2$ and $0<b,Q<1$.
Retain the two exact nonnegative centered excesses
\[
 E_{\rm cap}=4(a-a_*)D_uD_wN_y^2+4(\ell-\ell_*)D_uN_y^2.
\]
Define six separated roots
\[
\begin{aligned}
N_x&=X+uhZ,&N_y&=(Y-bwZ)/q,\\
\Gamma&=\sigma Z-uwW,&R&=wZ/q-uqW,\\
A_\theta&=q(t-x)W,&V&=uw(t-x)W.
\end{aligned}\tag{JR1}
\]
No test/source quotient is used in their definition. Their inverses
and radicals are pure source.

\paragraph{Exact quotient specification.}
The contact identities give the exact divisions
\begin{align*}
 \phi&=z(t-x)^2A+t(t-x)\delta B+t\delta^2C_r,\\
 L_x&=(t-x)(tA_x+zB_x)+t\delta C_x,\\
 L_y&=(t-x)A_y+\delta B_y,\\
 L_z&=(t-x)A_z+\delta B_z.
\end{align*}
All quotients are monic polynomial divisions in one testing variable.
Their complete support can equivalently be specified by
\begin{gather*}
 A=4\mu+\frac{2h^2(2D+1)}\sigma t,\qquad
 A_x=-\frac{4h(2D+1)}\sigma,\qquad B_x=8\mu x,\\
 A_y=8(p-l)+4z\{(1+2m)g+2Dh+2p\sigma\},\\
 S:=A_z/\sigma-B=S_0+
       \frac{2h^2(2D+1)}{\sigma^2}(t-x).
\tag{JR2}
\end{gather*}
The same monic divisions give the much shorter exact expression
\[
 S_0=\frac{2\{2h^2(2D+1)x+4Dh-2bp\sigma^2z+h^2\}}{\sigma^2}.
\tag{JR2a}
\]
Thus $S_0\sim2$ follows immediately from $h/\sigma\sim1$ and
$D/\sigma^2\sim1/2$. This removes any dependency on the longer R10
reduction modulo the original beta numerator. The source contact jets
are still needed for the quotient remainders to vanish.
For a denominator-free quotient specification, $C_r$ is obtained by
dividing $\phi$ twice by $\xi-\xi_0$ and then dividing the quotient
by $t$; $B$ is the quotient of
$t^{-1}\phi_\xi(t,\xi_0)$ by $t-x$; and $A_z$ is the quotient of
$L_z(t,\xi_0)$ by $t-x$. These divisions introduce no mixed inverse.

The remaining exact coefficients are most compactly written as
\begin{align*}
 K_r:=\sigma^2C_r={}&k_0+k_1t+k_2\xi,\\
 k_0={}&2b\{(2C+1)b+4D+4p\sigma
             +2(2C+2D+1)\sigma(bz-1)\},\\
 k_1={}&2(1+2m)g^2+8Dhb+8pb\sigma,\\
 k_2={}&2b^2(2C+2D+1),\\
 R_x:=\sigma C_x={}&-4\{2Db+(1+2m)g\},\\
 R_y:=\sigma B_y={}&4\{t[2Dh+2p\sigma+(1+2m)g]
      +b(2C+2D+1)(\xi+\sigma z)\\
 &\hspace{23mm} +(2C+1)(b-\sigma)+2D(1-\sigma)+2p\sigma\},\\
 E_r={}&2(C+D)(bz-1)+2p+bz-1,\qquad
 R_\gamma=-4(E_r+2pt),\\
 Z_0={}&\frac{4D+2t(4Dh+h^2)+2t^2h^2(2D+1)}{\sigma^2},\\
 Z_x={}&-\frac{4u}{\sigma}\{2D+h+th(2D+1)\}.
\tag{JR3}
\end{align*}
In particular all of these coefficients are bounded when $u$ is bounded
and $w$ infinitesimal. The $\sigma^{-1}$ and $\sigma^{-2}$ in $Z_0,Z_x$
are controlled by the \emph{proved pure-source} bounds on $h,D$.

\paragraph{The joint cancellation, including all first normals.}
The exact full normal block EN3--EN4 obeys
\begin{align*}
 H_z/\sigma^2+\xi C_r-(\xi/\sigma)B_z&=Z_0,\\
 B_z-2\sigma C_r&=bR_\gamma,\\
 \frac4\sigma\{2DT_{\rm test}+bD_w(1+2C)+tg(1+2m)\}-B_y&=R_\gamma,\\
 \frac{4u}{\sigma}\{-2DT_{\rm test}-hD_u-(1+2m)\xi g\}
                         -u\xi C_x&=Z_x.
\tag{JR4}
\end{align*}
Here $T_{\rm test}=1+th+b\xi$. These are polynomial identities after clearing
pure-source powers of $\sigma$. No coefficient is discarded: $K_r$
contains the full BR3 term, and the last three lines keep its radial
couplings. The cancellation explains why a scalar-only extraction of
BR1 followed by boundedness checking cannot see this certificate.

Put
\begin{align*}
 H_x&=2(1+t+\xi)+4Dt+4m\xi,\\
 H_y&=4D_u(a_*D_w+\ell_*)+4mtz+4CzD_w+(4D+2)\xi z+2D_uz,\\
 c_G&=4D_u\{pbD_w+lQ\},\qquad
 c_N=8D_u\{pD_w-l\}.
\end{align*}
The \emph{entire original parent}, without restriction to a state
subspace, is exactly $F_w=E_{\rm cap}+G$, where
\begin{align*}
 G={}&H_xN_x^2+H_yN_y^2+Z_0\Gamma^2
        +(c_G+zK_r)R^2+AA_\theta^2\\
 &-4uqw(2m+2D+1)N_xN_y+Z_xN_x\Gamma
       +qwR_\gamma N_y\Gamma+qwbR_\gamma R\Gamma\\
 &+uqwR_xN_xR+(c_N+zR_y)N_yR
       +qwB_xN_xA_\theta+uA_yN_yA_\theta\\
 &+uA_xN_xV+u\sigma B A_\theta R+S\Gamma V.
\tag{JR5}
\end{align*}
For example, the scalar radial conductor is
$t\delta^2W^2=\sigma^2(qR-w\Gamma/\sigma)^2$;
the radial $Z$ first normal is
$uw\delta W(\Gamma/\sigma)=qwR\Gamma-(\xi/\sigma)\Gamma^2$.
Their joint use gives the first two identities of JR4. The other
conductors are those of the verified large-$L$ proof with
$N_z=\Gamma/\sigma$. These identities prove JR5 exactly and include
every first normal and every hard quadratic coupling.

\paragraph{One bounded positive Gram for every bounded Alice test.}
All entries of JR5 are bounded separated tensors. Indeed
\[
 h/\sigma\sim1,\quad D/\sigma^2\sim1/2,\quad
 \mu\sim1/2,\quad C=O(z),\quad p,l\sim1/2,
 \quad p-l\in z\mathcal O_{\rm source},
\]
and $\sigma B=A_z-\sigma S\sim4$ uniformly for bounded $u$.
The displayed formula for $A_z$ obtained from $L_z$ has only the
bounded terms $h^2/\sigma$, $Dh/\sigma$, $p,l$, and bounded multiples
of $z$; hence this last assertion does not order any mixed rate.

For completeness, the possibly needed opening-uniform divisibility
$p-l\in z\mathcal O_{\rm source}$ follows directly from the exact
source formulas, not from a residue subtraction. Write $H=h/\sigma$.
The two numerators over their common unit denominator $2P$ are
\begin{align*}
 U_A={}&z(1+x\sigma H)(1-xH)+(1-bz)(1+x\sigma Hz),\\
 U_\ell={}&(1+\sigma bz)\big[xH(1+xbz)(1+\sigma z)
       +\{1+xbz-xz(1-\sigma H)\}(1-xH)\big].
\end{align*}
They are bounded polynomials in the source letters and both equal
one at $z=0$. Consequently $(U_A-U_\ell)/z$ is a bounded polynomial.
Neither numerator has a $b$ or $Q$ denominator. This proves the claimed
divisibility through both openings and is the source argument in EN6's
adjoining radial paragraph.

The exact syzygy
\[
 V^2=\sigma t A_\theta^2+u\sigma(t-x)A_\theta R-(t-x)V\Gamma
\tag{JR6}
\]
allows the addition of zero times the pure-testing function
\[
 \lambda(t)=\frac{2(t^2+t+1)}{(1+t)^2}>0.
\]
Specifically add
$\lambda[V^2-\sigma tA_\theta^2-u\sigma(t-x)A_\theta R+(t-x)V\Gamma]$
to JR5. There is no source/test denominator in $\lambda$.

At any real residue $u_0$, write $t_0=u_0^2$. The $N_y$ pivot is
$2(1+t_0)$. The $(A_\theta,R)$ block is
\[
 \begin{pmatrix}2&2u_0\\2u_0&2(1+t_0)\end{pmatrix},
 \qquad \det=4.
\]
The remaining $(N_x,\Gamma,V)$ block is
\[
 \begin{pmatrix}
 2(1+t_0)&-2u_0(1+t_0)&-2u_0\\
 -2u_0(1+t_0)&2(1+t_0+t_0^2)&1+t_0+\lambda(t_0)t_0/2\\
 -2u_0&1+t_0+\lambda(t_0)t_0/2&\lambda(t_0)
 \end{pmatrix}.
\tag{JR7}
\]
Its first two leading pivots are positive, with leading $2\times2$
determinant $4(1+t_0)$. The last Schur complement is exactly
\[
 \lambda(1+t_0^2)+\frac{2t_0^2}{1+t_0}
       -\frac12(1+t_0+\lambda t_0/2)^2
   =\frac{4t_0+3}{2(1+t_0)^4}>0.
\tag{JR8}
\]
Thus the residue Gram is strictly positive definite for every finite
real $u_0$, including zero and both signs. Crucially it does not depend
on the residue of $b$: the result remains strict through $b\to0$
and $Q\to0$. The finite bounded positive-Gram criterion already used
by the project then gives the exact legal tensor-SOS.

\paragraph{Boundaries and scope.}
JR1--JR6 are identities, including $u=0$ and $w=0$; no root is inverted
at a zero. Their sign coverage is literal. Apply the real identity to
the real and imaginary parts of the state for complex states.
The exact remote join is the common-parent theorem in
Section~\ref{component:p89fd3edb3edd},
Section~\texttt{sec:r6-c1-common-remote}. First, $A,\alpha$ are
unbounded on this source, so its adjacent-event theorem transported by
(CE9)--(CE10) covers every $v$ that is not projectively unbounded.
For projectively unbounded $v$ and projectively unbounded $u$, put
$p_t=1/u$ and apply the complete raw frame (CE20), whose strict
residue is (CE21). Its exact source premises (CE17), without the
unneeded $m$ lower bound in this $p_t$ chart, hold with $q=\sqrt z$:
$C\sim z/2$, $A\ge1/(4z)$, and $\alpha\ge1/(4z)$.
The $C$ lower bound may likewise be chosen as $z/4$. The unselected
$D,\beta$ events are retained. The identities include literal
$p_t=0$, $w=0$ and all signs. These two remote cases and JR1--JR8
exhaust all projective tests on the printed source face. No source
domain with $L$ a unit, or $\beta>0$, is covered by this proof.

\paragraph{Corroboration.}
 The
original six contact identities are explicitly identified imported
results, not silently replaced by numerical residue checks.



% Source formulas: p909aa5765721
\subsection{Bounded $H$ at unit Alice testing angle}

\label{component:p909aa5765721}

\label{p909aa5765721:sec:r13-beta-positive-extension}

Retain the same original C14.1 source as JR1--JR8, with $x,z,\sigma$
positive infinitesimals, but allow $\beta\ge0$. Write
\[
 b_\theta=\beta\theta,\qquad H=h/\sigma,\qquad
 M_\beta=b_\theta/\sigma^2.
\]
This section proves two precise extensions. First, if $M_\beta$ is
bounded, the \emph{entire source domain is covered for every test}.
This includes the beta-zero theorem. Second, if $H$ is bounded but
$M_\beta$ is unbounded, every bounded nonzero-unit $u$ is covered
at the residual infinitesimal-$w$ tests. On this second source domain
only $u$ infinitesimal remains. Neither statement assumes $d_0$ a
unit: the second includes $D=0$. The $H$-unbounded domain and C14.2
are not covered here.

\paragraph{The generic joint radial identity.}
The exact six-root identity JR5 continues to hold if its coefficients
are given the following beta corrections:
\begin{gather*}
 k_1\ \mapsto\ k_1+4b_\theta\theta,\qquad
 R_x\ \mapsto\ R_x+8b_\theta,\\
 Z_0\ \mapsto\ Z_0+\frac{4b_\theta\theta t}{\sigma^2},\qquad
 Z_x\ \mapsto\ Z_x+8ub_\theta/\sigma,\\
 S\ \mapsto\ S+4b_\theta\theta/\sigma^2.
\tag{JB1}
\end{gather*}
Here the uncorrected formulas on the right are the polynomials in
JR2--JR3 evaluated at the actual $h,D,m,p,l$, not at a beta-zero
source. The scalar $B$ remains its full monic quotient. The additional
retained centered excess is $4\beta D_wN_x^2$. The complete generic
coefficient is, explicitly,
\[
 S=\frac{2\{h^2(2D+1)(t+x)+4Dh-2bp\sigma^2z
                       +2b_\theta\theta+h^2\}}{\sigma^2}.
\tag{JB2}
\]
All four cancellations JR4 are exact with these changes. They were
checked symbolically with $b_\theta$ free. Thus no radial pole remains
unallocated in this extension.

\paragraph{The source estimates needed by the extension.}
The source equation BE3 gives exactly
\[
 H=H_0/B+\frac{2P}{B}\frac{b_\theta}{\sigma}.
\tag{JB3}
\]
Consequently $M_\beta$ bounded implies $H\sim1$ and $\theta\sim1$.
The sphere budget then implies $D/\sigma^2\sim1/2$, since
$xb_\theta/\sigma^2=xM_\beta$ is infinitesimal. Also $\mu\sim1/2$.

More generally, suppose only $H$ bounded. Then $\res H\ge1$;
no exact inequality $H\ge1$ is imposed. One has $\theta\sim1$,
and $b_\theta/\sigma$ is bounded. Write
\[
 d=\res(D/\sigma^2),\qquad
 c=\res(xb_\theta\theta/\sigma^2).
\]
BE4--BE6 give $d,c\ge0$ and $d+c=1/2$. Formula JB3 gives
$\res(b_\theta/\sigma)=(\res H-1)/2$. The bounded source formulas
also give $\mu\sim H/2$, $p,l\sim1/2$, $C=O(z)$, and
$p-l\in z\mathcal O_{\rm source}$. These statements retain both
source-sphere endpoints and both $b,Q$ openings.

\paragraph{An exact square joining the beta pole to its angular reserve.}
Suppose first $u$ is a bounded nonzero unit, so that $1/u$ is an
allowed bounded pure-testing inverse. Put
\[
 a_\beta=b_\theta\theta,\qquad
 C_\beta=4a_\beta/\sigma^2,\qquad T_c=t\Gamma+V,
 \qquad P_\beta=uZ-(x/\sigma)wW.
\]
The identity $T_c=\sigma uP_\beta$ has no $u$ denominator.
Adding $(C_\beta/t)$ times JR6's zero, and completing the resulting
$T_c$ square, gives the exact entire-parent identity
\[
 G=4a_\beta P_\beta^2+\widehat G,
\tag{JB4}
\]
where $G$ is the complete generic JR5 and
\begin{align*}
 \widehat G=G
 &-C_\beta t\Gamma^2
 -(C_\beta+C_\beta x/t)\Gamma V
 -C_\beta\sigma A_\theta^2\\
 &-C_\beta\sigma(t-x)A_\theta R/u.
\tag{JB5}
\end{align*}
Every coefficient of $\widehat G$ is bounded when $H$ is bounded
and $u$ a unit. The only apparent angular reserve subtraction is
$4a_\beta/\sigma$. Its exact source comparison is
\[
 \mu-a_\beta/\sigma
 =\kappa H_0/B+
    \frac{b_\theta}{\sigma}\left(\frac{2\kappa P}{B}-\theta\right),
\tag{JB6}
\]
where $\kappa$ is the BE8 source factor. The bracket is infinitesimal,
so this difference has residue $1/2$. Thus the beta pole is paid for
by the actual shared $m$ term. No leading coefficient is substituted
in the exact identity JB4.

\paragraph{Bounded $M_\beta$: all residual tests.}
When $M_\beta$ is bounded and $u$ a unit, the residue of $\widehat G$
is exactly the beta-zero residue in JR7 before the finite $V$ pivot:
$H\sim1$, $d=1/2$, $c=0$. Add the same pure-testing pivot
$\lambda(t)=2(t^2+t+1)/(1+t)^2$ and apply JR8. Restore the positive
square in JB4 and every centered excess. This covers every unit $u$.

\paragraph{Unbounded $M_\beta$, bounded $H$: all unit-Alice tests.}
For any real nonzero residue $u_0$, set $t_0=u_0^2$. In the notation
of JB3, and writing $H$ now for its finite real residue, the only
nontrivial $\widehat G$ residue block is
\begin{align*}
 Q={}&2(1+t_0)N_x^2+[4d+2H^2t_0(1+t_0)]\Gamma^2\\
 &-4u_0(1+Ht_0)N_x\Gamma-4Hu_0N_xV\\
 &+[2H^2(1+t_0)-4c/t_0]\Gamma V.
\tag{JB7}
\end{align*}
Its angular $(A_\theta,R)$ block is the strict matrix from JR7 and
its $N_y$ pivot is $2(1+t_0)$. All other residue cross entries vanish.
The beta square in JB4 equals $(C_\beta/t)T_c^2$, with
$C_\beta$ unbounded. On its exact residue kernel $V=-t_0\Gamma$,
JB7 reduces to
\[
 2(1+t_0)N_x^2-4u_0N_x\Gamma+2\Gamma^2,
 \qquad \det=4.
\tag{JB8}
\]
Here the cancellation uses exactly $d+c=1/2$. It remains strict at
$d=0$ and $c=0$. The row $T_c=V+t\Gamma$ has a bounded separated
right inverse, obtained by changing $V$ alone. The elementary finite
penalty lemma therefore chooses a positive real $K$ such that
$\widehat G+(K/t)T_c^2$ has strictly positive residue. This is possible
because $t_0>0$. The exact retained excess is
$((C_\beta-K)/t)T_c^2$: $C_\beta-K$ is a positive pure-source
coefficient, and its square root separates from the pure-testing
$1/\sqrt t$. No ordering of $C_\beta/t$ minus a real number is used.
This proves the entire
parent for all bounded nonzero-unit $u$, with infinitesimal or zero
$w$, all signs and complex states.

The new exact remaining obstruction on the bounded-$H$ source branch
is therefore $M_\beta$ unbounded and $u$ infinitesimal. JB5 has
$x/u$ and $x/u^2$ terms there, so its bounded-Gram argument cannot be
imported. The square $P_\beta$ itself is still legal and has no
$u$ inverse; a joint allocation with the remaining angular terms is
required to remove those coefficient denominators without comparing
independent source and testing rates. No such allocation is claimed.

\paragraph{Denominators and coefficient bounds.}
Every source inverse in the roots or exact coefficients is a pure-source
inverse. On the large-$L$, bounded-$H$ domains, $P,J,B$ and the original
geometry coefficient $M=1+xbz-xz\theta$ are units; $1-r$ and $\theta$
are units as well. Possible $b,Q$ inverses occur only in the original
source weights and the retained positive centered excesses, while their
opening-cleared products $p,l$ are bounded. The normalized quantities
$H$, $D/\sigma^2$, $b_\theta/\sigma$, and
$xb_\theta\theta/\sigma^2$ are bounded by JB3 and the BE sphere.
The corrections in JB5 therefore leave bounded coefficients: the
remaining $\Gamma^2$ terms involve $D/\sigma^2$, $Dh/\sigma^2$ and
$H^2$; the $\Gamma V$ coefficient involves those same bounded source
products and $(xb_\theta\theta/\sigma^2)/t$; the angular-radial
coefficient involves $b_\theta\theta/\sigma$ times $(t-x)/u$.
The factors $1/u,1/t$ are used only when $u$ is a testing unit.
All $u=0$ and infinitesimal-$u$ assertions for bounded $M_\beta$
instead use the unaltered generic six-root identity and its bounded
pure-source pivot $(1+2M_\beta)^2$. The everywhere bounded testing
denominator $1+t=1+u^2$ in JR8 is also pure testing. No mixed inverse,
mixed radical, or source/testing rate comparison occurs in either
covered domain or its remote join.



% Source formulas: pcc84202479d2
\subsection{Unbounded $H$ at unit Alice testing angle}

\label{component:pcc84202479d2}

\label{pcc84202479d2:sec:r18-sphere-unit-u}

Retain the entire original C14.1b source intersection and the exact
SP6 parent from the common identity.  Assume $u$ is a bounded nonzero unit
and $w$ is infinitesimal.  The argument below proves the whole parent
for every testing state, including the remaining
$W/Z$-unbounded states.  It does not cover infinitesimal $u$.
All source gaps, $D=0$, both source sphere endpoints and the literal
coefficient boundaries remain unchanged.

\paragraph{No state chart is needed.}
All roots below are homogeneous polynomials in the original testing
state, with pure-source coefficients.  In particular the proof
uses no state-coordinate inverse, no bound on $W/Z$, and no sign
restriction.  The only testing inverse is that of the unit $u$.
Write $U=u^2$, $s=x/\sigma$, and preserve the exact SP/BS coefficients
\[
 L_h=\lambda\sigma/h,\quad \rho=\lambda x/\sigma,
 \quad d=4D/\sigma^2,\quad K_A=4\kappa h/\sigma-\lambda.
\]
The source identities give $s\in\mathrm{Inf}$,
$\sigma/\sqrt h\in\mathrm{Inf}$, $L_h\sim2$, $\nu\sim1$, and
\[
 d+\rho\sim2,\quad K_A-(2-4h)\in\mathrm{Inf},\quad
 \chi-(4-4h)\in\mathrm{Inf},\quad
 (c_G+zK_r)-2(1+U)\in\mathrm{Inf}.
 \tag{SU2}
\]
The last infinitesimal is a finite separated expression for bounded
$u$ and infinitesimal $w$.  For the only new source assertion in
SU2, the sphere identities give exactly
\[
 d=\frac{2\mathcal H}{PJ}d_0^2,\qquad
 \rho=\frac{2xH_0}{\sigma B}
       +\frac{2\mathcal H}{B}b_0^2+\lambda_0s.
 \tag{SU3}
\]
Thus $d+\rho\sim2$ follows from $d_0^2+b_0^2=1$, with no inverse
of either sphere coordinate.  All other bounds follow from the exact source cancellations in
equations~\eqref{eq:c1-unification-nu}--\eqref{eq:c1-unification-KA} and the source estimates,
including $p-l\in zB_c$, in Section~\ref{sec:c1-unification-bounded-z-all-H}.

\paragraph{Shift the troublesome angular product into the normal block.}
Set $A_0=qW$, $A=A_\theta=(U-x)A_0$, and retain $R$.
There is an exact conductor:
\[
 A(A+uR)=u(U-x)wWZ-x(U-x)A_0^2.
 \tag{SU4}
\]
Consequently, collecting the radial terms in \eqref{eq:compression-full-parent} gives
\[
\begin{split}
 K_AA^2+\chi uAR+(c_G+zK_r)R^2
 ={}&2(A+uR)^2+2R^2-4hu(U-x)wWZ\\
 &+4hx(U-x)A_0^2+\delta_A A^2+
   \delta_\chi uAR+\delta_RR^2,
\end{split} \tag{SU5}
\]
where the three $\delta$ coefficients are the infinitesimals in
SU2.  This step keeps the signed product exactly; it is not an
inequality based on the sign of $U-x$.

\paragraph{A complete strict normal block.}
Introduce the legal separated roots
\[
 H=\sqrt h\,uZ,\qquad S=uwW-hu^2Z,\qquad N=N_x,
 \qquad Z_g=uwW=S+u\sqrt h\,H.
\]
Their exact allocations are
\[
\begin{split}
 \Gamma&=-S+\left(\frac{\sigma}{u\sqrt h}-u\sqrt h\right)H,\\
 P_h:=\sqrt h P_\beta&=(1-sh)H-\frac{s\sqrt h}{u}S,\\
 X_e&=N-\sqrt h P_h.
\end{split} \tag{SU6}
\]
Every coefficient is bounded; the factors $u^{-1}$ are pure testing
units.  The full scalar term $\mathscr S$ in \eqref{eq:compression-full-parent} also satisfies exactly
\[
 \mathscr P=L_hP_h^2+
      \rho(1-x/U)Z_g^2-\rho\sigma(U-x)A_0^2.
 \tag{SU7}
\]
The apparently signed factor $1-x/U$ is simply a bounded Gram
coefficient of residue one.  No mixed inverse is used.

Combine the $wWZ$ term in SU5 with every $N,\Gamma,P_h,Z_g$ term:
\[
\begin{split}
 \mathscr N={}&[2(1+w^2)+4mw^2]N^2
     +2(2D+1)U X_e^2+d\Gamma^2-4u\nu N\Gamma\\
 &+L_hP_h^2+\rho(1-x/U)Z_g^2
       -4\sqrt h\,\frac{U-x}{u}\,Z_gH .
\end{split} \tag{SU8}
\]
This is an exact equality for the entire indicated normal block.
Its coefficients in $(N,S,H)$ are bounded separated expressions.
Using SU2 and SU6, its residue is exactly
\[
 2(1+u_0^2)N^2+4u_0NS+2S^2+2H^2.
 \tag{SU9}
\]
The $N,S$ determinant is $4$, so this residue is strictly positive.
The $H$ cross terms cancel: the $X_e$ square cancels the leading
$N\Gamma$ coupling to $H$, and the angular shift in SU5 cancels
the remaining $S H$ and $hH^2$ terms.  This is where the sphere
sum $d_0^2+b_0^2=1$ is essential.

The remaining angular/radial block is
\[
\begin{split}
 \mathscr A={}&2\{(U-x)A_0+uR\}^2+2R^2\\
 &+(4hx-\rho\sigma)(U-x)A_0^2
   +\delta_A(U-x)^2A_0^2
   +\delta_\chi u(U-x)A_0R+\delta_RR^2.
\end{split} \tag{SU10}
\]
Its residue on $(A_0,R)$ has determinant $4u_0^4>0$.

\paragraph{Every other original-parent term.}
The unchanged hard $N_y$ diagonal has residue $2(1+u_0^2)$.
All non-Gamma cross terms of \eqref{eq:compression-full-parent} already have
infinitesimal coefficients on $(N,N_y,A_0,R)$, by SP's complete
opening-uniform coefficient bounds and $A=(U-x)A_0$.
The two $\Gamma$ terms have a factor $qw$ and become infinitesimal
crosses after the bounded allocation SU6.  Thus the complete Gram
on $(N,N_y,S,H,A_0,R)$ has strictly positive residue, the direct
sum of SU9, the $N_y$ entry and SU10.

The remaining exact nonnegative rows are retained once: both
centered cap excesses, $4\beta D_wN^2$, and
$4bpzU\Gamma^2$.  The $d\Gamma^2$ part and
$2(2D+1)(uX_e)^2$ have already been used in SU8 and are not
retained a second time.  All scalar, hard and first-normal terms
of the original SP6 parent are therefore included.

The bounded separated-Gram criterion proves tensor SOS for all real
states and both testing signs, including literal $Z=0$ and $w=0$.
Applying the real identity separately to real and imaginary parts
gives all complex states.  Neither source sphere endpoint is inverted.
The precise remaining C14.1b testing domain has $u$ infinitesimal
(with the earlier state and sign certificates still retained).



\subsection{Unit $L$: noninfinitesimal Alice tests}

\label{component:pb0d9d28d0849}

\label{pb0d9d28d0849:sec:r18-c142-remote}

\begin{proposition}
Keep the unchanged physical C14.2 source: $d=1$, $x>0$
infinitesimal, $L=x\zeta$ a positive bounded unit, $\eta$ unbounded,
all original strict gaps, nonnegative event weights, total-scale and
earlier-exclusion intersections.  Then the complete original parent
is finite tensor SOS for every noninfinitesimal Alice angle $u$,
including $u=\infty$, every independent Bob angle, and all real or
complex states.  Literal $D=0$ and $\beta=0$ are retained.  Thus
only infinitesimal or zero $u$, with arbitrary Bob type, remains in
C14.2.  No entire-source certificate is asserted.
\end{proposition}

\paragraph{Opening-safe source input.}
Use the source estimates in
Section~\ref{sec:c1-unification-bounded-z-all-H}, including
the exact positive source inequality~\eqref{eq:c1-unification-e-bound}
and its sphere and opening-cleared products.
Put $e_r=1-r$, $Q=1-Lr>0$ and define positive pure-source weights
\[
 C_0=\frac1{2Le_r},\qquad
 A_0=\frac{1+Le_r}{2Lr},\qquad
 a_0=\frac L{2Q}.
                                                               \tag{LA1}
\]
Their actual-weight ratios satisfy
$C/C_0,A/A_0,\alpha/a_0\sim1$ uniformly through $r$, $e_r$,
and $Q$ openings.  Each leading and actual weight has a positive
real lower bound, since $L$ is a positive bounded unit and
$0<r,e_r,Q<1$.  No upper bound for the individual weights is assumed.
The BE source sphere and both its endpoints remain unchanged; no
coefficient $D$, $\beta$, $d_0$ or $b_0$ is divided by.

\paragraph{The reciprocal Bob strip with a unit Alice angle.}
Let $u$ be a bounded nonzero unit and $w=1/v$ infinitesimal or zero.
Use four legal raw roots
\[
 P=uW,\quad X,\quad Y,\quad T=wZ,
 \qquad U=u^2,\quad D_u=1+U,\quad D_w=1+w^2.
\]
The original signed core is given in Section~\ref{component:p89fd3edb3edd}.
Written in these roots it has the exact split
\[
 K_w=K_0+w^2\{(P-Y)^2+X^2\}
                      +2uwX(T+P-Y),                         \tag{LA2}
\]
\[
 K_0=D_uX^2+D_uY^2+UT^2-2UTY-2PT.
\]
This is the complete core, including its sign; no state graph is used.
The selected positive events are
$CD_w(P-Y)^2$, $AD_uD_wY^2$, and
$\alpha D_u(Y-T)^2$ in the quarter-parent.  The whole original
$m,D,\beta$ terms are retained separately.

Direct expansion using LA1 gives the following exact identity for
the leading selected form at $w=0$:
\[
\begin{split}
 Q_0={}&C_0(P-Y)^2+A_0D_uY^2+a_0D_u(Y-T)^2+K_0/2\\
 ={}&\frac{D_u}{2}X^2
       +C_0\{P-Y-Le_rT\}^2
       +(A_0+a_0+1/2)\{Y-LrT\}^2\\
 &+U\{A_0Y^2+(a_0+1/2)(Y-T)^2\}.
\end{split}                                                   \tag{LA3}
\]
Every displayed coefficient is nonnegative pure source or a positive
pure-testing factor.  The source coefficients in LA3 can be clipped
at positive real lower bounds, retaining their exact positive square
excesses.  The resulting bounded Gram on $(P,X,Y,T)$ has strictly
positive residue: $X$ is controlled directly, the two $U$ terms
control $Y,T$ because $u$ is a unit, and the first contact square
then controls $P$.  The coefficients $Le_r,Lr$ in these roots are
bounded pure source, including all original openings.

\paragraph{Exact common-factor accounting for the actual weights.}
To transfer LA3 to the actual parent, define solely in the source field
\[
 \delta=|C/C_0-1|+|A/A_0-1|+|\alpha/a_0-1|,
 \qquad f=1-\delta.
\]
Then $\delta$ is nonnegative infinitesimal, $f>0$ is a unit,
and each of $C-fC_0$, $A-fA_0$, $\alpha-fa_0$ is nonnegative.
Retain their original event terms as exact squares.  The complete
remaining selected polynomial is, identically,
\[
\begin{split}
 S={}&f\{C_0D_w(P-Y)^2+A_0D_uD_wY^2
                          +a_0D_u(Y-T)^2\}+K_w/2\\
  ={}&fQ_0+fw^2\{C_0(P-Y)^2+A_0D_uY^2\}\\
    &+(1-f)K_0/2+(K_w-K_0)/2.
\end{split} \tag{LA4}
\]
The second line's $w^2$ terms are retained positive squares with
their complete actual source scales.  The third line, including the
signed $(1-f)K_0/2$ term, is a bounded infinitesimal Gram on the
four raw roots, by LA2 and $1-f=\delta\in\operatorname{Inf}$.
Clip the completed source-square coefficients of $fQ_0$ at positive
real lower bounds as above.  Its retained bounded Gram remains
strict under this entire infinitesimal perturbation.  The separated
bounded-Gram lemma gives tensor SOS; restoring all square excesses
proves the whole original parent.  This is exact polynomial
accounting in the tensor ring, not an inequality that orders mixed
source and testing quantities.

\paragraph{The other Bob types and the unbounded Alice endpoint.}
If $u$ is a bounded nonzero unit and $v$ is a bounded nonzero unit,
use the common noncommuting-residue theorem on the unchanged
physical parent.  If $u$ is a bounded nonzero unit and $v$ is
infinitesimal or zero, clip the original $C,A,\alpha$ at positive
real lower bounds $c_*,a_*,a'_*$ and retain every excess and the
whole original $m,D,\beta$ terms.  The complete selected residue is
\[
 (c_*+1/2)(uW-Y)^2+a_*D_uY^2
          +a'_*D_uZ^2+\tfrac12(X+uZ)^2\succ0.     \tag{LA5}
\]
It is strict on all four original states when $u$ is a unit; every
finite-angle correction has bounded infinitesimal coefficients.
This covers all bounded Bob types without invoking LA4 outside its
stated reciprocal strip.

For $u$ unbounded, put $p=1/u$ and use the original homogeneous
parent.  The exact CE4--CE5 raw-frame proof, with any positive real
caps for $C,A$, covers every bounded $v$.  Its strict selected
residue is
\[
 4c_*(1+v^2)W^2+4a_*(1+v^2)Y^2
                      +2(W+R)^2+2(B-vY)^2,
 \quad B=pX+Z,\quad R=vX.                         \tag{LA6}
\]
For unbounded $v$, the exact twice-homogeneous CE7--CE8 proof has
strict residue
\[
 4c_*W^2+4a_*Y^2+2X^2+2(T-Y)^2,
 \qquad T=wZ.                                    \tag{LA7}
\]
These two homogeneous proofs need only positive real lower bounds
for $C,A$ at these endpoints.  They do not use the stronger
``$C,A$ unbounded'' hypothesis needed for the finite-unit-angle part
of the adjacent-event theorem.  All unselected original terms and
positive excesses are retained exactly.  Literal $p=0$ and $w=0$
are the polynomial values of these original homogeneous parents.

The finite unit-Alice and homogeneous unbounded-Alice cases exhaust
all noninfinitesimal $u$.  Exact sign gauges and real/imaginary splitting
give both signs and complex states.  The possible additional $D=0$
null direction at $u=w=0$ is untouched, since the proposition does
not include the infinitesimal-Alice chart.  The sole remaining C14.2
testing obligation is now $u$ infinitesimal or zero, with arbitrary
Bob types and the unchanged full source intersection.



% Source formulas: pda4a21937f64
\subsection{Unit $L$: bounded Bob tests}

\label{component:pda4a21937f64}

\label{pda4a21937f64:sec:r18-c142-bounded-bob}

\begin{proposition}
Retain the unchanged full C14.2 source and every original intersection
as in the preceding noninfinitesimal-Alice proposition.  For every
infinitesimal or zero Alice test $u$ and every bounded Bob test $v$,
the entire original parent is finite tensor SOS for all real or
complex states.  Every source opening and both $D=0$, $\beta=0$
are included.  Together with that proposition this leaves exactly
$u$ infinitesimal and $v$ projectively unbounded, equivalently
$u,w=1/v$ infinitesimal, as the remaining C14.2 testing corner.
\end{proposition}

\paragraph{The same source weights and a pure-source strict matrix.}
Keep the opening-safe positive source weights from LA1,
\[
 C_0=\frac1{2L(1-r)},\quad
 A_0=\frac{1+L(1-r)}{2Lr},\quad
 a_0=\frac L{2(1-Lr)}.
\]
Let $c=(1+v^2)^{-1/2}$, $s=vc$, so $c^2+s^2=1$.
Since $v$ is bounded, $c$ is a positive testing unit.  Write
$s_0=\operatorname{res}s\in(-1,1)$ and $c_0^2=1-s_0^2>0$.
These are real numbers, not source/testing comparisons.  Consider
the pure-source quadratic form in three placeholder variables
\[
 M_0(P,Y,T)=
  (C_0+c_0^2/2)(P-Y)^2+(A_0+s_0^2/2)Y^2
          +a_0(s_0Y-T)^2-s_0PT.                    \tag{LB1}
\]
Its Gram has first leading minor $C_1=C_0+c_0^2/2>0$,
second leading minor $C_1[A_0+s_0^2(a_0+1/2)]>0$, and
\[
 \det M_0=
 \frac{c_0^2\{1+L(1-r)\}^2}
                {8Lr(1-r)(1-Lr)}>0.               \tag{LB2}
\]
This is an exact source identity, not just a limiting determinant.
Moreover $M_0^{-1}$ has bounded entries through every opening.
Indeed every cofactor is a sum of bounded multiples of
$C_0A_0,C_0a_0,A_0a_0,C_0,A_0,a_0,1$; terms involving a square
of a single divergent weight cancel.  Multiplication by
$8L^2r(1-r)(1-Lr)$ makes every cofactor polynomial in bounded
$L,r,s_0$.  Division by LB2 therefore leaves only the source-unit
denominators $L$, $c_0^2$, and $1+L(1-r)$.
Explicitly, with $e=1-r$ and $B_L=1+Le$,
\[
 M_0^{-1}=\frac1{c_0^2B_L}
 \begin{pmatrix}
 2L&2Lr&2s_0\\
 2Lr&2Lr(1-s_0^2e)&2s_0r(1+Le c_0^2)\\
 2s_0&2s_0r(1+Le c_0^2)&
       2(1-Lr c_0^2)(1+Le c_0^2)/L
 \end{pmatrix}.                                      \tag{LB2a}
\]

Choose any exact positive pure-source congruence
$M_0=N_1^2+N_2^2+N_3^2$.  Its inverse allocation of $P,Y,T$
on these three roots is bounded: each inverse-row norm squared is
the corresponding diagonal of $M_0^{-1}$.  This uses only a
pure-source Gram, with the fixed real number $s_0$ as parameter;
no mixed matrix is inverted or square-rooted.

\paragraph{Freeze the only divergent Bob-dependent event exactly.}
Use the legal raw roots
\[
 P=uW,\quad X,\quad Y,\quad
 T=cZ-(s-s_0)Y,
 \qquad d_s=s-s_0\in\operatorname{Inf}_{\rm test}.
                                                               \tag{LB3}
\]
In particular the entire alpha event satisfies
$sY-cZ=s_0Y-T$ exactly.  Thus no product of an unbounded alpha
weight with $d_s$ is ever estimated.  No inverse of $c$, $u$, or
a state coordinate is needed in these root definitions.

Let $B_0=c_0^2(P-Y)^2+s_0^2Y^2-2s_0PT$.
The complete original half-core after Bob normalization is
$K_B/2$, where
\[
 K_B=\{c(P-Y)+usX-sZ\}^2
       +(1+u^2)(sY-cZ)^2+(X+uZ)^2-(1+u^2)Z^2.
\]
It has the following exact, division-free allocation on LB3:
\[
\begin{split}
 K_B={}&X^2+B_0
       +d_sP\{-(s+s_0)P+2s_0Y-2T\}\\
       &+2ucX(sP-s_0Y+T)
       +u^2s^2X^2+u^2(s_0Y-T)^2.
\end{split} \tag{LB4}
\]
This keeps the full signed core and all $u$-dependent hard normals.

\paragraph{Exact actual-parent reassembly.}
Use exactly LA4's pure-source factor
$f=1-(|C/C_0-1|+|A/A_0-1|+|\alpha/a_0-1|)>0$.
Retain the actual original event excesses
$C-fC_0$, $A-fA_0$, $\alpha-fa_0$, and the whole $m,D,\beta$
terms.  On the finite Bob chart the original quarter-parent divided
by $1+v^2$ is then the sum of those positive terms and exactly
\[
 S=f\{C_0(P-Y)^2+A_0(1+u^2)Y^2
                         +a_0(1+u^2)(s_0Y-T)^2\}+K_B/2.
\]
Substitution of LB1 and LB4 gives
\[
\begin{split}
 S={}&fM_0+\tfrac12X^2
       +fA_0u^2Y^2+fa_0u^2(s_0Y-T)^2\\
    &+\tfrac12u^2s^2X^2+\tfrac12u^2(s_0Y-T)^2\\
    &+\tfrac12(1-f)B_0
       +\tfrac12d_sP\{-(s+s_0)P+2s_0Y-2T\}\\
    &+ucX(sP-s_0Y+T).
\end{split} \tag{LB5}
\]
The first two lines consist of a strict source-normalized Gram on
$(N_1,N_2,N_3,X)$ and exact retained positive squares.  All terms
in the last two lines have infinitesimal bounded coefficients on
that same frame: $1-f$ is source infinitesimal, $d_s,u$ are testing
infinitesimals, $c,s,s_0$ bounded, and LB2 supplies a bounded inverse
allocation for all three raw roots.  Since $f\sim1$, the resulting
complete Gram has residue $\operatorname{diag}(1,1,1,1/2)$.
The separated bounded-Gram lemma and restoration of every retained
term prove the full original parent.  The normalization factor
$1+v^2$ is pure testing and positive.

Both coefficient-boundary terms were retained rather than capped
or inverted.  The argument works literally at $u=0$ or $v=0$,
allows signed testing angles, and extends to complex states by
real/imaginary splitting.  It gives no claim for $c$ infinitesimal:
there $c_0=0$, the determinant in LB2 vanishes, and the reciprocal
two-small-angle corner remains a separate obligation.

The new focused checks establish LB2, opening-regular cofactors,
the complete original-core identity LB4 and the entire selected-parent
identity LB5.  They do not rerun prior source identities or assert
the remaining C14.2 corner closed.



% Source formulas: p98e371f5b846
\section{Contact gaps when $L$ vanishes}
Here the coefficient gaps, rather than $L$ alone, determine the remaining quadratic form. The gap $j$ first separates a strict block from its singular limit. The completed $J$ relation controls the remaining state normal when that block degenerates.

\subsection{Small $L$: the source openings and complementary tests}

\label{component:p6d2f659c0fde}

\label{p6d2f659c0fde:sec:r8-c1-small-L-bounded-eta}

Work in the original C1 source, with $d=1$, $x\in\Inf_{>0}$,
$\zeta$ noninfinitesimal, $L=x\zeta\in\Inf_{>0}$, and $\eta$
bounded. Native $r=1-a\eta$, $h=1-\theta$ remain strictly positive.
This section admits every original projective test except possibly
$u,v$ both infinitesimal. It does not admit the remaining two-soft
local chart.

\paragraph{Exact source restrictions, including the total-scale intersection.}
Put $\rho=r+\eta$ and $j=\eta J=\eta(1+xh)-1+r>0$.
Since $a<1+xh$, $\rho$ is a positive bounded unit. The exact beta
numerator is
\[
 2\beta\theta P_d=K\rho(xh+L)-\theta P_d.       \tag{SB1}
\]
It implies $\theta=O(x+L)=O(L)$, since $x/L=1/\zeta$ is bounded.
Consequently
\[
 h,M,E,Q,P_d,P_4\sim1,\qquad B_5\sim1.           \tag{SB2}
\]
For the last assertion, use the exact expression
$B_5=K(1+L\eta)\rho-jP_d$: its residue is
$r+\eta-(\eta-1+r)=1$. Thus original $D=0$ cannot occur on
this source face; literal $\beta=0$ remains included.
The original formulas, with openings canceled before taking residues,
give uniformly
\[
\begin{gathered}
 C(1-r)\sim\rho/2,\quad Ar/\eta\sim1/2,
 \quad Dj\sim1/2,\\
 mL/x^2\sim\rho/2,\quad
 \alpha/[L(L+x)]\sim\rho/2,\quad
 (2\beta+1)\theta/(x+L)\sim\rho.
\end{gathered}                                         \tag{SB3}
\]
In particular $C,D$ are strong, $m,\alpha$ weak, and $A$ is strong
precisely when $\eta$ is a unit. The four exact large-weight
alternatives are
\[
\begin{array}{c|c}
\text{weight unbounded}&\text{equivalent opening}\\\hline
C&1-r\in\Inf\\
D&j\in\Inf\\
A&r\in\Inf\\
\beta&\theta/(x+L)\in\Inf.
\end{array}                                                   \tag{SB4}
\]
For the $A$ equivalence, $r\in\Inf$ forces $\eta$ to be a unit
because $j>0$; if $r$ is a unit, $\eta/r$ is bounded.
All coefficients are bounded exactly when the four quantities in
the right column of (SB4) are positive bounded units. That entire
interior is proved by the positive-total-scale theorem.
After intersection with $\delta_{\rm tot}\in\Inf_{>0}$, at least
one of (SB4) therefore occurs. These are exhaustive necessary source
alternatives, not assertions of their independent realizability.

\paragraph{The unit-$\eta$ remote tests: two complete three-state blocks.}
Suppose first $\eta$ is a unit, with real residue $e>0$, and write
$r_0=\operatorname{res}r\in[0,1]$. For an interior opening,
the source residues are
\[
 C_0=\frac{r_0+e}{2(1-r_0)},\quad
 A_0=\frac e{2r_0},\quad D_0=\frac1{2(e+r_0-1)}.       \tag{SB5}
\]
Retain the original $m,\alpha,\beta$ terms throughout. Bounded
selected weights retain their actual coefficients; unbounded ones
are capped at sufficiently large finite real numbers and their
excesses retained as exact squares.

For $u$ a bounded nonzero unit and $v$ infinitesimal, use the full
raw frame
\[
 p=uW,\quad Y,\quad R=vZ,\quad B_t=X+uZ.
\]
The exact core is (CE2), and the $D$ event is
$J=p-Y+uvB_t-D_uR$. Every Gram coefficient is bounded after
capping, and no state normal is omitted. Set $U=\operatorname{res}u^2>0$.
The residue, divided by four, is $B_t^2/2$ plus
\[
\begin{split}
 f_U={}&C_0(p-Y)^2+A_0(1+U)Y^2
       +D_0\{p-Y-(1+U)R\}^2\\
 &+\tfrac12\{(p-Y)^2+(1+U)UR^2-2(1+U)pR\}.
\end{split}                                            \tag{SB6}
\]
Its first two principal minors and determinant are, respectively,
\[
\begin{gathered}
\Delta_1=\frac{e(e+r_0)}{2(1-r_0)j_0},\qquad
\Delta_2=\frac{e^2(1+U)(e+r_0)}{4r_0(1-r_0)j_0},\\
\Delta_3=\frac{Ue(1+U)^2(e+r_0)^2}{8r_0(1-r_0)j_0}>0,
\qquad j_0=e+r_0-1.
\end{gathered}                                      \tag{SB7}
\]

For $u$ infinitesimal and $v$ a bounded nonzero unit, use the full
frame $(p=uW,X,Y,R=vZ)$. The same exact parent has residue,
divided by four, $(1+V)X^2/2$ plus
\[
\begin{split}
 f_V={}&C_0(1+V)(p-Y)^2+A_0(1+V)Y^2
       +D_0(p-Y-R)^2\\
       &+\tfrac12\{(p-Y)^2+VY^2-2pR\},
 \qquad V=\operatorname{res}v^2>0.
\end{split}                                            \tag{SB8}
\]
The leading two minors are positive, and
\[
 \det f_V=\frac{V(1+V)(e+r_0)^2}{8r_0(1-r_0)j_0}>0. \tag{SB9}
\]

These arguments include every opening in (SB4). To see explicitly
why the finite cap remains valid, impose the kernels of the divergent
events and inspect the remaining two-state Gram. The determinants
are
\[
\renewcommand{\arraystretch}{1.7}
\begin{array}{c|c|c}
\text{opening}&\det(f_U|_{\ker})&\det(f_V|_{\ker})\\\hline
r_0=0,\ e>1&\dfrac{Ue^2(1+U)}{4(e-1)}&\dfrac{Ve}{4(e-1)}\\
r_0=1&\dfrac{U(1+U)^2(e+1)}4&\dfrac{V(e+1)}{4e}\\
j_0=0,\ 0<r_0<1&
 \dfrac{U(1+U)^2}{4r_0}&\dfrac{V(1+V)}{4r_0(1-r_0)}.
\end{array}                                                   \tag{SB10}
\]
The respective kernels are $Y=0$, $p=Y$, and
$p-Y=(1+U)R$ for $f_U$ or $p-Y=R$ for $f_V$.
All first pivots are positive. At the only double opening,
$r_0=j_0=0,e=1$, the two kernels leave the positive scalar
$U(1+U)R^2/2$, respectively $VR^2/2$. There are no other
multiple openings since $e>0$. The elementary finite-penalty
Schur lemma therefore supplies finite real caps simultaneously
for every divergent selected event.

For completeness the remaining projective ends also use original
homogeneous degree. If $u$ is unbounded, (CE4)--(CE5) for bounded
$v$, and (CE7)--(CE8) for unbounded $v$, are strict with any positive
finite $C,A$ caps. If $v$ is unbounded and $u$ bounded, put
$w=1/v$ and use $(p=uW,X,Y,Z)$. The exact $D$ event is
$w(p-Y)+uX-Z$. At $w=0$ the core is
$(1+u^2)(X^2+Y^2)$, while the three caps give
$C(p-Y)^2+A(1+u^2)Y^2+D(uX-Z)^2$; these are strict on
all four roots, including $u=0$. Each term suppressed in this
residue has a bounded infinitesimal coefficient in this same
full frame. Literal $w=0$ is therefore included. The both-unit
tests are covered by the common noncommuting-residue theorem.
Together these statements exhaust the complement of $u,v\in\Inf$.

\paragraph{The infinitesimal-$\eta$ remote tests.}
Here $r\sim1$, $C,D$ are unbounded, and $y=L\eta$ is infinitesimal.
Define the pure-source harmonic coefficient
$\gamma=yCD/(C+D)$. Uniformly through the $a,J$ openings,
\[
 \gamma\sim L/2,\qquad A/y\sim1/(2L),\qquad
 \gamma A/y\sim1/4,\qquad y/\gamma\sim2\eta\in\Inf. \tag{SB11}
\]
The strong-II remote proof (RR3)--(RR6), which uses only $C,D$
unbounded, covers every noninfinitesimal $v$ and every $u$.
For $u$ noninfinitesimal and $v$ infinitesimal, use (RR7)--(RR9)
with zero alpha cap and $k=-1$. In its notation let
$a=u/\sqrt{1+u^2}$, $b=1/\sqrt{1+u^2}$,
$\ell=vY$, $x_t=bX$, $R_0=x_t+a(Z-\ell)$, and $q=\sqrt yY$.
Then
\[
 M_{-1}=\tfrac12R_0^2+\tfrac12(b\ell)^2,\qquad
 v(Z-\ell)=avR_0-bv(b\ell)-b\sqrt y\,V_h,       \tag{SB12}
\]
where $V_h=v(aX-bZ)/\sqrt y$. Normalize the harmonic pair as
$P_h=\sqrt\gamma V_h$, $Q_h=\sqrt{A/y}\,q$. Its residue
determinant is $1-b_0^2>0$ by (SB11), where $b_0=\operatorname{res}b<1$;
the off-diagonal entry is $b_0$. The sole negative term in (RR9),
$-v^2(Z-\ell)^2/2$, is an infinitesimal Gram perturbation of
$(R_0,b\ell,P_h)$ by (SB12), since $\sqrt{y/\gamma}$ is
infinitesimal. Thus the full four-root Gram is strict. Retain the
original positive $\alpha(Z-\ell)^2$, the C1 beta event, and every
other unused term. The homogeneous endpoint $u=\infty$ has
$a=\pm1,b=0$ and is included; literal $v=0$ also causes no change.



% Source formulas: pc15b7caff0a3
\subsection{The small-$L$ native scalar and complete state normals}

\label{component:pc15b7caff0a3}

\label{pc15b7caff0a3:sec:r8-small-L-native-normal}

Normalize $d=1$, use the native C1 source $\lambda=n=0$, and keep
all other original event weights, including literal $\beta=0$.
Put $U=u^2,V=v^2,D_u=1+U,D_v=1+V$, $\theta=1-h$, and
\[
             X=-huZ+\xi,\qquad Y=ruW+\nu.
\]
Here $\nu$ denotes this state normal only.  Direct expansion of the
original event terms gives
\[
 \frac{F_g}{4}=A_wW^2+A_zZ^2-2uv(b_0+b_1U)WZ,
\]
with the complete graph coefficients
\[
\begin{aligned}
 A_w={}&[(C+D+\tfrac12)(1-r)^2+Ar^2]U+Ar^2U^2\\
 &+\{m+[2mr+C(1-r)^2+(A+\alpha+\tfrac12)r^2]U
                  +r^2(m+A+\alpha+\tfrac12)U^2\}V,\\
 A_z={}&\alpha+[\alpha+(m+\beta+\tfrac12)\theta^2]U\\
 &+\{D+[2Dh+\tfrac12h^2+\beta\theta^2]U
                     +(D+\tfrac12)h^2U^2\}V,\\
 b_0={}&m\theta+D(1-r)+\alpha r+\tfrac12,\\
 b_1={}&mr\theta+(D+\tfrac12)h(1-r)+(\alpha+\tfrac12)r.
\end{aligned}
\]
These identities are homogeneous in $W,Z$ and need no state-ratio
bound.  The first-normal and hard-block formulas that follow are
obtained by the same substitution, or by differentiating the original
parent twice in $X,Y$.
The entire original parent is
exactly
\[
\begin{split}
 F/4={}&F_g/4+2\xi(vWf_w+uZf_z)+2\nu(uWk_w+vZk_z)\\
 &+H_{xx}\xi^2+2H_{xy}\xi\nu+H_{yy}\nu^2.
\end{split}                                               \tag{SN1}
\]
where
\[
\begin{aligned}
 f_w={}&-m(1+rU)+(D+\tfrac12)(1-r)U,\\
 f_z={}&m\theta-(D+\tfrac12)V(1+hU)
                    +(\beta+\tfrac12)D_v\theta,\\
 k_w={}&mV(1+rU)-CD_v(1-r)+AD_uD_vr
             -(D+\tfrac12)(1-r)+(\alpha+\tfrac12)D_urV,\\
 k_z={}&-m\theta U+(D+\tfrac12)(1+hU)
                              -(\alpha+\tfrac12)D_u,\\
 H_{xx}={}&m+\beta D_v+(D+\tfrac12)UV+D_v/2,\\
 H_{yy}={}&mUV+CD_v+AD_uD_v+D+\tfrac12
                              +(\alpha+\tfrac12)D_uV,\\
 H_{xy}={}&-(m+D+\tfrac12)uv.
\end{aligned}                                                   \tag{SN2}
\]
There is no discarded normal or excess term in this identity.
Its two exact first-normal contact equations are
\[
 \zeta f_w(x,y)+f_z(x,y)=0,\qquad
 k_w(x,y)+\eta k_z(x,y)=0,\qquad y=x\eta\zeta.             \tag{SN3}
\]
They follow by substituting (NC4), not by setting a testing angle
equal to a source angle in the desired tensor assertion.

On the small-$L$, bounded-$\eta$ source, with $x,L$ infinitesimal
and $\zeta$ noninfinitesimal, the beta-opening identity gives:
\[
 2\beta\theta P_d=K(r+\eta)(xh+L)-\theta P_d,
 \qquad \beta\theta=O(L),\quad \beta\theta^2=O(L^2).       \tag{SN4}
\]
In particular, after the pure-source normalization
$\widehat\xi=\sqrt{1+\beta}\,\xi$, its cross coefficient is
$\beta\theta/\sqrt{1+\beta}=O(L)$, including every unbounded
beta opening.  The complete normal allocations are the retained
RO/JO/HO and DG/LD/CJ terms.



% Source formulas: pe3ef590c01d6
\subsection{The small-$L$ scalar and finite Taylor quotients}

\label{component:pe3ef590c01d6}

\label{pe3ef590c01d6:sec:r9-c1-small-interior}

Use the normalized native C1 parent, $d=1$, with $x=k^2>0$
infinitesimal, and set
\[
 \rho=\eta+r,\quad y=x\eta\zeta,\quad L=x\zeta,\quad
 h=1-\theta,\quad s=\sqrt{\eta\zeta},\quad
 \lambda=\rho/(2\eta).
\]
The reference interior has
\[
 \eta,\zeta,r,1-r,j_0=\eta+r-1
       \quad\text{positive bounded units}.                    \tag{IC1}
\]
The formulas below are the exact scalar and full-normal dictionary.
The opening-uniform RO proof supplies their normalized bounds when
$r$ or $1-r$ tends to zero.  The displayed reference $X$ normal
uses $\sqrt\beta$ when $\beta>0$; RO replaces it throughout by
$\sqrt{\beta+1/2}$, which includes literal $\beta=0$.
The full scalar graph coefficients $A_w,A_z,b_0,b_1$ are printed
explicitly in the native dictionary SN1--SN2.
For IC4--IC21, first assume $Z\ne0$ and normalize $Z=1$ inside
the testing field.  Homogenizing the resulting state-linear roots
restores every $Z$; the literal $Z=0$ state has the direct core
certificate.  The reciprocal identities below are already homogeneous.
\paragraph{A bounded-ratio frame retaining all three state normals.}
First suppose the pure-testing ratio $t=v/u$ is bounded, possibly
infinitesimal. Define
\[
 \begin{gathered}
 E=\eta uW-v,\quad
 b={r(1-r)\over\eta(\eta+2r-1)},\quad
 N=Y-ruW-bE,\\
 \xi=X+hu,\qquad N_x=\sqrt\beta\,\xi,\quad h=1-\theta.
                                                               \end{gathered} \tag{IC4}
\]
The state change is exact and uses only the pure-testing inverse $u$:
$W=(v+E)/(\eta u)$, $Y=r(v+E)/\eta+bE+N$, $X=-hu+\xi$.
Every normal vanishes at the actual source contact. 

\paragraph{The exact scalar, with the negative defect absorbed.}
Set $U=u^2,V=v^2$, $s=\sqrt{\eta\zeta}$, $k=\sqrt x$, and introduce
the positive separated roots
\[
 A_t=uv,\quad B_t={k v^2\over s u},\quad
 C_t=xs,\quad G=kv,\qquad F_v=V-y.
                                                               \tag{IC6}
\]
They satisfy $A_tB_tC_t=G^3$ and coincide at actual contact
$(u,t)=(k,s)$. The pure-testing pole in $B_t$ is inside the root;
none of the coefficient arrays below has a negative power of $t$.

Use $A_w,A_z,b_0,b_1$ from (SN1)--(SN2), including the beta term.
The scalar at $E=N=\xi=0$ is exactly
\[
 \mathcal S={VA_w\over\eta^2U}+A_z
                  -{2V\over\eta}(b_0+b_1U)
       =c_FF_v^2+f_AA_t^2+f_BB_t^2+f_CC_t^2-f_GG^2.    \tag{IC7}
\]
Here the complete coefficients are
\[
\begin{aligned}
 c_F={}&{C(1-r)^2+Ar^2+(\alpha+\tfrac12)r^2+2mr\over\eta^2},\\
 a_{UV}={}&{Ar^2\over\eta^2}+{h^2\over2}+2Dh+\beta\theta^2
                                     -{2b_1\over\eta},\\
 a_{UUV}={}&h^2(D+\tfrac12),\qquad
 a_{UVV}={r^2(m+A+\alpha+\tfrac12)\over\eta^2},\\
 q_V={}&{C(1-r)^2+Ar^2+(D+\tfrac12)(1-r)^2\over\eta^2}
                                     +D-{2b_0\over\eta},\\
 q_U={}&\alpha+(m+\beta+\tfrac12)\theta^2,\\
 f_A={}&a_{UV}+a_{UUV}U+a_{UVV}V,\qquad
 f_B={m\zeta\over\eta x},\\
 f_C={}&{\alpha-c_Fy^2+q_UU\over\eta\zeta x^2},\qquad
 f_G=-{q_V+2c_Fy\over x}.                              \end{aligned} \tag{IC8}
\]

\paragraph{Finite scalar Taylor arrays.}
Let $\epsilon_i=f_i-\lambda$ and $\epsilon_G=f_G-3\lambda$,
$h_c=\epsilon_A+\epsilon_B+\epsilon_C-\epsilon_G$.
Equations (IC8) show that $h_c=h_{00}+aU+cV$ with bounded
pure-source $a,c$. Exact contact value gives
\[
 h_c=a(u^2-k^2)+c(u^2t^2-k^2s^2).                     \tag{IC11}
\]
For $d_u=u-k,d_t=t-s$, its finite second Taylor identity is
\[
\begin{split}
 h_c={}&2k(a+cs^2)d_u+2k^2sc\,d_t
 +(a+ct^2)d_u^2+2kc(t+s)d_ud_t+k^2c\,d_t^2.            \end{split} \tag{IC12}
\]
The coefficient of $d_t^2$ is infinitesimal. For any affine function
$f=f_{00}+a_fU+c_fV$ among (IC8), the exact first difference is
\[
 f-f(k,s)=(u-k)(u+k)(a_f+c_ft^2)
                    +(t-s)k^2(t+s)c_f.                \tag{IC13}
\]
These are the required bounded $u$ arrays and infinitesimal $t$ arrays;
no appeal to unspecified Taylor membership is needed.
The two exact first-jet balances are
\[
 2(\epsilon_{A0}-\epsilon_{C0})+k h_{c,u0}=0,
 \quad 2(\epsilon_{B0}-\epsilon_{C0})+s h_{c,t0}=0.      \tag{IC14}
\]
They follow from the original stationary contact, and are also checked
directly using (NC4). The extra square $c_FF_v^2$ has zero value and
first jets there. Apply the accompanying three-circuit contact-lift
lemma with the exact conductors
\[
 G(u-k)=k(A_t-G),\qquad G(t-s)=s(B_t-G).                \tag{IC15}
\]

\paragraph{Explicit finite allocations of every first normal.}
Use $f_w,f_z,k_w,k_z$ in (SN2), and define the polynomials
\[
\begin{aligned}
 P_Y(U,V)&=2(k_w/\eta+k_z),\\
 P_E(U,V)&={2\over\eta^2}(A_w/U-mV/U)
              -{2\over\eta}(b_0+b_1U)+bP_Y(U,V).
\end{aligned}                                                   \tag{IC16}
\]
The quotient $A_w/U-mV/U$ is polynomial, as its displayed numerator
shows. For either polynomial $P$ use the exact divisions
\[
 P=P_{00}+UP_U+VP_V,\quad
 P_U={P-P(0,V)\over U},\quad
 P_V={P(0,V)-P(0,0)\over V}.                           \tag{IC17}
\]
The normalized bounds on these exact arrays are supplied by RO. The complete first normals
of the original $F/4$ in the frame (IC4) are exactly
\[
\begin{aligned}
 \mathcal L_N&=vP_Y,\\
 \mathcal L_E&=2mv t^2/\eta^2+vP_E,\\
 \mathcal L_x&={1\over\sqrt\beta}\left[
 2u\theta(m+\beta+\tfrac12)
 -{2m\over\eta}ut^2(1+rU)+2uVZ_X\right],\\
 Z_X&=(\beta+\tfrac12)\theta-(D+\tfrac12)(1+hU)
                         +(D+\tfrac12)(1-r)/\eta .     \end{aligned} \tag{IC18}
\]
These identities follow by differentiating the entire original parent
after the full state substitution (IC4).

Their actual root allocations are
\[
\begin{aligned}
 vP&={P_{00}\over k}G+u(P_U+t^2P_V)A_t,\\
 {2m\over\eta^2}vt^2&={2ms t\over\eta^2k}B_t,\\
 \mathcal L_x&={2u\theta(m+\beta+\tfrac12)
                    \over\sqrt\beta\,xs}C_t
 -{2ms(1+rU)\over\eta k\sqrt\beta}B_t
 +{2vZ_X\over\sqrt\beta}A_t.                          \end{aligned} \tag{IC19}
\]

The exact original contact makes each complete normal zero at
$(u,t)=(k,s)$. For a normal $\ell=\sum b_iA_i+b_GG$, put
$B=\sum b_i+b_G$, and use the explicit monic divisions
\[
 B=(u-k){B(u,t)-B(k,t)\over u-k}
       +(t-s){B(k,t)-B(k,s)\over t-s}.                 \tag{IC21}
\]

\paragraph{Reciprocal testing-ratio identities.}
Return to homogeneous state coordinates and set
\[
 p=uW,\ q=vZ,\ E=\eta p-q,\quad
 Y=rp+bE+N,\quad B=X+uZ=\xi+\theta uZ.
\]
Use the complete raw root list
\[
 E,\ N,\ N_x=\sqrt\beta\xi,\quad
 P=\sqrt m\,vW,\quad Q=\sqrt\alpha Z,\quad
 R_t=u^2W,\quad S_t=uvW.                              \tag{IC22}
\]
\[
 \sqrt x\,p=\sqrt{x/m}\,(u/v)P,\quad
 Y=rp+bE+N,\quad q=\eta p-E,                          \tag{IC23}
\]
For completeness, with $B=0$ the exact higher-angle part of the
$C,A,D$ terms and half core, after the quadratic base is removed, is
\[
\begin{split}
 &U\{AY^2-2Dq(p-Y-q)+q^2/2-pq\}\\
 &\quad+V\{C(p-Y)^2+(A+\tfrac12)Y^2\}
 +(A+\tfrac12)UVY^2+(D+\tfrac12)U^2q^2.               \end{split} \tag{IC24}
\]
The terms involving $B$ are equally explicit. Relative to $B=0$, the
$D$ term adds $2DuvB(p-Y-D_uq)+DUVB^2$, and the half core adds
$uvB(p-Y-D_uq)+(D_v+UV)B^2/2$.
Here
\[
 B={N_x\over\sqrt\beta}+{\theta u\over\sqrt\alpha}Q,
 \quad uvp=uS_t,
\]
The opening-specific proof bounds their coefficients.  Write
$M=vW-X+uvY-uZ$. The three remaining
original terms keep their entire roots:
\[
\begin{aligned}
 \sqrt m\,M&=P+\sqrt m\,uvY-\sqrt m\,B,\\
 \sqrt\alpha\sqrt{D_u}(vY-Z)
  &=\sqrt{D_u}\{\sqrt\alpha(rS_t+v(bE+N))-Q\},\\
 \sqrt\beta\sqrt{D_v}B
  &=\sqrt{D_v}\{N_x+\sqrt{\beta\theta^2/\alpha}\,uQ\}.
\end{aligned}                                                   \tag{IC26}
\]



% Source formulas: p3df6e6bdd8f4
\subsection{Unit positive $j$: all $r$ openings}

\label{component:p3df6e6bdd8f4}

\label{p3df6e6bdd8f4:sec:r13-c1-r-openings}

\begin{theorem}
Use the original physical native C1 source, normalized by $d=1$, and
retain all original intersections. Suppose that $x=k^2$ is positive
infinitesimal and
\[
 \eta,\quad\zeta,\quad j_0=\eta+r-1
       \quad\hbox{are positive bounded units}.                 \tag{RO1}
\]
Then the complete original parent is tensor SOS at every independent
projective test and every state. The original $0<r<1$ may approach
either endpoint at any positive rate, and the literal $\beta=0$ face
is included. No lower positive real bound on $r$, $1-r$, or $\beta$
is imposed. This extends the covered IC1 domain; it does
not cover an infinitesimal $j_0$, infinitesimal $\eta$, or unbounded
$\zeta$.
\end{theorem}

All scalar and Gram identities in this section use the normalized
parent $F/4$; multiplication of every square root by $2$ restores $F$.

The proof reuses the unchanged exact scalar and full-normal identities
(IC7)--(IC8), (IC16)--(IC19), and the three-circuit lemma. Its new
content is a uniform weighting of the two divergent event openings,
and a different $X$ normal that retains bounded and zero beta.

\paragraph{Opening-cleared source arrays.}
Put $\rho=\eta+r$, $y=x\eta\zeta$, $h=1-\theta$,
$c=C(1-r)$ and $a=Ar$. The beta numerator (SB1) gives
$\theta=x\pi$ with $\pi>0$ bounded. We do not require $\pi$
infinitesimal. The original source formulas give, uniformly for
all $0<r<1$,
\[
\begin{gathered}
 c=\rho/2+O(x),\quad a=\eta/2+O(x),\quad
 D=1/(2j_0)+O(x),\quad h=1+O(x),\\
 m/x\sim\rho/(2\zeta),\qquad
 \alpha/x^2\sim\rho\zeta(1+\zeta)/2,\qquad
 B_x\theta/x={K\rho(h+\zeta)\over2P_d},\quad
 B_x=\beta+\tfrac12\ge\tfrac12.                      
\end{gathered} \tag{RO2}
\]
These are finite canceled numerator identities, followed by residue
evaluation. Every canceled denominator at $x=0$ is a product of
positive units from $\eta,\zeta,j_0,j_0+r$ and $\rho$, with no
factor $r$, $1-r$, or $\pi$. In particular $D>0$ on this domain;
the original $D=0$ face has empty intersection here, as already proved
in (SB2). Both $\theta/\sqrt\alpha$ and
$B_x\theta^2/\alpha$ are bounded, not necessarily infinitesimal.

The exact source formulas for the two cleared strong weights are
\[
 c={K\rho E\over2P_d},\qquad a={\eta P_4\over2P_d}.
\]
Their order-$x$ differences in (RO2), and the following arrays, have
opening-free finite canceled formulas:
\[
 {P_{Y,00}\over x},\quad {P_{E,00}\over x},\quad
 {\theta(m+B_x)\over x},\quad
 {\alpha+(m+B_x)\theta^2\over x^2},\quad {q_V\over x}.
                                                               \tag{RO3}
\]
Here $P_Y,P_E,q_V$ are precisely (IC16) and (IC8). To check why no
unbounded event remains in these arrays, replace every occurrence of
$C(1-r)^2$ by $c(1-r)$, every $Ar^2$ by $ar$, every
$C(1-r)$ by $c$, and every $Ar$ by $a$ before division. No other
occurrence of $C$ or $A$ survives in the scalar or first-normal
polynomials.  This is a finite symbolic
criterion for boundedness on the stated field domain, not a numerical
endpoint approximation.

\paragraph{The normalized hard frame.}
Define the pure-source coefficients
\[
 b={r(1-r)\over\eta(j_0+r)},\quad
 H_0={\eta(j_0+r)\over2r(1-r)j_0},\quad
 \gamma={\rho\over2\eta(j_0+r)}.
\]
Both $1/\sqrt{H_0}$ and $1/\sqrt{B_x}$ are bounded; $H_0$
itself may be unbounded. Put
\[
 E_t=\eta uW-vZ,\quad
 N_y=\sqrt{H_0}(Y-ruW-bE_t),\quad
 N_x=\sqrt{B_x}(X+huZ).                              \tag{RO4}
\]
All three normals vanish at the actual source contact. For the
comparison weights
\[
 C_0={\rho\over2(1-r)},\quad A_0={\eta\over2r},\quad
 D_0={1\over2j_0},
\]
the exact quadratic base is
\[
 C_0(p-Y)^2+A_0Y^2+D_0(p-Y-q)^2+\tfrac12(p-Y)^2-pq
       =H_0(Y-rp-b(\eta p-q))^2+\gamma(\eta p-q)^2.  \tag{RO5}
\]
Thus the divergent normal has already been normalized before taking
residue. The cancellation factors needed for its entire hard block are
\[
 {C_0\over H_0}={r\rho j_0\over\eta(j_0+r)},\quad
 {A_0\over H_0}={(1-r)j_0\over j_0+r},\quad
 {D_0\over H_0}={r(1-r)\over\eta(j_0+r)}.             \tag{RO6}
\]
For example, after $p=(v+E_t)/\eta$ and $Z=1$, the $E_t$
coefficients in $p-Y$ and $Y$ are respectively
\[
 {1-r\over\eta}-b={(1-r)j_0\over\eta(j_0+r)},\qquad
 {r\over\eta}+b={r\rho\over\eta(j_0+r)}.            \tag{RO7}
\]
The strong $C,A$ factors therefore cancel even on those complete
first-state directions. For instance an error
$\delta_c(p-Y)^2/(1-r)$, with $\delta_c=O(x)$, has a
$N_y^2$ coefficient $\delta_c/((1-r)H_0)=O(x)$,
an $E_t^2$ coefficient $O(x)(1-r)$, and an $E_tN_y$
coefficient $O(x)/\sqrt{H_0}$. The $A$ error has the same
properties using $rH_0$ and the second formula in (RO7).
The original higher-angle factors have the same cleared coefficients,
multiplied by $U=u^2$ or $V=v^2$.

Suppose first that $u,v>0$ are infinitesimal and the pure-testing
ratio $t=v/u$ is bounded. After the state substitution (RO4), the
complete hard quadratic on $(E_t,N_y,N_x)$ consequently has bounded
entries and residue
\[
                         \operatorname{diag}(\gamma,1,1).       \tag{RO8}
\]
Here $\gamma$ is a positive bounded unit even when either $r$
opening is infinitesimal. The additional $m t^2$ contribution is
infinitesimal. The beta and half-core coefficients jointly give the
third pivot because $m/B_x$ is infinitesimal. These estimates concern
the full hard block, not only its restriction to the contact graph.

\paragraph{The complete scalar and all first normals.}
Normalize $Z=1$, homogenizing at the end as in IC1. Use the unchanged
positive separated roots
\[
 A_t=uv,\quad B_t={kv^2\over s u},\quad C_t=xs,\quad
 G=kv,\quad F_v=V-y,\qquad s=\sqrt{\eta\zeta}.
\]
The scalar at all three normals zero is exactly (IC7). Its coefficient
residues are still
\[
 c_F\sim{\rho\over2\eta^2},\quad
 f_A,f_B,f_C\sim\lambda={\rho\over2\eta},\quad
 f_G\sim3\lambda.                                   \tag{RO9}
\]
They do not depend on the residue of $\pi$. Every coefficient of
$U,V$ in (IC8) is bounded after the clearing described above;
for the additional term in $f_C$ use
$B_x\theta^2/x^2=(B_x\theta/x)\pi$. The exact contact
value and the two first-jet identities (IC11)--(IC14) are unchanged.
Their finite arrays therefore obey precisely the hypotheses of the
three-circuit lemma, with $s$ a unit. It gives the same strict finite
circuit Gram and the strict $F_v$ pivot. In particular the negative
soft defect has been allocated by the entire scalar (IC7), not
dropped by comparing only its quadratic base.

The complete first normals are (IC18), with the $N$ line divided
by $\sqrt{H_0}$ and the $X$ line divided by $\sqrt{B_x}$
instead of $\sqrt\beta$. Their exact root allocations are (IC19)
with those same replacements. Every coefficient remains infinitesimal:
the $N_y$ factors are bounded source factors times the old cleared
polynomials; the $X$ coefficient on $C_t$ is
\[
 {2u\theta(m+B_x)\over\sqrt{B_x}\,xs},
\]
which is infinitesimal by (RO3), and its coefficients on $B_t,A_t$
are respectively $O(k)$ and $O(v)$. In the latter coefficient
$Z_X=B_x\theta-(D+\tfrac12)(1+hU)
 +(D+\tfrac12)(1-r)/\eta$ is bounded. The $E_t$ line uses
$m/x,P_{E,00}/x$ and bounded polynomial arrays, as in IC1.
There is no negative power of $t$ in any coefficient.

Each coefficient is polynomial in $u,t$ with bounded pure-source
arrays; every term depending on $t$ carries $u,k$ or $m/k=O(k)$.
The exact monic differences (IC21) hence have bounded $u$ quotient
and infinitesimal $t$ quotient after evaluation at $u=k$. The
unchanged exact source contact cancels the bare $G$ term in each
complete normal. The conductors
$G(u-k)=k(A_t-G)$ and $G(t-s)=s(B_t-G)$ therefore allocate
all normals with infinitesimal coefficients on the strict circuit
roots. Joining this allocation to (RO8) proves the entire parent
on the bounded-ratio chart.

\paragraph{The reciprocal ratio chart.}
For $v/u$ unbounded use the homogeneous root list
\[
 E_t,\quad N_y,\quad N_x,\quad P_m=\sqrt m\,vW,
 \quad Q_\alpha=\sqrt\alpha Z,\quad R_t=u^2W,
 \quad S_t=uvW.                                      \tag{RO10}
\]
The error of the actual quadratic base relative to (RO5) is an
$O(x)$ quadratic in $(p,E_t,N_y)$ after the opening cancellations
(RO6)--(RO7). Every $\sqrt x\,p$ factor is allocated exactly by
\[
 \sqrt x\,p=\sqrt{x/m}\,(u/v)P_m,
\]
with a bounded source unit $x/m$ and an infinitesimal pure-testing
ratio $u/v$. Thus the entire error has zero Gram residue on (RO10).
All coefficients in the exact higher-angle identity (IC24) remain
bounded after $Y=rp+bE_t+N_y/\sqrt{H_0}$. This follows term
by term from (RO6), $b=O(r(1-r))$, $Ar=a$, and $C(1-r)=c$.
Its only nonzero scalar residues are the two positive diagonal entries
\[
                 \kappa_R=\rho\eta/2,\qquad
                 \kappa_S=\rho/2.                  \tag{RO11}
\]
No opening has been substituted before cancellation.

It remains essential to combine beta with its half-core part. Set
$B_t^{\rm state}=X+uZ$ and $J_0=p-Y-D_uvZ$. Relative to
$B_t^{\rm state}=0$, the exact $D$, beta and half-core increment is
\[
 2(D+\tfrac12)uvB_t^{\rm state}J_0
 +(D+\tfrac12)UV(B_t^{\rm state})^2
 +B_xD_v(B_t^{\rm state})^2.                          \tag{RO12}
\]
This whole-polynomial identity is checked in  Use
\[
 B_t^{\rm state}={N_x\over\sqrt{B_x}}
               +{\theta u\over\sqrt\alpha}Q_\alpha.
\]
The first two terms of (RO12) have zero Gram residue on (RO10);
$D$ is bounded, $1/\sqrt{B_x}$ and $\theta/\sqrt\alpha$
are bounded, and every coefficient carries $u$ or $v$.
The three remaining positive terms have their entire roots
\[
\begin{aligned}
 \sqrt m\,R_m&=P_m+\sqrt m\,uvY-\sqrt m\,B_t^{\rm state},\\
 \sqrt\alpha\sqrt{D_u}(vY-Z)
   &=\sqrt{D_u}\{\sqrt\alpha[rS_t+v(bE_t+N_y/\sqrt{H_0})]
                         -Q_\alpha\},\\
 \sqrt{B_x}\sqrt{D_v}B_t^{\rm state}
   &=\sqrt{D_v}\{N_x+\sqrt{B_x\theta^2/\alpha}\,uQ_\alpha\}.
\end{aligned}                                                    \tag{RO13}
\]
Each is its leading root plus an infinitesimal-coefficient combination
of (RO10). The square root in the last line is pure source and bounded;
its factor $u$ supplies the needed infinitesimal even when beta is
zero or bounded. The full parent therefore has a bounded Gram with
strict diagonal residue
\[
                 (\gamma,1,1,1,1,\kappa_R,\kappa_S).
                                                               \tag{RO14}
\]
This proves the remaining ratio chart.

\paragraph{All-tests join and exact scope.}
The earlier small-$L$, bounded-$\eta$ remote theorem includes (RO1)
with every original opening, so it covers the complement of
$u,v$ both infinitesimal. Literal axes have the elementary full-core
identities recorded in IC1. Its exact sign gauges apply unchanged.
Homogenizing the $Z=1$ identity gives polynomial state identities,
including $Z=0$; applying them to real and imaginary parts covers
complex states. All inverse and radical factors above are pure
source or pure testing. This completes the claimed full-parent
all-tests certificate.

The newly removed source region is the portion of (RO1) outside IC1,
namely either $r$ or $1-r$ infinitesimal, with all total-scale and
native source inequalities retained. Remaining C14.3 source
configurations still include infinitesimal $\eta$, unbounded
$\zeta$, and nonunit or nonpositive $j_0$ while the actual
$j=\eta(1+xh)-1+r$ remains positive. No conclusion on those
regions or on C14.4 follows from this theorem.

\paragraph{A local algebraic extension below beta zero.}
The preceding two local ratio-chart arguments have a stronger precise
scope that is useful when another event is restored later. Assume the
exact native stationary coefficient formulas (NC1)--(NC4),
$x>0$ infinitesimal, (RO1), $0<r,h<1$, and
\[
              B_x=\beta+\tfrac12\ge b_*>0
              \quad\hbox{for some real }b_* .                 \tag{RO15}
\]
Original source physicality and $\beta\ge0$ are not required for
this local assertion. Then the full C1-form parent is tensor SOS at
all independent $u,v$ both infinitesimal, including their literal
zero axes. In particular $-1/2<\beta<0$ is allowed whenever
$\beta+1/2$ is a positive bounded unit.

Indeed $P_d\sim1$ before any $h$ estimate, since all its source
arguments are bounded. The exact formula
$B_x\theta/x=K\rho(h+\zeta)/(2P_d)$ has bounded right side,
and (RO15) therefore implies $\theta/x$ bounded. All of (RO2)--(RO14)
follows, with $1/\sqrt{B_x}\le1/\sqrt{b_*}$. In particular
$m,\alpha>0$, $C,A>0$, and $D\sim1/(2j_0)>0$ on this
domain. No remaining use of $\beta\ge0$ occurs in the local
scalar, first-normal or hard-block proof: its complete contribution
is combined with the half-core as (RO12). At $u=0$ the signed
core contributes $D_v(X^2+Y^2)/2$, so its $X^2$ part plus
the beta term is $B_xD_vX^2$. At $v=0$ the same combination
is $B_x(X+uZ)^2$ and the remaining half-core is
$(uW-Y)^2/2$. Thus both axes are also exact SOS under (RO15).
This proves the local and literal-axis assertions under (RO15).



% Source formulas: p48da253a150f
\subsection{The positive $j$ opening}

\label{component:p48da253a150f}

\label{p48da253a150f:sec:r13-c1-j-opening}

\begin{theorem}
Use the original physical native C1 source with $d=1$, $x=k^2>0$
infinitesimal, and
\[
 \eta,\ \zeta,\ r,\ 1-r\quad\hbox{positive bounded units},\qquad
 j=\eta(1+xh)-1+r\quad\hbox{positive infinitesimal}.       \tag{JO1}
\]
Then the complete original parent is tensor SOS for every independent
projective test and every state. Literal $\beta=0$ is retained.
There is no assumption that $j_0=\eta+r-1$ is positive: $j_0=0$
and a slightly negative $j_0$ are included. Simultaneous infinitesimal
$r$ and $j$ are outside this theorem.
\end{theorem}

All scalar and Gram identities in this section use the normalized
parent $F/4$; multiplication of every square root by $2$ restores $F$.

The divergence $D\asymp1/j$ must be centered by its full event,
including its testing-angle dependence. A normal centered only by its
quadratic limit would lose control when $j\ll x$. The following exact
frame removes that problem before any residue is taken.

\paragraph{Source bounds and contact-exact coordinates.}
Write $\theta=1-h=x\pi$, where $\pi>0$ is bounded by (SB1),
$B_x=\beta+1/2\ge1/2$, and $y=x\eta\zeta$. Put
\[
 R_0=1+j-\eta=r+\eta xh,\qquad
 g(U)=R_0-\eta hU,\qquad U=u^2,\ V=v^2.
                                                               \tag{JO2}
\]
Thus $g(x)=r$ exactly. The native source formulas, with
$r=1+j-\eta(1+xh)$, imply
\[
 \begin{gathered}
 C,A=O(1),\quad Dj=\tfrac12+O(x),\quad
 m/x\sim1/(2\zeta),\quad
 \alpha/x^2\sim\zeta(1+\zeta)/2,\\
 B_x\theta/x=K(\eta+r)(h+\zeta)/(2P_d),\qquad
 \theta/\sqrt\alpha,\quad B_x\theta^2/\alpha=O(1).
 \end{gathered}                                                \tag{JO3}
\]
In particular $D>0$, so its square root and inverse are pure source
and legitimate. Set
\[
 \begin{gathered}
 p=uW,\quad q=vZ,\quad E_t=\eta p-q,\qquad
 N_d=\sqrt D\,(J_t+jp),\\
 J_t=uW+uvX-Y-vZ,\qquad
 N_x=\sqrt{B_x}(X+huZ),\quad
 n=N_d/\sqrt D,\quad \xi=N_x/\sqrt{B_x}.
 \end{gathered}                                                \tag{JO4}
\]
The exact inverse state formulas are polynomial and homogeneous:
\[
 Y_0=(1+j)p-q-n,\qquad
 Y=Y_0-hUq+uv\xi,\qquad X=-huZ+\xi,\qquad
 J_t=n-jp.                                             \tag{JO5}
\]
All three normals vanish at the actual source contact. At their zero
graph, $q=\eta p$ and $Y=g(U)p$, so the original native contact
graph is recovered at $U=x$ without changing any source contact jet.

For later use introduce positive reference weights
\[
 \widetilde C={1+j\over2(\eta-j)},\quad
 \widetilde A={\eta\over2R_0},\quad
 \widetilde D={1\over2j}.
\]
Since $R_0,1-R_0,\eta$ are units, exact source numerator division gives
\[
 C-\widetilde C=O(x),\qquad A-\widetilde A=O(x),\qquad
 Dj-\tfrac12=O(x).                                    \tag{JO6}
\]
The last assertion is deliberately about $Dj$, not $D$. The canceled
denominators of (JO3),(JO6), at $x=j=0$, contain only the unit
factors $\eta,\zeta,1-\eta$ and constants. No comparison of $j/x$
has been made.

\paragraph{The complete bounded hard block.}
The reference quadratic base has the exact identity
\[
 \begin{aligned}
 &\widetilde C(p-Y_0)^2+\widetilde A Y_0^2+
 \widetilde D(p-Y_0-q)^2+\tfrac12(p-Y_0)^2-pq\\
 &=\widetilde\gamma E_t^2+
 \widetilde H\{(1-\widetilde b)E_t-n\}^2,               \\
 \widetilde H&={\eta(j+R_0)\over2R_0(1-R_0)j},\quad
 \widetilde b={R_0(1-R_0)\over\eta(j+R_0)},\quad
 \widetilde\gamma={1+j\over2\eta(j+R_0)} .
 \end{aligned} \tag{JO7}
\]
No $p$ term survives. Furthermore
\[
 1-\widetilde b={j(1+j)\over\eta(j+R_0)},\qquad
 \widetilde H/D\sim1.                                 \tag{JO8}
\]
On $(E_t,N_d)$ this is a bounded hard Gram with diagonal residue
$(1/(2\eta(1-\eta)),1)$ and vanishing off-diagonal residue.
The actual-base error has coefficients $O(x)$ in $(p,E_t,N_d)$.
For its apparently singular part use
\[
 (D-\widetilde D)(n-jp)^2,\quad
 {D-\widetilde D\over D}=O(x),\quad
 {(D-\widetilde D)j\over\sqrt D}=O(x\sqrt j),\quad
 (D-\widetilde D)j^2=O(xj).                           \tag{JO9}
\]
These are consequences of (JO3),(JO6) and hold for every relative
source rate. All other weights in the base are bounded.

For bounded $t=v/u$ at $u,v>0$ both infinitesimal, put $Z=1$.
Then $p=(v+E_t)/\eta$ in (JO5). The higher-angle terms have
bounded infinitesimal hard coefficients, and the $m$ hard coefficient
$m t^2$ is infinitesimal. Combining beta with its half-core part gives
the third hard pivot on $N_x$. Thus the full hard residue is
\[
                 \operatorname{diag}(1/(2\eta(1-\eta)),1,1).
                                                               \tag{JO10}
\]
The exact higher-angle identity given below also proves these hard
bounds directly; no state normal is suppressed here.

\paragraph{An exact scalar three-circuit with bounded polynomial coefficients.}
Use $s=\sqrt{\eta\zeta}$ and the same separated roots
\[
 A_t=uv,\quad B_t={kv^2\over su},\quad C_t=xs,\quad
 G=kv,\quad F_v=V-y.
\]
On the zero graph of (JO4), the exact scalar is
\[
 \mathcal S=c_FF_v^2+f_A(U,V)A_t^2+f_BB_t^2+
                       f_C(U)C_t^2-f_GG^2,            \tag{JO11}
\]
where, with $g=g(U)$,
\[
 \begin{aligned}
 c_F&={2mR_0+C(1-R_0)^2+(A+\alpha+\tfrac12)R_0^2\over\eta^2},\\
 q_U&=\alpha+(m+B_x)\theta^2,\\
 q_V&={C(1-R_0)^2+AR_0^2+(D+\tfrac12)j^2\over\eta^2}
       -{2m\theta\over\eta}-{(2\alpha+1)R_0\over\eta}-\tfrac12,\\
 f_B&={m\zeta\over\eta x},\quad
 f_C={\alpha-c_Fy^2+q_UU\over\eta\zeta x^2},\quad
 f_G=-{q_V+2c_Fy\over x},                              \\
 f_A&={mV(g^2-2\eta h)\over\eta^2}-{2m\theta g\over\eta}
       +C(1+V)\{2h(1-R_0)/\eta+h^2U\}\\
 &\quad+\{A(1+V)+(\alpha+\tfrac12)V\}
                   \{g^2/\eta^2-2R_0h/\eta+h^2U\}\\
 &\quad-(2\alpha+1)(g/\eta-h)+B_x\theta^2-\tfrac12 .
 \end{aligned} \tag{JO12}
\]
In particular $f_A$ is an explicit bounded polynomial of degree at
most two in $U$ and one in $V$; the divergent $D$ is absent.
The only $D$ factor elsewhere is $Dj^2=O(j)$. These formulas
follow by substituting $W=v/(\eta u)$, $X=-hu$, $Y=g(U)v/\eta$
in the entire parent; 

The source bounds and (JO12) give
\[
 c_F\sim1/(2\eta^2),\qquad
 f_A,f_B,f_C\sim\lambda=1/(2\eta).
\]
Exact scalar contact gives
$f_G=f_A(x,y)+f_B+f_C(x)$, proving that $f_G$ is bounded
and $f_G\sim3\lambda$ without estimating a singular quotient
$q_V/x$. Every nonconstant coefficient in $f_A,f_C$ is bounded;
for $q_U/x^2$ use (JO3).

For clarity the required finite Taylor step works for these higher
polynomial coefficients too. Let $H(u,t)=f_A(u^2,u^2t^2)+f_B+
f_C(u^2)-f_G$. It is a finite polynomial in $u,t$ with bounded
source arrays and $H(k,s)=0$. Write $\Delta_u=u-k$,
$\Delta_t=t-s$. A finite identity, obtained by monic polynomial
division, is
\[
 H=H_u(k,s)\Delta_u+H_t(k,s)\Delta_t+
       H_{uu}\Delta_u^2+H_{ut}\Delta_u\Delta_t+
       H_{tt}\Delta_t^2,                              \tag{JO13}
\]
where one may take
\[
 \begin{aligned}
 H_{uu}&={H(u,t)-H(k,t)-(u-k)H_u(k,t)\over(u-k)^2},\\
 H_{ut}&={H_u(k,t)-H_u(k,s)\over t-s},\\
 H_{tt}&={H(k,t)-H(k,s)-(t-s)H_t(k,s)\over(t-s)^2}.
 \end{aligned}
\]
These displayed divisions mean their finite polynomial quotients,
not mixed inverses. The first two arrays are bounded and $H_{tt}$
is infinitesimal: every $t$-dependent monomial of $H(k,t)$
contains $k^2$. The coefficient first differences similarly have a
bounded $u$ quotient and infinitesimal $t$ quotient. Exact first
angular stationarity supplies the two unchanged contact-jet balances
of the three-circuit lemma. The chain rule is legitimate here because
$g(x)=r$ and the original first state derivatives vanish at contact.
Consequently (JO11) has the strict finite three-circuit Gram and the
strict $F_v$ pivot, including the complete soft defect.

\paragraph{Every complete first normal.}
The following bounded polynomials are explicit:
\[
 \begin{aligned}
 P_Y/2={}&m\{V(1+gU)/\eta-\theta U\}
       -C D_v(1-g)/\eta+A D_uD_vg/\eta\\
       &+(D+\tfrac12)j/\eta+
                  (\alpha+\tfrac12)D_u(gV/\eta-1),\\
 P_E={}&{2mV\over\eta^2}\{g+1+j+g(1+j)U\}
       -{2m\theta\over\eta}\{1+(1+j)U\}\\
       &-{2CD_vj(1-g)\over\eta^2}
        +{2AD_uD_vg(1+j)\over\eta^2}
        +{2(D+\tfrac12)j^2\over\eta^2}\\
       &+{2(\alpha+\tfrac12)D_u(1+j)\over\eta}
                                      (gV/\eta-1),\\
 Z_X={}&B_x\theta+P_Y/2-(D+\tfrac12)j/\eta .
 \end{aligned}                                                \tag{JO14}
\]
Their only $D$ factors are $Dj$ and $Dj^2$; $D$
cancels entirely from $Z_X$. Direct full-parent differentiation gives
\[
 \begin{aligned}
 \mathcal L_d&=-vP_Y/\sqrt D,\\
 \mathcal L_E&=2mvt^2/\eta^2+vP_E,\\
 \mathcal L_x&={1\over\sqrt{B_x}}
 \{2u\theta(m+B_x)-2mut^2(1+gU)/\eta+2uVZ_X\}.
 \end{aligned}                                                \tag{JO15}
\]
No normal is missing. For $P=P_Y,P_E$ use the exact finite
division $P=P_{00}+UP_U+VP_V$. Contact implies respectively
\[
 \begin{aligned}
 P_{Y,00}/x&=-P_{Y,U}(x,y)-\eta\zeta P_{Y,V}(y),\\
 P_{E,00}/x&=-2(m/x)\eta\zeta/\eta^2
                  -P_{E,U}(x,y)-\eta\zeta P_{E,V}(y).
 \end{aligned}                                                \tag{JO16}
\]
Thus these coquotients are bounded by finite formulas with no $j$
inverse. Allocate $vP$ by
$vP=(P_{00}/k)G+u(P_U+t^2P_V)A_t$; allocate the additional
$2mvt^2/\eta^2$ on $B_t$ with coefficient
$2mst/(\eta^2k)$. The $X$ line allocates exactly on $(C_t,B_t,A_t)$
with coefficients
\[
 {2u\theta(m+B_x)\over\sqrt{B_x}xs},\quad
 -{2ms(1+gU)\over\eta k\sqrt{B_x}},\quad
 {2vZ_X\over\sqrt{B_x}}.                               \tag{JO17}
\]
Every displayed coefficient is infinitesimal; for the first use
$\theta(m+B_x)/x=O(1)$, and for the others use $m/x=O(1)$.
The $N_d$ row has the additional bounded factor $1/\sqrt D$.
Each sum of coefficient functions vanishes at contact. Its finite
$u$ difference is bounded and its $t$ difference at $u=k$ is
infinitesimal, because every $t$ dependence carries $u$ or $k$.
The two unchanged conductors therefore place every whole normal on the
strict circuit roots with infinitesimal coefficients. Joining (JO10)
proves the full parent on the bounded-ratio chart.

\paragraph{The reciprocal testing ratio: a complete polynomial remainder.}
For $v/u$ unbounded, return to homogeneous state coordinates and use
\[
 E_t,\ N_d,\ N_x,\ P_m=\sqrt m\,vW,\ Q_\alpha=\sqrt\alpha Z,
                  \ R_t=u^2W,\ S_t=uvW.              \tag{JO18}
\]
The entire base error from (JO9) has zero Gram residue on this list:
its $p^2$ coefficient is $O(x)$ and
$\sqrt x\,p=\sqrt{x/m}(u/v)P_m$. Every coefficient is a product
of a bounded source factor and a bounded testing factor; $u/v$ is
pure testing.

Here is the full higher-angle identity used to control the other terms.
Write $T_s=-hUq+uv\xi$, so $Y=Y_0+T_s$, with $Y_0$ in
(JO5). After removing the actual quadratic base and the three entire
positive terms $mR_m^2$, $\alpha D_u(vY-Z)^2$ and
$B_xD_v(X+uZ)^2$, the exact remainder is
\[
 \begin{aligned}
 \mathcal R={}&C\{V(p-Y_0)^2-2D_v(p-Y_0)T_s+D_vT_s^2\}\\
 &+A\{(D_uD_v-1)Y_0^2+2D_uD_vY_0T_s+D_uD_vT_s^2\}\\
 &+\tfrac12D_uVY^2-UqY+(h-\tfrac12)Uq^2-uvq\xi .
 \end{aligned}                                                \tag{JO19}
\]
This is a quadratic in the full state, with bounded coefficients
$C,A,h$ and no $D$. It is checked as a whole polynomial. Insert
$q=\eta p-E_t$, $n=N_d/\sqrt D$, $\xi=N_x/\sqrt{B_x}$.
Every term contains $U,V$ or $uv$. Cross terms with $p$ have
infinitesimal coefficients on (JO18), using $Up=uR_t$,
$Vp=vS_t$ and $uvp=uS_t$. Hard quadratic terms have zero
residue. Terms $U^aV^bp^2$ are allocated on $R_t^2$ when
$a>0$, and on $S_t^2$ when $a=0,b>0$. Their two only nonzero
residues, verified from (JO19), are
\[
                         \kappa_R=\eta/2,\qquad
                         \kappa_S=1/2.               \tag{JO20}
\]
Every remaining such coefficient contains another $U$ or $V$.
These identities explicitly exhaust the scalar and every normal
coupling of (JO19).

The three retained positive terms have leading roots
$P_m,-Q_\alpha,N_x$. Their remaining root coefficients are
infinitesimal because $m,\alpha$ are infinitesimal and
$\theta/\sqrt\alpha$, $B_x\theta^2/\alpha$ are bounded:
\[
 \sqrt m\,R_m=P_m+\sqrt m\,uvY-\sqrt m(\xi+\theta uZ),
 \quad
 \sqrt{B_xD_v}(X+uZ)=
 \sqrt{D_v}\{N_x+\sqrt{B_x\theta^2/\alpha}\,uQ_\alpha\}.
\]
The alpha term is $\sqrt{D_u}(\sqrt\alpha\,vY-Q_\alpha)$.
Thus the complete parent has a bounded Gram on (JO18) with strictly
positive diagonal residue
\[
             (1/(2\eta(1-\eta)),1,1,1,1,\eta/2,1/2).
                                                               \tag{JO21}
\]
This proves the reciprocal-ratio chart and completes the local proof.

\paragraph{All-tests join and receipt.}
The common small-$L$, bounded-$\eta$ remote theorem applies to
(JO1), including the exact original source openings and literal beta
zero. It covers the complement of both-infinitesimal testing angles.
The original exact sign gauges, literal axes, homogeneous state
extension and real/imaginary decomposition apply as in IC1 and RO.
All coefficients and radicals used above are pure-field quantities;
no mixed inverse or comparison of source and testing rates occurs.

 They verify
the original full-parent substitution, every complete first normal,
the scalar identity, all newly used source denominator factors,
the reference hard cancellation, the whole reciprocal higher-angle
identity and both strict soft pivots. The finite polynomial bounds
and full-rates circuit allocation are the proof above, rather than
an inference from those checks. This admits precisely (JO1) and its
exact symmetry images. It does not remove the simultaneous
$r,j$ opening or any unbounded-$\zeta$ or infinitesimal-$\eta$ source.



% Source formulas: pbf6de1d7328a
\subsection{The simultaneous positive $r,j$ opening}

\label{component:pbf6de1d7328a}

\label{pbf6de1d7328a:sec:r13-c1-simultaneous-rj}

\begin{theorem}
Use the original physical native C1 source, with $d=1$ and $x=k^2>0$
infinitesimal. Suppose $\eta,\zeta$ are positive bounded units and
$r$ and the actual $j=\eta(1+xh)-1+r$ are both positive
infinitesimals. Then the complete original parent is tensor SOS for
every independent projective test and every state. Literal $\beta=0$
is retained. No relative rate among $x,r,j$ or between source and
testing parameters is imposed.
\end{theorem}

All scalar and Gram identities in this section use the normalized
parent $F/4$; multiplication of every square root by $2$ restores $F$.

The essential change is to combine the whole positive $A,D$ events
before choosing the scalar graph. The weighted $uvX$ term is retained
inside a hard root. Solving for $p=uW$, rather than for $q=vZ$,
then gives a polynomial inverse with a pure-source unit denominator.

\paragraph{Source bounds and the harmonic frame.}
Write $U=u^2$, $V=v^2$, $D_u=1+U$, $D_v=1+V$,
$\theta=1-h=x\pi$, $B_x=\beta+1/2$, and
\[
 T=1+xh,\qquad e=\eta T=1-r+j,\qquad y=x\eta\zeta .
 \tag{HO1}
\]
The bounded-$\eta$ source bound gives $\pi=O(1)$.
Here $e,\eta,h,1-r\sim1$. Exact native numerator cancellation gives
\[
 \begin{gathered}
 C-C_* =O(x),\qquad Ar-e/2=O(x),\qquad Dj-1/2=O(x),\\
 C_*={1+j\over2(1-r)},\qquad A_*={e\over2r},\qquad D_*={1\over2j},\\
 m/x\sim1/(2\zeta),\qquad
 \alpha/x^2\sim\zeta(1+\zeta)/2,\\
 B_x\theta/x=O(1),\qquad
 q_U/x^2=O(1),\quad q_U=\alpha+(m+B_x)\theta^2 .
 \end{gathered}
 \tag{HO2}
\]
 This proves boundedness at arbitrary simultaneous
source rates, rather than just nonvanishing at a selected boundary.
In particular $A,D>0$ and both are unbounded.
Also $\theta/\sqrt\alpha$ and $B_x\theta^2/\alpha$ are bounded,
including $\beta=0$.

Define pure-source positive quantities
\[
 \mathsf H=A+D,\qquad \rho={D\over\mathsf H},\qquad
 \mathsf K={AD\over\mathsf H},\qquad
 \delta={Ar+Dj\over\mathsf H}.
 \tag{HO3}
\]
Both $\mathsf H$ and $\mathsf K$ are unbounded,
$0<\rho<1$, and
$\mathsf K^{-1}=A^{-1}+D^{-1}=O(r+j)$.
Put $p=uW$, $q=vZ$, $J_0=p+uvX-q$, and use the full-event roots
\[
 \begin{aligned}
 E&=\sqrt{\mathsf K}\{J_0-(r-j)p\},\\
 N&=\sqrt{\mathsf H}\{Y-\rho J_0-\delta p\},\\
 N_x&=\sqrt{B_x}(X+huZ),\qquad
 \ell=E/\sqrt{\mathsf K},\quad n=N/\sqrt{\mathsf H},
 \quad \xi=N_x/\sqrt{B_x}.
 \end{aligned}
 \tag{HO4}
\]
All three vanish at the actual source contact. Their exact inverse is
\[
 \begin{aligned}
 p&={(1+hU)q-uv\xi+\ell\over e},&
 Y&=rp+\rho\ell+n,&X&=-huZ+\xi,\\
 J_t&=-jp+(1-\rho)\ell-n,&J_t&=p+uvX-Y-q .
 \end{aligned}
 \tag{HO5}
\]
The denominator $e$ is pure source. In particular no inverse of
$1+hU$ occurs, even if the testing and source rates are unrelated.
The useful exact weighted identities are
\[
 {A\rho^2\over\mathsf K}=\rho,\qquad
 {A\over\mathsf H}=1-\rho,\qquad
 {D(1-\rho)^2\over\mathsf K}=1-\rho,\qquad
 {D\over\mathsf H}=\rho .
 \tag{HO6}
\]

\paragraph{The complete hard block and its reference cancellation.}
The $A,D$ base has the exact identity
\[
 \begin{split}
 AY^2+DJ_t^2={}&N^2+E^2
 +2p\left\{{Ar+Dj\over\sqrt{\mathsf H}}N
                 +\sqrt{\mathsf K}(r-j)E\right\}\\
 &+(Ar^2+Dj^2)p^2 .
 \end{split}
 \tag{HO7}
\]
Both displayed cross coefficients and the last coefficient are
infinitesimal: the squared $E$ cross coefficient is $O(r+j)$
and the squared $N$ cross coefficient is $O(\mathsf H^{-1})$.
For bounded testing ratio $t=v/u$, (HO5) has bounded coefficients
on the hard roots, and the only potentially large $m$ coefficient is
$mt^2/\mathsf K$, which is infinitesimal.
The higher $A$ term is controlled by (HO6); for example
\[
 {Ar^2\over\mathsf K}=r^2+{Ar\over Dj}rj=O((r+j)^2),
 \qquad {Ar\rho\over\mathsf K}=r .
 \tag{HO8}
\]
Consequently its hard coefficients carry an additional $U$ or $V$.
The bounded $C$ and half-core terms have infinitesimal coefficients
on $E,N$, and combining the beta term with its half-core part gives
the third hard pivot. The full hard residue on $(E,N,N_x)$ is
$\operatorname{diag}(1,1,1)$.

A stronger exact base statement is needed for unbounded $t$.
Set $\widetilde q=q-uvX=ep-\ell$, so that
$J_t=p-Y-\widetilde q$. The reference quadratic form
\[
 \mathcal B_*=
 C_*(p-Y)^2+A_*Y^2+D_*(p-Y-\widetilde q)^2
                    +\tfrac12(p-Y)^2-p\widetilde q
 \tag{HO9}
\]
has no soft $p$ coordinate after substituting
$Y=rp+\rho\ell+n$ and $\widetilde q=ep-\ell$:
\[
 \mathcal B_*=
 (A_*+C_*+\tfrac12)(\rho\ell+n)^2
                 +D_*\{(1-\rho)\ell-n\}^2 .
 \tag{HO10}
\]
This is an exact contact cancellation, checked before taking residues.
Since $A_*/A=1+O(x)$ and $D_*/D=1+O(x)$, (HO6) shows that
its bounded hard Gram on $(E,N)$ has residue $\operatorname{diag}(1,1)$.
Replacing the reference weights by the actual $C,A,D$ creates a
quadratic error with coefficients $O(x)$ in $(p,E,N)$.
For the singular weights this follows from
$(A-A_*)Y^2=O(x)\,AY^2$ and
$(D-D_*)J_t^2=O(x)\,DJ_t^2$, using (HO6),(HO7);
neither an $x/r$ nor an $x/j$ estimate is used.

\paragraph{The exact bounded-ratio scalar circuit.}
At $E=N=N_x=0$ and $Z=1$, let
\[
 w(U)={1+hU\over e},\qquad
 p=w(U)v,\quad Y=rw(U)v,\quad X=-hu .
 \tag{HO11}
\]
Since $w(x)=1/\eta$, this graph agrees exactly with the native
source contact. The source contact value and all first state and
angular jets are therefore common by the chain rule.

For compact scalar coefficients set
$a=C(1-r)^2$, $b=Ar^2$, $d=(D+1/2)j^2$,
$c=(\alpha+1/2)r^2$, and
\[
 L_A(U)=1+2h+(2h+h^2)U+h^2U^2 .
 \]
The complete scalar is exactly
\[
 \mathcal S=c_F(V-y)^2+f_A(U,V)A_t^2+f_BB_t^2+
                         f_C(U)C_t^2-f_GG^2 ,
 \tag{HO12}
\]
on the separated roots
$A_t=uv$, $B_t=kv^2/(su)$, $C_t=xs$, $G=kv$,
where $s=\sqrt{\eta\zeta}$. The coefficients are
\[
 \begin{aligned}
 c_F&={2m(h+r)+a+b+c\over e^2},\\
 q_V&={a+b+d\over e^2}-{2m\theta\over e}
                       -{(2\alpha+1)r\over e}-\tfrac12,\\
 f_B&={m\eta\zeta\over e^2x},\qquad
 f_C={\alpha-c_Fy^2+q_UU\over\eta\zeta x^2},\qquad
 f_G=-{q_V+2c_Fy\over x},\\
 f_A&={mV\over e^2}
 \{(h+r)^2+2hr+2hr(h+r)U+h^2r^2U^2\}
       -{2m\theta\over e}(h+r+hrU)\\
 &\quad+{a\over e^2}(2h+h^2U)(1+V)
       +{b\over e^2}L_A(U)(1+V)
       +{d\over e^2}(2h+h^2U)+{c\over e^2}L_A(U)V\\
 &\quad-{(2\alpha+1)r\over e}(1+h+hU)
                              +B_x\theta^2-\tfrac12 .
 \end{aligned}
 \tag{HO13}
\]
The only opening weights here are $Ar^2$ and $Dj^2$, both
infinitesimal. Thus $f_A$ is a bounded polynomial of degree at most
two in $U$ and one in $V$, and all nonconstant coefficients of
$f_C$ are bounded by (HO2). Direct source reduction gives
\[
 c_F\sim\tfrac12,\qquad f_A,f_B,f_C\sim\tfrac12 .
 \tag{HO14}
\]
Exact scalar contact gives
$f_G=f_A(x,y)+f_B+f_C(x)\sim3/2$, without separately estimating
$q_V/x$.

Apply the finite three-circuit lemma with $\lambda=1/2$.
The polynomial Taylor arrays are the monic quotients in (JO13),
now applied to
$f_A(u^2,u^2t^2)+f_B+f_C(u^2)-f_G$.
Their $u$ and mixed arrays are bounded; the $t,t$ array and the
coefficient $t$ first differences are infinitesimal, since every
$t$ dependence at $u=k$ contains $k^2$.
The two exact angular contact jets give the two circuit balances.
This supplies the strict three-circuit Gram and the strict
$V-y$ pivot for the entire scalar, including its negative soft term.

\paragraph{All complete first normals.}
The following formulas come from differentiating the entire original
parent after (HO5), with $Z=1$. Put
$a_E=r/e+\rho$, $b_E=(1-r)/e-\rho$,
$c_E=1-\rho-j/e$. The bounded polynomials are
\[
 \begin{aligned}
 P_N/2={}&m\{wV(1+rU)-\theta U\}
 -CD_v(1-r)w+AD_uD_vrw\\
 &+(D+\tfrac12)jw
             +(\alpha+\tfrac12)D_u(rwV-1),\\
 P_E/2={}&{mV\over e^2}(h+r+hrU)-{m\theta\over e}
 +ma_E\{wV(1+rU)-\theta U\}\\
 &+CD_v(1-r)wb_E+AD_uD_vrwa_E
 -(D+\tfrac12)jwc_E
            +(\alpha+\tfrac12)D_u(rwV-1)a_E,\\
 Z_X={}&{m\theta(1+rU)\over e}+B_x\theta\\
 &-{CD_v(1-r)^2w+AD_uD_vr^2w+(D+\tfrac12)j^2w
                 +(\alpha+\tfrac12)D_ur(rwV-1)\over e}.
 \end{aligned}
 \tag{HO15}
\]
Their only opening-weight factors are $Ar,Dj,Ar^2,Dj^2$.
Every coefficient is bounded uniformly in the source rates and
$0<\rho<1$. The three complete first normals are
\[
 \begin{aligned}
 \mathcal L_N&={vP_N\over\sqrt{\mathsf H}},\\
 \mathcal L_E&={2mvt^2/e^2+vP_E\over\sqrt{\mathsf K}},\\
 \mathcal L_x&={1\over\sqrt{B_x}}\left\{
 2u\theta(m+B_x)-2mwut^2(1+rU)
           \left(1+{V(1+rU)\over e}\right)+2uVZ_X\right\}.
 \end{aligned}
 \tag{HO16}
\]
For either $P=P_N,P_E$, use its exact finite monic decomposition
$P=P_{00}+UP_U+VP_V$.
State contact gives
\[
 \begin{aligned}
 P_{N,00}/x&=-P_{N,U}(x,y)-\eta\zeta P_{N,V}(y),\\
 P_{E,00}/x&=-2(m/x)\eta\zeta/e^2
                         -P_{E,U}(x,y)-\eta\zeta P_{E,V}(y).
 \end{aligned}
 \tag{HO17}
\]
Thus these coquotients are bounded without any new opening inverse.
Allocate $vP=(P_{00}/k)G+u(P_U+t^2P_V)A_t$.
The extra $E$ term allocates on $B_t$ with coefficient
$2mst/(e^2k\sqrt{\mathsf K})$.
The three $X$ terms allocate respectively on $C_t,B_t,A_t$,
with coefficients
\[
 {2u\theta(m+B_x)\over xs\sqrt{B_x}},\qquad
 {-2msw(1+rU)\{1+V(1+rU)/e\}\over k\sqrt{B_x}},\qquad
 {2vZ_X\over\sqrt{B_x}} .
 \tag{HO18}
\]
Every displayed allocation coefficient is infinitesimal.
The coefficient sums vanish at contact by the exact state jets;
their finite $u$ first differences are bounded and their $t$
first differences at $u=k$ are infinitesimal.
This follows directly from the bounded polynomial arrays in
(HO15): every $t$ dependence carries $u$ or $k$.
The two unchanged circuit conductors therefore place every whole
normal on the strict circuit roots with infinitesimal coefficients.
Together with the complete hard block, this proves all bounded $t$,
including infinitesimal $t$, without a missing state normal.

\paragraph{The reciprocal ratio: retain both soft directions.}
If $v/u$ is unbounded, work homogeneously and use seven roots
\[
 E,\quad N,\quad N_x,\quad
 P_m=\sqrt m\,vW,\quad Q_\alpha=\sqrt\alpha Z,\quad
 R_t=uvZ,\quad S_t=uvW .
 \tag{HO19}
\]
The base error following (HO10) has zero Gram residue on this list:
use the exact separated identity
$\sqrt x\,p=\sqrt{x/m}(u/v)P_m$.
This absorbs every $O(x)$ coefficient in $(p,E,N)$, since $x/m$
is a source unit and $u/v$ is a pure-testing infinitesimal.

After removing the actual quadratic base (HO9), with actual $C,A,D$,
and the three entire positive terms
$mM_t^2$, $\alpha D_u(vY-Z)^2$,
$B_xD_v(X+uZ)^2$, the full remaining quadratic is exactly
\[
 \begin{split}
 \mathcal R={}&CV(p-Y)^2+A(U+V+UV)Y^2
                      +\tfrac12D_uVY^2-\theta UqY\\
 &+(h-\tfrac12)Uq^2+\tfrac12h^2U^2q^2
       -uv\xi\{Y+(1+hU)q\}+\tfrac12UV\xi^2 .
 \end{split}
 \tag{HO20}
\]
This identity retains the whole $D$ event in the quadratic base,
so no coefficient containing a bare divergent $D$ remains.
It is checked as one complete homogeneous polynomial.

Here is a finite allocation of every term. The exact identity
\[
 up={(1+hU)R_t-u^2v\xi+uE/\sqrt{\mathsf K}\over e}
 \tag{HO21}
\]
has bounded coefficients and uses no mixed inverse.
The $C$ term has leading root
$\sqrt C(1-r)S_t$; its other coefficients are
$-\sqrt C\,\rho v/\sqrt{\mathsf K}$ and
$-\sqrt C\,v/\sqrt{\mathsf H}$ on $E,N$, both infinitesimal.
All three $A$ higher terms have zero residue. For example,
\[
 \sqrt A\,uY=\sqrt A\,r(up)+\sqrt\rho\,uE+
                               \sqrt{1-\rho}\,uN,\qquad
 \sqrt A\,vY=\sqrt A\,rS_t+\sqrt\rho\,vE+
                               \sqrt{1-\rho}\,vN .
 \tag{HO22}
\]
Here $\sqrt A\,r=O(\sqrt r)$ and (HO8),(HO21)
control every coefficient; the $uvY$ term is $v$ times the first.
The half-$VY^2$ term is smaller by the same identities.
The $-\theta UqY$ term is $-\theta R_t(uY)$ and has zero residue.
The next two terms are $(h-1/2)R_t^2$ and
$h^2U R_t^2/2$.
Finally, the cross term is
$-\xi v(uY)-\xi v(1+hU)R_t$ and all its coefficients,
as well as those of $UV\xi^2/2$, are infinitesimal.
Thus this allocation exhausts the complete quadratic (HO20);
its only nonzero soft residues are $1/2$ on each of $R_t,S_t$.

The three retained positive terms have leading roots
$P_m,-Q_\alpha,N_x$:
\[
 \begin{aligned}
 \sqrt m\,M_t&=P_m+\sqrt m\,uvY-\sqrt m\,\xi
                  -\sqrt{m/\alpha}\,\theta uQ_\alpha,\\
 \sqrt{\alpha D_u}(vY-Z)
             &=\sqrt{D_u}(\sqrt\alpha\,vY-Q_\alpha),\\
 \sqrt{B_xD_v}(X+uZ)
             &=\sqrt{D_v}\{N_x+
                      \sqrt{B_x\theta^2/\alpha}\,uQ_\alpha\}.
 \end{aligned}
 \tag{HO23}
\]
Every other coefficient here is infinitesimal by (HO2),(HO21),
because $\theta\sqrt{m/\alpha}=O(\sqrt x)$.
The complete bounded Gram on (HO19) therefore has strictly positive
diagonal residue
\[
             \operatorname{diag}(1,1,1,1,1,1/2,1/2).
 \tag{HO24}
\]
This proves the reciprocal chart for the full parent and every state.

\paragraph{All-tests join and receipt.}
The common small-$L$, bounded-$\eta$ remote theorem covers the
complement of both-infinitesimal testing angles on this original
physical source domain. Its exact source openings and literal beta
zero remain available. The original exact sign gauges, literal axes,
homogeneous state extension and real/imaginary decomposition join
as in IC1, RO and JO. All roots, radicals and inverses above are
pure-field or separated; every bound is a product of a source bound
and a testing bound.

Those checks support the finite bounds and circuit allocation proved
above; they are not being used as a substitute for them.
This theorem removes precisely the simultaneous positive
infinitesimal $r,j$ opening with unit $\eta,\zeta$.

\begin{corollary}[The entire unit-$\eta,\zeta$ small-$L$ region]
Within the original physical native small-$L$ C1.4 source, suppose
$x>0$ is infinitesimal and $\eta,\zeta$ are positive bounded units.
Then the full parent is tensor SOS at every independent projective
test and every state, with all original source intersections retained.
\end{corollary}
\begin{proof}
The positive actual gap $j=\eta(1+xh)-1+r$ is bounded.
If $j$ is a unit, then $j_0=j-\eta xh$ is a positive unit,
so RO applies, including either $r$ opening. If $j$ is infinitesimal,
then $1-r\sim\eta$ is a unit. A unit $r$ is exactly the JO domain;
an infinitesimal $r$ is exactly the HO domain. These alternatives
exhaust the stated source, require no comparison of $j$ with $x$,
and each supplies its complete local, remote and endpoint join.
\end{proof}



% Source formulas: pa73a17051258
\subsection{The simultaneous small-$\eta$ parent identity}

\label{component:pa73a17051258}

\paragraph{Polynomial identity.}
This is an independently reconstructed exact identity for the complete
original C1 parent, not a certificate admitting C14.3b. It retains the
same stationary source, $\Lambda=n=0$, $\beta\ge0$, all $j$ and $1-r$
openings, independent source/testing fields, and the remaining local
$v/u$-infinitesimal testing chart. No mixed inverse or comparison of
source and testing rates is used below.

Write $\mathcal P=F/4$, initially at $Z=1$. Set
\[
 a=1/\zeta,\quad t=\eta,\quad y=tx/a,\quad U=u^2,\quad V=v^2,
 \quad B=\beta+\tfrac12,\quad D_p=D+\tfrac12,
\]
\[
 \begin{gathered}
 E=tuW-v,\quad Q_t=Y-rv/t,\quad P=E/t-Q_t,\\
 \xi=X+hu,\quad R_D=v(U-x),\quad T_D=R_D-uv\xi/h.
 \end{gathered}
\]
Thus $Q_t=rE/t+\nu$ in the R15 notation. Put
\[
 q=\alpha+(m+B)(1-h)^2,\quad\mu=m-a^2q,
 \quad\lambda=\mu/(tax),\quad H=avW-u-a\xi,
\]
\[
 F_0=V-y,\quad F_b=F_0+vE,\quad F_t=F_0+(t/r)vQ_t,
 \quad T_0=A+\alpha+\tfrac12,\quad T_m=T_0+m,
 \quad e_4=r^2T_m/t^2.
\]
Here $F_t=\widetilde F$ and $F_b=\overline F$. Define pure-source
coefficients
\[
 b=Ar/t-m(1-h)-\alpha-\tfrac12,\qquad
 b'=b+e_4yt/(2r),
\]
\[
 \kappa=2[B(1-h)-D_pj/t-mr/t],\qquad
 \gamma=2b'/h+2D_pj/r-2tB(1-h)/r.
\]
They satisfy $b'=O(x)$, $t\kappa\sim-1$ and $\gamma\sim1$.

\paragraph{Complete parent, with no omitted hard row.}
Let
\[
 L=UV+(ax/t)V^2/U+tx^2/a-3xV,
\]
\[
 \mathscr C=\lambda L+\frac{\mu V}{t^2U}E^2
                +2\left(\frac{\mu v^3}{t^2U}-\lambda xv\right)E.
\]
The exact full-parent identity is
\begin{align*}
 \mathcal P={}&qH^2+\mathscr C
 +CP^2+AQ_t^2+D_p(P-hT_D)^2+2b'Q_tT_D\\
 &+\mathscr R+e_4UF_t^2-e_4yu(H+a\xi)F_t\\
 &+AUQ_t^2+B(1+V)\xi^2
                 -2(\mu/a)H\xi-\mu\xi^2\\
 &+\kappa uF_t\xi+\gamma uvQ_t\xi,                 \tag{JF1}
\end{align*}
where the entire joint radial row is
\[
 \mathscr R=\frac C{t^2}(F_b-rF_t)^2
       +\frac{r^2T_0}{t^2}F_t^2
       +\left(\frac{2mr}{t^2}+e_4x\right)F_bF_t.    \tag{JF2}
\]
In particular all $mVEQ_t/t$ terms are included in JF2; they are not
separate errors. The $mUVQ_t^2$ contribution is included in the complete
$e_4UF_t^2$ row, so it cannot also be retained independently.

A short independent way of certifying JF1 is useful. Without imposing
any source formulas, the difference between the original parent and
its displayed right side is exactly a linear combination of only
\[
                       vE,\ vQ_t,\ u\xi,\ UV,\ V,\ 1.
\]
The displayed expression has zero value and all five first derivatives
at the original contact $U=x,V=y,E=Q_t=\xi=0$.
Its six coefficients therefore vanish by the six original stationary
contact balances. This support calculation and the six contact identities
prove the identity for the complete original parent, including every state normal.

\paragraph{One strict radial pivot, with its exact companion.}
Set $\rho=m+rT_mx/2$. Then JF2 has the pure-source completion
\[
 \mathscr R=\frac C{t^2}
       [F_b-r(1-\rho/C)F_t]^2
       +\frac{r^2}{t^2}[T_0+2\rho-\rho^2/C]F_t^2.  \tag{JF3}
\]
Here $\rho=O(x)$ and $T_0+2\rho-\rho^2/C\sim1/2$;
all denominators are pure source and this remains uniform at both
source openings. Both $F_t/t$ and $F_b/t$ have bounded allocations
on JF3. This does not provide an independent strict $F_0/t$ pivot.
The non-small remaining coefficient is precisely $\gamma uvQ_t\xi$.
In particular a source split making $x/\sqrt t$ bounded does not
remove it, because $\gamma\sim1$.

\paragraph{The extra circuit row participates in the leading cancellation.}
Use $c^2=a/t$, $k^2=x$, $p=ckv/u$, and homogenize now in $Z$:
\[
 J_0=uv-x/c,\quad R_c=kv(tcW-Z),\quad
 R=\sqrt\lambda\,tk(cv-u)W,
 \quad \tau=tcW-Z.
\]
The generic exact circuit completion is
\[
 \mathscr C=\lambda J_0^2Z^2+R^2+
      \underbrace{\frac{2\lambda k}{p}R_c^2}_{\mathcal T^2}
      -\lambda xE^2.
\]
The full-$D$ conductor is
\[
 R_D=uJ_0Z+kR_c-(x/c)(H+a\xi),\qquad
 H+a\xi=avW-uZ.                                  \tag{JF4}
\]
The additional exact polarization is
\[
 E(H+a\xi)=uv\tau^2-t(u-cv)^2WZ.                  \tag{JF5}
\]
Consequently the extra circuit row and the principal cross left after
using JF4 satisfy
\[
 \mathcal T^2+\kappa h(x/c)E(H+a\xi)
 =\frac{x}{c}(2\lambda+\kappa h)uv\tau^2
      -\kappa ht\frac{x}{c}(u-cv)^2WZ.             \tag{JF6}
\]
This cancels the previously non-small $E$--angular Gram coupling.
JF6 is used as an exact identity, rather than as a separate positive
term. In particular, its first coefficient has not been allocated
to a strict positive row. The sign of $WZ$ must be retained. The second summand is a
literal SOS on the pure-testing sign chart $WZ\ge0$ because
$-\kappa>0$, but not on the opposite sign chart.

\paragraph{An exact allocation of the negative small extra-row fraction.}
There is nevertheless a useful conductor:
\[
 x\frac{k}{p}R_c^2
          =U R_c^2-kuJ_0R_c\tau.                   \tag{JF7}
\]
If the testing state is normalized with $Z$ a unit and $W$ bounded,
then $\tau/Z=tc(W/Z)-1$ is bounded. Thus the right side is an
infinitesimal Gram on $\sqrt\lambda R_c$ and
$\sqrt\lambda J_0Z$, with coefficients $U$ and
$ku(\tau/Z)$. Since
$\sqrt\lambda R_c=R+\sqrt\lambda kE$, its additional coefficient
on $E/\sqrt t$ is $\sqrt{\lambda t}\,k=O(k)$.
Therefore a negative $O(x)$ fraction of $\mathcal T^2$ has a bounded
infinitesimal allocation on this state chart. No bound on $p$ is used.
The companion identity
\[
 F_b/t=(W/Z)J_0Z+(x/a)\tau                          \tag{JF8}
\]
may supply a strict $xZ$ radial pivot in a simultaneous state-chart
reassembly, since $\tau/Z\sim-1$ there. All divisions in JF7--JF8
are pure testing. Neither identity is an authorization to assume a
strict $E/\sqrt t$ pivot after the leading source block has become
semidefinite. That final simultaneous allocation is still required.

\paragraph{A direct conductor on the normalized $Z$-unit state chart.}
There is a simpler route on this state chart that does not perform JF4--JF6.
Put $e_E=E/\sqrt t$ and retain the original strict source completion
of the first line of JF1. Since $Q_t=E/t-P$ and $vZ=tuW-E$,
\[
 uvQ_t\xi=\sqrt t\,U(W/Z)e_E\xi
             -u(\xi/Z)e_E^2-uvP\xi.                \tag{JF9}
\]
When the testing state is normalized with $Z$ a unit, $W/Z$ and
$X/Z$ are bounded and $\xi/Z=X/Z+hu$ is bounded. Thus all three
coefficients in JF9 are infinitesimal in the bounded tensor ring.
The last uses the bounded pure-source allocation of $P$ on the
completed source roots. The first line of JF1 has source Schur pivot
$tS\sim1/2$, and $E/\sqrt t$ has a bounded allocation there,
including after the full $D_p$ shift. Consequently JF9 removes the
remaining $\gamma$ coupling on this state chart without producing
a semidefinite source frame.
The accompanying complete state-chart certificate, if used, must still
list the retained positive rows and the strict angular block; JF9
alone does not cover states with normalized $Z$ infinitesimal.

\paragraph{Use of the full identity.}
The complete JF1 identity is the input to DG3/LD9 and the completed $J$-row conductor in Section~\ref{component:pe3bbfd3a8151}. Those arguments allocate the remaining normal and retain the full angular and radial terms.



% Source formulas: peea9e9c546c5
\subsection{The small-$\eta$ angular cancellation}

\label{component:peea9e9c546c5}

Use the small-$\eta$, unit-$\zeta$ source: $x,\eta>0$
infinitesimal, $\zeta$ a positive bounded unit, and
$\beta\ge0$.  Put
\[
 a=1/\zeta,\quad c^2=a/\eta,\quad k^2=x,\quad
 B_x=\beta+\tfrac12,\quad q=\alpha+(m+B_x)(1-h)^2,
 \quad\mu=m-a^2q,\quad\lambda=\mu/(\eta ax).
\]
For homogeneous testing states use the exact roots
\[
\begin{gathered}
 \xi=X+huZ,\quad E=\eta uW-vZ,\quad
 D_0=u-cv,\quad J=(uv-x/c)Z,\\
 R_D=v(u^2-x)Z,\quad R_c=kv(\eta cW-Z),\quad
 \mathcal Q=\sqrt q(avW-uZ-a\xi).
\end{gathered}
\]
Direct multiplication gives
\[
 R_D-uJ=(x/c)D_0Z,\qquad
 \mathcal Q/\sqrt q=-D_0Z+(c/k)R_c-a\xi.
\]
Hence the complete principal $\xi$ row cancels exactly to
\[
 -2\eta\lambda ckR_c\xi+2\eta\lambda c(R_D-uJ)\xi
 =-2\eta\lambda x\mathcal Q\xi/\sqrt q
                    -2\eta\lambda xa\xi^2.             \tag{SC5}
\]
Together with the singular rank's $+\mu\xi^2$, the last term
leaves $-\mu\xi^2$, because $\mu=\eta\lambda ax$.
The latter coefficient is infinitesimal relative to $B_x\xi^2$.
The complete angular block on $(\mathcal Q,\sqrt{B_x}\xi)$
has strictly positive residue: the original source expressions give
\[
 B_xq/x^2\sim
 \frac{B_x\zeta(1+\zeta)}2+\frac{(1+\zeta)^2}4
       >\frac14\sim(\eta\lambda)^2.
\]
For unbounded $B_x$ the normalized off-diagonal tends to zero.
For bounded $B_x\ge1/2$, the displayed difference is bounded
below by $\zeta(3+2\zeta)/4>0$.  The source substitutions are
the regular DG/LD dictionary; the direct derivative and margin
calculation is also printed in the refinement proof.
All inverses and radicals in these coefficient expressions are
pure source.  The full radial and remaining normal allocations are
the retained JF/DG/LD/CJ terms; SC5 supplies their angular block.



% Source formulas: p60eca5442075
\subsection{Small $L$, unbounded $\zeta$ and unit $\eta$}

\label{component:p60eca5442075}

\label{p60eca5442075:sec:r17-lz-unit-eta}

\begin{theorem}
Retain the original physical C1 source and every original contact,
strict-gap, nonnegative-weight, total-scale and earlier-exclusion
intersection. Suppose $x=k^2,L=x\zeta$ are positive infinitesimals,
$\zeta$ is unbounded and $\eta$ is a positive bounded unit.
Then the complete original parent is finite tensor SOS for every
independent projective test and every real or complex state.
No positivity or noninfinitesimal lower bound is imposed on
$j_0=\eta+r-1$: only the actual original gap
$j=\eta(1+xh)-1+r>0$ is used. Both $r$ openings, simultaneous
$r,j$ openings and literal $\beta=0$ are included.
\end{theorem}

This theorem treats the whole $\eta$-unit source portion of C14.3a.
Every identity below uses the normalized quarter-parent $\mathcal P=F/4$.
The complete parent identity JF1--JF3 from R16 is algebraic in the
original stationary source; none of its identities assumes unit
$\zeta$. We retain that identity and change its source normalization.

\paragraph{Opening-cleared source bounds.}
Put
\[
 a=1/\zeta,\quad x=aL,\quad y=\eta L,\quad
 \theta=1-h=L\pi,\quad B=\beta+\tfrac12,\quad D_p=D+\tfrac12,
 \quad \rho_0=r+\eta.
\]
The original beta numerator gives bounded $\pi>0$;
$h,M,E,Q,P_d,P_4\sim1$ as in SB1--SB3. In particular
$\rho_0$ is a positive bounded unit. All formulas here preserve
actual $j>0$, rather than replacing it by $j_0$.
Define, exactly as in JF1,
\[
\begin{gathered}
 q=\alpha+(m+B)\theta^2,\quad \mu=m-a^2q,\quad
 \lambda=\mu/(\eta a^2L),\\
 T_0=A+\alpha+\tfrac12,\quad T_m=T_0+m,\quad
 e_4=r^2T_m/\eta^2,\\
 b'=Ar/\eta-m\theta-\alpha-\tfrac12+rT_mL/2,\\
 \kappa=2[B\theta-D_pj/\eta-mr/\eta],\qquad
 \gamma=2b'/h+2D_pj/r-2\eta B\theta/r.
\end{gathered}                                                    \tag{EU1}
\]
Finite canceled source division gives
\[
\begin{gathered}
 C(1-r)\sim\rho_0/2,\quad Ar\sim\eta/2,\quad
 D_pj\sim\rho_0/2,\quad m/(a^2L)\sim\rho_0/2,\\
 \alpha/L^2\sim\rho_0/2,\quad B\theta/L\sim\rho_0/2,
 \quad q/L^2-\rho_0(1+\pi)/2\in\operatorname{Inf},\\
 \mu/(a^2L)\sim\rho_0/2,\quad rT_m\sim\rho_0/2,
 \quad b'/L\sim-\rho_0/4,\quad \eta\kappa\sim-\rho_0.
\end{gathered}                                                    \tag{EU2}
\]
The canceled denominators at $a=L=0$ are only $1,2,4$;
there is no $r$, $1-r$, $j$, $\theta$ or $j_0$ factor.
Thus these are uniform through every opening in the theorem.
In particular $C,A,D_p$ are bounded below by positive real constants,
although any of them may be unbounded. Also $q/L^2$ and $\lambda$
are positive bounded units, $\mu>0$, and $b'=O(L)$.
The positive lower bound for $D_p$ uses bounded positive actual $j$
and the third estimate, not positivity of $j_0$.

\paragraph{Use the complete homogeneous parent.}
For nonnegative testing angles $u,v$ put $U=u^2,V=v^2$ and define
\[
\begin{gathered}
 E_t=\eta uW-vZ,\quad Q_t=Y-rvZ/\eta,\quad
 P=E_t/\eta-Q_t,\quad \xi=X+huZ,\\
 R_D=v(U-x)Z,\quad T_D=R_D-uv\xi/h,\quad
 H=avW-uZ-a\xi,\\
 F_b=(V-y)Z+vE_t,\qquad
 F_t=(V-y)Z+(\eta/r)vQ_t.
\end{gathered}                                                    \tag{EU3}
\]
All inverses in these roots are pure source. The exact full parent is
\[
\begin{split}
\mathcal P={}&qH^2+\mathscr C
 +CP^2+AQ_t^2+D_p(P-hT_D)^2+2b'Q_tT_D+\mathscr R\\
 &+e_4UF_t^2-e_4yu(H+a\xi)F_t
 +AUQ_t^2+B(1+V)\xi^2-2(\mu/a)H\xi-\mu\xi^2\\
 &+\kappa uF_t\xi+\gamma uvQ_t\xi .
\end{split}                                                       \tag{EU4}
\]
This is the complete JF1 identity, including every hard row and every
first normal. Set
\[
 \rho=m+rT_mx/2=O(x),\qquad
 K_r=T_0+2\rho-\rho^2/C.
\]
Then $K_r/T_m\sim1$ and $K_r$ is bounded below by a positive
real constant. The complete radial row is exactly
\[
\mathscr R=R_b^2+R_t^2,\qquad
 R_b={\sqrt C\over\eta}[F_b-r(1-\rho/C)F_t],\qquad
 R_t={r\sqrt{K_r}\over\eta}F_t.                         \tag{EU5}
\]
It follows that
\[
 {F_b\over\eta}={R_b\over\sqrt C}
               +{1-\rho/C\over\sqrt{K_r}}R_t          \tag{EU6}
\]
is a bounded pure-source allocation even when $r$ is infinitesimal.
No independent bare $(V-y)Z$ pivot has been inserted.

\paragraph{A strict source frame without either gap inverse bound.}
Use the three original weighted triangle roots
\[
 N_C=\sqrt C\,P,\quad N_A=\sqrt A\,Q_t,\quad
 N_D=\sqrt{D_p}(P-hT_D).                               \tag{EU7}
\]
Their positive squares are exactly the three original source-triangle
terms in EU4. They give the bounded pure-source allocations
\[
 P={N_C\over\sqrt C},\quad Q_t={N_A\over\sqrt A},\quad
 T_D={N_C/\sqrt C-N_D/\sqrt{D_p}\over h},\quad
 E_t=\eta(N_C/\sqrt C+N_A/\sqrt A).                    \tag{EU8}
\]
Consequently the entire $2b'Q_tT_D$ term is an infinitesimal Gram
on this strict diagonal frame. This does not cap or remove a divergent
original event and does not assert boundedness of its coefficient.

\paragraph{The complete two-normal cancellation.}
The source definitions in EU1 give the exact identity
\[
             {\eta\kappa\over r}+\gamma=2b'/h-2m.
\]
Since $F_t=F_b-vE_t+(\eta/r)vQ_t$, both complete normal rows in EU4
therefore satisfy
\[
\boxed{\quad
 \kappa uF_t\xi+\gamma uvQ_t\xi
 =\kappa uF_b\xi
     +uv\{(2b'/h-2m)Q_t-\kappa E_t\}\xi .\quad}       \tag{EU9}
\]
This is the crucial simultaneous allocation. It cancels the potentially
unbounded $1/r$ coefficients before any estimate is made. It neither
estimates the two original rows separately nor uses their contact
vanishing as a substitute for the identity.

Set $N_x=\sqrt B\xi$ and $\mathcal Q=\sqrt q H$.
By EU2, EU6 and EU8, every coefficient of EU9 on
$(N_C,N_A,R_b,R_t,N_x)$ is infinitesimal: the first uses
$\eta\kappa=O(1)$ and an extra $u$; the second uses
$b',m\in\operatorname{Inf}$ and an extra $uv$; the third uses
$\eta\kappa=O(1)$ and an extra $uv$.
The factors $1/\sqrt B$, $1/\sqrt A$, $1/\sqrt C$ and
$1/\sqrt{K_r}$ are bounded. No testing-state normalization occurs.

\paragraph{Every remaining scalar and hard term.}
The complete angular block on $(\mathcal Q,N_x)$ is
\[
 \mathcal Q^2-{2\mu\over a\sqrt{qB}}\mathcal QN_x
                         +(1-\mu/B)N_x^2.
\]
Its off-diagonal coefficient is $O(a)$ and $\mu/B$ is infinitesimal:
\[
 {\mu\over a\sqrt{qB}}
   =a\,{\mu/(a^2L)\over\sqrt{q/L^2}\sqrt B}.
\]
Thus its residue is the identity, including literal $\beta=0$.
The radial-angular cross has the exact normalized expression
\[
 -e_4yu(H+a\xi)F_t
 =-{rT_mL\over\sqrt{K_r}}u
       \left({\mathcal Q\over\sqrt q}
                         +{aN_x\over\sqrt B}\right)R_t.          \tag{EU10}
\]
Both coefficients are infinitesimal because $rT_m=O(1)$,
$L/\sqrt q=O(1)$ and $u$ is infinitesimal.

The coupled circuit $\mathscr C$ is exactly the sum of the following
three whole positive squares, followed by $-\lambda xE_t^2$:
\[
 \sqrt\lambda(uv-x/c)Z,\qquad
 \sqrt\lambda\,\eta k(cv-u)W,\qquad
 \sqrt{2\lambda/c}\,k\sqrt{uv}(\eta cW-Z),
 \qquad c^2=a/\eta,\quad k^2=x.                       \tag{EU11}
\]
The last negative term is an infinitesimal Gram on EU8 since
$\lambda$ is bounded and $x$ infinitesimal. All radicals in EU11
are separated pure-source and pure-testing radicals.

Retain EU11's three squares, $e_4UF_t^2$, $AUQ_t^2$ and
$BV\xi^2$ as whole positive terms. Every other term of EU4 has now
been assigned explicitly to the seven-root frame
\[
                 (N_C,N_A,N_D,R_b,R_t,\mathcal Q,N_x).             \tag{EU12}
\]
The Gram on this frame has identity residue and bounded coefficients.
The bounded-Gram lemma proves finite tensor SOS for the complete
parent whenever both testing angles are infinitesimal. Every estimate
is a product of a source bound and a testing bound; there is no
source/testing order comparison and no inverse of a mixed tensor.

\paragraph{All tests, endpoints and states.}
The roots and the whole-parent identity are homogeneous in all four
state coordinates. In particular no $Z$ inverse or state restriction
was used. The positive circuit rows EU11 are defined on literal
$u=0$ or $v=0$ and their squared identity is polynomial there.
The common exact sign gauges give all local sign charts.
The unchanged small-$L$, bounded-$\eta$ remote theorem covers the
complement of both infinitesimal testing angles on this exact source
domain, retaining all its projective endpoints and openings.
Applying the real homogeneous identity to real and imaginary state
parts proves the complex-state assertion. All of the source unit-$\eta$
portion of C14.3a is therefore all-tests covered.

\paragraph{The complementary source.}
The remaining C14.3a sources have $\eta$ infinitesimal, hence $r\sim1$, with all original $1-r,j$ openings retained. Their complete parent is treated by the LD identities and the completed $J$-row conductor in Sections~\ref{component:pc720c1e06aca} and~\ref{component:pe3bbfd3a8151}.



% Source formulas: p024c821bfb59
\subsection{Small $\eta$ and unit $\zeta$: the damped positive block}

\label{component:p024c821bfb59}

\begin{theorem}
Normalize $d=1$ and retain the unchanged physical C14.3b source, with $x,\eta>0$
infinitesimal, $\zeta$ a positive bounded unit, every original $j$
and $1-r$ opening, $\beta\ge0$ including literal $\beta=0$, and all common intersections.
On the additional \emph{pure-source} subchart
\[
                         x/\eta\text{ bounded},              \tag{DG1}
\]
the entire original parent is finite tensor SOS for every independent
projective angle and state test. No relation between a source rate
and a testing rate is assumed. The proof below treats the whole local
corner; the common remote theorem gives its complement. The source
subchart $x/\eta$ unbounded is not covered by this theorem.
\end{theorem}

\paragraph{The sign reduction and the complete frame.}
Use the sign gauges to take $u,v\ge0$. For real states with $WZ\le0$
the independent half-core sign lemma proved in this revision gives a
direct tensor-SOS decomposition of the original signed half core;
all remaining original event terms are positive source multiples
of squares. Thus only $WZ>0$ needs the following construction. No
normalization or inverse of a state coordinate is used here.

Use all homogeneous JF1 definitions, and abbreviate $Q=Q_t$,
$T=T_D$, $B=B_x$, $a=1/\zeta$, $x=k^2$, $a=\eta c^2$. In particular
\[
 E=\eta uW-vZ=\eta(P+Q),\quad H=avW-uZ,
 \quad \tau=\eta cW-Z,\quad D_0=u-cv,
\]
\[
 J=(uv-x/c)Z,\quad
 R=-\sqrt\lambda\,\eta kD_0W,\quad
 f_t=F_t/\eta,\quad \mathcal Q=\sqrt q(H-a\xi),
 \quad N_x=\sqrt B\xi.
\]
The full radial term $\mathscr R$ is JF2, with the strict two-root
completion JF3. Its roots will be denoted $R_b,f_t$; both normalized
pivots have positive unit residue. The angular term is exactly
\[
 \mathscr A=\mathcal Q^2
       -\frac{2\mu}{a\sqrt{qB}}\mathcal QN_x
       +(1-\mu/B)N_x^2.                              \tag{DG2}
\]
Its normalized Schur pivot is strict by the unchanged R15/R16 source
margin, including $\beta=0$ and unbounded $B$. Those previously
proved angular and radial blocks are used once, as their complete
terms.

\paragraph{Damp the remaining normal and complete both circuit companions.}
Set the positive source unit
\[
 G_0=2\eta\lambda-2\eta\sim1,\qquad
 \epsilon=\gamma-G_0/h,
\]
\[
 \widehat R=R+\frac{G_0k}{2\sqrt\lambda}Q,
 \qquad \widehat J=\sqrt\lambda J+
                         \frac{G_0u}{2\sqrt\lambda}Q,
 \qquad c_U=A-\frac{G_0^2}{4\lambda}.
\]
Use the exact full-$D$ conductor and polarization displayed in Chapter~\ref{ch:mechanisms} to conduct only
$(G_0/h)uvQ\xi$, leaving $\epsilon uvQ\xi$ unaltered. Complete
\emph{both} $R$ and $J$, incorporating the entire native $AUQ^2$
row into the latter completion. Since
$G_0/(2\lambda)=\eta-\eta/\lambda$, the complete original parent
has the following exact identity:
\[
\begin{split}
 F/4={}&\mathscr A+\mathscr R+S_0
    -\lambda x(\eta P+\eta Q/\lambda)^2
    +\widehat R^2+\widehat J^2\\
 &+c_UUQ^2+e_4UF_t^2+BV\xi^2
    +\frac{2x}{c}uv\tau^2
    +\frac{G_0x}{c}D_0^2WZ\\
 &-e_4yuHF_t+\kappa uF_t\xi
                 +\frac{G_0x}{c}PH+\epsilon uvQ\xi,       
\end{split}\tag{DG3}
\]
where
\[
 S_0=CP^2+AQ^2+D_p(P-hT)^2+2b_0QT,
 \qquad b_0=b'-G_0/2.                                 \tag{DG4}
\]
This is a simultaneous completion of the entire remaining normal and
the coupled circuit. It retains every full-$D$, full-$Y$ and higher
normal of JF1. There is no second bare radial pivot and no extracted
uncentered $m$ event. The coefficient of the extra $uv\tau^2$ row
is exactly $2x/c>0$, because $2\lambda-G_0/\eta=2$.

\paragraph{Two strictly positive source deficits replace the failed soft pivot.}
Put $K=CD_p/(C+D_p)$, $s=x+\eta$, and
\[
                         \Delta_0=h^2AK-b_0^2.
\]
The original source formulas imply exact finite divisions
\[
 \Delta_0=x d_x+\eta d_\eta,\qquad
 \frac{c_U}{\eta}=x e_x+\eta e_\eta,                \tag{DG5}
\]
with bounded source coefficients having the following residues:
\[
 \begin{array}{c|c}
 d_x&\pi/2+1/4+1/(4a)\\
 d_\eta&\sigma/2+3/4\\
 e_x&\pi+1+1/(2a)\\
 e_\eta&\sigma+3/2
 \end{array}
 \qquad h=1-x\pi,\quad r=1-\eta\sigma.             \tag{DG6}
\]
All four residues are strictly positive. Here $\pi,\sigma\ge0$
are bounded by the unchanged physical source estimates. Thus
\[
 \Delta_0/s\text{ is a positive bounded unit},\qquad
 c_U/(\eta s)\text{ is a positive bounded unit}.     \tag{DG7}
\]
These assertions compare source quantities only and do not use DG1.

For clarity, the divisions in DG5 are opening-safe. Write
$C_n=(1-r)C$, $D_j=jD_p$ and $\Theta_B=B(1-h)$, all given by their
unchanged unit-denominator source formulas. The identity
$\eta\Theta_B-D_j=-C_n$ gives
$b'-\gamma h/2=-hC_n/r$. The gap-cleared harmonic coefficient is
\[
 \eta K=\frac{C_nD_j}
 {C_n(1+xh-\sigma)+D_j\sigma}.
\]
Its denominator is a positive unit, even at either opening. The
quantities $\eta\lambda$, $b'$, $A/\eta$ and $\gamma h$ also have
unit-denominator expressions after their displayed monomial factors
are canceled. Both functions on the left of DG5 vanish at
$x=\eta=0$. Define their $\eta$ coefficients by their exact $x=0$
values divided by $\eta$, and then divide the remainder by $x$.
This gives DG5 with only source-unit denominators; differentiation at
the origin gives DG6. No assumption on $x/\eta$ enters these exact
finite divisions.

\paragraph{The source roots control the complete soft direction.}
There is the exact pure-source completion
\[
 S_0=N_s^2+N_T^2+N_Q^2,
\]
\[
 N_s=\sqrt{C+D_p}\left(P-\frac{hD_p}{C+D_p}T\right),
 \quad N_T=h\sqrt K\left(T+\frac{b_0}{h^2K}Q\right),
 \quad N_Q=\sqrt{\frac{\Delta_0}{h^2K}}Q.            \tag{DG8}
\]
As $\eta K\sim1/2$, DG7 gives
$N_Q/\{\sqrt{\eta s}\,Q\}$ a positive source unit. The inverse
identities show
\[
 T=\frac{N_T}{h\sqrt K}-\frac{b_0}{h^2K}Q,
 \qquad
 P=\frac{N_s}{\sqrt{C+D_p}}
                      +\frac{hD_p}{C+D_p}T.         \tag{DG9}
\]
Consequently $P$ has bounded coefficients on $(N_s,N_T,N_Q)$:
the first two coefficients are $O(\sqrt\eta)$ and the last is
$O(\sqrt{\eta/s})$, which is bounded because $\eta\le s$.
This is the complete $P$ control after the full $D$ shift; it is not
an assumption that the old $E/\sqrt\eta$ pivot remains strict.

The negative square in DG3 is an infinitesimal Gram on these three
roots. Its root is
\[
 \sqrt{\lambda x}(\eta P+\eta Q/\lambda)
   =\eta\sqrt{\lambda x}\,P+\eta\sqrt{x/\lambda}\,Q.
\]
Since $\eta\lambda\sim1/2$, the first coefficient is
$O(\sqrt{\eta x})$ on the bounded allocation of $P$; the second,
after normalizing $Q$ by $N_Q$, is
$O(\eta\sqrt{x/s})$. Both are infinitesimal, without DG1.

\paragraph{Every remaining complete normal is small on this frame.}
The angular-radial term in DG3 is
\[
 -e_4y\eta u\left(\mathcal Q/\sqrt q
                             +aN_x/\sqrt B\right)f_t.
\]
Its coefficients are $O(u)$ and $O(xu)$, using
$e_4\eta^2\sim1/2$ and $q/x^2$ a positive bounded unit.
The term $\kappa uF_t\xi$ has coefficient
$(\eta\kappa)u/\sqrt B$, again infinitesimal.
The new $PH$ term has coefficients
\[
 \frac{G_0x}{c\sqrt q}\quad\hbox{and}\quad
 \frac{G_0ax}{c\sqrt B}
\]
on $P\mathcal Q$ and $PN_x$; these are respectively
$O(\sqrt\eta)$ and $O(x\sqrt\eta)$, so DG9 allocates the whole
term infinitesimally.

Finally $\epsilon=O(s)$ by the same gap-free source expressions.
The complete, still unconducted normal has coefficient
\[
 \frac{\epsilon uv}{\sqrt{\eta s}\sqrt B}
         =O\!\left(\sqrt{s/\eta}\,uv\right)          \tag{DG10}
\]
on $N_QN_x$. DG1 is used \emph{exactly here}: it makes $s/\eta$
bounded, and hence the coefficient infinitesimal for independent
infinitesimal $u,v$. No testing comparison is imposed.

Thus the remainder on
\[
                  (N_s,N_T,N_Q,R_b,f_t,\mathcal Q,N_x)
\]
has a bounded Gram with strictly positive residue. The bounded-Gram
lemma makes it finite tensor SOS. Add the whole positive rows
$\widehat R^2$, $\widehat J^2$, $c_UUQ^2$, $e_4UF_t^2$,
$BV\xi^2$, $(2x/c)uv\tau^2$ and
$(G_0x/c)D_0^2WZ$. The last is a literal tensor square on the current
pure-testing sign chart; all radicals are source-only or testing-only.
This proves the complete local parent, not merely its scalar
restriction.

\paragraph{Equality, endpoints and the global join.}
At original contact $P,Q,T,J,R,\tau,F_t,\xi$ vanish, so all newly
completed rows vanish too. DG3 preserves all contact normals exactly.
The row formulas have no angle ratio or state-coordinate inverse;
$\sqrt{uv}$ is defined as zero on a literal axis. Thus they include
both literal local axes. The common gauges cover every angle sign.
For a complex state, split the real-coefficient Hermitian parent into
its real and imaginary evaluations and apply either the $WZ\le0$
core lemma or this $WZ>0$ proof separately to each real state. There
is no assumption that both real and imaginary state components have
a unit $Z$ or $W$. Combine the local proof with the unchanged remote
all-state theorem to obtain the all-tests assertion on DG1.

The hypothesis DG1 is used only in the final allocation DG10. The exact DG3 identity and preceding source bounds also apply without it; the completed $J$-row conductor in Section~\ref{component:pe3bbfd3a8151} supplies the remaining allocation for arbitrary $x/\eta$.



% Source formulas: pc720c1e06aca
\subsection{Small $\eta$ and unbounded $\zeta$: the damped positive block}

\label{component:pc720c1e06aca}

\label{pc720c1e06aca:sec:r17-large-zeta-damped}

\begin{theorem}
Normalize $d=1$ and retain the unchanged physical C14.3a source,
its actual stationary contact, every strict source gap,
$\beta\ge0$ including literal $\beta=0$, the original total-scale
intersection and every earlier exclusion. Suppose $x,L=x\zeta,\eta$
are positive infinitesimals, $\zeta$ is unbounded, and
\[
                         L/\eta\text{ is bounded}.       \tag{LD1}
\]
Then the entire original parent is finite tensor SOS at every independent
projective angle and state test, real or complex. This is a whole-state,
all-tests source-subchart certificate. No condition is imposed on the
testing angle ratio or on the source ratio $a/\eta$, where $a=1/\zeta$.
The source subchart with $\eta$ infinitesimal and $L/\eta$ unbounded
is not covered here.
\end{theorem}

The proof uses the exact complete damped identity DG3, but its source
scale is $L+\eta$ rather than $x+\eta$. This change must be proved:
one cannot invoke the unit-$\zeta$ estimates by substituting an
infinitesimal $a$ into their displayed $1/a$ coefficients. All formulas
below use the actual physical source and the original contact.

\paragraph{Source variables and opening-cleared coefficients.}
Put
\[
 a=1/\zeta,\quad x=aL,\quad y=\eta L,\quad
 h=1-L\pi,\quad r=1-\eta\sigma,\quad B=\beta+\tfrac12.
\]
The original source estimates give bounded $\pi,\sigma\ge0$, and
actual $j>0$ says
$0<\sigma<1+xh$, $j=\eta(1+xh-\sigma)>0$.
Neither opening has a positive-residue lower bound.
The exact native source formulas displayed below give
\[
 \frac{m}{a^2L}\sim\frac12,\quad
 \frac{\alpha}{L^2}\sim\frac12,\quad
 \frac{q}{L^2}-\frac{1+\pi}{2}\in\operatorname{Inf},\quad
 \frac{\mu}{a^2L}\sim\frac12,
\]
\[
 \eta\lambda\sim\frac12,\quad
 \eta K\sim\frac12,\quad
 \eta\kappa\sim-1,\quad \eta^2e_4\sim\frac12,
 \quad b'=O(L),                                     \tag{LD2}
\]
where $q=\alpha+(m+B)(1-h)^2$, $\mu=m-a^2q$,
$\lambda=\mu/(\eta a^2L)$ and $K=CD_p/(C+D_p)$.

Here is a completely regular specification of the new source
expressions. Write
\[
 K_x=1+x,\ R_s=1+xr,\ T_s=1+xh,\ M_s=R_s-aL\pi,
 \ E_s=1-x\eta h,\ Q_s=1-Lr,
\]
\[
 P_d=1+2xh+rx^2h+rxL+L
                  +\eta\{K_xL+x^2hL\pi\},\qquad
 P_4=T_sE_s+(1-r)(xh+L).
\]
Define the opening-cleared quantities
\[
 C_n=\frac{K_x(r+\eta)E_s}{2P_d},\qquad
 D_j=\frac{K_x(1+y)(r+\eta)}{2P_d},\qquad
 A_\eta=\frac{P_4}{2rP_d},
\]
\[
 m_n=\frac{K_x(r+\eta)h(1+y)}{2M_sP_d},\quad
 \alpha_n=\frac{(r+\eta)\{aR_sh(1+y)+M_sE_s\}}
                         {2Q_sP_d},\quad
 B_n=\frac{K_x(r+\eta)(1+ah)}{2P_d}.                 \tag{LD3}
\]
Thus $C_n=(1-r)C$, $D_j=jD_p$, $A_\eta=A/\eta$,
$m_n=m/(a^2L)$, $\alpha_n=\alpha/L^2$ and
$B_n=B(1-h)/L$. In particular
\[
 q_n:=q/L^2=\alpha_n+m\pi^2+B_n\pi,\qquad
 \Lambda:=\eta\lambda=m_n-Lq_n,
\]
\[
 \eta K=\frac{C_nD_j}
                  {C_n(T_s-\sigma)+D_j\sigma}.       \tag{LD4}
\]
There is no $a^{-1}$ in (LD3)--(LD4). Their denominators at
$a=L=\eta=0$ have residues $1$, $2$ or $1/2$, independently
of $\pi,\sigma$. The remaining definitions are
\[
 T_m=A+\alpha+\tfrac12+m,\quad
 b'=A_\eta r-mL\pi-\alpha-\tfrac12+rT_mL/2,
\]
\[
 G_f:=\gamma h=2b'+2D_jh/r-2\eta LB_nh/r,
 \quad G_0=2\Lambda-2\eta\sim1,\quad b_0=b'-G_0/2.
                                                               \tag{LD5}
\]
These are the actual native-source expressions used in the completion.

\paragraph{Positive source deficits at the new scale.}
Set $s=L+\eta$ and retain DG's definitions
\[
 \Delta_0=h^2AK-b_0^2,
 \qquad c_U=A-\frac{G_0^2}{4\lambda},
 \qquad\epsilon=\gamma-G_0/h.
\]
The regular source formulas above give exact finite divisions
\[
 \Delta_0=Ld_L+\eta d_\eta,\qquad
 \frac{c_U}{\eta}=Le_L+\eta e_\eta,\qquad
 G_f-G_0=Lf_L+\eta f_\eta,                           \tag{LD6}
\]
with bounded coefficients and residues
\[
\begin{array}{c|cc}
 &L\text{ coefficient}&\eta\text{ coefficient}\\\hline
 \Delta_0&\pi/2+1/4&\sigma/2+3/4\\
 c_U/\eta&\pi+1/2&\sigma+3/2\\
 G_f-G_0&\pi+1/2&\sigma+2
\end{array}                                                     \tag{LD7}
\]
(as usual the bounded $\pi,\sigma$ on the right can be replaced
by their real residues). To justify the finite divisions explicitly,
each expression on the left of (LD6) vanishes at $L=\eta=0$
\emph{for arbitrary $a$}. Subtract its exact value at $L=0$, divide
that difference by $L$, and divide the $L=0$ value by $\eta$.
Clearing the source-unit denominators proves each divisibility in a
polynomial ring. Evaluation of these quotients at $a=L=\eta=0$
gives (LD7). There is consequently no hidden pure-$a$ remainder.

All four coefficients for the first two rows have strictly positive
residue. Hence
\[
 \Delta_0/s\text{ and }c_U/(\eta s)
       \text{ are positive bounded source units},\qquad
 \epsilon=O(s).                                      \tag{LD8}
\]
This source calculation itself does not use (LD1).

\paragraph{The same complete damped identity, without a state chart.}
Use DG's unshifted $H=avW-uZ$; the strict angular root is
$\mathcal Q=\sqrt q(H-a\xi)$, so this $H$ differs from the
shifted letter $H$ in EU3. Retain every other homogeneous JF/DG
definition. In particular
\[
 E=\eta(P+Q),\quad N_x=\sqrt B\xi,\quad
 \tau=\eta cW-Z,\quad D_0=u-cv,\quad c^2=a/\eta.
\]
After the pure-testing sign gauges, the half-core lemma supplies all
$WZ\le0$ states. On $WZ>0$, the exact identity DG3 is valid for
this source without any change:
\[
\begin{split}
 F/4={}&\mathscr A+\mathscr R+S_0
       -\lambda x(\eta P+\eta Q/\lambda)^2
       +\widehat R^2+\widehat J^2\\
 &+c_UUQ^2+e_4UF_t^2+BV\xi^2
       +\frac{2x}{c}uv\tau^2+\frac{G_0x}{c}D_0^2WZ\\
 &-e_4yuHF_t+\kappa uF_t\xi
       +\frac{G_0x}{c}PH+\epsilon uvQ\xi.             
\end{split}\tag{LD9}
\]
Here $S_0=CP^2+AQ^2+D_p(P-hT)^2+2b_0QT$,
$\widehat R=R+G_0\sqrt xQ/(2\sqrt\lambda)$ and
$\widehat J=\sqrt\lambda J+G_0uQ/(2\sqrt\lambda)$.
The identity conducts only $G_0uvQ\xi/h$ and completes both
circuit companions exactly, retaining every first normal and hard
coupling. Its derivation uses only $x=k^2$, $c^2=a/\eta$ and the
homogeneous JF conductor, not a bounded reciprocal of $a$ or $c$.

\paragraph{Strict source, radial and angular frames.}
Use the exact source roots
\[
 N_s=\sqrt{C+D_p}\left(P-\frac{hD_p}{C+D_p}T\right),
 \quad N_T=h\sqrt K\left(T+\frac{b_0}{h^2K}Q\right),
 \quad N_Q=\sqrt{\frac{\Delta_0}{h^2K}}Q.            \tag{LD10}
\]
Their squares sum to $S_0$, and
$N_Q/[\sqrt{\eta s}\,Q]$ is a positive source unit.
The inverse source identities give bounded coefficients for $P$
on this frame: the coefficients of $N_s,N_T$ are $O(\sqrt\eta)$
and that of $N_Q$ is $O(\sqrt{\eta/s})$. This uses only
$\eta\le s$, so no old $E/\sqrt\eta$ pivot is imported.

The complete radial term has its common strict roots $R_b,f_t$,
with $f_t=F_t/\eta$. Indeed $\rho=m+rT_mx/2=O(x)$ and
$K_r\sim1/2$. The complete angular block is
\[
 \mathscr A=\mathcal Q^2-
       \frac{2\mu}{a\sqrt{qB}}\mathcal QN_x
                  +(1-\mu/B)N_x^2.
\]
It has identity residue, since
\[
 \frac{\mu}{a\sqrt{qB}}
       =\frac{a[\mu/(a^2L)]}{\sqrt{q/L^2}\sqrt B}=O(a),
 \qquad\mu/B\in\operatorname{Inf}.                 \tag{LD11}
\]
This includes unbounded $B$ and literal $\beta=0$.

\paragraph{The entire remainder has infinitesimal coefficients.}
The negative-square root in (LD9) has, on the source frame,
the same coefficient bounds as DG:
\[
 \eta\sqrt{\lambda x}\,P+
           \eta\sqrt{x/\lambda}\,Q,
 \qquad O(\sqrt{\eta x}),\quad
           O(\eta\sqrt{x/s}).                       \tag{LD12}
\]
They are infinitesimal; here $x=aL$ and no comparison of $a$ with
$\eta$ is needed. For the radial-angular term use
$H=\mathcal Q/\sqrt q+aN_x/\sqrt B$ and $F_t=\eta f_t$.
Since $e_4y\eta=(e_4\eta^2)L=O(L)$ and $\sqrt q\asymp L$,
its two coefficients are $O(u)$ and $O(xu)$.
The $\kappa$ term has coefficient
$(\eta\kappa)u/\sqrt B$, which is also infinitesimal.

The new $PH$ term improves at this source end:
\[
 \frac{G_0x}{c\sqrt q}=O(\sqrt{a\eta}),\qquad
 \frac{G_0ax}{c\sqrt B}=O(x\sqrt{a\eta}).            \tag{LD13}
\]
Thus the bounded $P$ allocation makes both of its angular coefficients
infinitesimal. Finally, the complete unprocessed normal is
\[
 \frac{\epsilon uv}{\sqrt{\eta s}\sqrt B}N_QN_x
   \quad\hbox{up to a bounded positive source-unit factor}.
\]
Its coefficient is
\[
                  O\!\left(\sqrt{1+L/\eta}\,uv\right),          \tag{LD14}
\]
which is infinitesimal under the bounded-source-ratio
hypothesis (LD1). This is the only use of that hypothesis.

Consequently the complete remainder on
$(N_s,N_T,N_Q,R_b,f_t,\mathcal Q,N_x)$ has a bounded strict-residue
Gram. The bounded-Gram lemma gives its finite tensor SOS. Every other
term in (LD9) is a whole positive square: in particular $c_U>0$,
$G_0>0$, and $uv,WZ$ are nonnegative pure-testing factors on this
sign chart. No mixed inverse or radical is taken.

\paragraph{All tests and exact remaining source end.}
No state normalization, state-coordinate inverse or testing ratio
occurs in this proof. The division-free circuit rows and
$\sqrt{uv}=0$ include both literal local axes. The half-core exit
and the exact angle sign gauges cover every real state; splitting
the real-coefficient Hermitian form covers complex states. The
unchanged small-$L$, bounded-$\eta$ remote theorem covers all other
projective testing pairs, including literal infinities. This proves
the theorem with every original source opening retained.

The restriction LD1 is used only in LD14. The complete LD9 identity, source bounds and all preceding normal allocations are inputs to the completed $J$-row conductor in Section~\ref{component:pe3bbfd3a8151}, which treats the full infinitesimal-$\eta$ source at arbitrary $L/\eta$. The unit-$\eta$ source is the EU theorem of Section~\ref{component:p60eca5442075}.



% Source formulas: p8cfada778490
\subsection{The direct large-$X/Z$ state bound}

\label{component:p8cfada778490}

In any local chart with $u$ infinitesimal, suppose $Z\ne0$ and
the pure-testing ratio $|X/Z|$ is unbounded.  Then
\[
 \frac{1+v^2}{2}\{(X+uZ)^2-(1+u^2)Z^2\}>0.             \tag{DS3}
\]
Indeed, after division by $Z^2>0$, the square $(X/Z+u)^2$
is unbounded and $1+u^2$ is bounded.  The factor $1+v^2$
is positive.  This expression is the complete $X$ square and the
negative $Z$ term of the signed quarter-core.  Its other two
terms are $(uW+uvX-Y-vZ)^2/2$ and
$(1+u^2)(vY-Z)^2/2$, both nonnegative.  Thus the signed core
itself is a nonnegative pure-testing scalar and has a square root
in the testing field.  Every original event remains a square with
its nonnegative source coefficient.  This proves a direct separated
SOS certificate for the complete parent, with arbitrary $W,Y$.

The state chart $Z=0$ and the sign chart $WZ\le0$ have the
direct core certificates already stated in the geometric reduction.
Consequently the C14.3 local certificate only needs the remaining
case $Z\ne0$, $WZ>0$, $X/Z$ bounded.  The completed-J
conductor CJ treats that whole case without bounding $W/Z$ or
$Y/Z$.  Apply the appropriate real chart to each real and imaginary
part to include complex states.



% Source formulas: pe3bbfd3a8151
\subsection{The completed $J$ relation for the entire remaining normal}

\label{component:pe3bbfd3a8151}

\label{pe3bbfd3a8151:sec:r18-completed-j}

\paragraph{Scope and unchanged inputs.}
Retain the complete actual-parent identity DG3, equivalently LD9,
with $d=1$, every original source opening and residual intersection,
and $\beta\ge0$ including literal $\beta=0$.  On C14.3b use
$s=x+\eta$ and the unit $a=1/\zeta$; on C14.3a use $s=L+\eta$
with $x=aL$ and $a$ infinitesimal.  In both cases $s$ and $\eta$
are positive infinitesimals.  The exact DG5--DG9 and LD3--LD10 source identities give
\[
 \Lambda=\eta\lambda\sim\tfrac12,\qquad
 \epsilon=O(s),\qquad
 N_Q=\sqrt{\frac{\Delta_0}{h^2K}}Q,
 \qquad
 Q=\frac{\omega}{\sqrt{\eta s}}N_Q,                 \tag{CJ1}
\]
where $\omega$ is a bounded positive source unit.  The complete strict
frame is
\[
 \mathcal F=(N_s,N_T,N_Q,R_b,f_t,\mathcal Q,N_x),
 \qquad N_x=\sqrt B\,\xi,\quad B=\beta+\tfrac12.
                                                               \tag{CJ2}
\]
All DG3/LD9 residuals other than $\epsilon uvQ\xi$ already have
infinitesimal Gram coefficients on this frame without any bound on
$s/\eta$.  The two residuals in the exact split CJ6a below have that same
property whenever $Z\ne0$ and $X/Z$ is bounded.  Thus the only
expression to be allocated on that state chart is
\[
 N=\epsilon U(W/Z)E\xi,
 \qquad E=\eta(P+Q)=\eta uW-vZ.                    \tag{CJ3}
\]
There is no restriction on $Y/Z$ or $W/Z$ below.

\paragraph{An exact identity from the already retained completed row.}
Let $\phi=\xi/Z=X/Z+hu$, and retain exactly DG's completed square
\[
 \widehat J=\sqrt\lambda\,(uv-x/c)Z
       +\frac{G_0u}{2\sqrt\lambda}Q,
 \qquad G_0=2\eta\lambda-2\eta,\qquad c^2=a/\eta.
\]
The defining identity $vZ=\eta uW-\eta(P+Q)$ gives
\[
 \boxed{\quad
 \frac{\widehat J}{\sqrt\lambda}
 =\eta U W-\frac{x}{c}Z
                -\eta u\left(P+\frac Q\lambda\right).
 \quad}                                                       \tag{CJ4}
\]
Multiply CJ4 by $\epsilon\phi E/\eta$.  The complete remaining
normal has the exact conductor
\[
 \boxed{\quad
 N=\frac{\epsilon\phi}{\eta\sqrt\lambda}E\widehat J
       +\epsilon\phi uE\left(P+\frac Q\lambda\right)
       +\frac{\epsilon x}{\eta c}E\xi .
 \quad}                                                       \tag{CJ5}
\]
The only state inverse here is the pure-testing $Z^{-1}$.  No inverse
of $u$, $v$, $W$, or a mixed expression occurs.  No radial row, source
pivot, or completed square is duplicated.

\paragraph{Bounded allocations and infinitesimal coefficients.}
Define two dependent roots
\[
 \mathcal E=\sqrt{s/\eta}\,E,
 \qquad \mathcal P=P+Q/\lambda.
\]
Both have bounded allocations on the existing source roots
$(N_s,N_T,N_Q)$.  Indeed the exact DG9/LD10 inverse identities give
coefficients $O(\sqrt\eta)$, $O(\sqrt\eta)$ and
$O(\sqrt{\eta/s})$ for $P$.  Hence the coefficients for
$\sqrt{\eta s}\,P$ are bounded (in fact infinitesimal), whereas
\[
 \mathcal E=\sqrt{\eta s}\,P+\omega N_Q,
 \qquad
 Q/\lambda=\frac{\omega}{\Lambda}\sqrt{\eta/s}\,N_Q.
                                                               \tag{CJ6}
\]
Here $\eta/s\le1$ is a pure-source inequality.  These are bounded
allocations of dependent roots, not new positive pivots.
For completeness, the other two terms of the same normal have
the particularly simple allocations
\[
 \epsilon uvQ\xi-N
   =-\frac{\epsilon}{s}u\phi\mathcal E^2
     -\frac{\epsilon uv}{\sqrt B}PN_x.               \tag{CJ6a}
\]
Their coefficients are infinitesimal, since $u,v$ are local
infinitesimals and $P$ has a bounded source-root allocation.

In these roots CJ5 is exactly
\[
\begin{split}
 N={}&\frac{\epsilon}{s}\frac{\sqrt s}{\sqrt\Lambda}
                  \phi\mathcal E\widehat J\\
 &+\epsilon\sqrt{\eta/s}\,u\phi\mathcal E\mathcal P
   +\frac{\epsilon}{s}
       \frac{x\sqrt s}{\sqrt a\sqrt B}\mathcal E N_x .
\end{split}                                                    \tag{CJ7}
\]
All three coefficients are infinitesimal bounded tensor elements.
For the first two use $\epsilon/s$ bounded, $s$ infinitesimal,
$\eta/s\le1$, and bounded
$\phi=X/Z+hu$.  The last coefficient is also infinitesimal:
\[
 \frac{x\sqrt s}{\sqrt a}=
 \begin{cases}
 O(x\sqrt s),&a\text{ a bounded positive unit},\\
 \sqrt a\,L\sqrt s,&x=aL,\ a\text{ infinitesimal}.
 \end{cases}                                                   \tag{CJ8}
\]
Also $1/\sqrt B$ is bounded, even at literal $\beta=0$ and when
$B$ is unbounded.  The apparent inverse $\sqrt a$ is canceled in
the large-$\zeta$ case before taking bounds.  Every coefficient is
a finite separated expression; $\phi$ is expanded as $X/Z+hu$.
No source rate is compared with any testing rate.

\paragraph{Full local reassembly.}
Add the already present row $\widehat J^2$ to the strict frame
CJ2 and allocate CJ7 to it and the source roots.  This produces an
infinitesimal Gram perturbation of the previous strict-residue
matrix plus a unit $\widehat J$ diagonal entry.  The separated
bounded-Gram lemma therefore applies to the full remainder.
Retain each other positive row from DG3/LD9 exactly once:
\[
 \widehat R^2,\quad c_UUQ^2,\quad e_4UF_t^2,\quad BV\xi^2,
 \quad(2x/c)uv\tau^2,\quad(G_0x/c)D_0^2WZ.
\]
The last two are literal tensor squares on the unchanged sign chart
$u,v\ge0$, $WZ>0$.  The complete angular and radial terms remain
inside CJ2; all full-$D$, full-$Y$ and higher normals are those of
the exact original-parent identity, not truncated expansions.

Consequently the entire local original parent is tensor SOS when
$Z\ne0$ and $X/Z$ is bounded, with no condition on $s/\eta$,
$u^2W/Z$, or $Y/Z$.  The proof is valid on literal local angle axes
because it uses no angle inverse.  At original contact $E,P,Q,\xi$
and $\widehat J$ vanish, so this reassembly preserves equality and
all original contact normals exactly.

\paragraph{The local state and complementary testing cases.}
The reviewed pure-testing half-core exit covers $WZ\le0$, including
$Z=0$.  When $Z\ne0$ and $X/Z$ is unbounded, DS3 in Section~\ref{component:p8cfada778490} is a direct
original-parent tensor-SOS exit.  The above argument covers every
remaining real state on the local angle chart.  Applying these real
certificates to real and imaginary state parts gives the full complex
parent.  Thus the shared local C14.3a/b obligation is closed, with
all source gaps and $\beta=0$ retained.  The SB remote theorem in Section~\ref{component:p6d2f659c0fde} and the unit-$\eta$ theorem in Section~\ref{component:p60eca5442075} supply the other testing and source charts. The complete source partition is printed in Section~\ref{component:pfd0dd7bb10cf}.



% Source formulas: p66922bb7f235
\section{Vanishing $L$ with unbounded \texorpdfstring{$\eta$}{eta}}
Both $L$ and $\sigma=1/\eta$ now vanish. With $S=L+\sigma$, the three source ratios $\kappa=x/S$, $\nu=\sigma/S$ and $\ell=L/S$ satisfy $\nu+\ell=1$. Their exact relations organize the remaining contacts, including the endpoints at which two of these ratios vanish.

\subsection{The small-$L$, unbounded-$\eta$ source sphere}

\label{component:p66922bb7f235}

\label{p66922bb7f235:sec:r13-small-L-eta-large-bounds}

This is a source theorem, not an additional all-tests certificate.
Use the original physical C1 source with $x,L=x\zeta>0$
infinitesimal, $\zeta$ noninfinitesimal and $\eta$ unbounded.
Put $\sigma=1/\eta$, $S=\sigma+L$, and retain the original
$0<r,h<1$, $\theta=1-h$, $M>0$, and $D,\beta\ge0$.
The native gap gives $E=1-xh/\sigma\ge1/18$. Set
\[
 e={xh\over\sigma}=1-E,\quad
 P=\sigma P_d,\quad J=1+xh-\sigma(1-r),\quad
 H_0=1+L(1-r),\quad B=P+x(J-\sigma),
\]
\[
 \mathcal H=R H_0+(1-r)xh(2+xr)
 -e\{2+x(1+h+r\theta)\}+Ke^2,
 \quad R=1+xr,\quad K=1+x.                           \tag{ES1}
\]
No $E\sim1$ assumption is made here. The original source formulas
give the two exact identities
\[
 2\theta P\beta=hB-\sigma H_0,\qquad
 2P(JD+x\theta\beta)=\sigma^2\mathcal H.             \tag{ES2}
\]
They are the beta-shared identities with the reciprocal-$L$ factors
cleared. In particular their use does not require bounded $1/L$.

The positive denominator is exactly
\[
 P=\sigma(1+L+2xh+rxL+rx^2h)+KL+x^2h\theta.
                                                               \tag{ES3}
\]
Since $x/L=1/\zeta$ is bounded and $x^2h=x\sigma e$,
all terms beyond $\sigma+L$ in (ES3), divided by $S$, are
infinitesimal. Consequently
\[
 {P\over S}\sim1,\qquad J\sim1,\qquad
 {B\over S}-(1+x/S)\in\Inf,\qquad
 H_0\sim1,\qquad \mathcal H-E^2\in\Inf.             \tag{ES4}
\]
Here $x/S$ is bounded. Both $B/S$ and $\mathcal H$ are
positive bounded units, uniformly through all original openings;
for the latter use the actual bound $E\ge1/18$.

Define nonnegative pure-source sphere coordinates
\[
 d_0^2={2PJD\over\sigma^2\mathcal H},\qquad
 b_0^2={2Px\theta\beta\over\sigma^2\mathcal H}.
\]
Then
\[
 d_0^2+b_0^2=1,\qquad
 D={\sigma^2\mathcal H\over2PJ}d_0^2,\qquad
 x\theta\beta={\sigma^2\mathcal H\over2P}b_0^2.
                                                               \tag{ES5}
\]
Thus $D=O(\sigma^2/S)=O(\sigma)$ and
$x\theta\beta=O(\sigma^2/S)$, including either literal coefficient
boundary. The first identity of (ES2) becomes
\[
 eB=xH_0+\sigma\mathcal H b_0^2.                     \tag{ES6}
\]
Put $\kappa=x/S$ and $\nu=\sigma/S$. Multiplying
$e(B/S)=\kappa H_0+\nu\mathcal H b_0^2$ by $e$ and using
$e\nu=\kappa h$ gives
\[
 e^2(B/S)=\kappa\{eH_0+h\mathcal H b_0^2\}.
                                                               \tag{ES7}
\]
All factors in braces are bounded and nonnegative. Hence
$e^2=O(x/S)$, and (ES6) also gives $x/S=O(e)$.
In particular $\zeta$ unbounded implies $e$ infinitesimal,
since $x/S\le1/\zeta$. The unit-$\zeta$ cases may have
$e$ a positive bounded unit.

The other original coefficients have the opening-cleared bounds
\[
 {CS(1-r)\over E}\sim\tfrac12,\qquad
 {ASr\over E}\sim\tfrac12,\qquad
 m={K(1+\sigma r)eS\over2M\zeta P},                  \tag{ES8}
\]
\[
 \alpha={(1+\sigma r)\{LR eS+L^2ME\}\over2QP}
       =O(L),\qquad Q=1-Lr\sim1.                   \tag{ES9}
\]
For $A$ use $P_4=(1+xh)E+(1-r)(xh+L)\sim E$.
The bound on $\alpha$ follows because $0<M<R$ is bounded and
$L/S\le1$. Thus $C,A$ are unbounded while $D,\alpha$ are
infinitesimal on the entire stated source domain. No bound on $m$
is asserted through $M=0$, and no bound on beta is asserted through
$\theta=0$.

Finally put $\mu=K(1+\sigma r)S/(2M\zeta P)>0$. The
original $m$ term has the exact shared-sphere split
\[
 m=m_0+m_1b_0^2,\qquad
 m_0={\mu xH_0\over B},\qquad
 m_1={\mu\sigma\mathcal H\over B}.                  \tag{ES10}
\]
The remaining source opening $M$ is displayed explicitly. These
identities provide the weights and boundedness estimates for C14.4.
Section~\ref{component:pfd0dd7bb10cf} assembles the full-parent
certificates for the remaining infinitesimal-$u$ tests and the
complementary testing charts. The inequalities above use only the
order in each individual field.



% Source formulas: p8c30238a5f7d
\subsection{The unmatched Bob residue}

\label{component:p8c30238a5f7d}

\label{p8c30238a5f7d:sec:recovery-c144-small-e}

This is a partial-domain theorem. It retains the original physical C14.4
source, all strict gaps, nonnegative event weights, the original total-scale
and native source inequalities, and independent source and testing fields.
Its source has $x,L>0$ infinitesimal, $\zeta=L/x$ noninfinitesimal and
$\eta$ unbounded. Write $\sigma=1/\eta$, $S=\sigma+L$ and
\[
 e=\frac{xh}{\sigma},\quad E=1-e,\quad
 \kappa=\frac{x}{S},\quad \nu=\frac{\sigma}{S},\quad
 \ell=\frac{L}{S},\quad \nu+\ell=1.
 \tag{SE1}
\]
The letter $\kappa$ in this section denotes $x/S$, not the different
coefficient denoted by $\kappa$ in the R19 reciprocal-$L$ chart.
All source inequalities are taken in the coefficient field. Real residues
are compared only after taking residues separately in the two fields.

\paragraph{The exact source input.}
Use ES1--ES10 from the retained R13 source proof. In particular
$1/18\le E<1$, $0<h,r<1$, $\theta=1-h$, and
\[
\begin{gathered}
 J=1+xh-\sigma(1-r),\qquad M=1+xr-\theta/\zeta>0,\\
 Q=1-Lr,\quad H_0=1+L(1-r),\quad B=P+x(J-\sigma),\\
 P=\sigma(1+L+2xh+rxL+rx^2h)+(1+x)L+x^2h\theta,\\
 \mathcal H=(1+xr)H_0+(1-r)xh(2+xr)
       -e\{2+x(1+h+r\theta)\}+(1+x)e^2.
\end{gathered} \tag{SE2}
\]
The opening-safe sphere identities are
\[
\begin{gathered}
 2\theta P\beta=hB-\sigma H_0,\qquad
 d_0^2+b_0^2=1,\quad d_0,b_0\ge0,\\
 D=\frac{\sigma^2\mathcal H}{2PJ}d_0^2,\qquad
 x\theta\beta=\frac{\sigma^2\mathcal H}{2P}b_0^2,\\
 e\frac BS=\kappa H_0+\nu\mathcal H b_0^2,\qquad
 e^2\frac BS=\kappa(eH_0+h\mathcal H b_0^2).
\end{gathered}\tag{SE3}
\]
Here $P/S\sim1$, $J,Q,H_0\sim1$, $B/S-(1+\kappa)$ is
infinitesimal, and $\mathcal H-E^2$ is infinitesimal. In particular
$B/S$ and $\mathcal H$ are positive bounded units. These estimates
are the proved ES bounds; they do not use bounded $\eta$.

\paragraph{Source dichotomy.}
On the entire C14.4 source,
\[
 e\text{ is infinitesimal}\quad\Longleftrightarrow\quad
 \kappa=x/S\text{ is infinitesimal}.\tag{SE4}
\]
Indeed the first nonnegative sum in SE3 gives $\kappa=O(e)$,
and the second gives $e^2=O(\kappa)$. Both implications are
ordered-field statements entirely inside the source field.

Suppose $e$ is infinitesimal. If $\zeta$ is unbounded, then
$\theta/\zeta$ is infinitesimal and $M\sim1$ directly. If $\zeta$
is a positive bounded unit, then the exact relation
\[
                  \ell=\zeta\kappa\tag{SE5}
\]
gives $\ell$ infinitesimal and $\nu\sim1$. Consequently
$B/\sigma\sim1$. Since $\beta\ge0$, the first identity of SE3
implies $hB\ge\sigma H_0$, whence $h\sim1$ and $\theta$ is
infinitesimal. Thus $M\sim1$ in this case as well. We have proved,
uniformly on the whole small-$e$ source,
\[
                    E\sim1,\quad M\sim1,\quad m\in\operatorname{Inf}_{>0}.
                    \tag{SE6}
\]
For the last statement use the exact original coefficient
\[
 m=\frac{(1+x)(1+\sigma r)eS}{2M\zeta P};\tag{SE7}
\]
$1/\zeta$ is bounded and every other coefficient apart from $e$ is
bounded, including the reciprocal of $M$ just proved to be a unit.
No assumption on the ratio of a source infinitesimal to a testing
infinitesimal occurs.

The two nonnegative terms in the first sum in SE3 imply
$\nu b_0^2$ infinitesimal. Hence the exact sphere formula and the
exact formula for $\alpha$ yield
\[
 \frac DS-\frac{\nu^2}{2}\in\operatorname{Inf},\qquad
 \frac\alpha S-\frac{\ell^2}{2}\in\operatorname{Inf},\qquad
 SC(1-r)\sim\tfrac12,\qquad SAr\sim\tfrac12.\tag{SE8}
\]
For clarity the second formula is obtained without dividing by $\ell$:
\[
 \frac\alpha S=
 \frac{(1+\sigma r)\{\ell(1+xr)e+\ell^2ME\}}
      {2Q(P/S)}.\tag{SE9}
\]
These are additive bounded residue estimates, including $D=0$ whenever
that boundary has points in the stated domain. They do not assert a
relative estimate at a vanishing pivot.

\paragraph{Theorem: every unmatched Bob residue on the small-$e$ source.}
Let $e$ be infinitesimal. In the testing field choose normalized homogeneous
Bob coordinates $(c,s)$ with $c^2+s^2=1$. For every independent projective
test and every real or complex state satisfying
\[
                         \overline{s^2}\ne\overline\ell,\tag{SE10}
\]
the complete original C14.4 parent is a finite sum of squares in
$K_c\otimes_{\mathbb R}K_t$.

\paragraph{Proof with the complete parent retained.}
Noninfinitesimal Alice tests use the unchanged adjacent-$C,A$ theorem:
SE8 and $S$ infinitesimal imply that both $C$ and $A$ are unbounded.
For infinitesimal or zero $u$, apply the exact EB3--EB7 homogeneous
frame from the retained R18 proof. Those formulas are algebraic identities;
their only source-limit inputs to the bounded-Gram proof are precisely
$m$ infinitesimal and SE8, now established on the larger source domain.
To make the imported implication explicit, put
\[
 H=S(C+A),\quad \rho=\frac C{C+A},\quad a_h=\frac{SCA}{C+A},
 \qquad p=\frac{uW}{\sqrt S},\ q=\sqrt S Z,\ y=\frac Y{\sqrt S}.
\]
Then $1/\sqrt H$ is bounded, $0<\rho<1$, $a_h\sim1/2$, and
the harmonic normal $N_y=\sqrt H(y-\rho p)$ can be substituted
by $y=\rho p+N_y/\sqrt H$ with bounded source coefficients.
The remaining normals are the EB4 $N_m,N_b,N_r$, with only pure-field
radicals and the positive testing denominators $1+u^2$ and $1+u^2s^2$.
The exact full quarter-parent is EB7,
\[
 N_y^2+a_hp^2+ASu^2y^2+M_B^2+D_B^2+A_B^2+\beta(X+uZ)^2
                         +\tfrac12 K_B.\tag{SE11}
\]
All EB5 core terms and all EB6 event amplitudes are retained. In particular
no state graph is imposed and neither first-normal coupling is removed.
The two positive rows $ASu^2y^2$ and $\beta(X+uZ)^2$ are left whole.
On the other six roots $(N_y,N_m,N_b,N_r,p,q)$ the bounded Gram has
residue equal to $\operatorname{diag}(1,1,1/2,1/2)$ plus the scalar block
\[
 \frac12\begin{pmatrix}
 1&-\bar s\bar c\\
 -\bar s\bar c&\bar\nu^{\,2}\bar s^{\,2}
                            +\bar\ell^{\,2}\bar c^{\,2}
 \end{pmatrix}.
\tag{SE12}
\]
Its determinant is $(\overline{s^2}-\bar\ell)^2/4$, so SE10
makes this finite coefficient Gram strictly positive definite. The
separated bounded-Gram lemma supplies tensor SOS; restoring the two
whole rows gives the complete parent. The exact homogeneous formulas
allow $c=0$, $s=0$ and $u=0$. Signs and real/imaginary state parts are
handled exactly as in EB, and the original positive source and testing
normalizations are pure-field factors. This proves the stated domain.

\paragraph{New part, and exact retained complement.}
Unbounded $\zeta$ already forces $\kappa\le1/\zeta$ infinitesimal,
so the new theorem contains the entire R18 unmatched-residue certificate.
Its additional domain is $\zeta$ a positive bounded unit, $e$
infinitesimal, and Bob sine square with nonzero residue; this includes
Bob infinity. In this added domain $\ell$ has zero residue, so the
remaining Bob chart is exactly finite infinitesimal $v$ (including zero,
which may be handled by independent literal-axis exits).

After the unmatched-residue certificates, the complementary C14.4 domain is contained in the union of (i) $e$ infinitesimal and matched Bob residue, and (ii) $\zeta$ a positive bounded unit and $e$ a positive bounded unit. All original intersections remain in both. In case (ii), SE7 shows that $m$ has a positive real lower bound, while its upper bound may diverge through $M=0$. This is a source estimate. The complete certificates for both complementary domains, with their exact source partition, are joined in Section~\ref{component:pfd0dd7bb10cf}.



% Source formulas: p1a6cc124f5db
\subsection{The homogeneous unmatched-Bob quadratic}

\label{component:p1a6cc124f5db}

\label{p1a6cc124f5db:sec:r18-c144-bob-residue}

\paragraph{Precise certificate.}
Retain the entire original C14.4 source: $x,L>0$ infinitesimal,
$\eta$ unbounded, $\zeta=L/x$ noninfinitesimal, all original physical
intersections, every positive opening, and either permitted literal
$D=0$ or $\beta=0$. This section additionally assumes $\zeta$
unbounded. Put
\[
 \sigma=1/\eta,\qquad S=\sigma+L,\qquad
 \nu=\sigma/S,\qquad \ell=L/S=1-\nu.
\]
For a real projective Bob test choose its normalized homogeneous
coordinates $(c,s)$, so $c^2+s^2=1$; on the finite chart these are
$c=(1+v^2)^{-1/2}$, $s=vc$. The complete original parent is tensor SOS
at every independent test and state satisfying
\[
                         \overline{s^2}\ne\overline\ell.       \tag{EB1}
\]
Bars mean real residues of bounded elements in their respective fields.
This compares real residues only, not independent infinitesimal rates.
The adjacent-$C,A$ theorem already covers noninfinitesimal Alice
tests. Hence only $u$ infinitesimal is proved below. No all-tests
closure of C14.4 is asserted: the matched residue and bounded-unit
$\zeta$ source remain.

\paragraph{Source limits, including both sphere endpoints.}
Use ES1--ES10 without changing their source entry. Since
$x/S\le x/L=1/\zeta$ is infinitesimal, ES7 gives
$e=xh/\sigma$ infinitesimal, and hence $E\sim1$.
Moreover $M=1+xr-\theta/\zeta\sim1$. ES8 then gives $m$ infinitesimal.
The exact equation ES6, divided by $S$, is
\[
 e(B/S)=(x/S)H_0+\nu\mathcal H b_0^2.
\]
Both summands are nonnegative; $B/S$ and $\mathcal H$ tend to one.
Consequently $\nu b_0^2$ is infinitesimal. Therefore ES5 yields
\[
 \frac DS-\frac{\nu^2}{2}\in\operatorname{Inf},\qquad
 \frac\alpha S-\frac{\ell^2}{2}\in\operatorname{Inf},
 \qquad m\in\operatorname{Inf}_{>0}.                         \tag{EB2}
\]
For the second formula, divide the exact ES9 expression by $S$:
\[
 \frac\alpha S=
 \frac{(1+\sigma r)\{\ell R e+\ell^2ME\}}
 {2Q(P/S)}.
\]
All displayed denominator factors tend to one. These are additive
residue assertions, valid also when $\nu$ or $\ell$ is infinitesimal.
In particular no division by either sphere coordinate is used.

The opening-cleared ES8 bounds give
\[
 c_*=SC(1-r)\sim\tfrac12,\qquad
 a_*=SAr\sim\tfrac12.
\]
Define the pure-source harmonic quantities
\[
 H=S(C+A),\qquad \rho=C/(C+A),\qquad
 a_h=SCA/(C+A)=\frac{c_*a_*}{c_*r+a_*(1-r)}.
                                                               \tag{EB3}
\]
Then $1/\sqrt H$ is bounded, $0<\rho<1$, and $a_h\sim1/2$,
uniformly through both $r$ openings. No upper bound on $C$ or $A$
has been assumed.

\paragraph{An exact all-state frame for every Bob chart.}
Let $D_u=1+u^2$ and $A_u=1+u^2s^2$, both positive testing units.
Use the separated roots
\[
 p=\frac{uW}{\sqrt S},\quad q=\sqrt S Z,\quad
 y=\frac Y{\sqrt S},\quad N_y=\sqrt H(y-\rho p),
 \quad N_m=\sqrt m\,sW,
\]
\[
 N_b=\sqrt{A_u}\left(X+\frac{uc^2}{A_u}Z\right),
 \qquad N_r=\frac{\sqrt{D_u}\,usc}{\sqrt{A_u}}Z.
                                                               \tag{EB4}
\]
Only pure-source and pure-testing radicals or inverses occur.
In particular no inverse of $c$, $s$, $u$, or a state coordinate is
needed. The inverse harmonic allocation
$y=\rho p+N_y/\sqrt H$ has bounded coefficients.

Write $b=X+uZ$. The exact original signed core $K/D_v$ in finite
Bob coordinates equals the following homogeneous expression:
\[
\begin{split}
 K_B={}&N_b^2+N_r^2-2D_usc\,pq
       +Sc^2(p-y)^2+SD_us^2y^2\\
 &+\frac{2\sqrt S\,usc}{\sqrt{A_u}}(p-y)N_b
  +\frac{2\sqrt S\,us^2\sqrt{D_u}}{\sqrt{A_u}}(p-y)N_r.
\end{split}                                                    \tag{EB5}
\]
Here $D_us$ denotes $D_u\,s$, and similarly in the other products.
One derivation starts with CE2 divided by $D_v$, sets
$R=usZ$, and completes its entire $(b,R)$ block:
\[
 A_ub^2-2D_us\,bR+D_uR^2=N_b^2+N_r^2.
\]
The cross term from
$[\sqrt S\,c(p-y)+usb]^2$ is allocated by
$cb=cN_b/\sqrt{A_u}+\sqrt{D_u}sN_r/\sqrt{A_u}$.
This proves EB5 without deleting a first state normal.

The original positive event roots have the exact bounded expressions
\[
 \begin{split}
 M_B={}&N_m-\frac{\sqrt m\,c}{\sqrt{A_u}}N_b
       -\frac{\sqrt{mD_u}\,s}{\sqrt{A_u}}N_r
       +\sqrt{mS}\,us\,y,\\
 D_B={}&\sqrt{DS}\,c(p-y)
       +\frac{\sqrt D\,us}{\sqrt{A_u}}N_b
       -\sqrt{D/S}\,\frac{D_us}{A_u}q,\\
 A_B={}&\sqrt{\alpha D_uS}\,s y
       -\sqrt{\alpha D_u/S}\,c q.
 \end{split}                                                   \tag{EB6}
\]
They are, respectively, the normalized original $m,D,\alpha$ event
amplitudes. The formulas are valid at $D=0$ as well. In the $D$ row,
the identity
$usb-D_usZ=usN_b/\sqrt{A_u}-D_usZ/A_u$
avoids an apparent reciprocal-$c$ pole.

Consequently the full original quarter-parent, divided by $D_v$ on
the finite chart, is exactly
\[
 \mathcal P_B=N_y^2+a_hp^2+ASu^2y^2
              +M_B^2+D_B^2+A_B^2+\beta b^2+\tfrac12K_B.
                                                               \tag{EB7}
\]
The two explicit positive rows $ASu^2y^2$ and $\beta b^2$ can be
retained unchanged. All other rows in EB7 use EB4 once. The formula
is a full four-state identity and is not a restriction to a state ray.

\paragraph{Strictness precisely off the matched residue.}
Every coefficient of EB5 except the printed $pq$ term is either
its fixed positive pivot or infinitesimal after EB3. In EB6, the
$m$ row is $N_m$ plus infinitesimal coefficients; the $D$ and
$\alpha$ rows have only their $q$ coefficients at residue.
Use EB2, $A_u,D_u\sim1$, and EB3. The resulting residue Gram on
$(N_y,N_m,N_b,N_r,p,q)$ is the direct sum of
$\operatorname{diag}(1,1,1/2,1/2)$ and
\[
 \frac12
 \begin{pmatrix}
 1&-\bar s\bar c\\
 -\bar s\bar c&\bar\nu^2\bar s^2+\bar\ell^2\bar c^2
 \end{pmatrix}.                                                \tag{EB8}
\]
The determinant of this last block is
\[
             \frac14(\bar s^2-\bar\ell)^2.                    \tag{EB9}
\]
Thus EB1 makes the entire bounded finite Gram strictly positive.
The separated bounded-Gram lemma gives tensor SOS, and restoring the
two retained positive rows proves the complete parent. The roots are
allowed to be dependent or even identically zero at an endpoint; the
claim concerns the coefficient Gram, so this does not affect the
argument.

The same homogeneous formulas allow $c=0$ and $s=0$ literally.
They include every Bob projective endpoint, every sign, and every
local Alice axis. Nonlocal Alice charts use the common adjacent
$C,A$ theorem. The real and imaginary state parts are treated
separately. Every normalization is pure source or pure testing.

\paragraph{The exact remaining obstacle.}
The leading scalar has the useful two-square identity
\[
 \tfrac12p^2-sc\,pq
 +\tfrac12(\nu^2s^2+\ell^2c^2)q^2
 =\tfrac12(p-scq)^2+\tfrac12(s^2-\ell)^2q^2.                  \tag{EB10}
\]
It does not complete the matched-residue chart: there the second
pivot collapses, and the actual source corrections must be allocated
through exact contact conductors. Arbitrary small positive $D$ or
$\beta$ openings do not supply a missing fixed pivot. Unit $\zeta$
also requires new source limits because $e$ can have nonzero residue.
These two obligations are retained, not declared empty.



% Source formulas: pe55ec4890781
\subsection{The common finite matched-Bob angular relation}

\label{component:pe55ec4890781}

\label{pe55ec4890781:sec:new-c144-matched-budget}

This exact auxiliary identity supplies the angular helper budget in FM.
Work in the separated tensor ring $K_c\otimes_{\mathbb R}K_t$.
Let $S,\kappa,\nu,\ell$ be positive source quantities, with
$\nu+\ell=1$, and put
\[
 x=\kappa S,\qquad j=\sqrt{\nu\ell},\qquad
 a_B=\sqrt{\ell/\nu}.
 \tag{MB1}
\]
For independent testing variables $u,v,W,Z$ define
\[
\begin{gathered}
 T=\kappa W,\quad U=uZ,\quad P=uW/\sqrt S,\quad Q=\sqrt S Z,\\
 N=T-jU,\quad G=P-jQ,\quad \delta=v-a_B,\\
 B=\delta Q,\quad B_u=\delta U.
\end{gathered}
 \tag{MB2}
\]
The monomial relation $TU=\kappa PQ$ is exact. It gives
\[
 \delta^2N^2+2j\kappa\delta GB+2j^2\kappa B^2-j^2B_u^2
       =(\delta T)^2.
 \tag{MB3}
\]
In particular, if $\nu,\ell$ are positive bounded units,
$\kappa$ is infinitesimal and $\delta$ has zero real residue,
then any fixed real positive diagonal on $B_u$ can be allocationd
at an infinitesimal perturbation of the Gram matrix on $(N,G,B)$.
Indeed, for any fixed $c>0$ the exact SOS remainder is
\[
 \frac{c}{j^2}\delta^2N^2+
       \frac{2c\kappa}{j}\delta GB+2c\kappa B^2-cB_u^2
       =\frac{c}{j^2}(\delta T)^2.
 \tag{MB4}
\]
Every displayed coefficient is a finite sum of bounded separated
factors, with zero real residue except the constant $c$ in the
desired $B_u$ pivot. The only inverses are of positive source units.
The factor $\delta$ is a difference of a pure testing and a pure
source quantity; no inverse or radical of that difference is taken.
The right side has the displayed separated SOS root
$\sqrt c\,\delta T/j$.

Two exact radial conductors and one cross conductor are
\[
 \begin{split}
 (u^2-x)W/\sqrt S&=uG-\sqrt S N,\\
 (u^2-x)Z&=(\kappa\sqrt S/j)G-(u/j)N,\\
 (u^2-x)\delta Z^2&=(\kappa/j)GB-(1/j)NB_u.
 \end{split}
 \tag{MB5}
\]
These identities require no bound on the state coordinates or on
any ratio of source and testing infinitesimals. In the local matched
chart with $u$ infinitesimal their first two right sides have only
infinitesimal coefficients on $(N,G)$.



% Source formulas: p3f02c51074be
\subsection{Unit $\kappa,\nu,\ell$: complementary Bob charts}

\label{component:p3f02c51074be}

\label{p3f02c51074be:sec:new-c144-unit-kappa-all-tests}

\paragraph{Exact scope and conclusion.}
Retain the original physical C14.4 source, its strict source gaps,
and its native source inequalities. Assume
\[
 S=\sigma+L\in\operatorname{Inf}_{>0},\qquad
 \kappa=x/S,\quad\nu=\sigma/S,\quad\ell=L/S
 \quad\hbox{are positive bounded units},\qquad\nu+\ell=1.
 \tag{FU1}
\]
Together with the matched theorem UKR1--UKR21, the argument below
proves that the complete original parent is a finite separated tensor
SOS for every projective test and every state on FU1. It includes
arbitrary openings of $M,\theta,r,1-r$, literal $D=0$ and $\beta=0$,
and independent source and testing infinitesimals. It makes no claim
at the $\nu$- or $\ell$-infinitesimal source endpoints.

All normalizations are pure-source or pure-testing operations.
All real Schur complements below prove strictness of a real residue
matrix; no such inverse is used as a tensor-ring coefficient. The
strict bounded-Gram lemma is the one in the surviving
Lemma~\ref{p25af2c58da7e:lem:gram}.

\paragraph{The full parent and opening-safe source bounds.}
Write $D_u=1+u^2,D_v=1+v^2$, and divide the parent by $4$:
\[
\begin{split}
 F/4={}&m(vW-X+uvY-uZ)^2+CD_v(uW-Y)^2+AD_uD_vY^2\\
 &+(D+\tfrac12)(uW-Y+uvX-vZ)^2
 +(\alpha+\tfrac12)D_u(vY-Z)^2\\
 &+(\beta+\tfrac12)D_v(X+uZ)^2-\tfrac12D_uD_vZ^2.
\end{split}\tag{FU2}
\]
Keep the exact FM2--FM3 source formulas, with
$c_0=SC(1-r)$, $a_0=SAr$, $d_0=D/S$, and $a_1=\alpha/S$.
These four coefficients and $\beta\theta$ are bounded, with no
inverse of the possibly small $M$ or $\theta$ in their normalized
formulas. The $D/S$ numerator has its factor $S$ canceled exactly.

The following letters refer only to real residues:
\[
 A_* =\nu-\kappa h=\nu E,\quad R=1-A_*,\quad
 \nu=A_*+\kappa h,\quad\ell=1-A_*-\kappa h.
 \tag{FU3}
\]
The exact source yields
\[
\begin{gathered}
 E\ge1/18,\qquad h-A_*\ge0,\qquad
 A_*^2-\kappa(h-A_*)\ge0,\qquad R-\kappa\ge0,\\
 c_0=a_0=E/2,\qquad
 d_0=\{A_*^2-\kappa(h-A_*)\}/2,\\
 a_1=\frac{\ell\{\kappa h+(R-\kappa)A_*\}}{2\nu},
 \qquad 2\beta\theta=h-A_*.
\end{gathered}\tag{FU4}
\]
In particular $h\ge\nu/(1+\kappa)$, so $h,E,a_1,c_0,a_0$
have positive real lower bounds. The actual source has $m$ bounded
below by a positive real, while $C,A$ are unbounded. Here $m$
can be unbounded, and so can $\beta$. They cannot both be unbounded:
$m$ unbounded forces $R=\kappa$, whence
$\ell=\kappa\theta$ makes $\theta$ a unit; $\beta$ unbounded
forces $\theta=0$, whence $R-\kappa=\ell$ makes $M$ a unit.
For bounded $m,\beta$ their real values are
\[
 m=\frac{\kappa^2h}{2\nu(R-\kappa)},\qquad
 \beta=\frac{h-A_*}{2(1-h)}.
 \tag{FU5}
\]
These formulas are used only when the indicated real denominators
are positive. Every estimate below for the actual parent uses
bounded normalized source coefficients before taking residues.

\paragraph{A complete finite-Bob frame, including the state normals.}
Assume $u$ is infinitesimal and $v$ is finite. We may cap $m$ and
$\beta$ by finite pure-source values $m_c>0,\beta_c\ge0$ and
retain the complete excesses in their original event squares.
When a weight is bounded, keep it unchanged. When it is unbounded,
the sufficiently large real cap is chosen below. Define
\[
 B_0=X+uZ,\quad N_y=\frac{Y-ruW}{\sqrt{Sr(1-r)}},\qquad
 T=\kappa W,\ U=uZ,\ P=\frac{uW}{\sqrt S},\ Q=\sqrt S Z.
 \tag{FU6}
\]
All six rows are legal separated expressions, and $TU=\kappa PQ$
holds exactly. No restriction is placed on their values or on state
ratios. Set $a=u^2$, $s_y=\sqrt{Sr(1-r)}$, and
\[
\begin{split}
 M_0&=(v/\kappa)(1+ra)T-B_0,\\
 A_P&=D_v\{c_0(1-r)+a_0r(1+a)\}\\
 &\hspace{5mm}+S\{(D+\tfrac12)(1-r)^2
       +(\alpha+\tfrac12)(1+a)r^2v^2\}.
\end{split}\tag{FU7}
\]
The entire capped parent is $F_0+2L_yN_y+H_yN_y^2$, where
\[
\begin{split}
 F_0={}&m_cM_0^2+A_PP^2+(1+a)\{d_0(1+a)v^2+a_1\}Q^2
                   +\tfrac12(1+a)v^2U^2\\
 &-2(D+\tfrac12)(1+a)v^2B_0U
 -(1+a)v\{1+2D(1-r)+2\alpha r\}PQ\\
 &+\{(\beta_c+\tfrac12)D_v+(D+\tfrac12)av^2\}B_0^2
 +2(D+\tfrac12)(1-r)uv\sqrt S\,B_0P,
\end{split}\tag{FU8}
\]
\[
\begin{split}
 L_y={}&m_cuvs_yM_0
 +\sqrt{r(1-r)}D_v\{-c_0+a_0(1+a)\}P\\
 &+S\sqrt{r(1-r)}\{-(D+\tfrac12)(1-r)
          +(\alpha+\tfrac12)(1+a)rv^2\}P\\
 &-(D+\tfrac12)uvs_yB_0
       +\sqrt{r(1-r)}v(1+a)(D-\alpha)Q,\\
 H_y={}&rc_0D_v+(1-r)a_0(1+a)D_v\\
 &+Sr(1-r)\{m_cu^2v^2+D+\tfrac12
                      +(\alpha+\tfrac12)(1+a)v^2\}.
\end{split}\tag{FU9}
\]
These are exact identities in the two unrestricted original state
coordinates $X,Y$; $B_0$ is not set to zero. Their coefficients are
bounded across all source openings. In $L_y$ every coefficient has
zero residue, since $c_0,a_0$ have the same residue, and the other
terms contain $u$ or $S$. The $N_y$ pivot has residue $ED_v/2>0$.
Thus it suffices to make the remaining five-row real residue strict.

\paragraph{Finite nonzero unmatched Bob residue.}
Take $v>0$ with a positive finite residue different from
$b=\sqrt{\ell/\nu}$. The complete real residue of FU8--FU9 is
\[
\begin{split}
 &m_c\{(v/\kappa)T-B_0\}^2+qD_vB_0^2-v^2B_0U
       +\tfrac12v^2U^2-vPQ\\
 &\hspace{8mm}+\tfrac E2D_vP^2+(d_0v^2+a_1)Q^2
             +\tfrac E2D_vN_y^2,\qquad q=\beta_c+\tfrac12.
\end{split}\tag{FU10}
\]
For the moment use the uncapped finite values $m,\beta$ from FU5.
Eliminating $B_0$ in this real matrix, and using $TU=\kappa PQ$
only to name the adjusted cross coefficient, gives
\[
 a_TT^2+a_UU^2+a_PP^2+a_QQ^2-2fPQ,
 \tag{FU11}
\]
where, for $c=v^2$ and $H=m+q(1+c)$,
\[
\begin{gathered}
 a_T=\frac{mq(1+c)c}{\kappa^2H},\quad
 a_U=\frac{c(2H-c)}{4H},\quad a_P=E(1+c)/2,\quad
 a_Q=d_0c+a_1,\\
 f=\frac{v}{2}(1+mc/H).
\end{gathered}\tag{FU12}
\]
The replacement of the real Schur cross term $TU$ by $\kappa PQ$
is implemented in the actual bounded Gram by a real multiple of
the exact zero $TU-\kappa PQ$. The elimination of $B_0$ itself
is only a real strictness test. Every diagonal is positive,
including when $D$ or $\beta$ is zero.
The reduced determinant in a formal positive real variable $y$ is
\[
 \Delta(y,c)=(a_T\kappa^2+a_Py)(a_Uy+a_Q)-f^2y.
 \tag{FU13}
\]
It obeys $4H\Delta/(1+c)=\mathcal P(y,c)$ with
\[
 \mathcal P=p_{22}y^2c^2+p_{21}y^2c+p_{12}yc^2
              +p_{11}yc+p_{10}y+p_{02}c^2+p_{01}c,
 \tag{FU14}
\]
\[
\begin{array}{ll}
 p_{22}=E\beta,&p_{21}=E(m+q),\\
 p_{12}=2m\beta+2qEd_0,&
 p_{11}=-m-q+2E\{(m+q)d_0+qa_1\},\\
 p_{10}=2(m+q)Ea_1,&p_{02}=4mqd_0,\quad p_{01}=4mqa_1.
\end{array}\tag{FU15}
\]
All coefficients except possibly $p_{11}$ are nonnegative, and
$p_{21},p_{10},p_{01}$ are strictly positive. Direct substitution of
FU3--FU5 gives the exact three contact identities
\[
 \mathcal P(\kappa,\ell/\nu)=
 \mathcal P_y(\kappa,\ell/\nu)=
 \mathcal P_c(\kappa,\ell/\nu)=0.
 \tag{FU16}
\]
These finite rational identities can also be checked without
radicals. For positive real $y,c$,
\[
 \frac{\mathcal P(y,c)}{yc}
 =p_{22}yc+p_{21}y+p_{12}c+p_{11}
                     +p_{10}/c+p_{02}c/y+p_{01}/y.
 \tag{FU17}
\]
As a real function of $(\log y,\log c)$, this is strictly convex:
each nonconstant summand has nonnegative Hessian, and
$p_{21}y+p_{10}/c$ alone has positive definite Hessian. FU16 says
its value and gradient vanish at the contact. It is consequently
positive at every other positive point. This real argument is only
a way to determine a finite residue matrix, not a use of logarithms
or mixed inverses in the tensor ring.

Since $c\ne\ell/\nu$, FU13 is positive for all $y>0$.
Minimizing its quotient by $y$ gives
\[
 f<\kappa\sqrt{a_Ta_U}+\sqrt{a_Pa_Q}.
 \tag{FU18}
\]
Indeed the minimum of
$a_Pa_Uy+a_T\kappa^2a_Q/y$ is
$2\kappa\sqrt{a_Ta_Ua_Pa_Q}$, attained at a positive $y$;
equality in FU18 would give a zero of FU13. Choose a real
$\lambda$ with $|\lambda|<\sqrt{a_Ta_U}$ and
$|f-\kappa\lambda|<\sqrt{a_Pa_Q}$. Adding the exact zero
$-2\lambda(TU-\kappa PQ)$ makes both $2\times2$ blocks in
FU11 strictly positive. Restoring the positive real $B_0$ pivot
and the $N_y$ pivot proves strictness of the complete six-row
residue. All actual Gram entries were bounded in FU8--FU9, so
the strict bounded-Gram lemma gives the tensor SOS.

\paragraph{Unbounded $m$ or $\beta$ in the unmatched chart.}
The preceding strictness persists under a sufficiently large
finite cap in the one possible unbounded weight. This assertion
has an opening-safe direct verification. If $m\to\infty$,
divide FU14--FU15 by $m$ before taking the real limit. Then
\[
 (p_{22},p_{21},p_{12},p_{11},p_{10},p_{02},p_{01})/m
 \longrightarrow
 (0,E,2\beta,-1+2Ed_0,2Ea_1,4qd_0,4qa_1).
 \tag{FU19}
\]
At $R=\kappa$, the three contact identities FU16 remain exact
for this limit polynomial. The strict convexity proof FU17 still
applies because its $y$ and $1/c$ coefficients are positive.
The Schur coefficients tend to
\[
 a_T\to q(1+c)c/\kappa^2,\qquad
 a_U\to c/2,\qquad f\to v(1+c)/2,
 \tag{FU20}
\]
with $a_P,a_Q$ unchanged. Therefore FU18 is strict at the limit,
and continuity makes it strict for a sufficiently large finite
real $m_c$.

If $\beta\to\infty$, divide instead by $q$ before taking the
limit at $h=1$. The corresponding coefficient vector is
\[
 (E,E,2m+2Ed_0,-1+2E(d_0+a_1),
                         2Ea_1,4md_0,4ma_1),
 \tag{FU21}
\]
again with the three exact contact identities and the two strictly
positive convexity coefficients. Now
$a_T\to mc/\kappa^2$, $a_U\to c/2$, $f\to v/2$.
A sufficiently large finite $\beta_c$ therefore works.
All the other normalized actual coefficients have their real
residues with bounded unit denominators; no $O(S)$ term was
divided by $M$ or $\theta$. Restore the complete excesses
$(m-m_c)(vW-X+uvY-uZ)^2$ and
$(\beta-\beta_c)D_v(X+uZ)^2$. These caps are used only on the
unmatched chart and need not vanish at the actual contact.

\paragraph{Infinitesimal or zero Bob coordinate.}
If $v$ is infinitesimal, use FU8--FU9 with the six rows
\[
 (B_0,T_v,U_v,P,Q,N_y),\qquad T_v=vT,\quad U_v=vU.
 \tag{FU22}
\]
Cap $m$ at any sufficiently small fixed positive real $m_c\le m$,
and discard no excess. One may set $\beta_c=0$ and retain its
whole event. The exact coefficients remain bounded: the $B_0U$
coefficient is now a multiple of $vB_0U_v$, and $M_0$ becomes
$(1+ra)T_v/\kappa-B_0$. Their real residue is
\[
 m_c(T_v/\kappa-B_0)^2+\tfrac12B_0^2+\tfrac12U_v^2
                      +\tfrac E2(P^2+N_y^2)+a_1Q^2>0.
 \tag{FU23}
\]
This also covers literal $v=0$. The rows need not be independent
coordinates on the state space.

\paragraph{Infinitesimal reciprocal Bob coordinate.}
Put $w=1/v$ and use the homogeneous parent $w^2F$ at $w=0$ as
well as for nonzero infinitesimal $w$. The complete signed core
in the quarter-parent is $K_w/2$, with
\[
\begin{split}
 K_w={}&u^2H^2-2u^2HY-2uHW+2uwHX+u^2w^2W^2\\
 &+2u^2wWX-2uw^2WY+(1+u^2+w^2)(X^2+Y^2)-2uwXY,
 \qquad H=wZ.
\end{split}\tag{FU24}
\]
This is the exact CE6 full core, so all state normals are present.
FU4 gives the real alternative
\[
 d_0+\kappa\beta\theta=(\nu E)^2/2>0.
 \tag{FU25}
\]
Thus either $D/S$ has a positive real lower bound, or $\beta$
has one. Both cases allow independent rates for $u,w,S$.

If $\beta$ is strong, use the four rows $(W,X,Y,U)$, $U=uZ$,
in FU24 by replacing $uH$ with $wU$. Every coefficient is bounded
and the core residue is $X^2+Y^2$. The exact $m$ event is
$(W-wX+uY-wU)^2$, and the $\beta$ event is
$D_w(X+U)^2$. Positive finite caps give the strict residue
\[
 m_cW^2+\tfrac12(X^2+Y^2)+\beta_c(X+U)^2>0.
 \tag{FU26}
\]
All other event terms and both excesses remain whole.

If $D/S$ is strong, choose positive real $c,d,m_c$ with
$c\le c_0,a_0$, $d\le D/S$, and $m_c\le m$. Select exactly
\[
 C_c=c/[S(1-r)],\qquad A_c=c/(Sr),\qquad D_c=dS.
 \tag{FU27}
\]
Retain all excesses, including the whole $A_cu^2D_wY^2$.
Writing $Y=ruW+\sqrt{Sr(1-r)}N_y$, the other two selected
$C,A$ terms have the exact harmonic identity
\[
 C_cD_w(uW-Y)^2+A_cD_wY^2=cD_w(P^2+N_y^2).
 \tag{FU28}
\]
Use six rows $(W,V,X,P,N_y,Q)$ with $V=uwZ$ and the exact
monomial identity $VW=wPQ$. Replacing $uH$ by $V$ in FU24,
the potentially non-small cross term $-2VW$ is exactly $-2wPQ$.
Every other off-diagonal core coefficient is infinitesimal after
$Y=\sqrt S\{rP+\sqrt{r(1-r)}N_y\}$; its residue is $V^2+X^2$.
The selected $m$ and $D$ events are the complete squares
\[
 m_c(W-wX+uY-V)^2,\qquad
 d\{uw\sqrt S W+u\sqrt S X-w\sqrt S Y-Q\}^2.
 \tag{FU29}
\]
The full selected residue is strictly positive:
\[
 m_c(W-V)^2+\tfrac12(V^2+X^2)+c(P^2+N_y^2)+dQ^2>0.
 \tag{FU30}
\]
The strict bounded-Gram lemma applies to the exact complete
forms, and all unselected positive terms are restored. Both
reciprocal proofs include literal $w=0$ and $u=0$.

\paragraph{Exhaustion of tests and exact joins.}
If $u$ is noninfinitesimal, $C,A$ are unbounded by FU4. The
two-adjacent-events theorem CE1--CE10 in
Section~\ref{component:p89fd3edb3edd}
therefore applies to the original full parent and all states.
Its hypotheses require only nonnegative event weights and these
two unbounded coefficients; no small-$\kappa$ source statement
is imported.

For infinitesimal $u$, every positive projective $v$ is either
infinitesimal, a finite unit, or has infinitesimal reciprocal.
FU22--FU30 cover the two end cases. A finite unit has either the
matched residue, covered by UKR1--UKR21, or a different residue,
covered by FU10--FU21. The transformations
\[
 (u,Y,Z)\mapsto(-u,-Y,-Z),\qquad
 (v,X,Z)\mapsto(-v,-X,-Z)
 \tag{FU31}
\]
preserve FU2 exactly and handle both signs. Literal axes and
projective endpoints were included in the displayed raw frames.
Complex states split into real and imaginary parts. No proof in
this exhaustion bounds, normalizes, or relates the testing-state
ratios $W/Z$, $X/Z$, or $Y/Z$.

Consequently every testing chart is covered on exactly FU1,
with every original source predicate retained. The remaining
source endpoints are separate from this conclusion.



% Source formulas: pd26b89cdb3f0
\subsection{Unit $\kappa,\nu,\ell$: the matched Bob residue}

\label{component:pd26b89cdb3f0}

\label{pd26b89cdb3f0:sec:new-c144-unit-kappa-matched}

\paragraph{Scope.}
This theorem retains the entire original physical C14.4 source,
including its strict gaps and native source inequalities. Assume
$S\in\operatorname{Inf}_{>0}$ and $\kappa=x/S,\nu=\sigma/S,
\ell=L/S$ are positive bounded units, with $\nu+\ell=1$.
The testing chart is $u\in\operatorname{Inf}$ and finite $v$
whose absolute value has the same residue as
$b=\sqrt{\ell/\nu}$. All real or complex states are allowed.
The unmatched Bob tests on this source are proved in Section~\ref{component:p3f02c51074be}; the source endpoints are separate cells in Section~\ref{component:pfd0dd7bb10cf}. The present proof is confined to its stated matched-Bob chart and uses no mixed inverse.

The dictionary FM2--FM3 remains an exact source identity on this
larger domain, but its small-$\kappa$ estimates FM4 are replaced
below. Set $j=\sqrt{\nu\ell}$ and retain the rows
\[
 N_x=X+huZ,\quad N_y=(Y-ruW)/\sqrt{Sr(1-r)},\qquad
 T=\kappa W,\ U=uZ,\ P=uW/\sqrt S,\ Q=\sqrt S Z,
\]
\[
 N=T-jU,\quad G=P-jQ,\quad \delta=v-b,\quad
 B=\delta Q,\quad B_u=\delta U,\quad R_Z=(u^2-x)Z.
 \tag{UKR1}
\]
The radial and cross identities FM8 hold unchanged. Their
coefficients on $R_Z$ remain infinitesimal because $S,u$ are
infinitesimal and $\kappa,1/j$ are bounded.

\paragraph{A sharper exact angular budget.}
For any fixed real $0<\rho<1$, the monomial identity $TU=\kappa PQ$
gives
\[
 \frac{\delta^2}{j^2(1-\rho^2)}N^2+
 \frac{2\kappa\delta}{j(1+\rho)}GB+
 \frac{2\kappa}{1+\rho}B^2-B_u^2
   =\frac{(\delta T-\rho jB_u)^2}{j^2(1-\rho^2)}.
 \tag{UKR2}
\]
Its first two coefficients have zero residue. Thus a fixed
$cB_u^2$ pivot costs $q\kappa cB^2$ in residue, where
$q=2/(1+\rho)>1$ can be chosen arbitrarily close to one.
Every inverse in UKR2 is a pure source unit or a positive real.
There is no inverse of $\delta$.

\paragraph{The actual source through both coefficient openings.}
The source has $1/18\le E=1-\kappa h/\nu<1$ and
$\kappa$ is bounded because $1/\zeta$ is bounded. The exact
$\beta\theta$ identity gives a positive real lower bound on $h$.
Consequently $m$ has a positive real lower bound, although it may
be unbounded. Let
\[
 R_s=1+xr,\quad p=\frac{\kappa\theta}{\ell R_s}\in(0,1),\qquad
 \mu_c=m(1-p)=
 \frac{\kappa^2h(1+x)(1+\nu Sr)}{2\nu\ell P_0R_s}.
 \tag{UKR3}
\]
The quantity $\mu_c$ is a positive bounded unit through $M\to0$.
All of $SC(1-r),SAr,D/S,\alpha/S,\beta\theta$ are bounded
without inverses of $r,1-r,\theta,D,\beta$; this follows directly
from FM3 after the explicit factor $S$ in the $D/S$ numerator
is canceled. Their limits are evaluated below.

Only in the real residue calculation use the abbreviations
\[
 A=\nu-\kappa h=\nu E,\quad R=1-A,\quad d=h-A.
 \tag{UKR4}
\]
Here $A,R,h,\nu,\ell,\kappa$ are positive units, and the
original nonnegative weights and positive $M$ imply
\[
 d\ge0,\quad A^2-\kappa d\ge0,\quad R-\kappa\ge0,
 \qquad \nu=A+\kappa h,\quad\ell=1-A-\kappa h.
 \tag{UKR5}
\]
Indeed in residue $2\beta\theta=d$, $2D/S=A^2-\kappa d$,
and $\ell M=R-\kappa$. These are additive residue identities;
no relative error estimate is asserted at a vanishing opening.

\paragraph{Retain unbounded weights in actual contact-zero rows.}
Define
\[
 V=v(W+uY),\quad B_0=X+uZ,\quad N_m=B_0-pV.
 \tag{UKR6}
\]
At the actual source contact $u=\sqrt x$, $v=b$,
$W/Z=j\sqrt{S/\kappa}$, $X=-h\sqrt x Z$, $Y=r\sqrt x W$,
both $N_m$ and $N_x$ vanish. This is an identity at the source
contact, not an equation imposed on independent tests.

If $m$ is bounded, set $m_c=m$. If it is unbounded, choose a
positive real $m_c$ sufficiently large, as justified below.
Similarly keep $\beta_c=\beta$ when $\beta$ is bounded, and
choose a sufficiently large positive real $\beta_c$ when it is
unbounded. An unbounded positive field element exceeds either
fixed real cap. Retain the entire positive excesses
\[
 (m-m_c)N_m^2+(\beta-\beta_c)D_vN_x^2.
 \tag{UKR7}
\]
The remaining parent is exact, with its $m$ part replaced by
\[
 m_0 B_0^2-2m_1B_0V+m_2V^2,\qquad
 m_0=m_c,\quad m_1=\mu_c+m_cp,\quad
 m_2=\mu_c(1+p)+m_cp^2.
 \tag{UKR8}
\]
This follows from the identity
$m(V-B_0)^2-(m-m_c)(B_0-pV)^2$
and $\mu_c=m(1-p)$. Every coefficient of UKR8 is bounded.
The residual $\beta$ terms are exactly
\[
 \beta_cD_vN_x^2+2(\beta\theta)D_vN_xuZ
                         +(\beta\theta)\theta D_vu^2Z^2.
 \tag{UKR9}
\]
All coefficients of UKR9 are bounded. The complete cross terms
are retained, and both original excesses in UKR7 are restored
once at the end.

\paragraph{The complete bounded scalar and normal frame.}
Apply the FM9--FM19 polynomial construction to this capped
parent with the following explicit replacements: in $A_T$ replace
$m/\kappa^2$ by $m_2/\kappa^2$; in $A_U$ replace its
$m\theta^2$ term by $m_0\theta^2$; in $K$ replace its
$m\theta(1+ra)$ term by $m_1\theta(1+ra)$. In the complete
first-normal formulas replace
\[
\begin{array}{c|c}
 \text{entry}&\text{replacement}\\ \hline
 H_x\text{ first term}&m_0\\
 H_y\text{ term }mu^2v^2&m_2u^2v^2\\
 H_{xy}\text{ term }m&m_1\\
 L_x\text{ coefficient }-m/\kappa&-m_1/\kappa\\
 L_U\text{ terms }m\theta-(mj/\kappa)v(1+ra)&
                  m_0\theta-(m_1j/\kappa)v(1+ra)\\
 Y_P\text{ term }mv^2(1+ra)&m_2v^2(1+ra)\\
 Y_Q\text{ term }-m\theta a&-m_1\theta a.
\end{array}\tag{UKR10}
\]
Include $\beta_cD_v$ in $H_x$ and omit the separately retained
$\beta D_vN_x^2$ term of FM15. All other entries are unchanged.
These replacements follow by expanding UKR8--UKR9 in the two
unrestricted state normals. They are not a scalar-only operation.

The graph scalar and both complete first normals still vanish to
the orders required by FM11 and FM19: the subtracted rows in UKR7
vanish at the actual source contact. Their squares therefore
change neither the value nor any first derivative there.
The same finite monic divisions supply the exact bounded FM13
scalar and FM18--FM19 full normal conductors. The $Y$ normal has
zero-residue cross coefficients, since $SC(1-r)$ and $SAr$
both have residue $E/2$ and the remaining $Y$ coefficients
contain $S$ or $u^2$. Its diagonal residue is $E/(2\nu)>0$.

The only nontrivial real residue uses $(N_x,N,G,B,B_u)$.
Writing
\[
 H_* = h-\tfrac12+\tfrac\theta2(\ell+\kappa h),\qquad
 r_* = H_*-j^2m_2/\kappa^2,
 \tag{UKR11}
\]
its exact entries are
\[
\begin{gathered}
 H_c=\begin{pmatrix}
 m_0+(\beta_c+1/2)/\nu&-(m_1/\kappa)b\\
 -(m_1/\kappa)b&(m_2/\kappa^2)b^2
 \end{pmatrix},\qquad A_G=E/(2\nu),\\
 z=\begin{pmatrix}-j(m_1/\kappa+E)\\-r_*/\nu\end{pmatrix},\qquad
 g_*=E\ell-m_1\theta-\tfrac12+2m_2\ell/\kappa
                                        +\kappa r_*/\nu,\\
 H_B=\nu(m_1\theta+\tfrac12).
\end{gathered}\tag{UKR12}
\]
Thus the residue form is
$ (N_x,N)H_c(N_x,N)^t+A_GG^2+2g_*GB+H_BB^2
 +2z^t(N_x,N)^tB_u$ before the UKR2 pivot is added.
In particular the $\kappa r_*/\nu$ contribution to $g_*$ is
retained; it is not infinitesimal in the present source sector.

\paragraph{The exact real contact curvature and all source faces.}
First take an interior real source, $\theta>0$ and $R>\kappa$,
without capping, so
\[
 m=\frac{\kappa^2h}{2\nu(R-\kappa)},\qquad
 \beta=\frac{h-A}{2\theta},\qquad
 m_0=m_1=m_2=m.
 \tag{UKR13}
\]
The two hard pivots and the angular pivot obey the exact identities
\[
\begin{split}
 (N_x,N)H_c(N_x,N)^t
   &=\frac m{\kappa^2}(bN-\kappa N_x)^2
                                  +\frac R{2\nu\theta}N_x^2,\\
 J_*&=H_B-\frac{g_*^2}{A_G}-\kappa z^tH_c^{-1}z\\
 &=\frac{AR+\kappa(h-A)}{2AhR}
                  \{4A^2hR-\kappa(Rh-A)^2\}.
\end{split}\tag{UKR14}
\]
All inverses in UKR14 are used only in a real finite matrix
calculation to prove strictness. They are not inverses in the
tensor ring. The crucial bracket has the positive decomposition
\[
\begin{split}
4A^2hR-\kappa(Rh-A)^2={}&A^3R(4-A)+A^4(R-\kappa)\\
 &+A^2R(h-A)(3+A+2\kappa)\\
 &+R^2(h-A)\{A^2-\kappa(h-A)\}.
\end{split}\tag{UKR15}
\]
Its first term is a positive unit and all remaining terms are
nonnegative by UKR5. This includes the literal $D=0$ and
$\beta=0$ faces and both limiting openings.

For completeness the strictness is uniform toward those openings.
Put $\mu=m/\kappa^2$ and retain $H_*$ from UKR11 with uncapped
coefficients. Then the two entries of $H_c^{-1}z$ are
\[
 z_x=-\frac{2\nu\theta}{R}
                   \left(jE+\frac{\kappa H_*}{\nu b}\right),
 \qquad z_N=\frac\kappa b z_x+\nu-\frac{H_*}{\mu\ell}.
 \tag{UKR16}
\]
They stay bounded. Also
\[
 g_*=-\frac{2A^2+A h\kappa+A\kappa-A-h\kappa}{2\nu}
 \tag{UKR17}
\]
stays bounded, while $A_G$ and $J_*$ have positive real lower
bounds near any fixed allowed residue point. UKR14 gives a
positive real lower bound on $H_c$ because $\mu$ and
$R/(2\nu\theta)$ have positive lower bounds and $b$ is a unit.
Completing the real source quadratic form with the bounded
coefficients UKR16--UKR17 therefore gives a uniform positive
lower bound on
\[
 Q_{\rm phys}=(N_x,N)H_c(N_x,N)^t+A_GG^2+2g_*GB+H_BB^2
                         +2\sqrt\kappa\,z^t(N_x,N)^tB.
 \tag{UKR18}
\]
This is an auxiliary real matrix used to choose a diagonal budget.
It does not impose $B_u=\sqrt\kappa B$, or any relation between
source and testing angles, in the tensor certificate.

\paragraph{Why finite caps preserve a strict real residue.}
When $m,\beta$ are bounded no cap is needed and UKR18 is already
strict. If $m$ is unbounded, $R-\kappa=0$ in residue; since
$\ell$ is a unit and $\ell=R-\kappa h$, $\theta$ is then a unit,
so $\beta$ is bounded. If $\beta$ is unbounded, $\theta=0$ in
residue, and $R-\kappa=\ell>0$, so $m$ is bounded. Thus at
most one unbounded penalty is present.

For the $m$ opening, the real contact-zero linear form in UKR18 is
\[
 n_m=N_x-(pb/\kappa)N-(pj/\sqrt\kappa)B,
 \tag{UKR19}
\]
and $p\to1$. Its $N$ coefficient is a unit. UKR8 shows that
subtracting the whole $mn_m^2$ leaves a bounded coefficient
matrix with a well-defined limit at the opening. That limit is
strict on $\ker n_m$: approach the residue point by real
$\kappa<R$ with $A,h$ fixed, adjust any kernel vector continuously
to the exact kernel using the unit $N$ coefficient, and apply the
uniform lower bound from UKR18. The parameter inequalities UKR5
are preserved when $\kappa$ is decreased. A real symmetric matrix
strict on the kernel of a nonzero linear form becomes strictly
positive definite after adding a sufficiently large fixed real
multiple of the form's square; this follows by a real block Schur
complement. Choose $m_c$ that large in UKR7.

For the $\beta$ opening the same argument uses the exact penalty
$(\beta/\nu)N_x^2$ and its kernel $N_x=0$. Approach $h=1$ from
below with $A,\kappa$ fixed, which also preserves UKR5, and choose
$\beta_c$ sufficiently large. This is a calculation on real
residue matrices, not a mixed-field limiting certificate. The
actual capped coefficients have those residues because UKR3,
UKR8 and UKR9 have bounded denominators through the original
independently varying $S,M,\theta$ openings. In particular no
estimate divides an $O(S)$ error by $M$ or $\theta$.

We have proved that the actual caps can be chosen so that the
capped matrix UKR18 is strictly positive definite. It follows
that $H_c\succ0$, $A_G>0$, and
\[
 H_B-g_*^2/A_G-\kappa L>0,\qquad L=z^tH_c^{-1}z\ge0.
 \tag{UKR20}
\]

\paragraph{Complete separated SOS allocation.}
Choose $\rho<1$ sufficiently close to one and then a positive
real $c$ so that
\[
 L<c<\frac{H_B-g_*^2/A_G}{q\kappa},\qquad q=2/(1+\rho).
 \tag{UKR21}
\]
Such choices exist by UKR20. Use $c$ times UKR2 to introduce
the $cB_u^2$ pivot. The full real residue on
$(N_x,N,G,B,B_u)$ is strictly positive: eliminate $H_c$ to
leave $(c-L)B_u^2$, and eliminate $A_G$ to leave
$(H_B-q\kappa c-g_*^2/A_G)B^2$.
The $N_y$ pivot is $E/(2\nu)>0$ and its cross coefficients
have zero residue. Every other full normal and radial term has
already been allocated by the exact bounded FM constructions
and UKR10. The strict separated bounded-Gram lemma supplies a
finite tensor SOS. Restore the square on the right of UKR2 and
both complete excess terms UKR7. This is the full original
parent, not its restriction to a state graph.

The original sign transformations of FM5 handle negative $u,v$,
and the proof includes $u=0$. Complex states split into real and
imaginary parts. Arbitrary state ratios and independent source
and testing rates are allowed throughout. Unmatched unit-$\kappa$
Bob tests and the $\nu$- or $\ell$-infinitesimal source endpoints
remain outside this theorem.



% Source formulas: p578719565e9d
\subsection{Unit $\kappa,\ell$ and infinitesimal $\nu$}

\label{component:p578719565e9d}

\label{p578719565e9d:sec:new-c144-nu-endpoint}

\paragraph{Scope.}
Retain every original physical C14.4 source predicate, strict gap,
nonnegative event coefficient, and original native source inequalities.
This theorem assumes
\[
 S=\sigma+L\in\operatorname{Inf}_{>0},\quad
 \kappa=x/S\text{ a positive bounded unit},\quad
 \nu=\sigma/S\in\operatorname{Inf}_{>0},\quad\ell=L/S=1-\nu.
 \tag{NU1}
\]
It proves a finite separated tensor SOS for the complete original
parent at every projective test and every real or complex state.
The sector with both $\kappa,\nu$ infinitesimal is not included.
No relative source/testing rate, state-ratio bound, or mixed inverse
is used. The original literal $D=0$ and $\beta=0$ faces are included.

\paragraph{Opening-safe source parameters and the two sphere masses.}
Use the exact FM2--FM3 native dictionary, and write
\[
\begin{gathered}
 q=\nu/\ell,\quad \zeta=\ell/\kappa=L/x,\quad
 e=\kappa h/\nu,\quad E=1-e,\\
 c_0=SC(1-r),\quad a_0=SAr,\quad d_0=D/S,\quad a_1=\alpha/S.
\end{gathered}
 \tag{NU2}
\]
In NU1, $q$ is infinitesimal and $\zeta$ is a positive bounded
unit. The exact source sphere SE3 gives
$e(B/S)=\kappa H_0+\nu\mathcal H b_0^2$.
Since $B/S\sim1+\kappa$, $H_0\sim1$, and
$\mathcal H\sim E^2$, its real residues satisfy
\[
 e=\frac\kappa{1+\kappa},\quad E=\frac1{1+\kappa},\quad
 h\in\operatorname{Inf},\quad\theta\sim1,\quad
 c_0=a_0=a_1=E/2.
 \tag{NU3}
\]
The two equalities for $e,E$ in this display are residue
equalities; the definition $e=\kappa h/\nu$ remains exact.
In particular $e,E,c_0,a_0,a_1$ are positive bounded units.

Define actual nonnegative source coefficients
\[
 d_s=\frac{D}{Sq^2},\qquad n_s=\frac\beta{q^2}.
 \tag{NU4}
\]
These are bounded even when either is zero. More explicitly the
SE3 sphere identities give
\[
 d_s=\frac{\ell^2\mathcal H}{2P_0J}d_0^{\,\mathrm{sph}\,2},
 \qquad
 n_s=\frac{\ell^2\mathcal H}{2\kappa P_0\theta}
                                   b_0^{\,\mathrm{sph}\,2},
 \qquad
 Jd_s+\kappa\theta n_s=\frac{\ell^2\mathcal H}{2P_0}.
 \tag{NU5}
\]
Here $(d_0^{\mathrm{sph}})^2+(b_0^{\mathrm{sph}})^2=1$;
the superscript distinguishes sphere coordinates from $D/S$.
Every denominator in NU5 is a positive unit. At least one of
$d_s,n_s$ is a positive unit. No sphere coordinate is inverted.

The sole possibly unbounded event coefficient needed locally is
$m$. Remove its pole by its actual contact-zero event, not by
dividing an error by $M$. With $R_s=1+xr$ set
\[
 p=\frac{\kappa\theta}{\ell R_s}\in(0,1),\qquad
 \mu_c=m(1-p)=
 \frac{\kappa e(1+x)(1+\nu Sr)}{2\ell P_0R_s}.
 \tag{NU6}
\]
Thus $\mu_c$ is a positive bounded unit and $m\ge\mu_c$.
Choose a fixed positive real $\epsilon<m$. On the reciprocal
Bob chart define $V=W+uY$ and $B_c=w(X+uZ)$. Retain the whole
positive term $(m-\epsilon)(B_c-pV)^2$. The remaining $m$
quadratic is exactly
\[
 m_0B_c^2-2m_1B_cV+m_2V^2,\qquad
 m_0=\epsilon,\quad m_1=\mu_c+\epsilon p,\quad
 m_2=\mu_c(1+p)+\epsilon p^2.
 \tag{NU7}
\]
All three coefficients are bounded through $M\to0$ with arbitrary
rate. Their residues are
\[
 \mu_c=\frac{\kappa^2}{2(1+\kappa)},\quad p=\kappa,
 \qquad m_2=\kappa^2(\epsilon+\tfrac12)>0.
 \tag{NU8}
\]
The source $M>0$ also implies $0<\bar\kappa\le1$.

\paragraph{Local reciprocal chart and exact angular budgets.}
Assume $u,w$ are infinitesimal, including literal zero. Use
the homogeneous reciprocal parent $F_w/4$ from P4, with its
$m$ term replaced by NU7 and all other terms intact. Put
\[
\begin{gathered}
 X_e=X+ewW,\quad N_y=(Y-ruW)/\sqrt{Sr(1-r)},\quad
 a=u^2,\quad c=w^2,\quad \delta=c-q,\\
 P=uW/\sqrt S,\quad Q=\sqrt S\,wZ,\quad
 N=W-\zeta uwZ,\quad G=P-\ell Q,\\
 B=\sqrt S\,\delta Z,\quad B_u=u\delta Z,\quad
 R=w(a-x)Z=\frac{\sqrt S}{\zeta}G-\frac u\zeta N.
\end{gathered}\tag{NU9}
\]
Every row is a finite separated expression. No state coordinate
is divided by another. Since $W(uwZ)=PQ$, two exact identities are
\[
\begin{split}
 \delta^2N^2+2\zeta\delta wGB+2\zeta\ell w^2B^2
                  -\zeta^2(wB_u)^2&=(\delta W)^2,\\
 \delta^2G^2+2\ell\delta wNB_u+2\ell\zeta w^2B_u^2
                  -\ell^2(wB)^2&=(\delta P)^2.
\end{split}\tag{NU10}
\]
Consequently, from either positive sphere pivot in NU5, any fixed
positive diagonal budgets on both $wB$ and $wB_u$ can be obtained
at infinitesimal bounded-Gram cost. To see this without a hidden
division, if $d_s$ is a unit, replace $B$ by
$\sqrt{d_s}B/\sqrt{d_s}$ in the first identity, divide by
$\zeta^2$, and use $wB=w(\sqrt{d_s}B)/\sqrt{d_s}$ for the
other budget. All cost coefficients are infinitesimal. If $n_s$
is a unit, use the second identity divided by $\ell^2$ and
$wB_u=w(\sqrt{n_s}B_u)/\sqrt{n_s}$. The factor
$\zeta w^2$ is infinitesimal because $\zeta$ is a unit in
NU1. This is the precise restriction preventing a silent extension
to infinitesimal $\kappa$.

\paragraph{The entire scalar graph before normal restoration.}
Write $D=Sd_0$, $\alpha=Sa_1$, and define bounded coefficient
polynomials in the independent testing squares $a,c$:
\[
\begin{split}
 A_P={}&(1+c)\{c_0(1-r)+a_0r(1+a)\}\\
 &+S\{(D+\tfrac12)(1-r)^2c
                  +(\alpha+\tfrac12)(1+a)r^2\},\\
 A_N={}&m_2(1+ra)^2+2m_1ec(1+ra)+m_0e^2c^2
                  +(\beta+\tfrac12)e^2c(1+c)\\
 &-(D+\tfrac12)e\{2(1-r)-e\}ac,\\
 K={}&m_1(1+ra)+m_0ec+(D+\tfrac12)(1-r-e)\\
 &+(\alpha+\tfrac12)(1+a)r+(\beta+\tfrac12)e(1+c),\\
 F_Z={}&D+\alpha c+\beta a+(m_0+\alpha+\beta+\tfrac12)ac,\\
 J_1={}&K-\ell A_P,\qquad L_1=J_1-\zeta A_N,\\
 H={}&F_Z-\ell^2ScA_P+ac(\zeta^2A_N-2\zeta J_1).
\end{split}\tag{NU11}
\]
For $X=-ewW$, $Y=ruW$, the exact graph scalar is
\[
 F_0=A_NN^2+A_PG^2-2uwZL_1N+HZ^2.
 \tag{NU12}
\]
This graph is used to compute coefficients; the unrestricted
normals are restored explicitly below. The actual source contact
has $a=x,c=q,W=\zeta uwZ$, so $N=G=X_e=N_y=0$.
The retained row $B_c-pV$ also vanishes there. Direct source
substitution in NU11 yields
\[
 L_1(x,q)=0,\qquad H(x,q)=H_a(x,q)=H_c(x,q)=0,
 \qquad H(a,0)=D+\beta a.
 \tag{NU13}
\]
These are exact rational source identities; no independent
testing angle is identified with a source angle in the certificate.

All divisions in the next display are finite monic polynomial
divisions at the bounded pure-source contacts $x,q$. They define
bounded coefficient polynomials $L_a,L_c,H_{20},H_{11},H_{02}$:
\[
\begin{split}
 L_1(a,c)&=(a-x)L_a(a,c)+(c-q)L_c(c),\\
 H(a,c)&=H(x,c)+(a-x)H_a(x,c)+c(a-x)^2H_{20}(a,c),\\
 H_a(x,c)&=(c-q)\{-\beta/q+cH_{11}(c)\},\\
 H(x,c)&=S(c-q)^2\{d_s+\kappa n_s+cH_{02}(c)\}.
\end{split}\tag{NU14}
\]
Here the factor $c$ in the quadratic $a$ remainder follows from
$H(a,0)=D+\beta a$. The two constant terms follow by putting
$c=0$: $H_a(x,0)=\beta$ and $H(x,0)=D+\beta x$.
In particular the displayed apparent divisions by $q$ are
bounded: $\beta/q=n_sq$. Finally $H(x,c)/S$ is a bounded
coefficient polynomial directly from NU11, because
$x=\kappa S$, $D/S=d_0$, and $\alpha/S=a_1$ are bounded.
Thus dividing it by the monic $(c-q)^2$ introduces no unbounded
coefficient. This verifies all asserted coefficient bounds without
using relative orders of $S,q,u,w$.

The useful mixed identity is
\[
 c(a-x)(c-q)Z^2=\frac1\zeta\{G(wB)-N(wB_u)\}.
 \tag{NU15}
\]
Substituting NU14 into NU12 gives the complete exact scalar
allocation
\[
\begin{split}
 F_0={}&A_NN^2+A_PG^2+d_sB^2+n_sB_u^2
       +H_{02}(wB)^2+H_{20}R^2\\
 &+\frac{H_{11}-n_s}{\zeta}\{G(wB)-N(wB_u)\}
       -2uL_aNR-2L_cN(wB_u).
\end{split}\tag{NU16}
\]
The residues of $A_N,A_P$ are respectively
$\kappa^2(\epsilon+1/2)$ and $E/2$, both positive.
The row $R$ has infinitesimal coefficients on $(N,G)$ by NU9.
Every remaining signed term is a bounded quadratic term on
$N,G,wB,wB_u$. The angular rows are retained as the complete
nonnegative squares $(\sqrt{d_s}B)^2$ and
$(\sqrt{n_s}B_u)^2$; neither is prematurely discarded.

\paragraph{Both complete state normals.}
Put $s_y=\sqrt{Sr(1-r)}$ and $r_y=\sqrt{r(1-r)}$.
Define
\[
\begin{split}
 A_x={}&-m_0ec-m_1(1+ra)+(D+\tfrac12)a(1-r-e)
                          -e(\beta+\tfrac12)(1+c),\\
 T_x={}&\zeta cA_x+\beta-D+c(m_0+\beta+\tfrac12),\\
 Y_P={}&(1+c)\{-c_0+a_0(1+a)\}\\
 &+S\{m_2(1+ra)+m_1ec-(D+\tfrac12)c(1-r-e)
                             +(\alpha+\tfrac12)(1+a)r\},\\
 Y_Q={}&-m_1a+(D+\tfrac12)-(\alpha+\tfrac12)(1+a),\qquad
 Y_C=\ell Y_P+Y_Q.
\end{split}\tag{NU17}
\]
The entire capped quarter-parent, before any allocation, is
\[
 F_0+H_xX_e^2+H_yN_y^2+2H_{xy}X_eN_y+2L_xX_e+2L_yN_y,
 \tag{NU18}
\]
where
\[
\begin{split}
 H_x={}&m_0c+(D+\tfrac12)a+(\beta+\tfrac12)(1+c),\\
 H_y={}&rc_0(1+c)+(1-r)a_0(1+a)(1+c)\\
 &+Sr(1-r)\{m_2a+(D+\tfrac12)c+(\alpha+\tfrac12)(1+a)\},\\
 H_{xy}={}&-uws_y(m_1+D+\tfrac12),\\
 L_x={}&wA_xN+uZT_x,\qquad
 L_y=r_y\{Y_PG+Y_CQ\}.
\end{split}\tag{NU19}
\]
These formulas retain every quadratic state coefficient and both
first normals. Their real hard residues are
$H_x=1/2$, $H_y=E/2$, $H_{xy}=0$.

Exact source substitution gives $T_x(x,q)=Y_C(x,q)=0$.
Moreover $T_x(a,0)=\beta-D$, and its $a$ dependence has a
factor $c$. Thus bounded monic divisions give
\[
\begin{split}
 T_x(a,c)&=c(a-x)X_a+(c-q)\{(D-\beta)/q+cX_c(c)\},\\
 Y_C(a,c)&=(a-x)Y_a(a,c)+(c-q)Y_c(c).
\end{split}\tag{NU20}
\]
No coefficient here needs an inverse of a sphere coordinate.
Using NU4 and NU9, the full normal conductors become
\[
\begin{split}
 L_x={}&wA_xN+uwX_aR
       +d_s\sqrt S\,q u B-n_sqB_u+wX_c(wB_u),\\
 L_y={}&r_y\{Y_PG+\sqrt S\,Y_aR+Y_c(wB)\}.
\end{split}\tag{NU21}
\]
The two potentially weak terms are face-safe:
$d_s\sqrt S\,quB=(\sqrt{d_s}\sqrt S\,qu)(\sqrt{d_s}B)$
and $n_sqB_u=(\sqrt{n_s}q)(\sqrt{n_s}B_u)$.
Their coefficients are infinitesimal, even when either sphere
mass is zero. Also $Y_P$ has zero residue by NU3 and NU17.
All other terms in NU21 have infinitesimal coefficients on
$N,G,R$ or bounded coefficients on the two angular helper rows.

\paragraph{Strict complete Gram and restoration.}
Begin with the six rows
\[
 (N,G,X_e,N_y,\sqrt{d_s}B,\sqrt{n_s}B_u).
 \tag{NU22}
\]
Their real diagonal residue in NU16--NU21 is
\[
 \operatorname{diag}\bigl(
 \kappa^2(\epsilon+\tfrac12),E/2,1/2,E/2,1,1\bigr)>0.
 \tag{NU23}
\]
Every cross coefficient among these rows is infinitesimal after
the displayed $R$ substitution. Literal zero rows are harmless:
the coefficient matrix may be strict even when the chosen rows
are linearly dependent or vanish.

Use NU10 to add sufficiently large fixed real pivots for the
two further rows $wB,wB_u$, at infinitesimal cost to NU22.
The signed scalar terms in NU16 and the remaining normal terms
in NU21 have bounded coefficients on these helper rows.
A sufficiently large fixed pair of diagonals therefore gives
a strictly positive real eight-row Gram, by an ordinary real
Schur complement. All entries are finite sums of bounded separated
factors. The separated bounded-Gram lemma supplies a finite SOS.
Restore the exact positive square remainders in NU10 and the
whole excess $(m-\epsilon)(B_c-pV)^2$ in NU7. This gives the
entire original parent, with both unrestricted state normals.

\paragraph{All projective testing joins on NU1.}
Noninfinitesimal Alice tests use CE1--CE10 because $C,A$ are
unbounded. For infinitesimal $u$, the local proof above covers
every infinitesimal reciprocal $w$, including zero. It remains
to check finite Bob coordinates $v$.

For $v$ infinitesimal use the exact FU22--FU23 frame; its source
inputs here are only $m$ strong, $c_0,a_0$ with common positive
residue $E/2$, and $a_1$ positive, all proved in NU3--NU6.
For a finite positive unit $v$ use the complete FU6--FU12
bounded parent and real Schur calculation, capping the original
$m$ event at a finite real only if $m$ is unbounded. In the
present source residue $\beta=d_0=0$, $a_1=E/2$, and, if
$\kappa<1$,
\[
 m=\frac{\kappa^2}{2(1-\kappa^2)},\qquad E=\frac1{1+\kappa}.
 \tag{NU24}
\]
The determinant polynomial FU14 becomes exactly
\[
 \mathcal P(y,c)=p_{21}c(y-\kappa)^2+p_{10}y>0,
 \qquad p_{21}=E(m+\tfrac12),\quad
 p_{10}=(m+\tfrac12)E^2,\quad y,c>0.
 \tag{NU25}
\]
It gives the strict four-term circuit inequality FU18. If
$\bar\kappa=1$, then $m$ is unbounded. The normalized limit is
\[
 \mathcal P(y,c)/m\longrightarrow
              \tfrac12c(y-1)^2+\tfrac14y>0.
 \tag{NU26}
\]
The Schur coefficients have finite positive limits, so a
sufficiently large finite cap preserves strictness, exactly as
in FU19--FU20. The whole positive original event excess is
restored. These finite-Bob arguments use no lower bound on
$\nu$, and do not import the source scope of the FU theorem.

The exact sign transformations FU31 cover negative angles;
all displayed frames include literal axes and projective
endpoints. Complex states split into real and imaginary parts.
No chart normalizes a state coordinate or bounds a state ratio.
This completes the all-tests theorem on precisely NU1.



% Source formulas: p2edb55cd1f69
\subsection{Infinitesimal $\kappa$ and unit $\nu,\ell$}

\label{component:p2edb55cd1f69}

\label{p2edb55cd1f69:sec:new-c144-finite-matched}

\paragraph{Scope.}
This theorem treats the source sector stated in FM1: $S,\kappa$ are positive infinitesimals and $\nu,\ell$ are positive bounded units. The source endpoints $\nu\in\operatorname{Inf}$ and $\ell\in\operatorname{Inf}$, and the unit-$\kappa$ source, are separate cells in the complete partition of Section~\ref{component:pfd0dd7bb10cf}. All original physical C14.4 source conditions, strict gaps, nonnegative event coefficients, total-scale conditions and native source inequalities are retained.

\paragraph{Theorem.}
Let the original native C14.4 source have
\[
 S=\sigma+L>0,\qquad x=\kappa S,\qquad
 \nu=\sigma/S,\quad\ell=L/S,\quad\nu+\ell=1,
 \qquad S,\kappa\in\operatorname{Inf}_{>0},
 \tag{FM1}
\]
with $\nu,\ell$ positive bounded units. Then the complete original
parent is a finite SOS in $K_c\otimes_{\mathbb R}K_t$ for every
independent projective test and every real or complex state.
The argument below supplies the missing finite matched-Bob chart;
the other tests use the explicitly scoped SE theorem.

\paragraph{Exact opening-safe source dictionary.}
Write $\theta=1-h$, $0<r,h<1$, $e=\kappa h/\nu$, $E=1-e$,
and set
\[
\begin{split}
 J&=1+xh-\nu S(1-r),\quad
 M=1+xr-\kappa\theta/\ell,\quad Q_s=1-\ell Sr,\\
 P_0&=1+S\{\nu\ell+2\nu\kappa h+\kappa\ell+
                    \kappa^2h\theta\}
             +S^2\{\nu r\kappa\ell+\nu r\kappa^2h\}.
\end{split}\tag{FM2}
\]
These are the actual native source quantities, with $P=SP_0$.
The converted P1 coefficients, without applying P1's different
unit-$L$ asymptotic hypotheses, are exactly
\[
\begin{split}
 c_0=SC(1-r)&=\frac{(1+x)(1+\nu Sr)E}{2P_0},\\
 a_0=SAr&=\frac{(1+xh)E+(1-r)(\ell S+xh)}{2P_0},\\
 d_0=D/S&=\frac{(1+x)(1+\nu Sr)-JP_0}{2JSP_0},\\
 a_1=\alpha/S&=
 \frac{(1+\nu Sr)\{\ell(1+xr)e+\ell^2ME\}}{2Q_sP_0},\\
 b_1=\beta\theta&=
 \frac{(1+x)(1+\nu Sr)(\ell+\kappa h)-\theta P_0}{2P_0},\\
 \mu=m/\kappa^2&=
 \frac{(1+x)(1+\nu Sr)h}{2M\ell\nu P_0}.
\end{split}\tag{FM3}
\]
The numerator in the $d_0$ formula has a factor $S$, so its displayed
$S$ denominator cancels in the source polynomial ring. Thus all six
coefficients belong to the bounded source algebra generated by
$S,\kappa,h,r,\nu,\ell$ and inverses of $J,M,Q_s,P_0,\nu,\ell$.
None of these inverses is of an opening $r,1-r,\theta,D$ or $\beta$.
The original source gives
\[
\begin{gathered}
 J,M,Q_s,P_0,E\sim1,\qquad c_0,a_0\sim\tfrac12,\qquad
 d_0\sim\nu^2/2,\quad a_1\sim\ell^2/2,\\
 b_1-(h-\nu)/2\in\operatorname{Inf},\qquad
 \mu-h/(2\nu\ell)\in\operatorname{Inf},\qquad
 \bar h\ge\bar\nu>0.
\end{gathered}\tag{FM4}
\]
The last inequality follows from $b_1=\beta\theta\ge0$.
Hence $\mu$ is a positive bounded unit. These estimates are uniform
through $r,1-r,\theta\to0$; no comparison with a testing quantity
has been used. In this sector $D/S$ is a positive unit, so $D=0$
has empty intersection; $\beta=0$ is allowed.

\paragraph{Exact full parent and unrestricted normals.}
On the finite Bob chart write $D_u=1+u^2$, $D_v=1+v^2$.
The complete quarter-parent is
\[
\begin{split}
 F/4={}&m(vW-X+uvY-uZ)^2+C D_v(uW-Y)^2+A D_uD_vY^2\\
 &+(D+\tfrac12)(uW-Y+uvX-vZ)^2
 +(\alpha+\tfrac12)D_u(vY-Z)^2\\
 &+(\beta+\tfrac12)D_v(X+uZ)^2-\tfrac12D_uD_vZ^2.
\end{split}\tag{FM5}
\]
No state restriction is made. Put
\[
 N_x=X+huZ,\qquad N_y=\frac{Y-ruW}{\sqrt{Sr(1-r)}}.
 \tag{FM6}
\]
The source square root is nonzero even when its size is arbitrarily
small. It is used only to define a permitted tensor root; every
coefficient below is bounded without its inverse.

After the unmatched-residue and adjacent-$C,A$ exits it is
enough to take $u$ infinitesimal and $v$ a bounded positive unit with
the same residue as $b=\sqrt{\ell/\nu}$. Set
\[
\begin{gathered}
 a=u^2,\quad j=\sqrt{\nu\ell},\quad T=\kappa W,\quad U=uZ,\\
 P=uW/\sqrt S,\quad Q=\sqrt S Z,\quad N=T-jU,\quad G=P-jQ,\\
 \delta=v-b,\quad B=\delta Q,\quad B_u=\delta U.
\end{gathered}
 \tag{FM7}
\]
Every state coordinate remains arbitrary, including zero and
unbounded coordinates. The rows satisfy the exact MB identities
\[
\begin{gathered}
 R_Z=(a-x)Z=(\kappa\sqrt S/j)G-(u/j)N,\\
 (a-x)\delta Z^2=(\kappa/j)GB-(1/j)NB_u,\\
 \delta^2N^2+2j\kappa\delta GB+2j^2\kappa B^2-j^2B_u^2
                  =(\delta T)^2.
\end{gathered}\tag{FM8}
\]
These polynomial identities hold in the separated tensor ring.

\paragraph{The full graph scalar.}
Let $F_0$ denote FM5 with $N_x=N_y=0$; it is used only as the
constant coefficient of the full normal expansion. Define bounded
source/testing polynomials
\[
\begin{split}
 A_T={}&\mu v^2(1+ra)^2,\\
 A_P={}&(1+v^2)\{c_0(1-r)+a_0r(1+a)\}\\
       &+S\{(D+\tfrac12)(1-r)^2+
                   (\alpha+\tfrac12)(1+a)r^2v^2\},\\
 A_Q={}&d_0v^2+a_1,\\
 A_U={}&m\theta^2+\alpha+b_1\theta+\theta^2/2\\
       &+v^2\{2Dh+h-\tfrac12+b_1\theta+\theta^2/2\}
              +(D+\tfrac12)h^2v^2a,\\
 K={}&m\theta(1+ra)+(D+\tfrac12)(1-r)(1+ha)
                         +(\alpha+\tfrac12)(1+a)r,\\
 H={}&vK-\kappa jA_T,\quad R_U=A_U-j^2A_T,\\
 g={}&jA_P-H,\quad R_Q=A_Q+j^2A_P-2jH,\\
 J_Z={}&SR_Q+aR_U.
\end{split}\tag{FM9}
\]
Direct expansion of every term of FM5 gives
\[
 F_0=A_TN^2+A_PG^2+2gGQ+J_ZZ^2.
 \tag{FM10}
\]
For example the only relation used to obtain FM10 from the
five monomials $T^2,U^2,P^2,Q^2,PQ$ is $TU=\kappa PQ$.

The actual source stationarity identities, also verified directly
from FM2--FM3, are
\[
 g(x,b)=0,\qquad
 J_Z(x,b)=\partial_aJ_Z(x,b)=\partial_vJ_Z(x,b)=0.
 \tag{FM11}
\]
These are coefficient-field identities at the actual source
contact; they impose no equation on independent tests.

\paragraph{Finite polynomial division.}
All following quotients mean monic polynomial division in formal
$a$ or $v$, performed before evaluation in the tensor ring. No
mixed inverse is introduced. There are bounded coefficient
polynomials $g_a,g_v,H_{20},H_{11},H_{02}$ defined exactly by
\[
\begin{split}
 g(a,v)&=(a-x)g_a(a,v)+(v-b)g_v(v),\\
 J_Z(a,v)&=(a-x)^2H_{20}(a,v)+(a-x)(v-b)H_{11}(v)
                                      +S(v-b)^2H_{02}(v).
\end{split}\tag{FM12}
\]
For clarity, first divide $g(a,v)-g(x,v)$ by $a-x$ and then
$g(x,v)$ by $v-b$. For the second line, first subtract the
constant and linear Taylor terms in $a-x$ and divide the remainder
by $(a-x)^2$. Divide $\partial_aJ_Z(x,v)$ by $v-b$ for $H_{11}$.
Finally divide the bounded polynomial
$R_Q(x,v)+\kappa R_U(x,v)$ by $(v-b)^2$ for $H_{02}$.
FM11 makes every division exact. In particular $H_{02}$ is
bounded without dividing a bounded error by $S$.

Substituting FM8 into FM10 now yields the exact scalar identity
\[
\begin{split}
 F_0={}&A_TN^2+A_PG^2+2g_vGB+H_{02}B^2
          +H_{11}\{(\kappa/j)GB-(1/j)NB_u\}\\
       &+H_{20}R_Z^2+2\sqrt S\,g_aGR_Z.
\end{split}\tag{FM13}
\]
Thus every coefficient is bounded and all remaining radial
contributions are infinitesimal matrix entries on $(N,G)$.

Here and below a displayed residue formula uses the real residues
of its bounded source parameters. FM3 gives explicitly
\[
 \bar A_T=\frac h{2\nu^2},\quad
 \bar A_P=\frac1{2\nu},\quad
 \bar g_v=\frac{\ell-\nu}{2},\quad
 \bar H_{02}=\frac\nu2,\quad \bar H_{11}=-\theta j.
 \tag{FM14}
\]
For example at $S=\kappa=0$ and $a=0$ the three relevant
polynomials are
$g=j(1+v^2)/2-v/2$,
$R_Q=\nu(v-b)^2/2$, and
$R_U=\theta(\ell-\nu v^2)/2$.
These establish FM14 by the finite divisions just specified.

\paragraph{Both complete first-normal terms.}
Write the full expansion of FM5 as
\[
 F/4=F_0+H_xN_x^2+H_yN_y^2+2H_{xy}N_xN_y
                  +2N_xL_x+2N_yL_y+\beta D_vN_x^2.
 \tag{FM15}
\]
The exact bounded diagonal and cross coefficients are
\[
\begin{split}
 H_x={}&m+(D+\tfrac12)u^2v^2+D_v/2,\\
 H_y={}&r c_0D_v+(1-r)a_0D_uD_v\\
       &+Sr(1-r)\{mu^2v^2+D+\tfrac12
                            +(\alpha+\tfrac12)D_uv^2\},\\
 H_{xy}={}&-uv\sqrt{Sr(1-r)}(m+D+\tfrac12).
\end{split}\tag{FM16}
\]
In particular $\bar H_x=\bar H_y=1/(2\nu)$ and
$\bar H_{xy}=0$. The last term of FM15 is retained whole;
it allows $\beta$ to be arbitrarily large or zero.

Define
\[
\begin{split}
 L_U={}&m\theta-(mj/\kappa)v(1+ra)
 -(D+\tfrac12)v^2(1+ha)+(b_1+\theta/2)(1+v^2)\\
       &+(D+\tfrac12)(1-r)vjS,\\
 Y_P={}&(1+v^2)\{-c_0+a_0(1+a)\}
       +S\{mv^2(1+ra)-(D+\tfrac12)(1-r)
                        +(\alpha+\tfrac12)(1+a)rv^2\},\\
 Y_Q={}&-m\theta a+(D+\tfrac12)(1+ha)
                                  -(\alpha+\tfrac12)(1+a),\qquad
 Y_C=jY_P+vY_Q.
\end{split}\tag{FM17}
\]
The expression $mj/\kappa=\kappa\mu j$ is bounded and
infinitesimal. The complete original normal terms are exactly
\[
\begin{split}
 L_x={}&-\kappa\mu v(1+ra)N
       +(D+\tfrac12)(1-r)vu\sqrt S\,G+L_UU,\\
 L_y={}&\sqrt{r(1-r)}\{Y_PG+Y_CQ\}.
\end{split}\tag{FM18}
\]
The actual source gives $L_U(x,b)=Y_C(x,b)=0$. Perform the same
finite monic divisions to write
\[
 L_U=(a-x)l_a+(v-b)l_v,\qquad
 Y_C=(a-x)y_a+(v-b)y_v.
 \tag{FM19}
\]
Then $L_UU=u l_aR_Z+l_vB_u$ and
$Y_CQ=\sqrt S\,y_aR_Z+y_vB$. All their coefficients are
bounded. Moreover $\bar l_v=-j$, while $Y_P$ and $y_v$ have
zero residue. This follows directly from the source-limit
identities
$L_U=(\ell-\nu v^2)/2$ at $a=0$, and
$Y_P=Y_Q=0$ at $a=0$.
Consequently $L_x+jB_u$ is an infinitesimal bounded combination
of $(N,G,B_u)$ and $L_y$ is an infinitesimal bounded combination
of $(N,G,B)$. This accounts for both unrestricted first-normal
terms, including all $\beta\theta$ terms and the independent
$r,1-r,\theta$ openings.

\paragraph{A strict residue for the full six-row Gram.}
Use FM8 with a fixed real $c=4$ to introduce a $cB_u^2$ pivot
at infinitesimal Gram cost. More explicitly subtract
\[
 (c/j^2)\delta^2N^2+(2c\kappa/j)\delta GB+2c\kappa B^2
 \tag{FM20}
\]
from the bounded part of FM15, add $cB_u^2$ to that part, and
retain the exact positive square $(\sqrt c\,\delta T/j)^2$.
Their sum is unchanged by FM8. The residue on
$(N_x,N_y,N,G,B,B_u)$ is therefore
\[
 \frac{N_x^2+N_y^2}{2\nu}+\frac h{2\nu^2}N^2
 +\frac{G^2}{2\nu}+(\ell-\nu)GB+\frac\nu2B^2
       +\theta NB_u-2jN_xB_u+4B_u^2.
 \tag{FM21}
\]
The $(G,B)$ block has determinant $\nu\ell>0$ and positive
diagonal. The remaining block has positive diagonal pivots
$1/(2\nu),1/(2\nu),h/(2\nu^2)$, and its final scalar Schur
pivot is
\[
 4-2\nu j^2-\frac{\theta^2\nu^2}{2h}>0.
 \tag{FM22}
\]
Indeed $j^2\le1/4$, $0<\nu\le1$, $h\ge\nu$ and
$0\le\theta\le1$ in residue make the two subtracted terms
at most $1/2$ each. All entries before taking residues are finite
sums of bounded separated factors. The strict
bounded-Gram lemma therefore supplies a finite tensor SOS.
Restoring the single square from FM20 and the whole positive
$\beta D_vN_x^2$ row gives the entire FM5. No mixed matrix inverse
or mixed square root has been used.

\paragraph{Testing, sign, and boundary joins.}
The argument includes $u=0$ and every real state without a state
normalization or a sign condition. Negative $u$ is changed to
$-u$ together with $(Y,Z)\mapsto(-Y,-Z)$; negative $v$ is
changed to $-v$ together with $(X,Z)\mapsto(-X,-Z)$.
FM5 is exactly invariant under these transformations. On FM1,
$e=\kappa h/\nu$ is infinitesimal and $C,A$ are unbounded.
The SE10--SE12 theorem covers every unmatched
homogeneous Bob residue, including Bob zero and infinity,
and invokes the unchanged adjacent-$C,A$ theorem for
noninfinitesimal Alice tests. The only remaining Bob residue
has finite nonzero $v$ and $\bar v^2=\bar\ell/\bar\nu$,
which is the matched chart just proved. Thus the union covers
all projective tests on FM1. Real and imaginary state parts
are treated separately for complex states. All divisions and
normalizations used in these joins are pure source or pure
testing operations; all original source predicates remain.

\paragraph{Dependencies and limits.}
The analytical inputs are the native source formulas
P1 converted as FM2--FM3, the full original parent, the strict
separated bounded-Gram lemma in Lemma~\ref{p25af2c58da7e:lem:gram},
and Section~\ref{component:p8c30238a5f7d}
with its preserved EB and adjacent-$C,A$ dependencies.
The exact budget is independently printed in
Section~\ref{component:pe55ec4890781}. The source endpoint
and unit-$\kappa$ sectors require separate complete arguments;
this theorem makes no claim about them.



% Source formulas: pb96c9bfcc12d
\subsection{Infinitesimal $\kappa,\ell$ and unit $\nu$}

\label{component:pb96c9bfcc12d}

\label{pb96c9bfcc12d:sec:c144-ell-endpoint-rederived}

\paragraph{Theorem and exact scope.}
Retain the original native C14.4 physical source, all strict gaps,
nonnegative event weights, total-scale conditions and native source inequalities.  Let
\[
 S=\sigma+L,\quad x=\kappa S,\quad
 \nu=\sigma/S,\quad\ell=L/S,\quad\nu+\ell=1,
 \quad\zeta=L/x=\ell/\kappa.
\]
Suppose $S,\ell>0$ are infinitesimal and $\zeta$ is
noninfinitesimal.  Then the complete original parent is a finite
sum of squares in the separated ring
$K_c\otimes_{\mathbb R}K_t$ for every independently chosen
projective test and every real or complex state.  The proof includes
$\beta=0$ and unrestricted openings of $r,1-r,\theta$.
The intersection with $D=0$ is empty by the source estimate below.
No earlier missing R21 endpoint proof is used.

\paragraph{Source estimates at this endpoint.}
Use the exact actual-source formulas FM2--FM3 and SE1--SE3, with
$0<r,h<1$, $\theta=1-h$, $e=\kappa h/\nu$, $E=1-e\ge1/18$,
$U_s=(1+x)(1+\nu Sr)$, and
\[
\begin{gathered}
 J=1+xh-\nu S(1-r),\quad
 M=1+xr-\kappa\theta/\ell>0,\quad Q_s=1-\ell Sr,\\
 a_0=SAr,\qquad c_0=SC(1-r),\qquad P=SP_0.
\end{gathered}
\]
Put $g=\kappa/\ell=1/\zeta$, which is bounded.
Then $\nu\sim1$ and $\kappa\in\operatorname{Inf}$.
Since
\[
 0\le\beta\theta
  =\frac{U_s(\ell+\kappa h)-\theta P_0}{2P_0},
\]
we have $\theta=O(\ell+\kappa)=O(\ell)$ and
$\beta\theta=O(\ell)$, hence $h\sim1$,
$\beta\theta^2=O(\ell^2)$ and $M\sim1$.
The exact source formulas give
\[
\begin{gathered}
 J,Q_s,P_0,E\sim1,\quad c_0,a_0\sim\tfrac12,
 \quad D/S\sim\tfrac12,\quad m\in\operatorname{Inf}_{>0},\\
 \frac{qm}{\kappa^2}
   =\frac{U_s h}{2M\nu^2P_0}\sim\tfrac12,
 \qquad q=\ell/\nu,\\
 \frac{\alpha}{S\ell^2}
   =\frac{(1+\nu Sr)\{(1+xr)e/\ell+ME\}}{2Q_sP_0},
 \qquad
 \frac{\alpha}{S\ell^2}-\frac{1+g}{2}\in\operatorname{Inf}.
\end{gathered}                                                   \tag{EL0a}
\]
The last difference is uniform for bounded $g$, because
$e/\ell=gh/\nu$.  All statements are inside the source field;
none compares a testing rate to a source rate.  In particular
$D/S$ is a positive unit, while $\beta$ may be zero, bounded,
or unbounded.

\paragraph{The full parent and local testing chart.}
Noninfinitesimal Alice tests and unmatched Bob residues will be joined
at the end.  The remaining chart has $u,v\ge0$ infinitesimal,
including literal zero axes.  Write $a=u^2$, $c=v^2$,
$b=\sqrt q$, $j=\nu b$, and $d=D+1/2$.
The complete original quarter-parent is
\[
\begin{aligned}
 F/4={}&m(vW-X+uvY-uZ)^2+C(1+c)(uW-Y)^2\\
 &+A(1+a)(1+c)Y^2+d(uW-Y+uvX-vZ)^2\\
 &+(\alpha+\tfrac12)(1+a)(vY-Z)^2
 +(\beta+\tfrac12)(1+c)(X+uZ)^2\\
 &-\tfrac12(1+a)(1+c)Z^2.
\end{aligned}                                                     \tag{EL0b}
\]
All state coordinates are arbitrary elements of the testing field.
No state coordinate is divided by or bounded below.
Put $R_s=1+xr$, $p=\kappa\theta/(\ell R_s)$, $\mu=m(1-p)$, and
\[
 V=v(W+uY),\qquad N=X+uZ-pV,\qquad
 Y=ruW+\sqrt{Sr(1-r)}N_y.                         \tag{EL1}
\]
These are pure-source rescalings of polynomial state normals.
The original $m$ event is $(1-p)V-N$, the beta event is $N+pV$,
and the core event is
$uW-Y+puvV-(1+a)vZ+uvN$; thus no term has been dropped.
On $N=N_y=0$ define
\[
\begin{aligned}
 L_0&=\mu(1-p)+(\beta+\tfrac12)p^2,
 \qquad L_1=(\beta+\tfrac12)p^2,\\
 B_c&=1-r+pc(1+ra),\\
 \lambda(a,c)&=\frac q{\kappa^2}(1+ra)^2(L_0+L_1c),\\
 A_P&=(1+c)[c_0(1-r)+a_0r(1+a)]
       +S[dB_c^2+(\alpha+\tfrac12)(1+a)r^2c],\\
 K&=(1+a)[dB_c+(\alpha+\tfrac12)r],\\
 F_Z&=(1+a)[\alpha+Dc+dac].
\end{aligned}                                                    \tag{EL2}
\]
Here $c_0=SC(1-r)$ and $a_0=SAr$.  The complete graph scalar is
\[
 F_0=c(1+ra)^2(L_0+L_1c)W^2+A_PP^2-2vKPQ+F_ZZ^2,
 \quad P=uW/\sqrt S,\quad Q=\sqrt S Z.             \tag{EL3}
\]
Let $T=\kappa W,U=uZ$ and use the exact four-term circuit
\[
 \Psi=(\nu/\ell)cT^2+\nu^2cU^2+\kappa\nu\ell Q^2
       +\kappa P^2-4\kappa\nu vPQ.                \tag{EL4}
\]
Its four positive monomial square roots are
$A_1=(\kappa/b)vW$, $A_2=\nu vuZ$,
$A_3=\nu b\sqrt xZ$, and $A_4=\sqrt\kappa P$.
Their product is a square.  The usual two paired AM--GM identities
make $\Psi$ a separated SOS, including $WZ<0$ by sign excess.
The multiplier $\lambda$ itself is a sum of separated monomial
squares with positive pure-source coefficients; no mixed radical
of $\lambda$ is taken.

Set $\Gamma=P-\nu vQ$, $\Lambda=vP-\ell Q$, and
\[
 \widehat A=A_P-\kappa\lambda,\quad
 \widehat K=K-2\kappa\nu\lambda,\quad
 \widehat F=F_Z-\lambda(\nu^2ac+\kappa\nu\ell S).
\]
Then the following equality is purely polynomial and exact:
\[
 F_0-\lambda\Psi
   =\widehat A\Gamma^2-2vJ\Gamma Q+HZ^2,
 \quad J=\widehat K-\nu\widehat A,
 \quad H=\widehat F+\nu^2Sc\widehat A-2\nu Sc\widehat K.
                                                               \tag{EL5}
\]
The exact coefficient bounds are
$\lambda_0=qL_0/\kappa^2\sim1/2$ and
$\lambda_1=qL_1/\kappa^2=O(\ell)$.
Indeed the exact opening-safe formulas are
\[
 \lambda_1=\frac{\theta(2\beta\theta+\theta)}
                    {2\ell\nu R_s^2},\qquad
 \lambda_0=\frac{U_shM}{2\nu^2P_0R_s^2}+\lambda_1.
\]
Here $\theta/\ell$ and $(\beta\theta)/\ell$ are bounded, so
the formula for $\lambda_1/\ell$ has only bounded source factors.
No comparison between $S$ and $\ell$ is made.  The mixed term
arising from $-\nu^2\lambda_1c^2a$ is allocated in EL18 below
using an actual row of the same circuit.

\subsubsection*{Exact helper rows and the complete two-normal conductor}
The following exact budgets and complete normal conductors use a
positive real fraction of $\Psi$ and strict pivots on
$(\Gamma,\Lambda)$.  The scalar calculation below supplies precisely
these pivots; it does not use the state normals to obtain them.
Put $\delta=v-b$, $G_0=P-jQ$, $B_o=\delta Q$, and
\[
\begin{gathered}
 E_{12}=A_1-A_2=(v/b)(T-jU),\quad V_h=v\delta U,\\
 R_v=v(a-x)Z,\quad B=\sqrt S(c-q)Z,\quad
 R_0=\ell\sqrt S(a-x)Z.
\end{gathered}
                                                               \tag{EL6}
\]
The exact angular relations are
\[
 (v+b)G_0=\Lambda+b\Gamma,\qquad
 \nu(v+b)B_o=\Lambda-v\Gamma=\nu B.             \tag{EL7}
\]
In particular $B$ has a bounded linear expression in the two strict
angular rows.  Multiply MB3 by $v^2/j^2$ and use
$\eta=\kappa(v+b)/j\in\mathrm{Inf}$.  The exact SOS inequality is
\[
 V_h^2\preceq_{\rm SOS}
  \frac{\delta^2}{\nu^2}E_{12}^2
   +\eta(\Lambda+b\Gamma)^2
   +\frac{\eta+2\kappa}{\nu^2}(\Lambda-v\Gamma)^2.
                                                               \tag{EL8}
\]
For a direct verification, set $R_1=vG_0,R_2=vB_o$.
The mixed-term replacement uses the legal SOS
\[
 \eta(R_1^2+R_2^2)-2\kappa\delta R_1R_2/j
 =\frac{\kappa v}{j}(R_1-R_2)^2
  +\frac{\kappa b}{j}(R_1+R_2)^2.
\]
Replacing $v^2$ by $(v+b)^2$ leaves
$(2vb+b^2)G_0^2$ and $(2vb+b^2)B_o^2$, with positive
separated multipliers.  The remaining MB3 square is
$(v\delta T/j)^2$.  Every cost in EL8 is infinitesimal, since
$\kappa/j=(\kappa/\ell)b$ is infinitesimal.

Two further exact relations are
\[
 R_v=\frac{\kappa v\sqrt S}{\nu b}\Gamma
       -\frac u\nu E_{12}+\frac{xv(v-b)}bZ,          \tag{EL9}
\]
\[
 \frac{x^2}{qS}B^2-
       \left(\frac{xv(v-b)}bZ\right)^2
   =\frac{x^2}{q}(2vb+b^2)(v-b)^2Z^2.              \tag{EL10}
\]
Thus $R_v$ is free at infinitesimal Gram cost because
$x^2/(qS)=\kappa S(\kappa/q)$ is infinitesimal.  Finally
\[
 2\nu^2ScR_v^2+2\nu^2(a-x)^2B^2-R_0^2
       =[\nu\sqrt S(2c-q)(a-x)Z]^2               \tag{EL11}
\]
supplies $R_0$ with infinitesimal costs.  None of EL6--11 bounds or
divides a state coordinate or a mixed source/testing sum.

For the exact complete $N$ half derivative set
\[
\begin{aligned}
 B_0(c)&=[-\mu+(\beta+\tfrac12)p(1+c)]/\kappa,\\
 F_a(a,c)&=r\kappa B_0(c)+d(1-r)+dpc(1+ra),\\
 D_t(a,c)&=d(1+a)-\nu S F_a(a,c).
\end{aligned}                                                     \tag{EL12}
\]
All their coefficients are bounded, including at $\beta=0$ or
unbounded $\beta$: $\mu/\kappa$ and
$(\beta+1/2)p/\kappa=(\beta\theta+\theta/2)/(\ell R_s)$
are bounded.  Write the monic polynomial differences
$D_t-D_t(x,q)=(a-x)D_a+(c-q)D_c$ and
$B_0(c)-B_0(q)=(c-q)B_{01}$.  Source contact gives
$D_t(x,q)=\nu B_0(q)$.  The entire half derivative is
\[
\begin{aligned}
 L_N={}&bB_0(c)E_{12}+uv\sqrt S F_a\Gamma-uvD_aR_v\\
 &+[-D_t(x,q)+\nu bB_{01}(v+b)-v(v+b)D_c]V_h.
\end{aligned}                                                     \tag{EL13}
\]
Hence EL8--10 allocate every $N$ cross term without a state bound.
The normal $\widehat N=\sqrt{1+\beta}N$ has diagonal at least
$(\beta+1/2)/(1+\beta)\ge1/2$; dividing the cross coefficients
by $\sqrt{1+\beta}$ only improves the displayed bounds.

The complete other half derivative is
\[
 L_y=\sqrt{r(1-r)}[Y_P\Gamma+vY_KQ],             \tag{EL14}
\]
where
\[
\begin{aligned}
 Y_P={}&(1+c)[-c_0+a_0(1+a)]\\
 &+S\{c(1+ra)(L_0+L_1c)+d(-1+pac)B_c
                   +(\alpha+\tfrac12)(1+a)rc\},\\
 Y_K={}&\nu Y_P+(1+a)(D-\alpha-dpac).
\end{aligned}
\]
The polynomial coefficients are bounded, $Y_P$ is infinitesimal,
and actual source contact gives $Y_K(x,q)=0$.  Monic division gives
$Y_K=(a-x)Y_a+(c-q)Y_c$ with bounded coefficients, so
\[
 L_y=\sqrt{r(1-r)}[Y_P\Gamma+\sqrt S Y_aR_v+vY_cB].
                                                               \tag{EL15}
\]
The exact quadratic normal coefficients are
\[
\begin{aligned}
 H_N={}&m+dac+(\beta+\tfrac12)(1+c),\\
 H_y={}&rc_0(1+c)+(1-r)a_0(1+a)(1+c)\\
 &+Sr(1-r)[d(-1+pac)^2+(\alpha+\tfrac12)(1+a)c
                         +ac(L_0+L_1c)],\\
 H_{Ny}={}&\sqrt{Sr(1-r)}uv[\kappa B_0(c)-d+dpac].
\end{aligned}                                                     \tag{EL16}
\]
After the source-only normalization of $N$, the first two diagonal
residues are respectively in $[1/2,1]$ and $1/2$, and the cross
coefficient is infinitesimal.  Equations EL13--16 include all terms
of the full original parent; they require no $W/Z$ or testing/source
rate restriction.  The scalar allocation and exact strict-Gram assembly are given next.

\subsubsection*{Complete scalar allocation, including the mixed radial term}
Both $\widehat A$ and the coefficients of $J,H$ in EL5 are
bounded polynomials in $a,c$: their source coefficients are built
from bounded $p,d,\alpha,c_0,a_0,\lambda_0,\lambda_1,\kappa,S,\nu$.
In fact $p=O(\kappa)$,
$\mu/\kappa=O(\kappa/\ell)$, and
$(\beta+1/2)p^2=O(\kappa^2)$.
Consequently $\widehat A\sim1/2$ on the local chart.
Actual source contact, checked in the original parent, gives
\[
 J(x,q)=0,\qquad H(x,q)=H_a(x,q)=H_c(x,q)=0.        \tag{EL17}
\]
These are exact identities, not residue equalities.  They also follow
by evaluating the actual contact state in EL1: all four circuit
monomials agree, $\Gamma=0$, and both graph normals vanish.
The vanishing first derivatives of the original contact parent
then give the three scalar jets.  

The additional actual circuit row
\[
 E_{23}=A_2-A_3=\nu(uv-\sqrt{xq})Z
\]
resolves the otherwise unbounded mixed coefficient.  Direct expansion
using $x=\kappa S$ gives the exact identity
\[
 c(a-x)(c-q)Z^2
 =\frac{c-q}{\nu^2}E_{23}^2
  +\frac{2\sqrt{\kappa q}}{\nu}E_{23}B-\kappa B^2.
                                                               \tag{EL18}
\]
Every coefficient on the right is infinitesimal.  To see the equality
without division by a testing coordinate, use
\[
 (ac-xq)Z^2=E_{23}^2/\nu^2+2\sqrt{xq}\,ZE_{23}/\nu,
 \qquad c(a-x)=ac-xq-x(c-q).
\]
All square roots in this calculation are pure-source roots.

For completeness, the single circuit supplies positive fixed pivots
on both $E_{12}$ and $E_{23}$.  If $WZ\ge0$, use the two AM--GM
pairings having respectively these differences and average them;
then $\Psi-\tfrac14(E_{12}^2+E_{23}^2)$ is separated SOS.
If $W,Z$ are both negative, first change the common state sign;
the parent and all squared differences are unchanged.
If $WZ<0$, the negative product term of $\Psi$ is nonnegative,
and
\[
 \sum_{i=1}^4 A_i^2-\tfrac14[(A_1-A_2)^2+(A_2-A_3)^2]
\]
is a real positive-definite quadratic form in the $A_i$.
For example its leading $3\times3$ block has diagonal
$(3/4,1/2,3/4)$ and successive leading minors
$3/4,5/16,3/16$; the fourth diagonal is $1$.
This supplies a finite real Gram factorization.  The remaining
$-4\kappa\nu uvWZ$ is a separated square on this pure-testing sign
chart.  The zero-state boundaries use the same identities.
Since $\lambda_0\sim1/2$ and $\lambda_1\ge0$,
$\lambda-1/4$ is a separated SOS multiplier:
\[
 \lambda-\tfrac14=(\lambda_0-\tfrac14)
       +\lambda_0(2ra+r^2a^2)+\lambda_1c(1+ra)^2.
\]
Thus $\lambda\Psi$ leaves strict real diagonal $1/16$ on each of
$E_{12},E_{23}$ and an explicitly retained separated SOS remainder.
No mixed square root of $\lambda$ has been used.

Perform monic polynomial division in $a$ and then in $c$.
By EL17 there are bounded coefficient polynomials such that
\[
\begin{aligned}
 J&=(a-x)J_a+(c-q)J_c,\\
 H&=(a-x)^2H_{20}(a,c)
       +(a-x)(c-q)H_{11}(c)+(c-q)^2H_{02}(c).
\end{aligned}                                                     \tag{EL19}
\]
The division is by monic polynomials; it introduces no inverse of
$a-x$, $c-q$, or any mixed expression.  The following four particular
coefficient identities are exact:
\[
\begin{aligned}
 H_{20}(a,0)&=-\kappa\nu\ell S\lambda_0r^2,\\
 H_{11}(0)&=-\alpha/q+2\kappa\nu^2S\lambda_0r(1+rx),\\
 H_{02}(0)/S
   &=\nu^2\left[(1+x)\frac\alpha{S\ell^2}
                 -g\nu\lambda_0(1+rx)^2\right]\sim\tfrac12,\\
 H_{02}(c)&=H_{02}(0)+\nu^2S^2dp^2(1+rx)^2c.
\end{aligned}                                                     \tag{EL20}
\]
The first three also follow immediately by evaluating
$H(a,0)=\alpha(1+a)-\kappa\nu\ell S\lambda_0(1+ra)^2$
and using EL17; the last follows from the $c^3$ coefficient.
In particular $H_{11}(0)/(S\ell)$ is bounded and
$H_{02}(c)/S\sim1/2$.  The $g$ terms cancel in the real residue
of the third line by EL0a.  This cancellation is why the independent
ratio $S/\ell$ never enters the estimate.

Write $H_{20}=H_{20}(a,0)+c\widetilde H_{20}$ and
$H_{11}=H_{11}(0)+c\widetilde H_{11}$, with bounded coefficients.
Substitution into EL5 gives the exact scalar assembly
\[
\begin{aligned}
 F_0={}&\lambda\Psi+\widehat A\Gamma^2
                 +\frac{H_{02}(c)}S B^2\\
 &-2\sqrt S J_a\Gamma R_v-2vJ_c\Gamma B
       -g\nu\lambda_0r^2R_0^2+\widetilde H_{20}R_v^2\\
 &+\frac{H_{11}(0)}{S\ell}R_0B\\
 &+\widetilde H_{11}
       \left[\frac{c-q}{\nu^2}E_{23}^2
             +\frac{2\sqrt{\kappa q}}\nu E_{23}B-\kappa B^2\right].
\end{aligned}                                                     \tag{EL21}
\]
This includes every scalar term, including the one involving
$\lambda_1c^2a$.  It supplies strict real pivots on
$(E_{12},E_{23},\Gamma,B)$; the two radial rows have arbitrary
fixed positive pivots available at infinitesimal cost by EL9--11.
There is no discarded source-dependent signed term.

\subsubsection*{Strict bounded Gram and all-test completion}
Use the main rows
\[
 E_{12},\ E_{23},\ \Gamma,\ B,\ \widehat N=\sqrt{1+\beta}N,
 \ N_y
\]
and the three auxiliary rows $V_h,R_v,R_0$.
EL21 and EL13--16 express the entire EL0b as a bounded quadratic
form in these nine rows, together with the retained circuit SOS
remainder.  The real diagonal of the main block is bounded below
by $(1/16,1/16,1/2,1/2,1/2,1/2)$.
Its off-diagonal entries are infinitesimal except for couplings to
the auxiliary rows.  These are bounded, as are the possibly signed
auxiliary diagonal entries.

Choose a sufficiently large fixed real positive diagonal on each
auxiliary row, and allocation it using EL8--11.  This is an exact
SOS operation: if $Q-R^2$ is a displayed SOS budget, replace the
quadratic form by its difference with a real multiple of $Q-R^2$
and retain that multiple as a separate SOS.  EL8 and EL9--10 cost
only infinitesimal entries on the main rows; EL11 is allocationd
last and costs an infinitesimal entry on $R_v$ as well as on $B$.
The real auxiliary diagonals can be chosen large enough that the
real Schur complement is positive definite.  This argument uses
only a finite real matrix, never a mixed-field matrix inverse.
The resulting bounded coefficient Gram has strictly positive real
residue.  The common strict separated bounded-Gram lemma in
Lemma~\ref{p25af2c58da7e:lem:gram} makes the complete form a finite
SOS in $K_c\otimes_{\mathbb R}K_t$.  Restoring the budget and circuit
remainders proves the original full parent.

There is no restriction on $W,Z$, their ratio, or either normal;
$\beta$ was never discarded or divided by.  The actual retained
normal diagonal handles an arbitrarily large $\beta$ before any
cross-term estimate.  The only inverse involving an arbitrarily
small opening is the pure-source scaling used to define $N_y$;
its inverse never appears in a Gram coefficient.  All helper roots
are products of source and testing factors or finite sums thereof.

The local proof includes $u=0$ and $v=0$, since it never divides
by either.  Negative $u$ is reduced by
$(u,Y,Z)\mapsto(-u,-Y,-Z)$ and negative $v$ by
$(v,X,Z)\mapsto(-v,-X,-Z)$; EL0b is invariant under each map.
Real and imaginary state parts are handled separately.
By EL0a both $C,A$ are unbounded, so the adjacent-$C,A$
proof covers noninfinitesimal Alice tests.  Since $\kappa$ and $e$
are infinitesimal, the SE10--SE12 theorem covers every
homogeneous Bob sine-square residue different from $\bar\ell=0$,
including Bob infinity.  The only remaining Bob chart is exactly
finite infinitesimal $v$, proved above.  Therefore all projective
tests and all states on the original $\ell$-infinitesimal source
intersection are covered.  All original source predicates are
retained; no numerical sampling or completion marker
is an input to this conclusion.



% Source formulas: p1de5b12d1927
\subsection{Exact quotients at the doubly infinitesimal source}

\label{component:p1de5b12d1927}

\label{p1de5b12d1927:sec:new-c144-small-nu-bounded-beta}

Keep the original C14.4 source, its full native contact, and all
original coefficient signs.  The common doubly small endpoint has
\[
 S=\sigma+L\in\operatorname{Inf}_{>0},\quad\kappa=x/S\in\operatorname{Inf}_{>0},
 \quad\nu=\sigma/S\in\operatorname{Inf}_{>0},\quad
 \ell=L/S=1-\nu,\quad q=\nu/\ell,\quad n_s=\beta/q^2.
                                                               \tag{SN1}
\]
Here $n_s$ may be bounded or unbounded.  The formulas below are
exact identities; the complete BN and UN proofs supply their distinct
normalized bounds on the bounded and unbounded $\tau$ domains.
Set
\[
 c_0=SC(1-r),\quad a_0=SAr,\quad d_s=D/(Sq^2),\quad
 a_1=\alpha/S,\quad\mu=m/\kappa^2,\quad
 e=\kappa h/\nu,\quad\tau=h/q.                         \tag{SN2}
\]
The actual source-sphere identities are
\[
 d_s=\frac{\ell^2\mathcal H d_0^2}{2(P/S)J},\qquad
 n_s=\frac{\ell^2\mathcal H b_0^2}{2\kappa(P/S)\theta},\qquad
 Jd_s+\kappa\theta n_s=\frac{\ell^2\mathcal H}{2(P/S)}.   \tag{SN3}
\]
The source symbols in SN3 have the definitions in FM2--FM3 and
SE1--SE3.  In particular $e=\kappa\tau/\ell$ exactly.
\paragraph{The complete reciprocal parent.}
For infinitesimal Alice coordinate $u$ and reciprocal Bob coordinate
$w$, including either literal zero, use $D_u=1+u^2,D_w=1+w^2$ and
\[
\begin{split}
 F_w/4={}&m(W-wX+uY-uwZ)^2+CD_w(uW-Y)^2+AD_uD_wY^2\\
 &+(D+\tfrac12)(uwW+uX-wY-Z)^2
 +(\alpha+\tfrac12)D_u(Y-wZ)^2\\
 &+(\beta+\tfrac12)D_w(X+uZ)^2-\tfrac12D_uD_wZ^2.
\end{split}\tag{SN5}
\]
Keep both unrestricted state normals
\[
 X_e=X+ewW,\qquad N_y=(Y-ruW)/\sqrt{Sr(1-r)}.
 \tag{SN6}
\]
Every inverse and radical in these expressions is pure-source or
pure-testing. Write $a=u^2,c=w^2,x=\kappa S,\delta=c-q$ and set
\[
\begin{gathered}
 N=\kappa W-\ell uwZ,\quad G=uW/\sqrt S-\ell\sqrt S\,wZ,\\
 B=\sqrt S\,\delta Z,\quad B_u=u\delta Z,\quad
 R=w(a-x)Z=\frac{\kappa\sqrt S\,G-uN}{\ell}.
\end{gathered}\tag{SN7}
\]
The last equality is the exact homogeneous radial conductor.

\paragraph{Two free helper pivots.}
The exact polynomial identity
\[
 \delta^2N^2+2\kappa\ell\delta wGB
       +2\kappa\ell^2w^2B^2-\ell^2(wB_u)^2
                  =(\kappa\delta W)^2
 \tag{SN8}
\]

is an exact SOS identity.  Whenever a current strict frame controls
$N,G,B$, it supplies a prescribed real pivot on $wB_u$ at
infinitesimal Gram cost.  BN and UN give the required ordered helper
constructions even when the original $d_sB^2$ pivot is absent.

\paragraph{Exact scalar dictionary and the $1/\kappa$ cancellations.}
Put $d=D+1/2$ and define the following polynomials in $a,c$:
\[
\begin{split}
 A_P={}&(1+c)\{c_0(1-r)+a_0r(1+a)\}
       +S\{d(1-r-e)^2c+(\alpha+\tfrac12)(1+a)r^2\},\\
 A_N={}&m(1+ra+ec)^2+(\beta+\tfrac12)e^2c(1+c),\\
 K={}&m(1+ra+ec)+d(1-r-e)
       +(\alpha+\tfrac12)(1+a)r+(\beta+\tfrac12)e(1+c),\\
 F_Z={}&D+\alpha c+\beta a+(m+\alpha+\beta+\tfrac12)ac,\\
 J_1={}&K-\ell A_P,\quad L_1=J_1-(\ell/\kappa)A_N,\\
 H={}&F_Z-\ell^2ScA_P
                 +ac\{(\ell/\kappa)^2A_N-2(\ell/\kappa)J_1\}.
\end{split}\tag{SN9}
\]
These are the exact graph coefficients at $X=-ewW,Y=ruW$;
the negative $dac,e\{2(1-r)-e\}$ contribution has been placed
in $A_P$, not left in $A_N$. The displayed combination is used intact before taking any
normalized bounds.

The actual source satisfies
$L_1(x,q)=H(x,q)=H_a(x,q)=H_c(x,q)=0$ and
$H(a,0)=D+\beta a$. Finite monic division therefore gives
\[
\begin{gathered}
 L_1=(a-x)L_a+\delta L_c,\\
 H=S\delta^2(d_s+\kappa n_s+cH_{02})
        +(a-x)\delta(-qn_s+cH_{11})+c(a-x)^2H_{20}.
\end{gathered}\tag{SN10}
\]
The quotients mean literal polynomial differences at $a=x,c=q$,
not division by a mixed tensor. For example $L_a$ is
$(L_1(a,c)-L_1(x,c))/(a-x)$, and $L_c$ is the quotient of
$L_1(x,c)-L_1(x,q)$ by $c-q$.

Individually $H_{20},H_{11},L_c/\kappa$ need not be bounded.
The complete scalar uses the following combined coefficients:
\[
 C_R=H_{20}+2\ell L_a/\kappa,\quad
 C_G=\kappa(H_{11}-n_s)/\ell,\quad
 C_N=(n_s-H_{11})/\ell-2L_c/\kappa.
 \tag{SN11}
\]
Substitution of $e=\kappa\tau/\ell$ gives the exact cancellations
\[
\begin{split}
 C_R={}&-\ell\mu r(a\ell r+2c\kappa\tau+2\ell),\\
 C_N={}&n_s/\ell-Sa_0\ell r
       +(\beta+m+\tfrac12)\tau^2(c/\ell+1/\ell^2)
                   +2\kappa\mu\tau-m\tau^2/\ell,\\
 H_{02}={}&\kappa\tau^2(\beta+m+\tfrac12).
\end{split}\tag{SN12}
\]
For further pole cancellation, write
$A_P=P_0(a)+cP_1(a)$. Direct differentiation in SN9 gives
\[
\begin{split}
 \kappa H_{11}={}&2\ell^2P_1(x)+\ell^2x a_0r
 +2\ell^2\kappa\mu e(1+2rx)
 +\frac{\ell^2(\beta+1/2)e^2}{\kappa}
 -2\ell(m+\beta+\tfrac12)e\\
 &+(c+q)\frac{\ell^2(m+\beta+1/2)e^2}{\kappa}.
\end{split}\tag{SN13}
\]
These are identities of the complete coefficient combinations.
Their respective normalized estimates are proved in BN and UN.

The entire scalar graph is now
\[
\begin{split}
 F_0={}&(A_N/\kappa^2)N^2+A_PG^2+d_sB^2+n_sB_u^2
       +H_{02}(wB)^2+C_RR^2-2\sqrt S\,L_aGR\\
 &+C_GG(wB)+C_NN(wB_u).
\end{split}\tag{SN14}
\]
For verification the mixed identity is
$c(a-x)\delta Z^2=\{\kappa G(wB)-N(wB_u)\}/\ell$.
Combining this with $uN=\kappa\sqrt S\,G-\ell R$ proves
SN14 directly from SN9--SN10.

\paragraph{Both complete state normals.}
Define
\[
\begin{split}
 A_X={}&-m(1+ra+ec)+da(1-r-e)-e(\beta+\tfrac12)(1+c),\\
 T_X={}&(\ell/\kappa)cA_X+\beta-D+c(m+\beta+\tfrac12),\\
 Y_P={}&(1+c)\{-c_0+a_0(1+a)\}
 +S\{m(1+ra+ec)-dc(1-r-e)
                    +(\alpha+\tfrac12)(1+a)r\},\\
 Y_C={}&\ell Y_P-ma+d-(\alpha+\tfrac12)(1+a).
\end{split}\tag{SN15}
\]
Their actual contact values satisfy $T_X(x,q)=Y_C(x,q)=0$.
Use finite polynomial quotients
\[
 T_X=c(a-x)X_a+\delta\{(D-\beta)/q+cX_c\},\qquad
 Y_C=(a-x)Y_a+\delta Y_c.
 \tag{SN16}
\]
Although $X_a,A_X/\kappa$ separately may be unbounded, their
complete conductor has
\[
\begin{gathered}
 \widehat A_X=A_X/\kappa-aX_a/\ell
 =-\kappa\mu-\frac\tau\ell
                  \{(\beta+\tfrac12)(1+c)+mc\},\\
 \widehat X_a=\kappa X_a=\ell\{d(1-r-e)-mr\},\qquad
 X_c=-\tau(m+\beta+\tfrac12).
\end{gathered}\tag{SN17}
\]
The exact quarter-parent is $F_w/4=F_0+H_XX_e^2+H_YN_y^2+
2H_{XY}X_eN_y+2L_XX_e+2L_YN_y$, with
\[
\begin{split}
 H_X={}&mc+da+(\beta+\tfrac12)(1+c),\\
 H_Y={}&rc_0(1+c)+(1-r)a_0(1+a)(1+c)
       +Sr(1-r)\{ma+dc+(\alpha+\tfrac12)(1+a)\},\\
 H_{XY}={}&-uw\sqrt{Sr(1-r)}(m+d),\\
 L_X={}&w\widehat A_XN+
        \frac{uw\sqrt S}{\ell}\widehat X_aG
        +d_s\sqrt S\,quB-n_sqB_u+wX_c(wB_u),\\
 L_Y={}&\sqrt{r(1-r)}\{Y_PG+\sqrt S\,Y_aR+Y_c(wB)\}.
\end{split}\tag{SN18}
\]

Every formula retains the original $r$ and $1-r$ openings and both
complete state normals.  In particular, the term $n_sqB_u$ may be
written $q\sqrt{n_s}(\sqrt{n_s}B_u)$ when that row is used;
literal $n_s=0$ needs no division.  BN and UN allocate these actual
rows, with their own weights, instead of treating separate uncanceled
coefficients as bounded.

\paragraph{Explicit finite divisions.}
All quotient symbols above can be defined directly in the polynomial
ring in $a,c$ over the source field.  In addition to $L_a,L_c$
specified after SN10, use
\[
\begin{aligned}
 H_{20}&=\frac{H(a,c)-H(x,c)-(a-x)H_a(x,c)}{c(a-x)^2},\\
 H_{11}&=\frac{H_a(x,c)/(c-q)+qn_s}{c},\\
 H_{02}&=\frac{H(x,c)/(S(c-q)^2)-d_s-\kappa n_s}{c},\\
 X_a&=\frac{T_X(a,c)-T_X(x,c)}{c(a-x)},\qquad
 X_c=\frac{T_X(x,c)/(c-q)-(D-\beta)/q}{c},\\
 Y_a&=\frac{Y_C(a,c)-Y_C(x,c)}{a-x},\qquad
 Y_c=\frac{Y_C(x,c)-Y_C(x,q)}{c-q}.
\end{aligned}
\]
These displays specify polynomial division, not evaluation of an
inverse of $a-x$, $c-q$ or $c$.  Divisibility follows from
$H(x,q)=H_a(x,q)=H_c(x,q)=0$, $H(a,0)=D+\beta a$,
$T_X(x,q)=Y_C(x,q)=0$, and the explicitly affine dependence of
$T_X$ on $a$.  Thus all exact coefficients are defined without an
external program.  Pure-source divisions such as $1/S$ or $1/q$
are retained until their cancellations in the printed combinations.



% Source formulas: p80b0c5b9879f
\subsection{The doubly infinitesimal source with bounded $\tau$}

\label{component:p80b0c5b9879f}

\label{p80b0c5b9879f:sec:new-c144-bounded-h-ratio}

\paragraph{Source scope and result.}
Retain the original entered physical C14.4 source, all strict gaps and
all native source inequalities. Assume
\[
 S=\sigma+L,\ \kappa=x/S,\ \nu=\sigma/S
       \in\operatorname{Inf}_{>0},\qquad
 \ell=1-\nu,\quad q=\nu/\ell,\quad
 \tau=h/q\ \hbox{is bounded}.
 \tag{BN1}
\]
The complete original parent is a finite separated tensor SOS for
every test and every real or complex state on BN1. The quantity
$\beta/q^2$ is unrestricted, including zero and unbounded values.
The proof below is a circuit completion of SN, not an import of
a missing endpoint proof.

\paragraph{Exact source estimates.}
Use the definitions of SN2, except that $n_s=\beta/q^2$ need not be
bounded. The exact sphere SN3 remains valid. Put $b=\beta/q$ and
$z=\kappa\beta/q^2$. Then
\[
\begin{gathered}
 h=q\tau\in\operatorname{Inf},\quad\theta\sim1,\quad
 e=\kappa\tau/\ell\in\operatorname{Inf},\quad M\sim1,\\
 b=\frac{\tau(B/S)-\ell H_0}{2\theta(P/S)},\qquad
 \mu\ell^2=\frac{\tau(1+x)(1+\sigma r)}{2M(P/S)},\\
 Jd_s+\theta z=\frac{\ell^2\mathcal H}{2(P/S)}\sim\tfrac12.
\end{gathered}\tag{BN2}
\]
In particular $b,z,d_s,\mu,\tau$ are bounded, $\tau$ and $\mu$
have positive real lower bounds, and at least one of $d_s,z$ is a
positive bounded unit. The exact inequality $\beta\ge0$ supplies
the lower bound on $\tau$. Also
$c_0,a_0,a_1\sim1/2$, $D,\alpha,m,\beta$ are infinitesimal.
Although the two terms below need not have small difference before
subtraction, their displayed bounded formulas in BN2 give
\[
                  A_{00}:=\mu\ell^2-b\sim\tfrac12.
 \tag{BN3}
\]
This argument takes residues only of bounded source coefficients and
never uses an estimate divided by $q$. It allows arbitrary relative
rates of $S,\kappa,q$.

\paragraph{A retained circuit and two actual positive rows.}
Use the full reciprocal parent SN5, the two unrestricted normals SN6,
and the shifted exact scalar dictionary SN9. Let $\zeta=\ell/\kappa$,
$a=u^2,c=w^2$, $x=\kappa S$, and assume $u,w\ge0$ by the original
sign gauges. Both may be zero. For every state define
\[
\begin{array}{ll}
 A_1=\sqrt\beta\,uZ,&A_2=\sqrt\beta\,W/(\zeta\sqrt q),\\
 A_3=\sqrt\beta\sqrt{x/q}\,wZ,
 &A_4=\sqrt\beta\,uwW/(\zeta q\sqrt x),
\end{array}\tag{BN4}
\]
\[
 \Psi=\sum_{i=1}^4 A_i^2-\frac{4\beta}{\zeta q}uwWZ,
 \qquad E_1=A_1-A_2,\quad E_2=A_3-A_4.
 \tag{BN5}
\]
Every monomial and its radical is a product of a pure-source and a
pure-testing element. This remains true if $x/q$ is unbounded.
The exact inequality meant here is a separated SOS certificate:
\[
                    \Psi-\tfrac12(E_1^2+E_2^2)\in\Sigma T^2.
 \tag{BN6}
\]
Indeed, if $WZ\ge0$, then $A_1A_2,A_3A_4\ge0$ and
\[
 \Psi=E_1^2+E_2^2+
              2(\sqrt{A_1A_2}-\sqrt{A_3A_4})^2.
\]
The radicals displayed here split into pure-source and pure-testing
radicals. If $WZ<0$, use instead
\[
 \Psi-\tfrac12(E_1^2+E_2^2)
 =\tfrac12(A_1+A_2)^2+\tfrac12(A_3+A_4)^2
                  +\frac{4\beta}{\zeta q}|uwWZ|.
\]
The last positive separated monomial is a square. Thus BN6 is valid
for all states, including $WZ=0$. At the actual source contact every
$A_i$ agrees, so $\Psi$ and its first contact derivatives vanish.

\paragraph{The exact scalar after circuit subtraction.}
Write $N_0=W/\zeta-uwZ=N/\ell$, where $N$ is SN7, and keep
$G,B,B_u,R$ from SN7. Thus
\[
 R=\frac{\sqrt S}{\zeta}G-uN_0.
 \tag{BN7}
\]
Subtract the entire circuit BN5 from the exact scalar graph in SN9.
Define
\[
\begin{gathered}
 A=\zeta^2 A_N-b,\qquad A_P'=A_P-zc/\ell^2,\\
 K'=K-2\beta/(\zeta q),\quad J'=K'-\ell A_P',
 \quad L'=J'-A/\zeta,\\
 H'=F_Z-\beta a-\beta xc/q-\ell^2ScA_P'
                          +ac(A-2\zeta J').
\end{gathered}\tag{BN8}
\]
Then $F_0-\Psi=A N_0^2+A_P'G^2-2\zeta uwZL'N_0+H'Z^2$.
All coefficients of $A$ and $A_P'$ are bounded:
\[
 A=\mu\ell^2(1+ra+ec)^2
       +(\beta+\tfrac12)\tau^2c(1+c)-b,
 \qquad \overline A=\overline{A_P'}=\tfrac12.
 \tag{BN9}
\]
Here and below the residue is taken at infinitesimal $u,w$.

The circuit subtraction has especially simple exact effects on the
SN10 quotient dictionary:
\[
\begin{gathered}
 L'_a=L_a,\qquad L'_c=L_c+\kappa n_s/\ell,\qquad
 H'_{20}=H_{20},\quad H'_{02}=H_{02},\quad
 H'_{11}=H_{11}-2n_s,\\
 H'=S\delta^2(d_s+cH_{02})
               +c(a-x)\delta H'_{11}+c(a-x)^2H_{20}.
\end{gathered}\tag{BN10}
\]
In particular no coefficient $n_s$ remains on its own in the final
conductors. With $C_R$ and $L_a$ from SN11--SN13 put
\[
 C_G'=\kappa(H_{11}-2n_s)/\ell,\qquad
 C_N'=-H_{11}-2\ell L_c/\kappa.
 \tag{BN11}
\]
These are bounded. The cancellation for the second coefficient is
explicit:
\[
 C_N'=-Sa_0\ell^2r+(\beta+m+\tfrac12)\tau^2(c+1/\ell)
                 +2\ell\kappa\mu\tau-m\tau^2.
 \tag{BN12}
\]
For $C_G'$ use the bounded formula SN13 and $\kappa n_s=z$.
The exact full scalar is consequently
\[
\begin{split}
 F_0={}&\Psi+A N_0^2+A_P'G^2+d_sB^2+H_{02}(wB)^2
       +C_RR^2-2\sqrt S\,L_aGR\\
 &+C_G'G(wB)+C_N'N_0(wB_u).
\end{split}\tag{BN13}
\]
This is an identity in unrestricted state coordinates, not a
restriction to a contact graph in the full parent.

\paragraph{Free helper rows across the whole source sphere.}
First reserve the positive rows $E_1,E_2$ from BN6. If $d_s$ is a
unit, $wB$ has infinitesimal cost on $\sqrt{d_s}B$. Otherwise
$z$ is a positive bounded unit. The exact identity
\[
 E_2=-\sqrt z\,\{(w/\ell)G+
                          \sqrt S\,w(w-\sqrt q)Z\}
\]
then gives
\[
 wB=-\frac{w+\sqrt q}{\sqrt z}E_2
                   -\frac{w(w+\sqrt q)}\ell G.
 \tag{BN14}
\]
Every coefficient in BN14 is infinitesimal, so any fixed real pivot
on $wB$ is free at infinitesimal Gram cost. Only the chosen pure-source
unit $\sqrt z$ is inverted. Literal $\beta=0$ uses the first branch.

After this pivot has been allocationd, the normalized SN8 identity is
\[
 \delta^2N_0^2+\frac{2\kappa\delta}{\ell}G(wB)
          +2\kappa(wB)^2-(wB_u)^2
                         =(\delta W/\zeta)^2.
 \tag{BN15}
\]
It supplies any fixed pivot on $wB_u$ at infinitesimal cost. The order of
these two allocations is deliberate; no weak $B$ diagonal is needed
in the second step.

\paragraph{Both full normals and their only additional row.}
Use the complete $H_X,H_Y,H_{XY},Y_P,Y_a,Y_c$ from SN18 and its
bounded $\widehat A_X,\widehat X_a,X_c$ from SN17. These formulas
only require $\tau,\mu,\beta$ bounded, not $n_s$ bounded.
The entire first-normal part is
\[
\begin{split}
 L_X={}&\ell w\widehat A_XN_0
       +\frac{uw\sqrt S}{\ell}\widehat X_aG
       +d_s\sqrt S\,quB-bB_u+wX_c(wB_u),\\
 L_Y={}&\sqrt{r(1-r)}\{Y_PG+\sqrt S\,Y_aR+Y_c(wB)\}.
\end{split}\tag{BN16}
\]
The apparent extra $B_u$ coupling is paid by the actual circuit row
$E_1$, using the exact, zero-face-safe identity
\[
 bB_u=-(w+\sqrt q)\sqrt b\,E_1-b(w+\sqrt q)N_0.
 \tag{BN17}
\]
All coefficients on the right have zero residue because $b$ is
bounded. No inverse of $b$ or $\beta$ is used. In the $B$ term in
$L_X$, write $d_sB=\sqrt{d_s}(\sqrt{d_s}B)$; this is likewise
face-safe when $D=0$. All direct normal coefficients are
infinitesimal, and $H_X,H_Y\sim1/2$, $H_{XY}$ is infinitesimal.
There is no missing $X$ or $Y$ normal and no imposed state ratio.

\paragraph{Allocation and all-tests conclusion.}
Keep a fixed positive fraction of each circuit row in BN6.
Use the strict base rows $(N_0,G,X_e,N_y,\sqrt{d_s}B,E_1,E_2)$,
omitting any zero row. Their main four diagonal residues are $1/2$;
the last three have fixed positive pivots wherever their rows are
nonzero. Terms containing $R$ are infinitesimal matrix perturbations
by BN7. Purchase sufficiently large fixed pivots on $wB,wB_u$ by
BN14--BN15, and use them for the bounded cross terms in BN13 and
BN16. Formula BN17 pays the remaining first-normal row at
infinitesimal cost. Every allocation is taken from an explicitly
retained positive square or from the exact SOS remainder BN6.
The resulting bounded Gram has strictly positive definite real
residue. The strict bounded-Gram lemma therefore gives
finite tensor SOS for SN5 with both unrestricted normals restored.

This proves all infinitesimal $u,w$, including the literal axes.
For noninfinitesimal Alice use the unchanged adjacent-$C,A$ theorem
CE1--CE10; $C,A$ are unbounded on BN1. For the other Bob charts use
the small-$e$ unmatched theorem SE10--SE12, since
$\overline\ell=1$. The sign gauges and real/imaginary state split
are the original exact ones. The local argument already handles
either sign of $WZ$, so no state normalization or state-chart
exclusion is required. Both literal $D=0$ and $\beta=0$ are covered
whenever they intersect BN1.

\paragraph{The other source cells.}
The unbounded-$h/q$ source and the $\ell$-infinitesimal endpoint are proved in Sections~\ref{component:peaa56c1a3c7a} and~\ref{component:pb96c9bfcc12d}. Their source partition with the present bounded-$h/q$ theorem is printed in Section~\ref{component:pfd0dd7bb10cf}.



% Source formulas: peaa56c1a3c7a
\subsection{The doubly infinitesimal source with unbounded $\tau$}

\label{component:peaa56c1a3c7a}

\label{peaa56c1a3c7a:sec:new-c144-unbounded-h-ratio}

\paragraph{Exact scope.}
Retain the full original entered C14.4 source, all its strict gaps and
native source inequalities. Suppose $S,\kappa,\nu$ are positive
infinitesimals, $\ell=1-\nu$, $q=\nu/\ell$, and
\[
 \tau=h/q\ \hbox{is unbounded},\qquad
 \theta=1-h>0,\quad e=\kappa h/\nu\in\operatorname{Inf},\quad
 \zeta=\ell/\kappa.
 \tag{UN1}
\]
This section proves finite separated tensor SOS for the complete
parent at every projective test and every state on UN1. The opening
$\theta\in\operatorname{Inf}$ is included. There is no state-ratio
restriction or source/testing relative-rate assumption.

\paragraph{Normalized source estimates and a canceled radial mass.}
Keep $c_0=SC(1-r)$, $a_0=SAr$, $d_s=D/(Sq^2)$,
$a_1=\alpha/S$, $P_0=P/S$, and the exact SE2--SE3 source sphere.
Define
\[
 d=D+\tfrac12,\quad
 n=\beta\theta^2+2Dh+\tfrac12h^2>0,\quad
 z_n=\kappa n/q^2,\quad A_{00}=m\zeta^2-n/q.
 \tag{UN2}
\]
Here $n$ is a scalar coefficient, not an omitted event weight.
It is bounded even when $\beta$ is unbounded. The exact source gives
\[
\begin{gathered}
 m\zeta^2=\frac{\tau(1+x)(1+\sigma r)}{2MP_0},\qquad
 m\zeta\in\operatorname{Inf},\quad D,\alpha,m\in\operatorname{Inf},\quad
 c_0,a_0,a_1\sim\tfrac12,\\
 e\tau=\frac{\ell\{eH_0+h\mathcal H b_0^2\}}{B/S}\ \hbox{is bounded},\\
 z_n=\frac{\ell^2\mathcal H\theta b_0^2}{2P_0}
 +\frac{\ell^2(eH_0+h\mathcal H b_0^2)}{2(B/S)}
 +2\kappa S d_s h.
\end{gathered}\tag{UN3}
\]
Consequently $d_s,z_n$ are bounded and
$d_s+z_n\sim1/2$; at least one is a positive bounded unit. This
includes literal $D=0$. Literal $\beta=0$ has no points on UN1,
because $\beta\ge0$ has the exact boundary equation
$\tau(B/S)=\ell H_0$ when $\beta=0$.

For completeness, the following finite source division proves the
bound on $A_{00}$ without dividing an error by $\theta$ or $q$.
Put $U=(1+x)(1+\sigma r)$ and
$M_*=\ell(1+xr)-\kappa\theta=\ell M$, and form the bounded
source polynomial
\[
 T=\ell JU+JM_*P_0-JM_*\kappa\theta(J-\sigma)-2M_*U.
\]
In this expression $P_0,J,U,M_*$ are their exact native polynomial
formulas in $\kappa,\nu,S,h,r$, with $\ell=1-\nu$. Thus
\[
\begin{gathered}
 T=\kappa T_\kappa+\ell\nu S T_\nu,\qquad
 T_\nu=r-1-\nu-\nu S(1-r)(r+\ell),\\
 T_\kappa=\{T-T|_{\kappa=0}\}/\kappa
             \quad\hbox{is a bounded source polynomial},\\
 A_{00}=\frac{\theta\ell H_0}{2P_0}
            +\frac{eT_\kappa+h\ell S T_\nu}{2JMP_0}.
\end{gathered}\tag{UN4}
\]
The quotient is a finite monic polynomial division in a pure source
variable. All denominator factors are positive bounded units.
Hence $A_{00}$ is bounded and $A_{00}-\theta/2$ is infinitesimal.
Moreover $n/q=m\zeta^2-A_{00}\sim\tau/2$ is positive unbounded.
If $\theta$ is infinitesimal, then $n\sim1/2$ and
$\kappa/q=e\ell/h$ is infinitesimal. These facts supply the two
source subcases used below.

\paragraph{The actual circuit and three reserved positive rows.}
Use the complete original reciprocal quarter-parent SN5. For
infinitesimal $u,w\ge0$, including zero, retain the unrestricted
normals
\[
 N_x=X+huZ,\qquad N_y=(Y-ruW)/\sqrt{Sr(1-r)}.
 \tag{UN5}
\]
Put $a=u^2,c=w^2,x=\kappa S,\delta=c-q$ and
\[
\begin{gathered}
 N_0=W/\zeta-uwZ,\quad G=uW/\sqrt S-\ell\sqrt S\,wZ,\\
 B=\sqrt S\delta Z,\quad B_u=u\delta Z,\quad
 R=w(a-x)Z=\sqrt S\,G/\zeta-uN_0,\quad
 G_h=h(a-xc/q)Z.
\end{gathered}\tag{UN6}
\]
Define $A_1,A_2,A_3,A_4$ by BN4 with $\beta$ replaced by the
positive $n$, and set
\[
 \Psi_n=\sum_iA_i^2-4nuwWZ/(\zeta q),\qquad
 E_1=A_1-A_2,\quad E_2=A_3-A_4,\quad E_3=A_1-A_3.
 \tag{UN7}
\]
The explicit separated AM--GM square decompositions in BN6 give
\[
             \Psi_n-\tfrac14(E_1^2+E_2^2+E_3^2)\in\Sigma T^2.
 \tag{UN8}
\]
To verify the simultaneous reserve, average BN6 with the pairing
$\Psi_n-(A_1-A_3)^2-(A_2-A_4)^2\in\Sigma T^2$.
For the latter pairing the products $A_1A_3,A_2A_4$ are nonnegative;
if $WZ<0$, the additional signed cross term is a positive separated
monomial. Thus UN8 holds for both signs of $WZ$, as well as its
zero face. No radical of a general mixed sum is used.

\paragraph{Exact scalar subtraction and radial cancellation.}
The complete scalar graph at $N_x=N_y=0$ is
\[
 F_0=m(1+ra)^2W^2+A_PP^2+F_ZZ^2-2KuwWZ,
 \quad P=uW/\sqrt S,
 \tag{UN9}
\]
where
\[
\begin{split}
 A_P={}&(1+c)\{c_0(1-r)+a_0r(1+a)\}
       +S\{d(1-r)^2c+(\alpha+\tfrac12)(1+a)r^2\},\\
 K={}&m\theta(1+ra)+d(1-r)(1+ha)
                         +(\alpha+\tfrac12)(1+a)r,\\
 F_Z={}&D+\alpha c+na+
          \{m\theta^2+\alpha+(\beta+\tfrac12)\theta^2\}ac
                                                 +dh^2a^2.
\end{split}\tag{UN10}
\]
Subtract $\Psi_n$, put
\[
\begin{gathered}
 A=m\zeta^2(1+ra)^2-n/q,\quad A_P'=A_P-z_nc/\ell^2,\\
 K'=K-2n/(\zeta q),\quad J'=K'-\ell A_P',\quad
 L'=J'-A/\zeta,\\
 H'=F_Z-na-nxc/q-\ell^2ScA_P'+ac(A-2\zeta J').
\end{gathered}\tag{UN11}
\]
The exact residual graph is $A N_0^2+A_P'G^2-
2\zeta uwZL'N_0+H'Z^2$. At the actual source contact
$L'(x,q)=H'(x,q)=H'_a(x,q)=H'_c(x,q)=0$.
Also $H'(a,0)=D+dh^2a^2$. Finite monic contact division gives
\[
\begin{gathered}
 L'=(a-x)L_a+\delta L_c,\\
 H'=(D+dh^2x^2)\delta^2/q^2
       +(a-x)\delta(-2dh^2x/q+cH_{11})
       +(a-x)^2(dh^2+cH_{20}).
\end{gathered}\tag{UN12}
\]
The otherwise possible third-$c$ quotient is exactly zero.
Define $P_1(a)$ to be the coefficient of $c$ in $A_P$. The complete
combined coefficients have short bounded formulas:
\[
\begin{split}
 L_a={}&m\theta r+dh(1-r)+(\alpha+\tfrac12)r
       -\ell\{(1+c)a_0r+S(\alpha+\tfrac12)r^2\}\\
       &-m\zeta r\{2+r(a+x)\},\\
 C_R={}&H_{20}+2\zeta L_a=-m\zeta^2r(2+ra),\\
 C_G={}&H_{11}/\zeta=2\ell P_1(x)+\ell x a_0r-2z_n/\ell,\\
 C_N={}&-H_{11}-2\zeta L_c=-Sa_0\ell^2r.
\end{split}\tag{UN13}
\]
Only $C_R$ may be unbounded separately. Its entire contribution
must be combined before estimating: $A+aC_R=A_{00}$ exactly.
Using $R=\sqrt S\,G/\zeta-uN_0$ therefore gives
\[
\begin{split}
 F_0={}&\Psi_n+A_{00}N_0^2+\widetilde A_PG^2
             +2u\sqrt S\,(L_a-C_R/\zeta)N_0G\\
 &+d_sB^2+dG_h^2+C_GG(wB)+C_NN_0(wB_u),\\
 \widetilde A_P={}&A_P'+SC_R/\zeta^2-2SL_a/\zeta\sim\tfrac12.
\end{split}\tag{UN14}
\]
Every coefficient in UN14 is bounded after this combination:
$C_R/\zeta=-m\zeta r(2+ra)$ is infinitesimal and
$C_R/\zeta^2=-mr(2+ra)$ is bounded. The $N_0G$ coefficient
is infinitesimal. No estimate of an unbounded isolated radial
coefficient is used.

\paragraph{Complete $X$ conductor and its exact positive completion.}
The full $N_x$ hard coefficient is
$H_X=mc+da+(\beta+1/2)(1+c)$. Its complete half-linear
coefficient, after using the actual contact identity, is
\[
\begin{split}
 L_X={}&-m\zeta wN_0+d\tau a(wN_0+B_u)
        +uw\sqrt S\{d(1-r-e)-mr\}G+t_0B_u,\\
 t_0={}&\frac{D+h/2-\beta\theta}{q}
       =D/q+\frac{\ell H_0-\kappa\tau(J-\sigma)}{2P_0}
          \sim\tfrac12.
\end{split}\tag{UN15}
\]
The exact source identity in the second line cancels the potential
$1/q$ pole. Here $D/q=S q\,d_s$ is infinitesimal. The single
unbounded $\tau$ term is removed with the actual square $dG_h^2$,
not estimated. Namely
\[
 G_h=-\tau u(wN_0+B_u)+ew\sqrt S\,G,\quad
 N_x+\tau(wN_0+B_u)=X+ewW,
\]
\[
\begin{split}
 &dG_h^2+daN_x^2+2d\tau a(wN_0+B_u)N_x\\
 &\quad=da(X+ewW)^2
       -2de\tau uw\sqrt S\,(wN_0+B_u)G+de^2Sw^2G^2.
\end{split}\tag{UN16}
\]
The remainder has bounded infinitesimal coefficients on $N_0,G,wB_u$
because $e\tau$ is bounded. Retain the entire positive square on the
right and also the whole term $\beta(1+c)N_x^2$. The remaining
$N_x$ diagonal is $mc+(1+c)/2\sim1/2$. The remaining $L_X$ is
UN15 with its $d\tau a$ term removed; all its direct coefficients
are infinitesimal except the bounded coefficient $t_0$ on $B_u$.
The construction thus includes arbitrarily large $\beta$.

\paragraph{The unrestricted $Y$ normal and cross coefficient.}
Put
\[
\begin{split}
 Y_P={}&(1+c)\{-c_0+a_0(1+a)\}
 +S\{m(1+ra)-d(1-r)c+(\alpha+\tfrac12)(1+a)r\},\\
 Y_C={}&\ell Y_P-m\theta a+d(1+ha)-(\alpha+\tfrac12)(1+a).
\end{split}\tag{UN17}
\]
The actual source has $Y_C(x,q)=0$, so finite monic division gives
$Y_C=(a-x)Y_a+\delta Y_c$ with bounded polynomial coefficients.
The complete remaining entries are
\[
\begin{split}
 L_Y={}&\sqrt{r(1-r)}\{Y_PG+\sqrt S\,Y_aR+Y_c(wB)\},\\
 H_Y={}&rc_0(1+c)+(1-r)a_0(1+a)(1+c)\\
 &+Sr(1-r)\{ma+dc+(\alpha+\tfrac12)(1+a)\}\sim\tfrac12,\\
 H_{XY}={}&-uw\sqrt{Sr(1-r)}(m+d)\in\operatorname{Inf}.
\end{split}\tag{UN18}
\]
In particular $Y_P$ is infinitesimal. Every coefficient is bounded
through $r=0$ and $r=1$ openings; no inverse of $r$ or $1-r$ remains
after the displayed pure-source normal scaling. UN14--UN18 account
for all terms of SN5, including both independent state normals.

\paragraph{Free pivots, including the $\theta$-infinitesimal opening.}
Reserve the three actual circuit rows in UN8. In either sphere
chart, $wB$ is free at infinitesimal cost: use $\sqrt{d_s}B$ when
$d_s$ is a unit, and otherwise use the exact identity
\[
 wB=-\frac{w+\sqrt q}{\sqrt{z_n}}E_2
                 -\frac{w(w+\sqrt q)}\ell G.
 \tag{UN19}
\]
If $\theta$ is a unit, $A_{00}$ already supplies a strict $N_0$
pivot. Since $q/n$ is infinitesimal, the exact identity
\[
 B_u=-(w+\sqrt q)\sqrt{q/n}\,E_1-(w+\sqrt q)N_0
 \tag{UN20}
\]
supplies any fixed $B_u$ pivot at infinitesimal cost. Then $wB_u$ is
free by multiplication by the pure-testing infinitesimal $w$.

If $\theta$ is infinitesimal, first supply the missing $N_0$ pivot
from the actual circuit, using
\[
 N_0=\sqrt{q/n}(E_2-E_1)
       -\frac{w-\sqrt q}{\sqrt n}E_3
       +\frac{\sqrt{\kappa/q}}\ell\,wG.
 \tag{UN21}
\]
All coefficients are infinitesimal: $n\sim1/2$ and $\kappa/q$
is infinitesimal by UN4. Identity UN21 holds without requiring
$z_n$ to be a unit; the apparent source factor in UN19 has
canceled exactly before estimating. After this allocation, UN20
supplies $B_u$, and multiplication by $w$ supplies $wB_u$. Thus the
order of allocations involves no circular use of a missing pivot.
The bounded, possibly signed $A_{00}$ in this source subcase
has zero residue and is absorbed after the strict $N_0$ pivot is
obtained. No inverse of $\theta$ is introduced.

\paragraph{Final allocation and unchanged testing joins.}
Use the strict rows $G,N_x,N_y$, the applicable $N_0$ pivot, and
fixed retained fractions of $E_1,E_2,E_3,\sqrt{d_s}B$ wherever
nonzero. Purchase sufficiently large real pivots on $B_u,wB,wB_u$
by UN19--UN21. Every allocation has infinitesimal Gram cost on
the preceding strict rows. All signed terms in UN14--UN18 are
then either infinitesimal matrix perturbations or bounded
couplings to those freely allocationd helper pivots. A fixed
sufficiently large choice makes the complete residue matrix
strictly positive definite. The separated strict bounded-Gram
lemma proves finite SOS. Restore the positive terms retained
in UN8 and UN16, including $\beta(1+c)N_x^2$; nothing from the
original parent is omitted.

This proves all infinitesimal $u,w$, their literal zero faces and
arbitrary states. Noninfinitesimal Alice tests use CE1--CE10 because
$C,A$ are unbounded. The other Bob tests use the SE10--SE12
unmatched certificate, since $e$ is infinitesimal and
$\overline\ell=1$. The original sign gauges and real/imaginary state
split are exact; the circuit argument treats both signs of $WZ$.
These joins give the claimed all-tests theorem on UN1.

The theorem just proved has precisely the source domain UN1. Its join with the bounded-$h/q$ theorem and the other C14.4 source cells is printed in Section~\ref{component:pfd0dd7bb10cf}.



\section{Completion of the C1 certificate}
The source partition in Section~\ref{sec:c1-source-partition} is
exhausted. For unbounded $\eta$ and bounded $z=1/(x\zeta)$,
Theorem~\ref{thm:c1-unification-bounded-z-all-H} gives the local
certificate for every $H=h/\sigma$; the unit-Alice and complementary
projective charts cover the other tests. When $L$ is infinitesimal,
the positive-$j$ completions, the completed $J$ relation and the five
normalized source domains give the remaining certificates. Each
represents the full state quadratic form, with its positive event
terms restored.

The real statement applies separately to the real and imaginary state
parts because the parent is a real quadratic form. Source and testing
homogenization factors are pure-field squares. Restoring the original
positive $d$ therefore proves the complete projective C1 assertion.
